% =====================================================================
%  Non-GPU AI Accelerators: Lecture Notes
%  Companion notes to the 37-slide deck "Non-GPU AI Accelerators".
%  One subsection per slide: what the slide shows, the details behind
%  it, definitions of every term, and the intuition.
%
%  Compile:  pdflatex, makeindex, pdflatex, pdflatex
%  Needs only standard packages (mathptmx, tcolorbox, imakeidx, ...).
% =====================================================================
\documentclass[11pt,letterpaper]{article}
\usepackage[T1]{fontenc}
\usepackage{mathptmx}
\usepackage[scaled=0.9]{helvet}
\usepackage{courier}
\usepackage[margin=1in,headheight=14pt]{geometry}
\usepackage{amsmath,amssymb,bm}
\usepackage{booktabs,array}
\usepackage{xcolor}
\usepackage{microtype}
\usepackage{enumitem}
\usepackage{titlesec}
\usepackage[most]{tcolorbox}
\usepackage{imakeidx}
\usepackage{fancyhdr}
\PassOptionsToPackage{hyphens}{url}
\usepackage[hidelinks,bookmarksnumbered]{hyperref}
\usepackage{xurl}

% ---------- palette (same as the deck) --------------------------------
\definecolor{ink}{HTML}{1F2937}
\definecolor{muted}{HTML}{6B7280}
\definecolor{tok}{HTML}{5F6F85}
\definecolor{emb}{HTML}{0CA678}
\definecolor{posc}{HTML}{7048E8}
\definecolor{attn}{HTML}{E8590C}
\definecolor{mlp}{HTML}{1C7ED6}
\definecolor{outc}{HTML}{D6336C}
\definecolor{trn}{HTML}{9C6644}
\definecolor{vc}{HTML}{1098AD}
\colorlet{stagec}{ink}
\newcommand{\setstage}[1]{\colorlet{stagec}{#1}}

% ---------- headings ---------------------------------------------------
\titleformat{\section}{\sffamily\bfseries\Large\color{ink}}{\thesection}{0.8em}{}[{\color{stagec}\titlerule[1.2pt]}]
\titleformat{\subsection}{\sffamily\bfseries\large\color{ink}}{\thesubsection}{0.7em}{}
\titleformat{\subsubsection}{\sffamily\bfseries\color{ink}}{\thesubsubsection}{0.6em}{}
\titlespacing*{\section}{0pt}{20pt}{10pt}
\titlespacing*{\subsection}{0pt}{16pt}{6pt}

\pagestyle{fancy}
\fancyhf{}
\fancyhead[L]{\sffamily\small\color{muted}\nouppercase{\leftmark}}
\fancyhead[R]{\sffamily\small\color{muted}\thepage}
\renewcommand{\headrulewidth}{0pt}

\makeindex[title=Index of terms,intoc,columns=2]

% ---------- slide-note machinery -------------------------------------
% \slide{number}{title}{subtitle}
\newcommand{\slide}[3]{%
  \subsection[Slide #1: #2]{Slide #1: #2}%
  {\color{muted}\itshape #3\par}\medskip}
\newcommand{\para}[1]{\medskip\noindent\textbf{#1}\quad}

\tcbset{
  termsbox/.style={breakable,enhanced,colback=stagec!4,colframe=stagec!60,
    boxrule=0.5pt,arc=2pt,left=6pt,right=6pt,top=4pt,bottom=4pt,
    title=Terms,fonttitle=\sffamily\bfseries\small,coltitle=white,colbacktitle=stagec!80,
    attach boxed title to top left={yshift=-2mm,xshift=4mm},boxed title style={arc=1.5pt,boxrule=0pt}},
  intuitionbox/.style={breakable,enhanced,colback=white,colframe=white,
    borderline west={3pt}{0pt}{stagec},left=8pt,right=2pt,top=2pt,bottom=2pt,arc=0pt,boxrule=0pt,
    title=Intuition,fonttitle=\sffamily\bfseries\small,coltitle=stagec,colbacktitle=white,
    attach boxed title to top left={yshift=-1mm,xshift=0mm},boxed title style={boxrule=0pt,left=8pt}},
  formulabox/.style={enhanced,colback=ink!4,colframe=ink!30,boxrule=0.4pt,arc=2pt,
    left=6pt,right=6pt,top=4pt,bottom=4pt}
}
\newenvironment{terms}{\begin{tcolorbox}[termsbox]\begin{description}[leftmargin=1.1em,itemsep=3pt,topsep=2pt,font=\normalfont]}{\end{description}\end{tcolorbox}}
\newcommand{\term}[2]{\item[\textbf{#1}]\index{#1} #2}
\newenvironment{intuition}{\begin{tcolorbox}[intuitionbox]}{\end{tcolorbox}}
\newenvironment{formulas}{\begin{tcolorbox}[formulabox]}{\end{tcolorbox}}

\setlist[itemize]{leftmargin=1.4em,itemsep=2pt}
\setlength{\emergencystretch}{1.5em}
\setlist[enumerate]{leftmargin=1.6em,itemsep=2pt}
\newcommand{\paper}[1]{\textit{#1}}
\newcommand{\passk}{\mathrm{pass}@k}
\newcommand{\E}{\mathbb{E}}
\DeclareMathOperator{\clip}{clip}






\begin{document}

\begin{titlepage}
\thispagestyle{empty}
\vspace*{2.5cm}
{\color{attn}\rule{3cm}{3pt}}\par\vspace{0.6cm}
{\sffamily\bfseries\color{muted}\large LECTURE NOTES \ $\cdot$ \ HARDWARE, 2026\par}
\vspace{0.4cm}
{\sffamily\bfseries\fontsize{30}{34}\selectfont\color{ink} Non-GPU\\AI Accelerators\par}
\vspace{0.3cm}
{\sffamily\fontsize{18}{22}\selectfont\color{ink} Companion notes to the slides\par}
\vspace{1.2cm}
{\large\color{muted}\parbox{0.82\textwidth}{The chips built to beat the GPU, 2016--2027: Google's TPU from v1 to Ironwood and TPU 8, Cerebras's wafer-scale engines, the dataflow and streaming designs of Groq, SambaNova, Graphcore and Tenstorrent, Korea's Rebellions and FuriosaAI, the hyperscalers' own silicon (Trainium, Maia, MTIA, the OpenAI--Broadcom XPU), Intel Gaudi, Qualcomm Cloud AI, IBM Spyre and Tesla Dojo, Huawei's Ascend and China's ASIC lines, in-memory, photonic and neuromorphic machines; the principles behind them (the memory wall, systolic arrays, SRAM against HBM, static scheduling), what MLPerf and third-party benchmarks measured, the deals of 2025--26, and why most of the programs failed. One section per slide, every term defined, and a reading list for each part.}\par}
\vspace{1.5cm}
{\color{muted}For advanced undergraduates and beginning graduate students\par}
\vspace{0.3cm}
{\color{muted}November 2026\par}
\vfill
{\small\color{muted}\parbox{0.9\textwidth}{These notes follow the deck \textit{Non-GPU AI Accelerators} (37 slides) slide by slide. Numbers quoted from a conference paper, a journal article or a filing were checked against that document; numbers from product datasheets, product pages, press releases, analyst notes and press coverage are marked $\dagger$; numbers marked ``derived'' were computed from quoted ones and the formulas are given. Peaks are dense and carry their precision unless the vendor published neither, in which case the figure is quoted in the vendor's own words (``AI FLOPS'') and marked. Where a vendor has not published a figure the notes say so (``n.p.'') rather than estimate it. The account runs to November 2026: Ironwood was generally available, Trainium 3 and Maia 200 were in production, Cerebras had been public for six months, and the Ascend 950DT and TPU 8 were announced but not shipping.}\par}
\end{titlepage}

\tableofcontents
\clearpage

\section*{How to use these notes}
\addcontentsline{toc}{section}{How to use these notes}
Each section of these notes corresponds to one section of the deck, and each subsection to one slide. A subsection has four parts:
\begin{itemize}
  \item \textbf{What the slide shows}: a description of the figure or table, so the notes can be read without the slides.
  \item \textbf{Details}: the material the slide compresses, with the arithmetic written out and the numbers given their sources.
  \item A \textbf{Terms} box defining every technical term the slide uses, the first time it appears. All terms are collected in the index at the end.
  \item An \textbf{Intuition} box: the picture to keep in mind, and the question to ask when the picture seems to fail.
\end{itemize}
Conventions. An \emph{accelerator} here is one device as sold: a chip or package (a TPU, a Trainium2, an SN40L), a wafer (a Cerebras WSE), or a card (a Groq LPU, a Qualcomm AI 100). A \emph{system} is the device in its node, rack or pod with its fabric and software. Peak FLOP/s counts a fused multiply-add as two operations and is \emph{dense} with its precision stated (BF16, FP8, INT8, FP4); where a vendor quotes only an unqualified ``AI FLOPS'' or ``peak AI performance'' the notes quote it in the vendor's words, in quotation marks, with a dagger, and do not compare it with a dense figure. Sparse figures (2:4 structured sparsity, twice dense) are named as sparse. Memory is \emph{on-chip SRAM}, \emph{HBM} or \emph{DRAM} (DDR, LPDDR), and the notes always say which, because the difference is the argument of half the lecture. MiB and GiB are binary (Groq's 220 MiB); MB and GB are decimal unless a vendor's figure is quoted. Bandwidths use decimal prefixes (1 TB/s $=10^{12}$ B/s). A dagger ($\dagger$) marks any number whose source is a datasheet, product page, press release, analyst note or press coverage rather than a paper, talk at a peer-reviewed venue, or filing; ``n.p.'' means not published by the vendor as of November 2026; ``derived'' marks arithmetic on quoted numbers.

A reading plan for a single lecture: Sections 1--2 (the overview and the principles: memory wall, taxonomy, systolic arrays, SRAM against HBM) take about a quarter of the session; Sections 3--4 (the TPU, Cerebras) a quarter; Section 5 (Groq, SambaNova, Graphcore, Tenstorrent, Korea) a fifth; Sections 6--8 (the hyperscalers' chips, China, the emerging designs) a fifth; Sections 9--10 (the measurements, the deals, why most failed) the rest. The exercises in Appendix~C are arithmetic on the deck's own numbers and can be set as homework. Appendix~D is the reading list, section by section, and matches the paper kit that accompanies these notes.
% =====================================================================
\setstage{ink}
\section{Overview: a decade of alternatives (slides 1--3)}

\slide{1}{Non-GPU AI Accelerators}{Systolic arrays, dataflow engines, wafers and SRAM machines, 2016--2027}

\para{What the slide shows} The cover carries the title, a one-paragraph description of the lecture, and a drawing of three ways to build an AI chip that is not a GPU. At the top left is a small square grid of cells with arrows running across it and down it, labeled ``systolic array'': the organization of Google's TPU matrix unit, in which weights sit still in a grid of multiply-add cells while activations flow across and partial sums flow down. To its right is a disk cut into a mesh of tiny squares, labeled ``one wafer, one chip'': Cerebras's Wafer-Scale Engine, which keeps an entire 300 mm wafer as a single processor rather than dicing it into hundreds of dies. Below both is a long horizontal block divided into vertical slices with arrows running east and west through them, labeled ``data streams through SRAM slices on a fixed schedule'': Groq's tensor streaming processor, in which on-chip memory and functional units are laid out as stripes across the die and the compiler decides, to the clock cycle, when each operand moves. The description paragraph lists what the lecture covers: Google's TPUs; Cerebras, SambaNova, Groq, Graphcore, Tenstorrent, Rebellions and Furiosa; the hyperscalers' own chips; China's Ascend, Kunlun and Cambricon; what is emerging, what the benchmarks measure, and why most of these designs did not survive.

\para{Details} This is the fourth lecture in a sequence on AI hardware. The first explained how a graphics chip became a throughput processor; the second followed NVIDIA's data-center GPUs from Volta (2017) to Vera Rubin (2026); the third compared NVIDIA's platforms with AMD's and Intel's. All three were about GPUs: processors that run many threads under hardware scheduling, with large register files, caches at every level, and off-chip high-bandwidth memory. This lecture is about the chips that were designed, from about 2015 on, on the premise that a neural network does not need most of what a GPU carries, and that a processor built only for tensor arithmetic could do the same work faster, cheaper or with less energy.

The premise was not wrong. Google's first TPU, deployed in 2015 and described at ISCA in 2017, ran Google's inference workloads 15 to 30 times faster than the contemporary NVIDIA K80 GPU and Intel Haswell CPU, and 30 to 80 times better per watt (slide 6). Cerebras's wafer holds 44 GB of SRAM at about 21 PB/s, several thousand times the memory bandwidth of any GPU's HBM (slide 13). Groq's chip serves a single user's tokens faster than a GPU could at any batch size, because its weights never leave on-chip memory (slide 16). What the premise left out is the subject of the last part of the lecture: the GPU did not stand still, it absorbed the best ideas (matrix units, then narrow number formats, then very large on-package memory), and its software ran each new model the week it was published. By November 2026 most of the merchant start-ups of 2016--2019 had been sold, folded or turned into licensing businesses, while the chips that thrived were the ones whose buyer also owned the workload and the compiler: Google's TPU, Amazon's Trainium, Huawei's Ascend (slide 35).

The three drawings on the cover are three different answers to one question, which slide 4 states precisely: where should the operands of a matrix multiply wait, and who decides when they move? In the systolic array the answer is ``in the array itself, moved by the wiring'': a weight loaded into a cell is reused by every activation that passes it, so the memory system has to supply only the edges of the array (slide 7). In the wafer the answer is ``in SRAM next to every core'': there is no off-chip memory on the critical path at all, and the problem becomes making a 46{,}225 mm$^2$ piece of silicon work despite manufacturing defects (slide 13). In the streaming processor the answer is ``in SRAM slices, moved by a schedule the compiler wrote'': nothing in the hardware decides anything at run time, so latency is a constant known before the program runs (slide 16). The other designs in the lecture are variations or mixtures: SambaNova's reconfigurable dataflow unit wires compute and memory units per model (slide 17); Graphcore's IPU and Tenstorrent's Tensix mesh are many small independent cores with local SRAM (slides 18 and 19); Trainium, MTIA, Maia and Ascend are systolic or cube-shaped matrix engines with HBM (slides 21--25).

The lecture is organized in ten parts. Slides 2 and 3 give the landscape. Slides 4 to 8 develop the principles needed for everything after: the memory wall, a five-way taxonomy, what a special-purpose chip gives up, the systolic array, and the energy and capacity of SRAM against HBM. Slides 9 to 12 follow the most successful non-GPU accelerator, Google's TPU, from v1 to the eighth generation. Slides 13 to 15 cover Cerebras. Slides 16 to 20 cover the dataflow and streaming start-ups: Groq, SambaNova, Graphcore, Tenstorrent, and Korea's Rebellions and FuriosaAI. Slides 21 to 24 cover the hyperscalers' own silicon and the established vendors (Intel's Gaudi, Qualcomm, IBM, Tesla). Slides 25 to 27 cover China. Slides 28 to 30 cover in-memory, hardwired, photonic and neuromorphic designs. Slides 31 to 34 compare what was published and what was measured, and follow the money. Slides 35 to 37 draw conclusions and list sources.

Two conventions run through the whole lecture and these notes. First, peaks are dense and carry their precision: ``275 TFLOP/s BF16'' means 275 trillion floating-point operations per second at 16-bit bfloat precision, counting a multiply-add as two operations, without the factor of two that 2:4 structured sparsity would add. Several vendors in this lecture publish only an unqualified ``AI FLOPS'' figure (Cerebras's 125 ``peak AI'' PFLOP/s for the WSE-3, for example); these are quoted in the vendor's words and marked, and never compared directly with a dense figure. Second, the distinction between on-chip SRAM, on-package HBM and off-package DRAM (DDR or LPDDR) is stated every time, because the difference between them, a factor of tens in bandwidth and of hundreds in capacity, is the argument of half the lecture. The reading list for each part is in Appendix~D and matches the paper kit that accompanies these notes.

\begin{terms}
\term{accelerator}{A processor designed to run one class of computation (here, neural-network training and inference) faster or more efficiently than a general-purpose CPU; in this lecture, any such chip that is not a GPU.}
\term{GPU}{Graphics processing unit: a throughput processor that runs tens of thousands of threads in groups (warps) under hardware scheduling, with large register files, caches and off-chip memory; the reference design every chip in this lecture was built to beat.}
\term{TPU}{Tensor Processing Unit: Google's family of AI accelerators, built around systolic matrix units (MXUs), deployed since 2015.}
\term{wafer-scale integration}{Building one processor from an entire silicon wafer instead of cutting the wafer into many separate dies; Cerebras's approach.}
\term{SRAM}{Static random-access memory: fast memory built from transistors on the same die as the logic; low latency and very high bandwidth, but tens of times less dense than DRAM, so capacity per chip is limited to hundreds of megabytes (tens of gigabytes on a wafer).}
\term{dense peak}{The highest throughput a vendor states for its matrix units without counting structured-sparsity speedups; the unit of comparison in this lecture.}
\term{AI FLOPS}{A vendor's peak figure stated without a precision or a dense/sparse qualifier; quoted in quotation marks and not compared with dense figures.}
\end{terms}

\begin{intuition}
Every chip in this lecture is an answer to the same question: where should the numbers wait between multiplies, and who decides when they move? The GPU answers ``in registers and caches, decided by hardware at run time''; the alternatives answer ``in an array, in SRAM, or on a wafer, decided by wiring or by the compiler.'' Keep the question in mind and each design becomes a point in a small space rather than a new thing to memorize.
\end{intuition}

% ---------------------------------------------------------------------
\slide{2}{A decade of accelerators that are not GPUs}{Dense 16-bit peak per device by year shipped; the GPU line is the reference every one of them was built to beat.}

\para{What the slide shows} A chart on the left plots, on a logarithmic vertical axis from 20 to 40{,}000 TFLOP/s, the dense FP16 or BF16 matrix peak of one device against the year it first shipped, from 2017 to 2027. Six series share the axes, identified by a key along the top. NVIDIA's data-center GPUs form a dashed gray line, the reference: V100 at 125 TFLOP/s (2017), A100 at 312 (2020), H100 at 989 (2022), B200 at 2{,}250 (2025) and Rubin at about 4{,}000 (2026; this is NVIDIA's published 4 PFLOP/s figure, stated for FP8/FP16, since a dense BF16 peak for Rubin has not been published$^\dagger$). Google's TPUs form a solid green line just below it: v2 at 46 (2017), v3 at 123 (2018), v4 at 275 (2021), v5p at 459 (late 2023), v6e Trillium at 918 (2024) and v7 Ironwood at 2{,}307 (2025). Purple squares mark the hyperscalers' other chips: Trainium1 at 190 (2022), Trainium2 at 667 (2024), Meta's MTIA v2 at 177 (2024) and Microsoft's Maia 200 at about 1{,}270 (2026; derived, since Microsoft publishes only about 10 PFLOP/s FP4 and 5 PFLOP/s FP8 and no BF16 figure). Blue triangles mark the start-ups: Graphcore's GC2 at 125, GC200 at 250 and Bow at 350; Groq at 188; SambaNova's SN30 at 688 and SN40L at 638; Tenstorrent's Blackhole at 372; Rebellions' ATOM at 32 and REBEL at 1{,}024; FuriosaAI's RNGD at 256. Two orange diamonds high above everything else are Cerebras's WSE-2 (about 7{,}500, 2021) and WSE-3 (about 15{,}000, 2024). Two brown stars are Huawei's Ascend 910 (256, 2019) and 910C (about 800, 2025). The panel on the right, ``How to read it,'' gives the unit of account, the two lines that track the GPU, and what the chart leaves out.

\para{Details} The chart reduces each device to one number so that a decade of very different machines can be placed on one page, and the first job is to be clear about what the number is. It is the dense matrix-multiply peak at 16-bit precision (FP16 or BF16, with wider accumulation), for one package as sold, at the vendor's stated clock, with any 2:4-sparsity multiplier removed. Sixteen-bit precision was chosen because it is the one format every device on the chart supports; most of the newer parts are faster at 8 or 4 bits (Ironwood reaches 4{,}614 TFLOP/s at FP8; Maia 200 states 10 PFLOP/s at FP4), and on those formats the GPU line would also move up by two or four times. The unit is one device, and the devices are not the same size: a TPU v4 is one chip of under 600 mm$^2$, an H100 is one reticle-sized die, a B200 or Rubin is two dies in one package, a Cerebras WSE is an entire wafer of about 46{,}225 mm$^2$. The panel says this in its last line: each Cerebras point is a whole wafer. Divided by area, the wafer's peak per square millimeter is comparable to a GPU's; its advantage is elsewhere (slide 13).

The Cerebras points need a second warning. Cerebras does not publish a dense FP16 peak. Its WSE-3 release gives ``125 petaflops of peak AI performance,'' with no precision and no statement about sparsity; the figure quoted for the WSE-2 by analysts (TechInsights, from Cerebras's material) is about 7.5 PFLOP/s of FP16, with no statement at the source of whether it is dense or sparse$^\dagger$, and the WSE-3 value of about 15 PFLOP/s is derived by scaling that figure for the WSE-3's doubled per-core SIMD width at a similar core count, so it inherits the same gap. Both points are therefore marked as derived and should be read as order-of-magnitude placements. Several of the start-up points carry similar caveats: Groq's 188 TFLOP/s FP16 is from the GroqChip product brief, not the ISCA paper (which gives 820 TOPS at INT8); Tenstorrent's Blackhole is shown at 372, the Hot Chips 2024 FP16 figure for the 140-core design (745 at FP8), while the p150 cards now ship with 120 cores enabled, for which Tenstorrent's product page lists 664 TFLOP/s BLOCKFP8 (slide 19); Maia 200's 16-bit figure is derived, not published, since Microsoft's January 2026 material gives only FP4 and FP8 peaks. These are listed in full, with their sources, on slide 31.

With that in hand, the chart says three things. First, two lines track the GPU. The TPU line has stayed within about a factor of two to four of NVIDIA's per-chip peak since 2017, usually a year or so behind, and in 2025 Ironwood (2{,}307 BF16) passed the B200's dense BF16 (2{,}250) for the first time. The hyperscalers' other chips follow the same course a step further down: Trainium2 at 667 against the H100's 989 in the same year. These chips do not need to beat the GPU per chip, because their makers sell compute by the pod and by the hour, not by the chip; slide 11 shows that Google competes on 9{,}216-chip pods and slide 21 that Amazon competes on 500{,}000-chip clusters and on price.

Second, the start-ups mostly sit below the GPU line, and that is not because they failed to build fast arithmetic. The panel lists what each sold instead of peak FLOP/s: latency (Groq, Cerebras: a single user's tokens per second, slide 33), memory tiers (SambaNova: SRAM, HBM and DDR on one chip, so a trillion parameters fit on eight sockets, slide 17), openness (Tenstorrent: RISC-V cores and an open-source software stack, slide 19), and locality (Rebellions and Furiosa in Korea, Huawei's Ascend in China: a domestic supplier for a market that wants or must have one, slides 20 and 25). The chart cannot show these, which is why the next slide and the rest of the lecture exist.

Third, the GPU line is steep. From the V100 in 2017 to Rubin in 2026, NVIDIA's 16-bit peak rose from 125 to about 4{,}000 TFLOP/s, a factor of 32 in nine years (derived; Rubin's end of the range is the published FP8/FP16 figure, and a dense BF16 value, if lower, would shorten it), roughly doubling every 1.8 years, with larger factors at 8 and 4 bits. A start-up that taped out a chip twice as fast as the current GPU had about two years before the GPU caught it, and the start-up's next chip needed another two to three years and another nine-figure sum (slide 35). The chart's flattest series, Graphcore's three points from 125 to 350 between 2018 and 2022, ends where the company's independence ended (slide 18).

What the chart leaves out, the panel says, matters as much as what it shows. It leaves out memory bandwidth, where the SRAM-resident designs win by factors of tens or thousands (slides 4--8); a chip with a fifth of the GPU's peak can generate one user's tokens many times faster if its weights sit in SRAM. And it leaves out software, which decided most of the fates in this lecture: the question a buyer asked was rarely ``how many TFLOP/s?'' and usually ``does my model run on it this week, and at what fraction of peak?'' (slide 35).

\begin{terms}
\term{TFLOP/s}{Trillions of floating-point operations per second; a multiply-add counts as two operations. PFLOP/s is a thousand TFLOP/s.}
\term{FP16}{IEEE half-precision floating point: 1 sign bit, 5 exponent bits, 10 mantissa bits.}
\term{BF16}{Brain floating point, 16 bits: 1 sign, 8 exponent, 7 mantissa bits; the same range as FP32 with less precision. Introduced on TPU v2, now universal for training.}
\term{FP8, FP4}{Eight- and four-bit floating-point formats used for inference and increasingly training; each halving of width roughly doubles the throughput of a fixed matrix unit.}
\term{2:4 structured sparsity}{A pattern in which two of every four consecutive weights are zero; hardware that exploits it doubles matrix throughput. Vendors that headline sparse peaks quote twice the dense figure.}
\term{reticle limit}{The largest area a lithography stepper exposes at once, about 830 square millimeters; a conventional die cannot be larger. Cerebras stitches 84 reticle fields across one wafer.}
\term{hyperscaler}{A company operating cloud data centers at the largest scale (Google, Amazon, Microsoft, Meta); all four design their own AI chips.}
\term{pod}{A group of accelerators joined by a dedicated high-bandwidth fabric so they can work on one model; a TPU v4 pod has 4{,}096 chips, an Ironwood pod 9{,}216.}
\end{terms}

\begin{intuition}
Read the chart as a race against a moving target. Anyone can beat the current GPU on one axis for one generation; the question is whether you can do it again in two years, and whether the extra speed reaches the user through software. The two lines that track the GPU belong to companies that can afford to be second per chip because they own the workload, the pod and the price.
\end{intuition}

% ---------------------------------------------------------------------
\slide{3}{Who is who: fifteen programs, 2016--2027}{First and latest chips, the design idea, and where each program stood in November 2026.}

\para{What the slide shows} The slide is one table of fifteen rows and five columns: the program, its first chip with a date, its latest chip with a date, the design idea in a phrase, and its status in November 2026. The rows run from the oldest and largest program to the newest: Google TPU; Cerebras; Groq; SambaNova; Graphcore; Tenstorrent; Rebellions; FuriosaAI; AWS (Inferentia and Trainium); Microsoft (Maia); Meta (MTIA); Huawei (Ascend); Intel Habana (Gaudi); Qualcomm (Cloud AI); and d-Matrix. A caption below the table says that dates are first shipments or general availability where known, otherwise announcements; that business facts come from filings and press, with sources on the section slides; and that several programs are omitted from the table and treated elsewhere: Tesla Dojo (paused 2025), Nervana and Wave Computing (2020), Untether AI (2025), the photonic and analog companies (slides 28--30), and China's Kunlun, Hanguang and Cambricon (slide 27).

\para{Details} The table is the lecture's index: each row is a section or a slide later on, and its last column is the reason the program is in the lecture at all. Reading down that column sorts the programs into four groups.

\emph{The captive programs}, whose buyer is also the designer: Google, AWS, Microsoft, Meta and, in its home market, Huawei. Google's TPU began as an inference coprocessor deployed in 2015 (published at ISCA in 2017), and its latest generation is v7 Ironwood (generally available November 2025), with an eighth generation announced in April 2026 in two variants, one for training and one for inference (slide 11). Its design idea, systolic matrix units joined by an optically switched torus, has been stable for a decade (slides 7 and 10). Its status line shows how far it has traveled from a private chip: Gemini is trained and served on it; Anthropic signed in October 2025 for access to up to one million TPUs and more than a gigawatt of capacity in 2026; Meta, a competitor, rents TPUs through Google Cloud (slide 12). AWS's line began with Inferentia in 2019 and reached Trainium3 in 2025; it uses systolic NeuronCores with HBM, and its status line reads ``${\sim}$1M Trn2; 500k$+$ Anthropic'': in December 2025 AWS said more than a million Trainium2 chips were in production across its whole fleet, and that Anthropic was using over 500{,}000 of them, most in Project Rainier, which went live with about 500{,}000 in October 2025$^\dagger$ (slides 21 and 34). Microsoft's Maia 100 (2024) used Microsoft's own microscaling number formats (MX) on a reticle-sized die with an Ethernet-based scale-up fabric; Maia 200 (January 2026) runs Azure inference, and Microsoft has rights to OpenAI's chip designs (slides 22 and 23). Meta's MTIA began as a 25 W chip for its recommendation models in 2023 and moved to HBM in later generations; the status ``buys TPUs too'' is the point: even a company with its own chip buys someone else's when its own does not yet cover the workload (slide 22). Huawei's Ascend 910 (2019) was the last Ascend made at TSMC before US export rules cut Huawei off in 2020; the 910B and 910C are made in China on SMIC's 7 nm-class processes, and the CloudMatrix 384 SuperPoD made the Ascend China's default AI platform (slides 25 and 26).

\emph{The merchant start-ups that survived as companies}: Cerebras, SambaNova, Tenstorrent, Rebellions and FuriosaAI. Cerebras shipped its first wafer-scale system, the CS-1, in 2019; its current wafer is the WSE-3 (2024), and in August 2026 it announced the CS-4 ``Nexus,'' three wafers in one system. It went public in May 2026 and has a contract with OpenAI for 750 MW of capacity (slides 13--15). SambaNova's first reconfigurable dataflow unit, the SN10, shipped in 2020; the SN40L (2024) added a three-tier memory of SRAM, HBM and DDR, and an SN50 was announced for the second half of 2026; it raised \$1 billion at an \$11 billion valuation in July 2026 after sale talks with Intel stalled (slide 17). Tenstorrent's Grayskull (the table's 2021 is the silicon; cards first shipped in March 2024) and Blackhole (2025) are meshes of Tensix cores, each a cluster of small RISC-V processors around a matrix engine, programmed through an open-source stack; the table's status gives its \$693 million round of December 2024, and in October 2026 it raised a further \$1.5 billion at a \$6.2 billion valuation; it licenses its RISC-V CPU designs to others (slide 19). Rebellions' ATOM (2023) was a 5 nm inference chip; its REBEL line (2025--26) adds chiplets and HBM3E, and it merged with SK Telecom's Sapeon in December 2024 to become Korea's national champion. FuriosaAI's Warboy (2021) and RNGD (2025) are built around tensor contraction rather than matrix multiplication; in early 2025 it declined an acquisition approach from Meta reported at about \$800 million, and LG adopted RNGD for its EXAONE models (slide 20).

\emph{The merchant programs that ended or changed form}: Groq, Graphcore and Intel's Habana. Groq's tensor streaming processor (2020) holds all weights in SRAM and is scheduled entirely by its compiler; the chip was renamed the LPU without new silicon and sold as a cloud service, GroqCloud. In December 2025 NVIDIA licensed Groq's technology and hired its leadership in a deal reported at about \$20 billion; Groq the company continues to run its cloud (slide 16). Graphcore's GC2 IPU (2018) and Bow (2022) were meshes of 1{,}472 independent cores with 900 MB of SRAM; SoftBank bought the company in July 2024 for a reported \$600 million or so, against a peak private valuation of \$2.77 billion (slide 18). Intel bought Habana in December 2019 for about \$2 billion; Gaudi (2019) to Gaudi 3 (2024) combined matrix engines, programmable vector cores and Ethernet on the chip; the table's ``line ended 2025'' means that the roadmap ended (Intel cancelled Falcon Shores in January 2025, so Gaudi 3 has no successor), while Gaudi 3 itself kept shipping (slide 24).

\emph{The established vendors and the newest entrants}: Qualcomm and d-Matrix. Qualcomm's Cloud AI 100 (2021) scaled up a phone neural-processing unit with large SRAM and LPDDR memory; its AI200 (2026) and AI250 (2027) bet on memory capacity per card, with a 200 MW order from Saudi Arabia's Humain (slide 24). d-Matrix's Corsair, launched in late 2024 and shipping in 2025, performs multiplication inside its SRAM arrays (``digital in-memory compute''); its successor Raptor is to stack DRAM on the logic$^\dagger$ (Hot Chips 2025 coverage names a DRAM-stacking test vehicle, Pavehawk); it raised \$275 million at a \$2 billion valuation in November 2025 (slide 28).

Two patterns are visible before any architecture is discussed. The first is in the ``latest chip'' column: every program that is still advancing has released a chip within the last two years, and every program that stopped (Graphcore at 2022, the original Groq silicon at 2020) stopped there. Silicon cadence is the visible sign of a program that has money and customers. The second is in the ``status'' column: the strongest positions are held by programs whose customer owns the workload (Google, AWS, Huawei), or which found a large, specific buyer (Cerebras with OpenAI, Qualcomm with Humain). Slide 35 turns these two patterns into an argument.

The omissions are deliberate. Tesla's Dojo was a training machine for one company's video models; it was cut in August 2025, when the team was disbanded, and Elon Musk said on January 20, 2026, that a Dojo3 had been restarted, now aimed at space-based AI compute, so the caption calls it ``paused'' (slide 24). Intel's Nervana, the first well-funded AI chip start-up, was cancelled in January 2020 in favor of Habana, and Wave Computing, a dataflow pioneer, went bankrupt the same year (slide 35). Untether AI, an at-memory inference start-up, folded in June 2025 and its team went to AMD (slide 28). China's Baidu Kunlun, Alibaba T-Head's Hanguang and PPU, and Cambricon are covered with the export-control background on slide 27.

\begin{terms}
\term{captive chip}{An accelerator designed by the company that also runs the workload and owns the data centers (Google, Amazon, Microsoft, Meta, Huawei), and not sold as a separate product on the open market.}
\term{merchant chip}{An accelerator sold to other companies as a product, as GPUs are; most start-ups in this lecture.}
\term{inference}{Running a trained model to produce outputs (for a language model, generating tokens); distinct from training, which adjusts the weights.}
\term{training}{Adjusting a model's weights by gradient descent on data; needs more memory per parameter and much more inter-chip communication than inference.}
\term{RISC-V}{An open instruction-set architecture that anyone can implement without a license fee; Tenstorrent, Meta's MTIA and others use small RISC-V cores to control their matrix and vector engines.}
\term{dataflow}{A style of computation in which an operation runs when its inputs arrive rather than when a program counter reaches it; in this lecture, chips whose compiler wires units together according to the model's graph.}
\term{CGRA}{Coarse-grained reconfigurable array: a grid of compute and memory units with programmable connections, configured per program at the level of whole arithmetic operations rather than individual gates as in an FPGA. SambaNova's RDU is one.}
\term{export controls}{Government rules restricting the sale of chips, chip-making equipment or chip-design services to particular countries or companies; the US rules from 2019 on shaped China's accelerators.}
\end{terms}

\begin{intuition}
The table is easiest to remember as four groups: those whose buyer is also the designer (they thrived), the start-ups that found a niche and kept taping out (they survived), the start-ups whose silicon stopped (they were bought or licensed), and the newcomers betting on memory. The architecture matters, but the column that predicts the future is the one about who pays.
\end{intuition}
% =====================================================================
\setstage{tok}
\section{Principles: the memory wall, dataflow and SRAM (slides 4--8)}

\slide{4}{The memory wall: why the GPU's critics had a point}{Peak compute grew 3$\times$ every two years, DRAM bandwidth 1.6$\times$, interconnect 1.4$\times$; decode is bandwidth-bound.}

\para{What the slide shows} The chart on the left plots three exponential curves on a logarithmic axis, each normalized to 1 in 2012 and running to 2026: peak floating-point throughput of server hardware growing 3.0$\times$ every two years (gray), DRAM bandwidth growing 1.6$\times$ every two years (green), and interconnect bandwidth growing 1.4$\times$ every two years (orange). At 2026 the curves are labeled with their cumulative factors: 2{,}187$\times$, 27$\times$ and 11$\times$. The chart title calls them ``three exponentials (published rates, illustrative)'': the rates are those fitted by Gholami and co-authors in ``AI and Memory Wall'' (IEEE Micro, 2024), and the curves are drawn from the rates rather than from individual data points. The panel on the right, ``The argument every non-GPU company made,'' states in five entries why the gap matters for language-model decoding, what the GPU does about it, and the two alternatives the rest of the lecture follows.

\para{Details} The phrase \emph{memory wall} was coined in 1995 by Wulf and McKee for CPUs: processor speed was growing faster than DRAM speed, so programs would eventually spend most of their time waiting for memory. Gholami, Yao, Kim, Hooper, Mahoney and Keutzer revisited the argument for AI hardware in 2024 by fitting growth rates to twenty years of published data. They found peak server compute growing about 3.0$\times$ every two years, DRAM bandwidth about 1.6$\times$, and interconnect bandwidth (the links between chips) about 1.4$\times$. Over the 14 years on the chart these compound to very different totals (all derived):
\begin{formulas}
\[ 3.0^{14/2} = 3^7 = 2{,}187, \qquad 1.6^{7} \approx 26.8, \qquad 1.4^{7} \approx 10.5. \]
\end{formulas}
The ratio between the first and second is about 80: a chip in 2026 can do about eighty times more arithmetic per byte it reads from DRAM than a chip in 2012. In the same paper the authors report transformer parameter counts growing about 410$\times$ every two years over the period they studied, so the models were outgrowing both.

Why should the ratio matter? Any computation has an \emph{arithmetic intensity}: the number of floating-point operations it performs per byte it moves from memory. A machine has a \emph{ridge point}: its peak FLOP/s divided by its memory bandwidth in bytes per second. A computation whose intensity is above the ridge point is limited by arithmetic (compute-bound); one below it is limited by memory (bandwidth-bound) and runs at intensity $\times$ bandwidth, whatever the peak. This is the \emph{roofline model} (Williams, Waterman and Patterson, 2009). For an H100 with 989 TFLOP/s of dense BF16 and 3.35 TB/s of HBM bandwidth, the ridge point is about 295 FLOP per byte (derived). As compute grows faster than bandwidth, the ridge point rises, and more of every workload falls below it.

Language-model inference contains a workload that is far below any ridge point: \emph{decode}, the generation of output tokens one at a time. To produce one token, the model must multiply its current hidden state by every weight matrix in every layer. When only one sequence is being generated (batch size 1), each weight is read from memory once and used for exactly one multiply-add, two floating-point operations per weight. At 8 bits per weight that is two FLOP per byte; at 16 bits, one. The panel works the consequence for a 70-billion-parameter model:
\begin{formulas}
\[ \text{bytes per token} \approx \text{weight bytes} = 70\times10^{9} \text{ parameters} \times 1 \text{ byte} = 70 \text{ GB}, \]
\[ \text{tokens/s} \le \frac{\text{memory bandwidth}}{\text{bytes per token}} = \frac{3.35\times10^{12}\ \text{B/s}}{70\times10^{9}\ \text{B}} \approx 48 \ \text{(derived)}. \]
\end{formulas}
An H100 cannot generate more than about 48 tokens per second for one user of such a model, ``whatever its tensor cores do'': at 48 tokens per second the tensor cores perform about $48 \times 2 \times 70\times10^9 \approx 6.7$ TFLOP/s, under 1\% of their dense FP8 peak. (In practice the model must also read its key--value cache, the stored attention state of the sequence, so the real bound is lower.) This is the single most important number in the lecture, because it is the argument with which most of the non-GPU companies raised their money.

The panel then lists the responses. \emph{The GPU answer} is to batch: if 64 users' sequences are decoded together, each weight read from memory serves 64 multiply-adds, the arithmetic intensity rises 64-fold, and throughput per chip rises nearly as much, at the cost of each user's latency. GPUs also get more bandwidth per generation (HBM3e, then HBM4, from 3.35 TB/s on the H100 to about 8 on the B300 and about 19 on Rubin) and pool memory across GPUs with NVLink. Batching makes the GPU efficient in throughput; it does not make one user's tokens arrive faster than the bandwidth bound.

\emph{The alternative} is to put the weights where bandwidth is abundant: in on-chip SRAM, whose aggregate bandwidth on a large chip is tens of terabytes per second (80 TB/s on Groq's chip) and on a wafer about 21 PB/s. Then the bound above moves up by one to three orders of magnitude, and the limit becomes something else: how many chips the model needs (slide 8), and how fast data moves between them. Capacity, not bandwidth, becomes the constraint.

\emph{Or change the math unit.} A GPU moves every operand from a register file to an arithmetic unit and back, and a register-file access costs more energy than the multiply it feeds. Systolic and dataflow designs pass operands directly from one arithmetic unit to the next, so most multiply-adds need no register or memory access at all (slides 5 and 7). This attacks the energy side of the wall rather than the bandwidth side; at the scale of a data center, energy per operation is what sets the cost of a token.

The panel's last line gives the source. The three rates are fits, not laws: the interconnect rate in particular has risen faster since 2020 than the 1.4$\times$ fit suggests (NVLink went from 600 GB/s per GPU in 2020 to 3.0 TB/s on Rubin in 2026). But the direction of the gap is not in dispute, and every design in the lecture is a response to it.

\begin{terms}
\term{memory wall}{The widening gap between how fast processors compute and how fast memory delivers data; named by Wulf and McKee (1995), quantified for AI hardware by Gholami et al. (2024).}
\term{arithmetic intensity}{Floating-point operations performed per byte moved from memory by a computation.}
\term{ridge point}{A machine's peak FLOP/s divided by its memory bandwidth; computations with lower arithmetic intensity are bandwidth-bound on that machine.}
\term{roofline model}{A plot of attainable FLOP/s against arithmetic intensity: a sloped line (bandwidth times intensity) meeting a flat line (peak) at the ridge point.}
\term{decode}{The phase of language-model inference that generates output tokens one at a time, each requiring a pass over all the model's weights; bandwidth-bound at small batch.}
\term{prefill}{The phase that processes the whole input prompt at once; many tokens share each weight read, so it is compute-bound.}
\term{batch size}{The number of independent sequences processed together; at batch $B$ each weight read serves $B$ multiply-adds.}
\term{KV cache}{The stored key and value vectors of all previous tokens in a sequence, read at every decode step by attention; grows with context length and batch.}
\term{HBM}{High-bandwidth memory: stacks of DRAM dies placed on the same package as the processor, connected by a very wide bus; 3.35 TB/s on the H100, 7.4 TB/s on Ironwood. Generations HBM2, HBM2e, HBM3, HBM3e, HBM4.}
\term{DRAM}{Dynamic random-access memory: dense memory that stores each bit as charge on a capacitor; HBM, DDR and LPDDR are all DRAM.}
\end{terms}

\begin{intuition}
Picture a kitchen where the chef can chop eighty times faster than ten years ago but the pantry door is only twice as wide. When each dish needs every ingredient in the pantry once (decode at batch 1), the chef waits at the door. The GPU's answer is to cook many orders from each trip; the alternatives' answer is to move the pantry onto the counter.
\end{intuition}

% ---------------------------------------------------------------------
\slide{5}{A taxonomy: five ways to organize a tensor machine}{Where the operands wait and who decides when they move: in hardware, in a compiler schedule, or in a program.}

\para{What the slide shows} Five small diagrams run across the top of the slide, each with a heading and a one-line caption. The first, ``GPU (for reference),'' shows a block of streaming multiprocessors beside a tall bar labeled register file and L1, with the caption ``SIMT warps, big register file, HBM; run-time scheduling.'' The second, ``systolic array,'' shows a grid of cells with weights held in place and activations flowing through: ``weights stationary, activations flow; TPU, Trainium, MTIA.'' The third, ``spatial dataflow (CGRA),'' shows a checkerboard of compute and memory units with a path wired through them: ``units wired per model by the compiler; SambaNova, Wave.'' The fourth, ``tensor streaming,'' shows vertical slices labeled with their function (memory, vector, matrix, switch) and streams running across them: ``one static program streams data across slices; Groq.'' The fifth, ``MIMD mesh, local SRAM,'' shows a grid of small cores each with its own memory: ``thousands of cores with local memory; Graphcore, Cerebras, Tenstorrent.'' Below the diagrams are two panels. The left one, ``How to read the five,'' answers two questions for each organization: who schedules, and where the weights sit. The right one, ``The trade each one makes,'' gives the strength and weakness of each.

\para{Details} A taxonomy is useful only if it predicts behavior, and this one is built on two questions whose answers do. The first is \emph{who decides when an operand moves}. There are three possible deciders. Hardware can decide at run time, as a GPU's warp schedulers do: each clock they pick, among dozens of resident warps, one whose operands are ready, and caches decide what to keep. A compiler can decide before the program runs, producing a fixed schedule that the hardware simply follows. Or the program itself can decide, as threads on a CPU do, synchronizing with each other at points the programmer chooses. The second question is \emph{where the operands, especially the weights, wait between uses}: in a register file, in the arithmetic units themselves, in on-chip SRAM, or in off-chip DRAM.

\emph{The GPU} (for reference) answers ``hardware, at run time'' and ``registers, caches and HBM.'' A modern GPU keeps tens of thousands of threads resident and switches among them every clock to hide memory latency; this requires a very large register file (256 KB per streaming multiprocessor on recent NVIDIA parts), warp schedulers, scoreboards to track dependencies, and several levels of cache. All of this is area and energy not spent on arithmetic, and all of it is what lets the GPU run any program reasonably well. The panel's verdict, ``best at nothing on paper, adequate at everything,'' is meant as a compliment: it is the reason the GPU won most contests in this lecture.

\emph{The systolic array} answers ``hardware, by wiring'' and ``in the array.'' Weights are loaded into a grid of multiply-add cells and stay there; activations enter at one edge and pass from cell to cell; partial sums flow in the perpendicular direction and leave at the other edge. No cell decides anything: the schedule is the geometry. Slide 7 works through the mechanism. Weights that do not fit in the array wait in HBM (TPU, Trainium) or in SRAM and LPDDR (Meta's MTIA). The trade, in the right panel: the best utilization of any organization on large, dense matrix multiplies, and poor utilization on small or ragged ones, because a matrix smaller than the array leaves cells idle.

\emph{Spatial dataflow} (a coarse-grained reconfigurable array, CGRA) answers ``the compiler, by wiring.'' The chip is a checkerboard of pattern compute units and pattern memory units joined by a switched network. The compiler takes the model's computation graph and lays it out in space: one group of units performs the first matrix multiply, the next group the activation function, the next the second multiply, and the intermediate results flow from group to group through the network without returning to off-chip memory. The academic ancestor is Plasticine (Prabhakar et al., ISCA 2017), from Kunle Olukotun's group at Stanford; SambaNova, co-founded by Olukotun, commercialized it, and Wave Computing pursued a related design until 2020. Weights wait in on-chip SRAM, and on SambaNova's SN40L also in HBM and DDR tiers (slide 17).

\emph{Tensor streaming} answers ``the compiler, to the clock cycle.'' Groq's chip has no caches, no branch prediction, no arbitration and no out-of-order anything. Its functional units are arranged as vertical slices across the die (memory, vector, matrix, data-switching), and data moves east and west across the slices in streams. The compiler knows the latency of every unit and every wire, so it can schedule each instruction for an exact cycle; the hardware runs that schedule, and a 576-chip system behaves like one deterministic machine (Abts et al., ISCA 2020 and 2022). All weights wait in on-chip SRAM: 220 MiB per chip, no DRAM at all (slide 16).

\emph{The MIMD mesh with local SRAM} answers ``the program, in bulk-synchronous steps.'' The chip is thousands of small independent cores, each with its own memory and its own instruction stream (multiple instruction, multiple data), joined by an on-chip network. Graphcore's IPU (1{,}472 tiles) follows Valiant's bulk-synchronous parallel model (1990) literally: all tiles compute on local data, then all synchronize at a barrier, then all exchange data in a compiler-scheduled phase, and repeat. Cerebras's wafer (about 900{,}000 cores) and Tenstorrent's Tensix mesh are variations: Cerebras cores fire when data arrives; Tensix cores are small clusters of RISC-V processors. Weights wait in local SRAM. The trade: the most flexible of the non-GPU organizations, good on irregular and sparse work, but capacity is limited to the SRAM on the chip or wafer.

Reading the left panel's second line together with the first gives the lecture's map. Organizations that schedule in the compiler can do so only because the compiler knows where the data is; that is easiest when the data is in SRAM with fixed latency, which is why the streaming, dataflow and mesh designs are also the SRAM-resident ones. Organizations that keep weights in HBM need some tolerance of variable memory latency, which the systolic array gets from double-buffering (loading the next tile of weights while the current one is in use) and the GPU gets from switching threads. The right panel's second line names the cost of compiler scheduling: ``each new operator needs compiler work.'' When a new attention variant or mixture-of-experts router appears, the GPU runs it the day someone writes a kernel; the static machines run it when their compiler team has mapped it.

The taxonomy is simplified, and real chips mix categories. Trainium has systolic arrays and also vector and scalar engines under software control. Tenstorrent's cores contain matrix engines that are small systolic arrays. Ascend's ``cube'' unit is a three-dimensional matrix engine much like a systolic array (slide 25). NVIDIA's tensor cores are small matrix units under SIMT control (slide 7). The categories describe where each design puts its weight, not a pure form.

\begin{terms}
\term{SIMT}{Single instruction, multiple threads: the GPU execution model in which groups of threads (warps of 32 on NVIDIA) execute the same instruction on different data, scheduled by hardware.}
\term{warp}{A group of 32 threads on an NVIDIA GPU that execute in lockstep; the unit of hardware scheduling.}
\term{register file}{The fast storage next to the arithmetic units that holds each thread's operands; on a GPU it is very large so that many threads can be resident at once.}
\term{systolic array}{A grid of identical processing cells, each connected only to its neighbors, through which data moves rhythmically; each cell multiplies and adds as data passes. Named by Kung and Leiserson (1978).}
\term{weight-stationary}{A dataflow in which each weight is loaded into a processing cell and held there while many activations pass through; the TPU's MXU uses it.}
\term{static schedule}{A complete assignment of operations to hardware units and clock cycles, made by the compiler before the program runs.}
\term{tensor streaming processor}{Groq's architecture: functional units arranged as slices across the die, data streams moving between them, and every instruction scheduled by the compiler to a fixed cycle; renamed LPU (language processing unit).}
\term{MIMD}{Multiple instruction, multiple data: many cores each running their own instruction stream on their own data.}
\term{bulk-synchronous parallel}{Valiant's 1990 model of parallel computation in supersteps: local computation, a global barrier, then communication. Graphcore's IPU implements it in hardware.}
\term{double-buffering}{Loading the next block of data into one buffer while computing on another, so that memory latency is hidden behind computation.}
\end{terms}

\begin{intuition}
Two questions place any accelerator: who decides when data moves (hardware at run time, wiring, a compiler, or the program), and where the weights wait (registers, the array, SRAM or DRAM). Run-time decisions cost area and energy but forgive surprises; compile-time decisions are cheap and fast but have to be right in advance, so every new kind of model is a new problem for the compiler.
\end{intuition}

% ---------------------------------------------------------------------
\slide{6}{Why not a GPU? What a special-purpose chip gives up}{TPU v1 (2017): 15--30$\times$ a K80 and 30--80$\times$ its perf/W, by dropping what a GPU keeps.}

\para{What the slide shows} Three schematic floorplans stand side by side on the left, each labeled with what its silicon is spent on; a caption above says the drawings are illustrative and that the TPU v1 shares come from Jouppi et al. The first, ``GPU (SIMT),'' is a grid of streaming multiprocessors, each subdivided in color into scheduler, registers and math units, above a strip labeled ``L2 cache, memory controllers, I/O''; its caption says each SM carries a scheduler, registers and math units, with caches at every level. The second, ``TPU v1 (2015),'' is dominated by two blocks: the matrix multiply unit (``MXU 256$\times$256, 24\%'') and the unified buffer (``24 MiB, 29\%''), with a small accumulator block (4 MiB) and a strip for DDR3, PCIe and the activation pipeline; its caption says control takes 2\% of the die, there are no caches and no speculation, and the chip has five main instructions. The third, ``Groq TSP (2020),'' is a set of vertical stripes: SRAM of 115 MB at each side and vector and matrix slices in between; its caption says the slices span the die and data streams east and west on a compile-time schedule. A paragraph below the drawings summarizes the contrast and quotes the TPU v1 paper's numbers: 92 TOPS INT8, 28 MiB on-chip, 40 W measured, 15 to 30 times the K80 and Haswell on inference, 30 to 80 times per watt. The panel on the right, ``Given up, and gained,'' lists the three things a special-purpose chip gives up, the four things it gains, and where the bill comes due.

\para{Details} The question in the title was the one Google's engineers faced in 2013, when they estimated that if users spoke to their phones for three minutes a day, speech recognition on the CPUs of the time would double the company's data-center capacity. They had three options: buy more CPUs, buy GPUs, or build a chip. The TPU v1 paper (Jouppi et al., ``In-Datacenter Performance Analysis of a Tensor Processing Unit,'' ISCA 2017) explains why they built one, and its floorplan is the clearest illustration of what a special-purpose chip gives up.

The TPU v1 was a 28 nm chip running at 700 MHz, attached to a server over PCIe as a coprocessor. Its heart was one 256$\times$256 systolic array of 8-bit multiply-accumulate cells, 65{,}536 in all, giving
\begin{formulas}
\[ 65{,}536 \text{ MACs} \times 2 \text{ ops} \times 0.7\times10^{9}\ \text{Hz} \approx 92\times10^{12} \text{ ops/s} = 92 \text{ TOPS (INT8)}. \]
\end{formulas}
Activations were held in a 24 MiB software-managed unified buffer, results in 4 MiB of accumulators (28 MiB on chip in all), and weights streamed from 8 GiB of DDR3 at 34 GB/s. The paper gives the area breakdown: the matrix unit about 24\% of the die, the unified buffer about 29\%, and control about 2\%. On a CPU or GPU of the time, control (instruction fetch and decode, scheduling, branch prediction, reorder logic) and caches took a far larger share. The chip had a handful of CISC-style instructions issued by the host (read host memory, read weights, matrix multiply, activate, write host memory were the main five) and executed them deterministically, so that its 99th-percentile latency was predictable: Google's services had strict response-time limits, and the paper argues that the GPU's throughput advantages disappeared once a latency limit was imposed on batch size. Its thermal design power was 75 W; the paper reports about 40 W measured busy power.

Against the NVIDIA K80 GPU and an Intel Haswell server CPU on Google's six production inference workloads (two multilayer perceptrons, two LSTMs, two CNNs, which together were 95\% of Google's inference demand in 2016), the TPU was 15 to 30 times faster and 30 to 80 times better in performance per watt. Two caveats belong with those numbers. The K80 was a 2014 GPU on a 28 nm process, without the tensor cores NVIDIA added in the V100 a year after the TPU paper; and the comparison was at inference with latency limits, where the GPU's batching was most constrained. Neither caveat removes the point: by removing what it did not need, a small team built a chip several times more efficient than the best general-purpose part of its day.

The panel lists what was removed. \emph{Dynamic control}: no warp schedulers, no speculation, no scoreboards; ``a program is a static schedule the compiler proves correct.'' \emph{The cache hierarchy}: in its place, software-managed buffers whose contents and timing are known in advance; a cache spends area on tags and replacement logic and energy on lookups, to guess what a program will need, and a compiler that knows the program does not need to guess. \emph{Generality}: a fixed set of large operators (matrix multiply, convolution, a few activation functions), and ``new ops wait for the compiler.''

What is gained follows directly. More multiply-accumulate units and more SRAM per square millimeter, because the area is not spent on control and caches. ``One operand fetch feeds hundreds of multiplies,'' because a systolic array reuses every operand it reads (slide 7). Predictable latency, which matters because ``serving fleets are sized to'' their tail latency: if the 99th-percentile response time is fixed and known, fewer machines are needed to guarantee it. And better performance per watt at a given process node, which at data-center scale is better cost per inference.

The Groq floorplan on the slide takes the same reasoning further. The tensor streaming processor (14 nm, 2020) devotes much of its die to SRAM, 220 MiB in all, drawn here as two large blocks of about 115 MB at the edges, and arranges its vector and matrix units as slices between them. With no DRAM at all, it removes not only dynamic control and caches but also the off-chip memory system; the compiler schedules every data movement to the cycle (slide 16). It is the purest example of the trade: maximum speed and determinism, at the price that the whole model must fit in the SRAM of however many chips it takes (slide 8).

The panel's last entry is the lecture's thesis in advance: ``the bill comes due in software.'' Removing dynamic control from the hardware moves the scheduling problem into the compiler, and the compiler must now map every model (every PyTorch graph a customer brings) onto a static schedule that keeps the arithmetic units busy. For a fixed set of Google's own models in 2015 that was tractable. For the merchant start-ups of 2017--2022, facing customers whose models changed every few months (convolutional networks, then transformers, then mixtures of experts and long contexts), it was the hardest part of the business; ``most companies on the next slides spent more on that than on silicon.'' The GPU, whose hardware schedules at run time and whose vendor had fifteen years of CUDA libraries, kept winning ``where models change faster than compilers.'' Slide 35 returns to this as the first of five reasons most of the start-ups failed.

\begin{terms}
\term{MXU}{Matrix multiply unit: the systolic array at the core of every TPU; 256$\times$256 8-bit cells on TPU v1, 128$\times$128 BF16 cells from TPU v2 on.}
\term{unified buffer}{The TPU v1's 24 MiB software-managed on-chip memory for activations; the compiler decides what it holds.}
\term{software-managed memory}{On-chip memory (a scratchpad) whose contents are placed and removed by explicit instructions from the program or compiler, rather than by cache hardware.}
\term{cache}{On-chip memory managed by hardware that keeps copies of recently or frequently used data, guessing what the program will need next; costs area for tags and energy for lookups.}
\term{speculation}{Executing instructions before it is known that they are needed (for example past a branch), discarding the results if the guess was wrong; common in CPUs, absent in the accelerators of this lecture.}
\term{scoreboard}{Hardware that tracks which registers have pending writes so that dependent instructions wait; part of a GPU's dynamic control.}
\term{TOPS}{Trillions of (integer) operations per second; the integer counterpart of TFLOP/s.}
\term{INT8}{Eight-bit integer arithmetic; used for inference after a model has been quantized.}
\term{tail latency}{The response time exceeded by only a small fraction of requests (the 99th percentile, for example); what interactive services are designed around.}
\term{TDP}{Thermal design power: the sustained power a chip and its cooling are designed for; measured power on a real workload is often lower.}
\end{terms}

\begin{intuition}
A GPU is a factory with a large planning office that decides, every minute, which machine runs which job. A TPU or a Groq chip is a factory where the plan was made the night before and printed on the floor; the planning office is replaced by more machines. That works beautifully as long as the products do not change faster than the night shift can write new plans.
\end{intuition}

% ---------------------------------------------------------------------
\slide{7}{Systolic arrays and the MXU: reuse by geometry}{16,384 multiply-adds per clock from one operand read per row; the TPU has kept the idea for seven generations.}

\para{What the slide shows} The left half draws a 4$\times$4 weight-stationary systolic array; a caption above it notes that the TPU's MXU is 128$\times$128 or 256$\times$256. Each cell shows its stored weight ($w_{ij}$) and the operation it performs ($\times$, $+$). Activations $x$ enter from the left edge, ``skewed in time'' (each row's input arrives one clock later than the row above), and move right one cell per clock. Partial sums move down one cell per clock and leave at the bottom ``to accumulators'' as $\sum_k w_{ik}x_k$. Along the top, a note says the weights are loaded once and held in place, $N^2$ registers. Below the drawing, a formula line gives the peak: arrays $\times N^2 \times 2 \times f$, worked for TPU v1 ($1\times256^2\times2\times0.7$ GHz $=$ 92 TOPS) and TPU v4 (2 TensorCores $\times$ 4 MXUs $\times128^2\times2\times1.05$ GHz $=$ 275 TFLOP/s BF16, derived), and states that operand traffic per multiply-add falls as $1/N$, because the array needs $2N$ inputs per clock for $N^2$ multiplies. The panel on the right, ``Why it works, and where it hurts,'' has four entries: reuse by wiring, utilization, loading weights, and the GPU's convergence.

\para{Details} H.~T. Kung and Charles Leiserson described systolic arrays in 1978, and Kung's 1982 paper ``Why Systolic Architectures?'' gave the argument that still applies. In a conventional processor every operand travels from memory (or a register file) to an arithmetic unit and back, so the memory system must supply operands at the rate the arithmetic consumes them. In a systolic array, data enters at the edges and is passed from cell to neighboring cell, each cell doing a little work on everything that passes; the memory system supplies only the edges. The name comes from the heart: data pulses through the array in rhythm.

Follow one matrix--vector product $y = Wx$ through the drawn array, with $W$ a 4$\times$4 matrix. Before computation, weight $w_{ij}$ is loaded into the cell in row $i$, column $j$ (the drawing labels cells so that each column corresponds to an output). At each clock, every cell takes the activation arriving from its left, multiplies it by its stored weight, adds the product to the partial sum arriving from above, passes the activation right and the new partial sum down. Because the activations enter skewed, each partial sum meets exactly the right activations as it descends, and after the pipeline fills, one complete dot product leaves the bottom of each column every clock. For a batch of input vectors streamed one after another, the array keeps all $N^2$ cells busy. The counting is what matters:
\begin{formulas}
\[ \text{MACs per clock} = N^2, \]
\[ \text{operands crossing the edges per clock} = N \text{ (activations in)} + N \text{ (partial sums out)} = 2N, \]
\[ \frac{\text{operand traffic}}{\text{MAC}} = \frac{2N}{N^2} = \frac{2}{N}. \]
\end{formulas}
For $N=128$, the array does 16{,}384 multiply-adds per clock (the subtitle's number) while reading one activation per row; for $N=256$, 65{,}536. The weights themselves are read once per tile and then reused for every vector that passes. This is ``reuse by geometry'': the reuse is a property of the wiring, not of a cache that has to guess.

The formula line turns the cell count into a peak. TPU v1 had one 256$\times$256 array at 0.7 GHz, $65{,}536 \times 2 \times 0.7\times10^9 \approx 92$ TOPS at INT8. TPU v4 has two TensorCores per chip, each with four 128$\times$128 BF16 MXUs, at 1.05 GHz:
\begin{formulas}
\[ 2 \times 4 \times 128^2 \times 2 \times 1.05\times10^{9} = 275.3\times10^{12} \approx 275 \text{ TFLOP/s (BF16; derived, matches Google's figure)}. \]
\end{formulas}

\emph{Why it works} (the panel's first entry): each weight is used by every activation that passes through its cell, and each activation by every weight in its row. A GPU feeding 128 arithmetic lanes must read their operands from a register file through a crossbar every clock, and a register-file read costs more energy than the multiply it feeds; in a systolic array those reads become nearest-neighbor wires a few micrometers long.

\emph{Utilization is the catch.} The array is efficient only when the work fills it. A layer with fewer than $N$ outputs leaves columns idle; a batch smaller than $N$ leaves the pipeline mostly filling and draining; matrix shapes that are not multiples of $N$ leave a ragged edge of idle cells. A 256$\times$256 array doing a 300$\times$300 multiply needs four tiles, the second row and column of which are mostly empty, so it runs at about 34\% utilization (derived: $300^2 / (4\times 256^2)$). The panel cites Google's account of this as the reason TPU v2 moved from one 256$\times$256 array to several 128$\times$128 arrays: in their Hot Chips 32 talk (2020) on the design of TPU v2 and v3, Norrie et al. reported that at constant peak the smaller arrays achieve about 1.6 times higher utilization on real model shapes, and that the 128$\times$128 size balanced utilization against the wiring and control cost of more arrays (the same team wrote the design up in IEEE Micro in 2021). Google's later ``Ten lessons from three generations'' paper (ISCA 2021, on TPU v4i) draws general design lessons from the same history, but the 1.6$\times$ figure is from the Hot Chips talk, not from that paper.

\emph{Loading weights} into an $N\times N$ array takes $N$ clocks (one row per clock down the columns), so weight tiles are double-buffered: the next tile is shifted in while the current one computes. This works when each weight tile is used for many activation vectors, which is the case in training and in large-batch inference. In decode at small batch (slide 4) and in attention over short sequences, each weight tile is used for only a few vectors, the load time is not hidden, and the array starves. Inference chips therefore pair their matrix arrays with large vector units for the operations that do not fill them (softmax, normalization, small matrix products), and the later TPUs, Trainium and MTIA all do so.

\emph{GPUs converged halfway.} NVIDIA added tensor cores in the V100 (2017), each a small matrix unit that multiplies small tiles of matrices per instruction, under the control of the SIMT scheduler. Each generation since has enlarged them, and from Hopper (2022) they read operands directly from shared memory rather than registers, and from Blackwell (2024) accumulate into a dedicated tensor memory. The GPU has moved most of the way to a systolic machine for its matrix arithmetic while keeping SIMT control for everything else, which is one reason the advantage of a pure systolic design narrowed over the decade. The TPU, in turn, has kept the systolic MXU for seven generations, from v1 in 2015 to Ironwood in 2025, adding vector units, SparseCores and wider formats around it rather than replacing it (slides 9--11).

\begin{terms}
\term{multiply-accumulate}{The operation $a \leftarrow a + w \times x$; counted as two floating-point operations in peak figures. Abbreviated MAC.}
\term{processing element}{One cell of an array accelerator, typically a multiplier, an adder and a few registers. Abbreviated PE.}
\term{partial sum}{An incomplete dot product $\sum_{k\le j} w_{ik}x_k$ carried from cell to cell until all terms have been added.}
\term{skew}{Delaying each row's input by one more clock than the row before, so that data meets in a systolic array at the right cell at the right time.}
\term{utilization}{The fraction of a chip's peak arithmetic actually performed on a workload; for a systolic array, the fraction of cells doing useful work.}
\term{TensorCore}{Google's name for the compute core of a TPU chip, containing MXUs, a vector unit and a scalar unit; TPU v4 has two per chip. Distinct from NVIDIA's tensor cores.}
\term{tensor core}{NVIDIA's matrix multiply unit inside each streaming multiprocessor, introduced with the V100 in 2017; a small matrix engine under SIMT control.}
\term{vector unit}{An arithmetic unit that applies one operation to a short vector of elements per instruction; used for activation functions, normalization and other non-matrix work.}
\end{terms}

\begin{intuition}
Think of a bucket brigade instead of a crowd running to the well: each bucket (activation) is handed down the line and every person (cell) uses it once on the way. The brigade moves water with very little running, as long as the line is full; when there are only a few buckets, most people stand idle. That is the systolic array's strength and its utilization problem in one picture.
\end{intuition}

% ---------------------------------------------------------------------
\slide{8}{SRAM against HBM: bandwidth is cheap on chip, capacity isn't}{An on-chip read costs a hundredth of a DRAM access; a 70B model on chip takes hundreds of chips or a wafer.}

\para{What the slide shows} A horizontal bar chart on the left, titled ``Horowitz, ISSCC 2014 (45 nm),'' plots on a logarithmic axis the energy of six operations in picojoules: an INT8 add (0.03 pJ), an FP16 multiply (1.1), an FP32 multiply (3.7), a read from an 8 KB SRAM (10), a read from a 1 MB SRAM (100) and a DRAM read (1.3 to 2.6 nJ, drawn at about 2{,}000 pJ). The axis label reads ``energy per operation (pJ, log; reads are 64-bit)'': the arithmetic bars are for the stated widths and the memory bars are for 64-bit reads. A caption below says a DRAM access costs about 500 FP32 multiplies, and that the figures are 45 nm and 2014 but the ratios have widened since (Dally, CACM 2020). The panel on the right, ``How many chips hold a 70B model? (derived),'' works through the capacity arithmetic for six designs and ends with a rule.

\para{Details} The previous slides argued that weights should sit close to the arithmetic. This slide puts numbers on both sides of that argument: how much energy is saved by keeping data on chip, and how much capacity it costs.

The energy figures come from Mark Horowitz's keynote at ISSCC 2014, ``Computing's Energy Problem (and what we can do about it),'' which tabulated the energy of basic operations in a 45 nm process. The table has been reproduced in nearly every accelerator talk since, including Han, Mao and Dally's ``Deep Compression'' (ICLR 2016) and Sze et al.'s survey (2017). Three ratios carry the argument. An 8-bit integer add costs 0.03 pJ and a 32-bit floating-point multiply 3.7 pJ, so narrowing the arithmetic saves one to two orders of magnitude (the motivation for INT8, FP8 and FP4). Reading 64 bits from a small local SRAM costs about 10 pJ, a few times an FP32 multiply; reading from a larger 1 MB SRAM costs about 100 pJ, because the wires inside it are longer. Reading 64 bits from off-chip DRAM costs 1.3 to 2.6 nJ, or
\begin{formulas}
\[ \frac{2{,}000 \text{ pJ}}{3.7 \text{ pJ}} \approx 540 \ \text{FP32 multiplies (derived; } 350\text{--}700 \text{ over the range)}, \]
\end{formulas}
which the caption rounds to ``about 500.'' The ratio depends on the width of the read. Horowitz's table also gives a 32-bit DRAM read at about 640 pJ, which is about 640/3.7 $\approx$ 170 FP32 multiplies (derived); the caption's 500 compares a 64-bit read with a 32-bit multiply, and the like-for-like 32-bit ratio is about three times smaller, still two orders of magnitude. The subtitle's ``a hundredth'' compares a local SRAM read with a DRAM read: 10 pJ against 1{,}300--2{,}600 pJ. In other words, energy is dominated by data movement, and the further the data travels the more it costs; the arithmetic is nearly free.

Dally, Turakhia and Han updated the picture for 14 and 28 nm processes in ``Domain-Specific Hardware Accelerators'' (CACM, July 2020): local 8 KB SRAM about 50 fJ per bit, a 100 MB on-chip memory about 0.7 pJ per bit, LPDDR4 about 4 pJ per bit, DDR4 about 20 pJ per bit, a chip-to-chip SerDes link about 10 pJ per bit, and on-chip wire about 100 fJ per bit per millimeter. Arithmetic energy fell with each process generation faster than wire and off-chip energy did, so the ratios have widened. The same paper notes that executing an add as an instruction on a general-purpose core costs over 250 pJ against tens of femtojoules for the add itself, the overhead argument of slide 6. HBM sits between on-chip SRAM and DDR: it is DRAM, but the stacks are millimeters from the processor on the same package, so the energy per bit is a few picojoules and the bandwidth several terabytes per second, at a capacity of tens to hundreds of gigabytes.

The other side of the trade is capacity. SRAM cells are built from six transistors each, DRAM cells from one transistor and a capacitor, and DRAM is made on processes optimized for density; per square millimeter of silicon, DRAM holds tens of times more bits. A large accelerator die holds hundreds of megabytes of SRAM; a stack of HBM3e holds 24 to 36 GB. The panel works out what that means for a model of 70 billion parameters, which needs 70 GB of weights at 8 bits or 140 GB at 16 bits, plus key--value cache (all chip counts derived; the chips' figures are from the sections that follow):
\begin{itemize}
\item \emph{H100 (80 GB HBM, 3.35 TB/s):} one GPU at 8 bits (with little room left for KV cache), two at 16 bits.
\item \emph{Groq LPU (230 MB SRAM, no DRAM):} at least $70\times10^9 / 230\times10^6 \approx 304.3$, so 305 chips at 8 bits, for the weights alone. Groq served Llama 2 70B from 576 chips in 2024 (8 racks of 9 nodes of 8 cards); the extra chips hold activations, KV cache and the pipeline's duplication. Each chip reads its slice of the weights at 80 TB/s, so the aggregate bandwidth across 576 chips is about 46 PB/s and bandwidth is no longer the bound; the bound becomes the time for activations to cross hundreds of chips (slide 16). (Groq's 220 MiB is 230.7 MB; the panel rounds to 230 MB.)
\item \emph{Graphcore Bow (900 MB SRAM):} $70 / 0.9 \approx 78$ chips at 8 bits. Graphcore's IPU-POD systems did not do this; they streamed weights from DDR memory attached to the host, giving up much of the SRAM's bandwidth advantage for capacity (slide 18).
\item \emph{Cerebras WSE-3 (44 GB SRAM per wafer):} two wafers at 8 bits, four at 16 bits; Cerebras's inference service ran Llama 3.1 70B at 16 bits across four CS-3 systems (slide 15). For training, one wafer instead streams weights in layer by layer from an external MemoryX store (slide 14).
\item \emph{SambaNova SN40L:} 520 MB of SRAM, 64 GB of HBM and 1.5 TB of DDR5 per chip, a three-tier hierarchy; eight sockets hold a trillion parameters of DDR-resident weights, swapped into HBM and SRAM as needed (slide 17).
\end{itemize}
The panel's rule summarizes: ``SRAM buys bandwidth and latency; capacity comes from chip count, stacking or tiers, each at a cost.'' Chip count costs money, power and the latency of crossing chips (Groq). Stacking costs manufacturing complexity (d-Matrix says its Raptor will stack DRAM on logic$^\dagger$, slide 28). Tiers cost the bandwidth of the slower tier whenever data must come from it (SambaNova, and every GPU with HBM). Wafer-scale integration (Cerebras) is a fourth answer: make the chip so large that its SRAM approaches HBM capacity.

The same arithmetic explains why SRAM-resident machines are fast and why they are expensive. One user's decode step on Groq or Cerebras reads every weight from SRAM, at an aggregate bandwidth hundreds or thousands of times an H100's; the bound of slide 4 rises from about 48 tokens per second to thousands, and the measured figures (hundreds to thousands of tokens per second per user) are on slide 33. But the hundreds of chips or several wafers are dedicated to holding one copy of the model; to serve many users the system must batch, and at large batch the GPU's HBM bandwidth is shared across many users while the SRAM machine's chip count is not reduced. SemiAnalysis's 2024 analysis of Groq concluded that it had a silicon-cost advantage per token for latency-optimized serving but that NVIDIA was about ten times better in performance per dollar for throughput-optimized serving (slide 33). Both statements are consequences of the numbers on this slide.

\begin{terms}
\term{picojoule}{$10^{-12}$ joule; a femtojoule (fJ) is $10^{-15}$ joule and a nanojoule (nJ) $10^{-9}$ joule. Energy per operation or per bit moved is quoted in these units.}
\term{energy per bit}{The energy to move one bit between two places (across a wire, out of a memory, over a link); the dominant cost in modern accelerators.}
\term{SerDes}{Serializer--deserializer: the circuit that sends data between chips over a few high-speed differential wires; costs about 10 pJ per bit (Dally et al., 2020).}
\term{LPDDR}{Low-power DDR DRAM, designed for phones; cheaper and lower-power per gigabyte than HBM, with less bandwidth. Used by Meta's MTIA, Qualcomm's AI200 and d-Matrix.}
\term{DDR}{Double-data-rate DRAM on standard modules; the cheapest capacity, with tens to hundreds of GB/s per channel. SambaNova's third memory tier.}
\term{memory tier}{A level in a hierarchy of memories that trade capacity against bandwidth and energy; data moves between tiers as needed.}
\term{MemoryX}{Cerebras's external weight store, which streams weights into a wafer layer by layer during training.}
\end{terms}

\begin{intuition}
On-chip memory is a desk: everything on it is within reach, but it holds a few books. HBM is the bookshelf beside the desk; DRAM on a module is the library down the hall. The SRAM machines buy hundreds of desks, or one enormous desk the size of a wafer, so that no one ever walks to the shelf; the GPU keeps a large shelf close by and serves many readers from each trip.
\end{intuition}
% =====================================================================
\setstage{emb}
\section{The Google TPU (slides 9--12)}

\slide{9}{Google TPU v1 to v3: the systolic array goes to the data center}{A 2015 inference coprocessor became, by 2018, a liquid-cooled 1,024-chip training pod; the pod is the design unit.}

\para{What the slide shows} The upper left of the slide is a block diagram of the first Tensor Processing Unit, labeled ``TPUv1 (2015): one chip, one array, no cache.'' Six boxes are joined by arrows. A host CPU, attached over PCIe 3 x16, writes into a 24 MiB \emph{unified buffer} of on-chip SRAM; the buffer feeds a large green box, the \emph{matrix unit} of 256 $\times$ 256 8-bit multiply-accumulators; the matrix unit writes into 4 MiB of accumulators; the accumulators feed an activation and pooling unit, which writes back into the unified buffer. A separate box at the lower left, ``weights: 8 GiB DDR3, 34 GB/s,'' sends a curved arrow straight into the matrix unit. A caption under the diagram states the design's essentials: weights are loaded into the array and stay there while activations flow through; the chip delivers 92 TOPS of INT8 at 75 W and was 15--30$\times$ faster than an NVIDIA K80 on Google's 2015 inference workload (Jouppi et al., ISCA 2017); and the paper's own roofline analysis showed that the 34 GB/s DDR3 interface, not the array, was the limit.

Below the diagram a table compares the first three generations. TPU v1 was deployed in 2015 on a 28 nm process, peaks at 92 TOPS INT8, holds its weights in 8 GiB of DDR3 at 34 GB/s, draws 75 W and works alone, one chip per server slot. TPU v2, deployed in 2017 on 16 nm, peaks at 46 TFLOP/s of BF16, carries 16 GiB of HBM at 700 GB/s, draws 280 W, and is built into 256-chip pods totaling 11.5 PFLOP/s. TPU v3, deployed in 2018 on the same 16 nm process, peaks at 123 TFLOP/s BF16, carries 32 GiB of HBM at 900 GB/s, draws 450 W with liquid cooling, and is built into 1{,}024-chip pods of ``100+'' PFLOP/s (1{,}024 $\times$ 123 $\approx$ 126 PFLOP/s, derived). The right panel, ``What Google learned, and told,'' lists four lessons: remove what inference does not need; determinism makes the tail latency predictable; v2 was a different machine, designed as a pod; and v3 doubled the matrix units, added liquid cooling and grew the pod to 1{,}024 chips, with the closing lesson that a fixed array needs a compiler that pads and tiles for it.

\para{Details} The TPU began as a defensive project. In 2013 Google estimated that if its users spoke to their phones for three minutes a day, serving the speech-recognition networks on CPUs would require doubling the company's data centers. The response was a coprocessor designed, built and deployed in about fifteen months, in service from 2015, disclosed at Google I/O in May 2016 and described in detail at ISCA in June 2017 (Jouppi et al., arXiv:1704.04760). That paper is the most cited document in this lecture and the one that restarted the field of special-purpose deep-learning hardware: it showed, with production workloads and measured power, that a chip doing one thing could beat the best general-purpose processors of its day by more than an order of magnitude.

What TPU v1 does is narrow. It executes eight-bit integer matrix multiplications for inference. The heart of it is a 256 $\times$ 256 systolic array, 65{,}536 multiply-accumulate units (slide 7), clocked at 700 MHz:
\begin{formulas}
\[
\text{peak} = 65{,}536 \text{ MACs} \times 2 \tfrac{\text{ops}}{\text{MAC}} \times 0.7\ \text{GHz} \approx 92\ \text{TOPS (INT8)}.
\]
\end{formulas}
The array is \emph{weight stationary}: a 256 $\times$ 256 tile of weights is shifted into the array and stays put while the activations of many inputs stream through from the left and partial sums flow down. The 24 MiB unified buffer holds the activations; together with the 4 MiB of accumulators, that is the 28 MiB of on-chip memory the paper reports. The host issues a dozen or so CISC-style instructions (read host memory, read weights, matrix multiply, activate, write host memory), each of which runs for many cycles, so the host interface is not the bottleneck.

The paper's most quoted numbers are the comparisons with an Intel Haswell server CPU and an NVIDIA K80 GPU on six production networks (two MLPs, two LSTMs, two CNNs) that together were 95\% of Google's inference load in 2016. The TPU was 15 to 30 times faster and delivered 30 to 80 times the performance per watt, measured on total power. Two observations in the paper are less quoted and more important for this lecture. First, the gains came less from the array than from what the chip left out: no caches, no branch prediction, no out-of-order execution, no multithreading, no speculative prefetch, all of which a CPU or GPU spends transistors and energy on to make an unknown program run fast. The TPU runs a known program. Second, its own roofline (slide 4) showed most of the six workloads were memory-bound: the MLPs and LSTMs reuse each weight only as many times as there are inputs in a batch, and the 34 GB/s DDR3 could not deliver weights fast enough to keep 65{,}536 multipliers busy. The authors wrote that with GDDR5 memory the chip would have been two to three times faster. Every later TPU, and every SRAM-heavy design in this lecture, is an answer to that sentence.

The panel's second point, determinism, is a business argument as much as an engineering one. A serving fleet is provisioned for the 99th-percentile latency, not the average: a search or a voice query must return in a fixed budget. Caches and dynamic scheduling make a processor fast on average and variable at the tail; a machine with software-managed buffers and a fixed instruction stream takes the same time every time. The TPU v1 paper reports that the K80 and Haswell had to run at small batch sizes, and so far below peak, to meet the latency limit, while the TPU met it at batch sizes that kept its array busy. Groq (slide 16) made determinism its whole product.

TPU v2, announced at Google I/O in May 2017 and offered on Google Cloud from February 2018, is a different machine (Norrie et al., Hot Chips 32, 2020, written up in IEEE Micro in 2021). Training needs floating point, needs to store activations for the backward pass, and needs many chips to cooperate. Each v2 chip has two \emph{TensorCores}; each TensorCore is a VLIW (very long instruction word) scalar unit that issues instructions to a vector unit and to one 128 $\times$ 128 matrix unit (MXU). Google chose 128 rather than 256 because at constant FLOP/s, four smaller arrays are easier to keep busy on real layer shapes than one large one: about 1.6 times the utilization, by the Hot Chips 2020 account (slide 7). The multipliers take \emph{bfloat16} (BF16) inputs and accumulate in FP32. BF16, introduced here, keeps FP32's eight-bit exponent and cuts the mantissa to seven bits, so it has FP32's range and needs no loss scaling; it is now the default training format on every accelerator in this lecture. DDR3 was replaced by HBM, 16 GiB at 700 GB/s per chip, twenty times v1's bandwidth. And each chip carries an inter-chip interconnect (ICI) with its own router and four links, so 256 chips wire directly into a 16 $\times$ 16 two-dimensional torus with no switches and no host involvement: a pod of 256 $\times$ 46 TFLOP/s $\approx$ 11.8 PFLOP/s, quoted by Google as 11.5. The design unit had become the pod.

TPU v3 (May 2018) doubled the MXUs per TensorCore to two, raised the clock from 700 to 940 MHz, doubled HBM to 32 GiB at 900 GB/s, and drew 450 W, which forced liquid cooling. Four pods' worth of v2 area became one v3 pod of 1{,}024 chips in a 32 $\times$ 32 torus, ``over 100 PFLOP/s'' in Google's words (1{,}024 $\times$ 123 $\approx$ 126 PFLOP/s peak; derived). In 2019 v3 pods posted MLPerf Training results that matched the largest NVIDIA systems (slide 32).

The closing lesson on the panel is the one that made the TPU a platform rather than a chip. A fixed 128 $\times$ 128 array is efficient only when the matrices fed to it are multiples of 128 and laid out the way the array wants. Somebody has to pad, tile, fuse and schedule. That somebody is XLA, the ahead-of-time compiler Google released with TensorFlow in 2017, which takes a whole computation graph and emits a program for the VLIW cores; from 2018 it is JAX, a NumPy-like front end whose transformations trace straight into XLA. The TPU's success is inseparable from that software, and its absence is the first item in the list of reasons other programs failed (slide 35).

\begin{terms}
\term{TPU}{Tensor Processing Unit: Google's family of matrix-multiply accelerators, first deployed in 2015.}
\term{unified buffer}{TPU v1's 24 MiB on-chip SRAM, managed by software, that holds activations between layers; there is no cache.}
\term{weight stationary}{A systolic dataflow in which each processing element holds one weight while inputs stream past it; the opposite choice, output stationary, keeps the partial sum in place.}
\term{CISC instructions}{Complex instructions that each run for many cycles; TPU v1's host sends a few such instructions per layer rather than one per operation.}
\term{TOPS}{Tera (10 to the 12th) operations per second, used for integer arithmetic; a multiply-accumulate counts as two operations.}
\term{tail latency}{The response time of the slowest requests, usually the 99th percentile; serving systems are sized to it.}
\term{TensorCore (TPU)}{The TPU's compute core from v2 on: a VLIW scalar unit driving a vector unit and one or more MXUs. Not to be confused with NVIDIA's tensor cores, which are matrix units inside an SM.}
\term{VLIW}{Very long instruction word: an instruction format in which the compiler packs several independent operations into one wide instruction, so the hardware needs no dynamic scheduling.}
\term{bfloat16}{Brain floating point: a 16-bit format with an 8-bit exponent and 7-bit mantissa, introduced on TPU v2; it has the range of FP32 at half the bits.}
\term{ICI}{Inter-chip interconnect: the TPU's direct chip-to-chip links, with a router on every chip, wired into a torus without external switches.}
\term{torus}{A mesh whose edges wrap around, so every chip has the same number of neighbors; TPU v2 and v3 pods are 2D tori, TPU v4 and later 3D.}
\term{pod}{A TPU supercomputer of chips wired by ICI into one torus: 256 chips for v2, 1,024 for v3, 4,096 for v4, 9,216 for Ironwood.}
\term{XLA}{Accelerated Linear Algebra: Google's ahead-of-time compiler for whole computation graphs; the compiler for every TPU since v2.}
\term{JAX}{A Python library for NumPy-style array programs with automatic differentiation and compilation through XLA; the native TPU front end since 2018.}
\end{terms}

\begin{intuition}
TPU v1 is a lesson in subtraction: take away everything a processor spends to run unknown programs fast, keep the multipliers and a buffer, and a known program runs thirty times more efficiently. TPU v2 is a lesson in addition: training needs floating point, HBM and a network, and once the chips talk to each other directly the machine you design is the pod. The first lesson is why special-purpose chips exist; the second is why most of them could not keep up.
\end{intuition}

% ---------------------------------------------------------------------
\slide{10}{TPU v4: optical circuit switches, SparseCores, and 4,096 chips}{Mirrors route light between 64-chip cubes, so the pod's topology is chosen per job and routed around failures.}

\para{What the slide shows} The top left draws a TPU v4 pod as a row of six green squares, each a cube of 4 $\times$ 4 $\times$ 4 chips, with a note that there are 64 cubes in all. Below them a long orange bar, ``48 Palomar MEMS optical circuit switches: 136 $\times$ 136 ports each,'' is joined to every cube by a vertical line. A caption explains the wiring: inside a cube the 64 chips are connected electrically in a 3D mesh; between cubes the links are optical and are re-patched in milliseconds into a 3D torus, optionally twisted; the optical switching layer is under 5\% of the system's cost and under 3\% of its power (Jouppi et al., ISCA 2023). Two stat tiles give the headline numbers: 275 TFLOP/s of BF16 per chip on a 7 nm process, and 1.1 EFLOP/s for 4{,}096 chips in a 3D torus. The right panel, ``Three ideas in one machine,'' explains optical circuit switching, the SparseCore, and the on-chip memory lesson of the TPUv4i inference chip, and closes with the chip's numbers and dates and the observation that one XLA programming model compiled for v2 through v4 while the pods changed around it.

\para{Details} TPU v4 was deployed inside Google from 2020, offered on Google Cloud in 2022 and described at ISCA in June 2023 (Jouppi et al., arXiv:2304.01433). Each chip has two TensorCores with four 128 $\times$ 128 MXUs each, runs at 1{,}050 MHz on 7 nm, and peaks at 275 TFLOP/s BF16 (and 275 TOPS INT8; v4 does not double the integer rate). It carries 32 GiB of HBM2 at 1.2 TB/s and 128 MiB of on-chip common memory (CMEM), with a thermal design power of 192 W, well under half of v3's 450 W at more than twice the performance. A full pod is 4{,}096 chips, $4{,}096 \times 275$ TFLOP/s $= 1.126$ EFLOP/s of BF16, which the slide rounds to 1.1 EF.

The pod's novelty is how the 4{,}096 chips are joined. The building block is a cube of 64 chips, $4 \times 4 \times 4$, on 16 hosts in one rack, wired to each other with ordinary electrical links. Each face of a cube has 16 links leaving it, six faces make 96 optical links per cube, and every one of them goes to an \emph{optical circuit switch} (OCS). Google's OCS, developed in-house under the name Palomar, is a box of two arrays of tiny MEMS mirrors: light from an input fiber hits a mirror that is tilted to steer it to a chosen output fiber. Each switch has 136 ports (128 in use plus spares), and 48 of them connect the 64 cubes. Because the switch moves light rather than packets, it has no buffers, adds no protocol, and is indifferent to data rate; a link through it behaves like a long fiber. Re-aiming the mirrors takes milliseconds, which is far too slow to switch packets but fast enough to rewire the machine between jobs.

That rewiring buys three things. First, \emph{shape}: a job can be given a slice of any size from one cube (64 chips) to the whole pod, in any rectangular shape, and the OCS wraps the ends of each dimension to make it a torus. The paper adds a \emph{twisted} torus, in which the wraparound links are offset, which raises bisection bandwidth for all-to-all patterns. Second, \emph{availability}: in an electrically wired 4{,}096-chip torus, one failed host breaks the ring, and the scheduler must find a contiguous working block. With an OCS the scheduler picks any set of healthy cubes and patches them together, routing around the failed hosts. The paper quantifies this as the host availability the system can tolerate: with the OCS, large slices keep good goodput when hosts are available 99.0 to 99.5\% of the time, whereas without it about 99.9\% host availability would be needed for the same result. The slide compresses this to ``slices stay usable at 99\% host uptime.'' Third, \emph{deployment}: a pod can enter service a cube at a time rather than waiting for every rack. The price is small. The paper puts the OCS layer, switches and optics together, at under 5\% of the system's capital cost and under 3\% of its power. The optical switch is now a fixture of Google's design (Ironwood uses it, slide 11) and has been copied in spirit by every vendor that has since put optics in the scale-up network (slide 30).

The second idea is the \emph{SparseCore}. Google's largest workloads in 2020 were not language models but recommendation and ranking models for ads, YouTube and search, and those models spend much of their time not multiplying matrices but looking up rows in enormous embedding tables: sparse, irregular, memory-bound gathers and scatters that an MXU cannot accelerate. TPU v4 adds four SparseCores per chip, small dataflow processors with their own scratchpad that perform embedding lookups, reductions and all-to-all exchanges across the ICI. They take about 5\% of the chip's area and power and speed embedding-heavy models by 5 to 7 times. The point for this lecture is that a special-purpose chip is shaped by its owner's workloads; the TPU is as much an ads chip as a language-model chip, and an outside buyer gets a machine optimized for Google's mix.

The third idea comes from a different chip. In 2021, before the v4 paper, Google described TPUv4i, a one-TensorCore, 175 W inference chip deployed in 2020 (Jouppi et al., ``Ten lessons from three generations shaped Google's TPUv4i,'' ISCA 2021). Its main lesson is that for inference, performance per total cost of ownership (TCO) was driven more by on-chip memory than by FLOP/s: v4i carries 128 MiB of CMEM, an on-chip SRAM between the vector memory and HBM, which keeps the working set of many models on chip and cut HBM traffic. TPU v4 kept the 128 MiB CMEM. This is the same argument that Groq, Graphcore and Cerebras took to its extreme (slide 8), and that TPU 8i returns to with 384 MB of SRAM (slide 11).

The panel's comparison with the A100 needs its qualifier. The paper reports TPU v4 at 1.2 to 1.7 times an A100 on Google's workloads at similar scale, using 1.3 to 1.9 times less power; these are Google's measurements of Google's models. The A100's dense BF16 peak is 312 TFLOP/s, slightly above v4's 275, so the advantage is in delivered efficiency, interconnect and power, not peak. The H100, shipping in volume by the time the paper appeared, was not compared.

The panel ends on the software. XLA compiled the same TensorCore programming model, scalar unit, vector unit, MXUs, for v2, v3 and v4. The MXU count, the clock, the memory and above all the network changed between generations, but a JAX program written for v2 ran on v4. This stability, which Google's 2026 retrospective of five generations calls ``architectural stability'' (Jouppi et al., IEEE Micro 2026), is the property most startups never had time to acquire.

\begin{terms}
\term{optical circuit switch}{A switch that steers light from input fibers to output fibers with movable mirrors, establishing fixed circuits rather than routing packets; reconfigured in milliseconds. Abbreviated OCS.}
\term{MEMS}{Microelectromechanical systems: microscopic mechanical structures, here tilting mirrors, built with chip-fabrication methods.}
\term{cube (TPU v4)}{A block of 64 TPU v4 chips, 4 by 4 by 4, wired electrically inside one rack; the unit the optical switches connect.}
\term{slice}{The set of TPU chips allocated to one job, which the OCS wires into its own torus.}
\term{twisted torus}{A torus whose wraparound links are offset by a shift, raising the bisection bandwidth available to all-to-all traffic.}
\term{bisection bandwidth}{The total bandwidth across the narrowest cut that divides a network into two equal halves; it bounds all-to-all communication.}
\term{SparseCore}{A dataflow processor on TPU v4 and later chips that accelerates embedding lookups and sparse exchanges; later generations also offload collectives to it.}
\term{embedding table}{A large table of learned vectors indexed by an ID (a word, a user, an ad); looking up rows is memory-bound and irregular.}
\term{CMEM}{Common memory: the 128 MiB on-chip SRAM of TPUv4i and TPU v4, between the cores' local memories and HBM.}
\term{TCO}{Total cost of ownership: purchase cost plus power, cooling, space and operations over a system's life.}
\term{TPUv4i}{Google's 2020 single-TensorCore inference TPU, 138 TFLOP/s BF16 at 175 W with 128 MiB of CMEM, described at ISCA 2021.}
\end{terms}

\begin{intuition}
An electrical torus is a building: once wired, its shape is fixed, and a broken room blocks the corridor. An optical circuit switch turns it into a set of rooms on rails that can be rearranged overnight for each tenant, and the rails cost a few percent of the building. TPU v4's other two ideas, SparseCores and on-chip memory, show the owner's workloads writing the specification.
\end{intuition}

% ---------------------------------------------------------------------
\slide{11}{v5e to TPU 8: an inference line, a training line, FP8 and FP4}{Trillium quadrupled v5e; Ironwood added FP8 and 9,216-chip pods; generation 8 splits training from inference.}

\para{What the slide shows} A table occupies the left half, with six columns for six chips: v5e, v5p, v6e (Trillium), v7 (Ironwood), 8t (Sunfish) and 8i (Zebrafish). Its rows give the announcement date (August 2023, December 2023, May 2024, April 2025, April 2026, April 2026), the number of TensorCores per chip (1, 2, 1, 2, 1, 2), the dense peaks (BF16: 197, 459, 918 and 2{,}307 TFLOP/s for the first four, n.p. for TPU 8; FP8: only Ironwood, 4{,}614 TFLOP/s; FP4: only TPU 8, 12.6 PFLOP/s for 8t and 10.1 for 8i), HBM capacity (16, 96, 32, 192, 216 and 288 GB), HBM bandwidth (0.8, 2.8, 1.6, 7.4, 6.5 and 8.6 TB/s), ICI bandwidth per chip (0.4, 1.2, 0.8, 1.2 and 2.4 TB/s, n.p. for 8i), the largest pod (256; 8{,}960; 256; 9{,}216; 9{,}600 chips; n.p. for 8i) and the topology (2D torus; 3D torus with OCS; 2D torus; 3D torus with OCS; 3D torus; ``Boardfly''). A caption gives the sources (Google Cloud documentation, the Hot Chips 2025 Ironwood talk, press coverage of Cloud Next in April 2026 for TPU 8) and two conventions: ICI is quoted as bidirectional bytes per second, as the current documentation does, and HBM usable capacity is slightly below the nominal figure. The right panel, ``Reading the generations,'' takes the four steps in turn: the fork into ``e'' and ``p'' parts in 2023, Trillium in 2024, Ironwood in 2025 and the split TPU 8 in 2026.

\para{Details} After v4 Google stopped building one TPU for everything. The fifth generation, in 2023, came in two parts with different jobs. TPU v5e (announced August 29, 2023) is the efficiency part: one TensorCore, 197 TFLOP/s BF16 (393 TOPS INT8), 16 GB of HBM at about 0.8 TB/s, and 256-chip pods wired as a fixed 16 $\times$ 16 2D torus without optical switches. It was sold for inference and for training models that fit a pod, at a low price per chip-hour. TPU v5p (announced December 6, 2023, alongside Gemini 1.0) is the training part: two TensorCores, four SparseCores, 459 TFLOP/s BF16, 95 GiB usable of 96 GiB HBM at 2.765 TB/s, 1.2 TB/s of ICI, and pods of 8{,}960 chips in an OCS-switched 3D torus. Google's documentation also lists 459 TFLOP/s of ``FP8'' for v5p, equal to its BF16 rate; v5p has no native FP8 datapath, and the figure should not be read as one. For jobs larger than a pod, \emph{Multislice} runs one program across several slices or pods, carrying the data-parallel traffic between them over the ordinary data-center network.

Trillium, TPU v6e (announced May 14, 2024; generally available December 11, 2024), is the next ``e'' part, and its gain is the largest of any single TPU step: 918 TFLOP/s BF16 per chip against v5e's 197, which the panel calls 4.7$\times$ ($918 / 197 = 4.66$; derived). Google attributes it to larger matrix units and a higher clock. HBM doubled to 32 GB at 1.6 TB/s, ICI doubled to 0.8 TB/s, but the pod stayed at 256 chips in a 2D torus. Trillium therefore depends on Multislice for large jobs, and its showcase was MLPerf Training v4.1 (November 2024), in which Google ran GPT-3 175B on several 256-chip Trillium pods and reported 99\% weak-scaling efficiency (Google's figure: the time per step stayed nearly constant as pods and batch grew together).

Ironwood, TPU v7 (announced April 9, 2025 at Cloud Next, presented at Hot Chips on August 26, 2025, generally available from November 2025), is the first TPU built from two compute dies in one package, and the first with native FP8. Per chip: two TensorCores and four SparseCores, 2{,}307 TFLOP/s BF16 and 4{,}614 TFLOP/s FP8 dense, 192 GB of HBM3e in eight stacks at 7.4 TB/s (the Hot Chips slides say 7.3; 192 GiB in the documentation), and 1.2 TB/s of bidirectional ICI over six links. The pod returns to the OCS-switched 3D torus and grows to 9{,}216 chips, 144 cubes of 64:
\begin{formulas}
\[
9{,}216 \times 4.614\ \text{PFLOP/s} \approx 42.5\ \text{EFLOP/s (FP8)},\qquad
9{,}216 \times 192\ \text{GB} \approx 1.77\ \text{PB of HBM}.
\]
\end{formulas}
Google calls the 1.77 PB ``directly addressable'' across the pod. Two caveats. The 42.5 EFLOP/s is FP8; the BF16 figure for the pod is half. And the comparison Google drew at launch, that one Ironwood pod has more than 24 times the compute of El Capitan, set an FP8 pod peak against a supercomputer's FP64 HPL result, which are different quantities. Ironwood's documented peaks are nonetheless close to NVIDIA's B200 (2{,}250 TFLOP/s dense BF16, 8 TB/s); a TPU generation was, for the first time, level with NVIDIA's current part on paper. Google's IEEE Micro retrospective puts Ironwood's performance per watt at about six times v5p's.

The ICI figures in the table need the caption's convention. Google's documentation quotes ICI in gigabytes per second, both directions summed (Ironwood: six links of 200 GB/s, 1{,}200 GB/s); older blog posts quote gigabits per second per direction (v5e's ``1,600 Gbps'' is the same 400 GB/s), and the April 2025 Ironwood announcement's ``1.2 Tbps'' was a units slip for 1.2 TB/s, which is 9.6 Tb/s. The deck uses bytes, bidirectional, throughout.

The eighth generation, announced April 22, 2026, splits training from inference as deliberately as v5 did, but with different chips rather than different sizes of one chip. TPU 8t (``Sunfish'') is the training part: one large TensorCore, 12.6 PFLOP/s of dense FP4, 216 GB of HBM3e in six stacks at 6.5 TB/s, twice Ironwood's ICI (19.2 Tb/s, 2.4 TB/s), and superpods of 9{,}600 chips in a 3D torus, which press coverage puts at 121 EFLOP/s FP4. TPU 8i (``Zebrafish'') is the inference part: two TensorCores, 10.1 PFLOP/s FP4, 288 GB of HBM3e at 8.6 TB/s, and 384 MB of on-chip SRAM, three times Ironwood's. Its network is not a torus but a low-diameter topology Google calls Boardfly, at most seven hops between any two chips in a domain of 1{,}152 against sixteen in an equivalent torus; the SparseCores are replaced by a collectives acceleration engine. Both choices aim at mixture-of-experts inference, where each token's activations go to a few experts on other chips and back: an all-to-all pattern that a torus serves poorly and that rewards both short paths and keeping hot state on chip. Google claims twice Ironwood's performance per watt. The TPU 8 figures come from press transcriptions of Google's presentation; the BF16 and FP8 peaks, power, process and ship dates are not published (n.p.), and the codenames appear only in press. Press reports Broadcom as the design partner for 8t and MediaTek for 8i; Google has not confirmed the split.

Read across the six columns, the table shows three trends that recur throughout this lecture. Per-chip compute rose about 60-fold from v5e's 197 TFLOP/s BF16 to 8t's 12.6 PFLOP/s FP4, most of it from narrower formats. HBM bandwidth rose about tenfold, from 0.8 to 8.6 TB/s. And the networks diverged: Google now builds a torus for training, where neighbors exchange gradients and activations in regular patterns, and a low-diameter network with more SRAM for inference, where the traffic is irregular and latency matters. That is the same diagnosis Groq, Cerebras and SambaNova made, arrived at from the other side.

\begin{terms}
\term{TPU v5e}{Google's 2023 cost- and inference-oriented TPU: one TensorCore, 197 TFLOP/s BF16, 256-chip 2D-torus pods.}
\term{TPU v5p}{Google's 2023 training TPU: two TensorCores, 459 TFLOP/s BF16, 8,960-chip OCS-switched 3D-torus pods.}
\term{Multislice}{Google's technique for running one training job across several TPU slices or pods, with data-parallel traffic carried over the data-center network.}
\term{Trillium}{TPU v6e (2024): 918 TFLOP/s BF16, 32 GB HBM, 256-chip pods.}
\term{Ironwood}{TPU v7 (2025): Google's first two-die TPU and first with native FP8, 4,614 TFLOP/s FP8 dense, 192 GB HBM3e, 9,216-chip pods.}
\term{FP8}{Eight-bit floating point, in E4M3 or E5M2 variants; twice the throughput of BF16 on hardware that supports it natively.}
\term{FP4}{Four-bit floating point, usually with a shared scale per block of values (as in MXFP4 or NVFP4); four times BF16 throughput.}
\term{weak scaling}{Growing the problem with the machine so that work per chip stays constant; perfect weak scaling keeps the time per step flat.}
\term{TPU 8t, TPU 8i}{Google's eighth-generation TPUs announced in April 2026: 8t for training (FP4, 3D torus, 9,600-chip superpods), 8i for inference (384 MB SRAM, Boardfly network).}
\term{Boardfly}{The low-diameter network topology of TPU 8i, at most seven hops between chips, designed for mixture-of-experts all-to-all traffic.}
\term{mixture of experts}{A model architecture in which each token is routed to a few of many expert sub-networks, so parameters grow faster than compute per token; it generates all-to-all communication. Abbreviated MoE.}
\term{all-to-all}{A collective in which every chip sends a distinct piece of data to every other chip; it stresses bisection bandwidth and latency.}
\end{terms}

\begin{intuition}
For ten years the TPU line was one design refined: bigger arrays, more HBM, larger tori. From 2023 Google admitted that training and serving want different machines, and by 2026 the two lines differ in kind: the training chip keeps the torus, the inference chip trades it for short paths and three times the SRAM. When the most experienced builder of special-purpose chips concludes that inference needs on-chip memory and low latency, the startups' premise, if not their execution, is vindicated.
\end{intuition}

% ---------------------------------------------------------------------
\slide{12}{The TPU as a business: Google's own, then everyone's}{For eight years the TPU was Google's own; in 2025--26 it became the first ASIC sold to others at gigawatt scale.}

\para{What the slide shows} A horizontal timeline crosses the top of the slide with eight dated points, labels alternating above and below the line. 2017: the TPUv1 paper at ISCA and the announcement of Cloud TPU (v2). 2018: JAX released, TPU v3 pods on Cloud, MLPerf Training v0.5--0.6 entries. 2021--22: TPU v4 pods of 4{,}096 chips in MLPerf v1.0 and Cloud general availability in 2022. November 2024: Trillium in MLPerf Training 4.1, generally available in December. October 2025: Anthropic agrees to use up to one million TPUs, with over a gigawatt in 2026. December 2025: Broadcom books \$21 billion of TPU orders; Google announces TorchTPU with Meta. February 2026: Meta rents TPUs through Google Cloud in a multibillion-dollar agreement. April 2026, in bold: TPU 8t and 8i announced, and Anthropic adds multiple gigawatts from 2027. Below the timeline two panels sit side by side. The left, ``The stack that made it portable,'' traces the software from XLA (2017) and JAX (2018) through Pathways, PyTorch/XLA, TorchTPU and Pallas, and notes that Google skipped MLPerf Training 5.1 (November 2025) and listed no TPU result in 6.0. The right, ``Why it matters for the rest of this lecture,'' states the TPU's role in the argument: proof that a non-GPU accelerator can carry frontier training, and the benchmark every startup had to beat along with NVIDIA.

\para{Details} For most of its life the TPU was not for sale. From 2015 to 2017 it existed only inside Google, serving search, translation, photos and speech; AlphaGo ran on it in 2016. Cloud TPU, announced with v2 in 2017 and in beta from February 2018, let outsiders rent it, but the customers were a few research groups and some Google Cloud accounts, and for years the main external users were academics on free research credits (the TPU Research Cloud). Google never sold a TPU as a box. The chip was a competitive advantage in Google's own cost of serving and training, and the company's public statements were about efficiency, not revenue.

Two things made it usable outside Google. The first is the software stack in the left panel. XLA compiles whole graphs to the TPU's VLIW cores; JAX, released in 2018, made XLA the default for a generation of researchers, first inside Google and DeepMind and then at many outside labs, including Anthropic; Pathways, Google's runtime for running one program across thousands of chips in several pods from a single controller, is how Gemini is trained. PyTorch users reached TPUs through PyTorch/XLA, a lazy-tensor bridge that never matched native PyTorch performance; from December 2025 Google and Meta have collaborated on TorchTPU to make PyTorch run natively, a project whose existence explains the Meta deal two months later. Pallas exposes a kernel language for the cases the compiler handles badly. The second is MLPerf. Google submitted to MLPerf Training from its first round in 2018 through v5.0 (June 2025), except v2.1 (November 2022) and v3.0 (June 2023), usually with the largest systems in the round, and these results were the only third-party evidence that a non-GPU accelerator could train large models at scale (slide 32). That record ends in the slide's footnote: Google did not submit to Training v5.1 in November 2025, and its v6.0 entry of June 2026 is an NVIDIA GB200 NVL72 system, not a TPU. Ironwood has no MLPerf result at all. Once the TPU had customers who could test it themselves, Google stopped needing the benchmark.

The change in 2025--26 is the subject of the subtitle. On October 23, 2025, Anthropic and Google announced that Anthropic would use up to one million TPUs, bringing ``well over a gigawatt'' of capacity online in 2026, in a deal both companies valued in the tens of billions of dollars$^\dagger$. This is capacity in Google Cloud, not a purchase of chips: ``Anthropic bought a million TPUs'' is a misreading. Two months later Broadcom, which co-designs the TPU and handles its physical implementation, disclosed that one customer had placed a \$10 billion order for TPU-based systems in its fiscal third quarter of 2025 and an additional \$11 billion for delivery in late 2026; on its fourth-quarter call Broadcom identified the customer as Anthropic$^\dagger$. The \$21 billion is what Broadcom booked, not Anthropic's total spending. In February 2026 Meta, whose own MTIA chips (slide 22) are for inference and ranking, signed a multibillion-dollar agreement to rent TPUs through Google Cloud for training, with purchases of hardware for its own data centers reported under discussion for 2027$^\dagger$. And in April 2026, with TPU 8, Anthropic and Google, together with Broadcom, announced further capacity of ``multiple gigawatts'' from 2027$^\dagger$.

The slide's claim that the TPU became ``the first ASIC sold to others at gigawatt scale'' rests on these announcements, and its wording is deliberately narrow. It does not say the TPU is the only non-NVIDIA chip bought at that scale: AMD signed 6 GW of GPUs with OpenAI in October 2025 (slide 23), but those are GPUs, not ASICs. Amazon's Trainium (slide 21) reached gigawatt scale for Anthropic in the same months, with about 500{,}000 Trainium2 chips live in Project Rainier in October 2025, so ``first'' is a matter of weeks and of counting; the difference is that Trainium serves one outside customer inside Amazon's own data centers, while the TPU was taken by Anthropic and then Meta. OpenAI's 10 GW with Broadcom (slide 23) is a commitment, not a deployment. Analysts' cost estimates suggest why buyers came. SemiAnalysis, as reported in the press, estimated that an Ironwood system costs Google about 44\% less to build than an equivalent GB200 system costs to buy, and might offer external customers about 30\% lower total cost of ownership$^\dagger$; these are estimates by an analyst firm, not disclosures. What the deals show with certainty is that the largest AI developers judged the TPU's software and scale mature enough to bet frontier models on, a judgment that only Trainium, with one customer, had also earned among non-NVIDIA chips.

The right panel turns the TPU into a yardstick for the rest of the lecture. It proves the possibility: Gemini is trained on TPUs, Anthropic and Meta buy capacity on them, and Broadcom's custom-accelerator (XPU) business, which by 2025 served five customers, was built on the TPU partnership. It also sets the terms. A competitor had to match a systolic matrix core with HBM, a scalable network with optical switching, a compiler that owns the whole graph, and a decade of patience funded by one of the world's most profitable companies. Every startup in the next three sections had to explain why its approach would beat this, and NVIDIA's GPUs, before it ran out of money.

\begin{terms}
\term{Cloud TPU}{Google Cloud's rental service for TPUs, announced with TPU v2 in 2017; TPUs have never been sold as hardware.}
\term{Pathways}{Google's runtime that runs one program across thousands of accelerators in several pods from a single controller; the system under Gemini training.}
\term{PyTorch/XLA}{The bridge that runs PyTorch programs on TPUs by recording operations lazily and compiling them with XLA.}
\term{TorchTPU}{A Google--Meta collaboration announced in December 2025 to make PyTorch run natively on TPUs.}
\term{Pallas}{A kernel language in JAX for writing custom low-level kernels for TPUs (and GPUs) when the compiler's output is not good enough.}
\term{MLPerf}{The benchmark suites of MLCommons, measuring time to train (Training) and throughput or latency of serving (Inference) on fixed models with peer-reviewed rules; see slide 32.}
\term{XPU}{Broadcom's term for a custom accelerator it designs and builds for one customer, who owns the architecture.}
\term{gigawatt scale}{A deployment measured by the electric power of its data centers, about 1 GW being of the order of half a million to a million accelerators with their cooling and networking.}
\end{terms}

\begin{intuition}
The TPU spent eight years as Google's private advantage and two years becoming a product, and the second phase needed the first: the software, the scale and the benchmark record were in place before any outside company bet a frontier model on it. That is the timescale against which the startups' five-year runways should be judged.
\end{intuition}
% =====================================================================
\setstage{attn}
\section{Cerebras and the wafer (slides 13--15)}

\slide{13}{Cerebras: the wafer is the chip}{84 reticle fields wired across the scribe lines, spare cores for yield, 44 GB of SRAM on one 46,225 mm$^2$ die.}

\para{What the slide shows} On the left is a drawing of a wafer: a pale orange disk filled with a 12 $\times$ 7 grid of darker rectangles, captioned ``12 $\times$ 7 reticle fields, each a GPU-sized die.'' Beside it is an enlarged view of one of the WSE-3's 900{,}000 cores, drawn as four boxes: 48 KB of SRAM with single-cycle access, an 8-wide FP16 SIMD unit doing fused multiply-accumulates, a ``dataflow trigger'' that starts computation when data arrives, and a router with five ports onto a 2D mesh that passes messages in hardware. Below, a table compares the three wafer-scale engines. WSE-1 (2019): TSMC 16 nm, 1.2 trillion transistors, 400{,}000 cores, 18 GB of SRAM at 9 PB/s, dense FP16 not published, in the 20 kW CS-1. WSE-2 (2021): 7 nm, 2.6 trillion transistors, 850{,}000 cores, 40 GB at 20 PB/s, 7.5 PFLOP/s in the dense-FP16 row (the source figure carries no dense/sparse qualifier; see below), in the 23 kW CS-2. WSE-3 (2024): 5 nm, 4 trillion transistors, 900{,}000 cores, 44 GB at 21 PB/s, about 15 PFLOP/s dense FP16 (derived), in the 23 kW CS-3. The right panel, ``Why wafer scale works, and what it costs,'' takes four points: yield, why the design helps, the core, and a warning about the vendor's ``125 PFLOP/s.''

\para{Details} Every chip in this lecture so far is limited in size by the \emph{reticle}: the mask a lithography scanner projects onto the wafer covers at most about 26 $\times$ 33 mm, so no single exposure can pattern a die larger than about 830 to 858 mm$^2$. NVIDIA's largest dies sit at that limit, and the industry's way past it is to package several dies together and connect them through an interposer or bridge (two dies in Blackwell and Ironwood). Cerebras, founded in 2016 by Andrew Feldman, Gary Lauterbach, Sean Lie, Michael James and Jean-Philippe Fricker (most from the server-chip company SeaMicro), took the opposite course. It leaves the whole wafer in one piece. A 300 mm wafer is exposed reticle by reticle, as usual, as 84 identical fields in a 12 $\times$ 7 grid; then, working with TSMC, Cerebras adds wires that cross the \emph{scribe lines}, the narrow lanes between fields that are normally left empty for the saw. The square cut from the wafer, 215 mm on a side, is one die of 46{,}225 mm$^2$, about 56 times the area of the largest GPU die. The idea is old. Gene Amdahl's Trilogy Systems tried wafer-scale integration in the early 1980s and abandoned it in 1984, defeated by defects; the panel names it as the precedent that failed.

The first problem is \emph{yield}. Every wafer has defects, scattered randomly at some density per square centimeter; a die of 46{,}225 mm$^2$ is certain to contain many. A GPU maker discards or downgrades the dies that catch a defect. Cerebras cannot discard the wafer, so it makes defects harmless. The die is a uniform mesh of very small identical cores, so a defect kills one core out of hundreds of thousands. About 1\% of the cores are spare (Cerebras gave the figure at Hot Chips 2019$^\dagger$), and the mesh has redundant links, so after testing, each wafer's map of bad cores is used to configure the fabric to route around them and present a perfect rectangular array to software. The ``4 trillion transistors'' of WSE-3 therefore counts spares and disabled logic, and ``900{,}000 cores'' is the number usable.

The second problem is \emph{power and heat}. A wafer drawing tens of kilowatts cannot be powered through its edges or cooled by a fan. Cerebras feeds power vertically, from a board pressed against the back of the wafer through thousands of connections, and cools the front with a cold plate and water. The system that houses one wafer, the CS-1, CS-2 or CS-3, is a 15U box drawing about 20 to 23 kW (vendor figures$^\dagger$), roughly a third of a rack of GPUs.

What the wafer buys is memory bandwidth (slide 8). Each core has 48 KB of SRAM next to its arithmetic, readable in one cycle. On WSE-3, 900{,}000 cores times 48 KB is about 44 GB, and the sum of their bandwidths is 21 PB/s, which Cerebras compares with an H100's 3.35 TB/s of HBM, a factor of about 7{,}000 ($21{,}000 / 3 \approx 7{,}000$; Cerebras's comparison$^\dagger$). The panel is careful to say what that number is not. It is the sum of 900{,}000 private memories each serving its own core; it is not shared bandwidth. A core reads another core's SRAM only by sending messages across the mesh, whose total bandwidth is about 214 Pb/s, some 27 PB/s, and which delivers a message from one corner of the wafer to the other in many hops. The 21 PB/s is available only when the data a core needs is already in that core's SRAM. For a neural network laid out with each core owning a slice of a layer's weights, that is exactly the case, and the comparison with HBM is fair; for an arbitrary program, it is not.

The core itself is tiny. On WSE-3 it has an 8-wide FP16 SIMD datapath (4-wide on WSE-2) doing fused multiply-adds, 48 KB of SRAM, and a router with five ports (four neighbors and the core). It has no cache and runs in order. Its distinctive feature is the \emph{dataflow trigger}: instructions fire when their input data arrives over the fabric, and a core that receives a zero can skip the multiply. This is the mechanism behind Cerebras's claims about sparsity, and the reason its headline FLOP/s include it. Transistor density makes the point that the wafer is mostly memory and wire: 4 trillion transistors on 46{,}225 mm$^2$ is about 87 million per mm$^2$, against about 130 million on a Blackwell die built on a similar process (derived). SRAM and long wires are less dense than logic.

The table's dense FP16 row and the panel's last paragraph deal with the most misquoted number in the company's literature. The WSE-3 launch release of March 2024 claims ``125 petaflops of peak AI performance,'' with no precision and no statement about sparsity$^\dagger$. Cerebras's own figure for a 2{,}048-system cluster, ``256 exaflops,'' is described in Next Platform's account as FP16 ``with sparse data.'' The dense rate can be estimated from the architecture:
\begin{formulas}
\[
\text{WSE-2: } 850{,}000 \times 4 \times 2\ \tfrac{\text{FLOP}}{\text{FMA}} \times {\approx}1.1\ \text{GHz} \approx 7.5\ \text{PFLOP/s FP16 dense},
\]
\[
\text{WSE-3: } 900{,}000 \times 8 \times 2 \times {\approx}1.1\ \text{GHz} \approx 15\text{--}16\ \text{PFLOP/s FP16 dense (derived)}.
\]
\end{formulas}
The WSE-2's 7.5 PFLOP/s is the FP16 figure analysts quote (TechInsights, from Cerebras's material$^\dagger$), and the source does not say whether it is dense or sparse; the first line above shows that a dense count gives about that number at about 1.1 GHz, but the clock is itself an estimate, so the dense reading is plausible rather than confirmed. Cerebras has not published a dense figure for WSE-3, and the WSE-3 estimate inherits the WSE-2's uncertainty. Dividing 125 by 15.6 gives 8, which matches Cerebras's statements that unstructured sparsity can give up to an eightfold gain. The lecture therefore quotes about 15 PFLOP/s dense FP16 for WSE-3, labeled derived, and puts ``125 PFLOP/s'' in quotation marks. Fifteen dense PFLOP/s on one wafer is about seven B200s' worth of dense FP16 (2{,}250 TFLOP/s each); the wafer's advantage is bandwidth and latency, not arithmetic.

\begin{terms}
\term{reticle}{The mask pattern a lithography scanner exposes in one shot, at most about 26 by 33 mm; it bounds the size of an ordinary die.}
\term{reticle field}{One exposure's worth of area on the wafer; WSE dies are made of 84 identical fields.}
\term{scribe line}{The narrow lane between dies on a wafer, normally left empty for the dicing saw; Cerebras runs wires across it.}
\term{wafer-scale integration}{Building one circuit from an entire wafer rather than cutting it into dies; tried by Trilogy in the 1980s and first made to work by Cerebras.}
\term{WSE}{Wafer-Scale Engine: Cerebras's processor, one 46,225 mm2 die per 300 mm wafer, in three generations (2019, 2021, 2024).}
\term{yield}{The fraction of manufactured dies (or cores) that work; defects scale with area, so large dies need redundancy.}
\term{spare cores}{Extra cores built in so that defective ones can be mapped out after test; about 1 percent on the WSE.}
\term{CS-1, CS-2, CS-3}{The Cerebras systems housing WSE-1, WSE-2 and WSE-3 respectively, each a 15U, 20 to 23 kW box with water cooling.}
\term{dataflow trigger}{Cerebras's core-level mechanism that starts an instruction when its input data arrives, and can skip work on zeros.}
\term{SIMD}{Single instruction, multiple data: one instruction applied to several data elements at once; the WSE-3 core is 8-wide for FP16.}
\term{unstructured sparsity}{Zeros scattered anywhere in a weight or activation tensor, without a fixed pattern; exploited by skipping work on zeros, unlike 2:4 structured sparsity.}
\end{terms}

\begin{intuition}
A GPU is a big room with a fast elevator to a warehouse (HBM). A wafer is 900,000 small rooms, each with its own shelf: anything already on a core's shelf is reached instantly, and the sum of all shelves' speeds is enormous, but fetching from a neighbor's shelf means walking. Wafer scale works when the program can be laid out so that each core's data stays on its own shelf, which is what a neural network allows.
\end{intuition}

% ---------------------------------------------------------------------
\slide{14}{Weight streaming: MemoryX, SwarmX and Condor Galaxy}{Activations stay on the wafer, weights stream in layer by layer from outside; clusters scale by data parallelism.}

\para{What the slide shows} The upper left draws a Cerebras training cluster. A box labeled MemoryX holds the weights and the optimizer state, 1.5 TB to 1.2 PB in DRAM and flash. An arrow leads to SwarmX, which broadcasts weights and reduces gradients. SwarmX fans out to three CS-3 systems, each marked ``activations resident,'' with a note that there can be up to 2{,}048. The caption explains: each CS-3 holds a full copy of the activations and receives the same weights; only the batch is split, so one program fits any cluster size (Lie, IEEE Micro 2023); Cerebras claims models of up to 24 trillion parameters. Below, a table lists the three Condor Galaxy supercomputers built with G42: CG-1 in Santa Clara (July 2023) and CG-2 in Stockton (November 2023), each 64 CS-2, ``4 AI exaflops'' with sparsity; CG-3 in Dallas (2024), 64 CS-3, 8. The right panel, ``What the design buys and gives up,'' lists one program, data parallelism only, the customers, the business catch of revenue concentration, and the absence of MLPerf.

\para{Details} A wafer holds 44 GB. That is enough for the weights of a model of about 20 billion parameters in FP16, and far too little for training anything larger, since training also needs gradients and optimizer state (with Adam in mixed precision about 16 bytes per parameter in all). Cerebras's first systems, the CS-1 and early CS-2, ran in \emph{layer-pipelined} mode, mapping a whole model's layers onto regions of the wafer, which capped models at what the SRAM could hold. From 2021 the company changed the execution model to \emph{weight streaming}, described in Sean Lie's IEEE Micro article of 2023.

Weight streaming turns the usual arrangement inside out. On a GPU the weights live in HBM next to the chip and the activations, which grow with batch size and sequence length, compete with them for space. On a Cerebras cluster the \emph{activations} live on the wafer and the \emph{weights} live outside, in MemoryX, a separate appliance of DRAM and flash that stores every parameter and the optimizer state. Training proceeds one layer at a time. MemoryX sends a layer's weights through SwarmX to the wafer; the wafer computes that layer's forward pass on its resident activations and discards the weights; the next layer's weights arrive. On the backward pass the wafer computes the gradients for each layer and sends them back out, and MemoryX applies the optimizer update. Because each layer's weights are used for a whole batch while they are on the wafer, the external bandwidth needed is modest: the arithmetic intensity is the batch size times the sequence length. The model's size is bounded by MemoryX, not by the wafer: 1.5 TB, 12 TB, and in the largest configuration 1.2 PB, which Cerebras says is enough for 24 trillion parameters$^\dagger$.

Scaling out follows. SwarmX is a tree of switches and servers that broadcasts the same weights to every CS-3 in a cluster and sums the gradients coming back. Every system computes every layer on its own share of the batch, so the cluster is purely \emph{data parallel}. There is no tensor parallelism, splitting a layer's matrices across chips; no pipeline parallelism, assigning layers to different chips; and no expert parallelism, placing mixture-of-experts experts on different chips. These are the three techniques that make training large models on GPUs a specialist's job, and Cerebras's pitch is that its users write none of them. The slide quotes Cerebras's illustration: a GPT-3-class model in 565 lines of PyTorch, against tens of thousands for a GPU framework that implements those parallelisms$^\dagger$. The same program runs on one system or 2{,}048.

The panel's second paragraph is the cost. Pure data parallelism has a well-known limit: the global batch grows with the number of systems, because each needs enough of the batch to keep its wafer busy, and beyond a \emph{critical batch size} larger batches stop reducing the number of training steps needed. A data-parallel cluster also sends the full weights and full gradients over SwarmX every step; that traffic does not shrink as the cluster grows. ``Efficiency is linear in theory'' is the honest wording: Cerebras has published near-linear scaling on its own measurements, but there are no third-party training results at large scale.

The Condor Galaxy machines are the largest Cerebras deployments for training, and all were built with G42, the Abu Dhabi technology group. Condor Galaxy 1, announced in July 2023 in Santa Clara, was 64 CS-2 systems; the first 32 cost about \$100 million (Next Platform$^\dagger$). Condor Galaxy 2 followed in Stockton, California, in November 2023, and Condor Galaxy 3, 64 CS-3 in Dallas, was announced with the WSE-3 in March 2024. The table's ``AI exaflops'' are Cerebras's sparse figures: 64 systems $\times$ about 62.5 ``AI PF'' for a CS-2 (derived from the totals), and 64 $\times$ 125 PF for a CS-3. With the dense estimates of slide 13, CG-3 is about $64 \times 15 \approx 1$ EFLOP/s of dense FP16 (derived). G42 and its partners trained several open models on them, including the Arabic-English Jais family.

The customer list on the panel explains Cerebras's early reputation. Argonne received the first CS-1, in late 2019, and Lawrence Livermore, Sandia, the Pittsburgh Supercomputing Center and the Edinburgh Parallel Computing Centre followed; Mayo Clinic and GSK were early commercial users. National laboratories took the first units of each generation, used them for science (the SC20 paper by Rocki et al. ran a memory-bound stencil code on a CS-1 at close to its peak), and published evaluations, which is how much of the independent knowledge of the machine was produced.

The business catch is concentration. The S-1 registration statement Cerebras filed in September 2024 showed that G42 accounted for 83\% of revenue in 2023 and 87\% in the first half of 2024, and the second S-1, filed in April 2026, showed that MBZUAI (62\%) and G42 (24\%) together still accounted for 86\% of revenue in 2025 (as read from the filing by Futurum$^\dagger$). The panel's ``86--87\% of revenue in 2024 and 2025'' is best read as ``most of revenue'': the 87\% is for the first half of 2024, the full-year 2024 share was not located in the filings, and only the 2025 figure of 86\% is confirmed. G42's investment in Cerebras triggered a review by the Committee on Foreign Investment in the United States (CFIUS), which delayed the first IPO attempt (slide 15). Last, the panel notes that Cerebras has never submitted to MLPerf (slide 32), so all its training claims are its own.

\begin{terms}
\term{weight streaming}{Cerebras's execution model since 2021: activations stay on the wafer and each layer's weights are streamed in from external memory, used for the whole batch, and discarded.}
\term{layer pipelining}{The earlier Cerebras mode in which all layers of a model are mapped onto regions of the wafer at once, limiting model size to on-wafer memory.}
\term{MemoryX}{Cerebras's external weight store, an appliance of DRAM and flash from 1.5 TB to 1.2 PB that holds parameters and optimizer state and applies updates.}
\term{SwarmX}{Cerebras's interconnect that broadcasts weights to all systems in a cluster and reduces their gradients.}
\term{data parallelism}{Every device holds the whole model and processes a different part of the batch; gradients are summed across devices each step.}
\term{tensor parallelism}{Splitting each layer's weight matrices across devices, which then exchange partial results within the layer.}
\term{pipeline parallelism}{Assigning consecutive layers to different devices and passing activations between them in micro-batches.}
\term{critical batch size}{The batch size beyond which larger batches no longer reduce the number of training steps proportionally; it limits pure data parallelism.}
\term{optimizer state}{The extra values an optimizer keeps per parameter, such as Adam's two moment estimates; with master weights, several times the size of the model.}
\term{Condor Galaxy}{A series of Cerebras supercomputers built with G42: CG-1 and CG-2 (64 CS-2 each, 2023) and CG-3 (64 CS-3, 2024).}
\term{S-1}{The registration statement a company files with the US Securities and Exchange Commission before an initial public offering.}
\term{CFIUS}{Committee on Foreign Investment in the United States, which reviews foreign investments for national-security risk.}
\end{terms}

\begin{intuition}
Weight streaming accepts that the wafer is too small to hold a large model and makes a virtue of it: the wafer becomes a fast, stateless layer engine and the model lives in cheap memory outside, so every system in a cluster runs the same simple program. The price is that the only way to add systems is to add batch, which works until the batch gets too large to help.
\end{intuition}

% ---------------------------------------------------------------------
\slide{15}{Cerebras Inference, the CS-4, and the IPO}{From August 2024 the wafer became an inference engine; speed per user, not training, took it public.}

\para{What the slide shows} On the left, a horizontal bar chart titled ``Cerebras's announced records'' gives the output tokens per second for one user on six occasions: Llama 3.1 70B at 450 (August 2024) and 2{,}100 (October 2024); Llama 3.1 405B at 969 (November 2024); Llama 4 Scout at 2{,}600 (April 2025); gpt-oss-120B at 3{,}000 (August 2025); and gpt-oss-120B on the CS-4 at 4{,}400 (August 2026). The caption says these are vendor announcements citing Artificial Analysis, with 16-bit weights at launch and lower precision later; that GPU clouds served Llama 3.1 70B at 30 to 100 tokens per second in 2024; and that the bars measure one user's speed, not throughput. The right panel is a timeline in five steps: the inference launch in August 2024; the first S-1 in September 2024 and its withdrawal in October 2025; the OpenAI contract and the second S-1 in January to April 2026; the IPO on May 14, 2026; and the CS-4 ``Nexus'' in August 2026.

\para{Details} Cerebras's turn to inference follows directly from the arithmetic of slide 4. When a language model generates text one token at a time for one user, every weight is read once per token and used once, so the speed is bounded by memory bandwidth divided by the model's size in bytes. Cerebras's launch post made the calculation itself: Llama 3.1 70B in 16-bit precision is about 140 GB, so 1{,}000 tokens per second requires reading 140 TB per second, some forty H100s' worth of HBM bandwidth for a single user.
\begin{formulas}
\[
\text{tokens/s per user} \;\le\; \frac{\text{weight bandwidth}}{\text{bytes per token}}
= \frac{B}{2\,\text{bytes} \times N_{\text{params}}},
\]
\[
\text{one H100's HBM: } \frac{3.35\ \text{TB/s}}{140\ \text{GB}} \approx 24\ \text{tok/s}.
\]
\end{formulas}
With weights spread across the SRAM of several wafers at 21 PB/s each, the bandwidth bound is far above anything achieved, and the limits become latency across the fabric and between wafers. Since 44 GB holds only about 22 billion 16-bit parameters, a 70B model at 16 bits spans four CS-3 systems (176 GB of SRAM), pipelined layer by layer, and a 405B model needs about 810 GB, some twenty wafers (derived; Cerebras has not published the count).

The service launched on August 27, 2024, with Llama 3.1 8B at 1{,}800 tokens per second and 70B at 450, both with 16-bit weights, which Cerebras described as twenty times the speed of GPU-based clouds$^\dagger$. The bars that follow are a record of software as much as hardware: the 70B figure rose to 2{,}100 tokens per second on October 24, 2024 with the same systems and a new software release, and Llama 3.1 405B reached 969 tokens per second in November 2024. Mistral's Le Chat, Perplexity's Sonar, Meta's Llama API (Llama 4 Scout at over 2{,}600 tokens per second) and OpenAI's open-weight gpt-oss models (3{,}000 tokens per second in August 2025) followed as customers or launch partners. Three cautions apply to every bar. They are Cerebras's announcements, citing Artificial Analysis where the bar says so, of measurements of a public API endpoint, not a chip benchmark. They measure the speed one user sees, not the system's total throughput, which is what determines cost. And the precision changed over time: 16-bit at launch, lower for the later models. Slide 33 compares them with the other SRAM machines and with GPUs.

The right panel tells the financial story. Cerebras filed its first S-1 on September 30, 2024. It showed revenue of \$78.7 million in 2023 and \$136.4 million in the first half of 2024, with G42 accounting for 83\% of 2023 revenue (the slide's figure) and 87\% of the half-year's, and losses of \$127.2 million in 2023. The listing stalled while CFIUS reviewed G42's stake, and in October 2025, after raising \$1.1 billion at an \$8.1 billion valuation, Cerebras withdrew the filing as stale. In January 2026 OpenAI signed a multi-year contract for 750 MW of Cerebras capacity through 2028, valued at over \$10 billion at signing and reported in April as more than \$20 billion over three years, with an option to 2 GW by 2030 and warrants that give OpenAI a stake if its spending reaches agreed levels$^\dagger$. The second S-1, filed April 17, 2026, showed 2025 revenue of \$510 million, still 86\% from MBZUAI and G42, and a backlog of \$24.6 billion, most of it OpenAI.

The IPO priced at \$185 a share, above its marketed range, and the shares began trading on Nasdaq as CBRS on May 14. The base offering of 30 million shares raised about \$5.5 billion (TechCrunch's figure; \$5.55 billion exactly), and with the overallotment option exercised, 34.5 million shares and \$6.38 billion: the slide's ``\$6.4 billion'' is the total with the greenshoe, per the research file R1. It was the largest US IPO of 2026 at pricing. The stock opened at \$350 and closed its first day at \$311.07. The slide's ``\$95B diluted'' is the fully diluted value at that price, scaled from the \$56.4 billion fully diluted value at the \$185 offer price (derived); the press quoted a market value of about \$66 billion at the close, on shares outstanding. By early November 2026 it traded below \$100 a share, under a third of the first close; the slide's ``${\sim}$\$96'' is a snapshot, and a price tracker showed about \$89 in the first days of November$^\dagger$. Cerebras reported revenue of \$193.4 million for the first quarter of 2026 and \$180.1 million (GAAP) for the second, with remaining performance obligations of \$25.4 billion; third-quarter results had not been released when these notes were written.

The last step is hardware. The CS-4, announced August 18, 2026 under the name ``Nexus,'' is not a new wafer. Its WSE-3T is the WSE-3 die, the same 4 trillion transistors, 900{,}000 cores and 44 GB, driven at roughly twice the power through a new power-delivery ``backpack'' mounted behind the wafer (Next Platform estimates the clock rose from about 1.4 to 2.8 GHz; Cerebras has not confirmed it). Cerebras quotes 250 ``AI PFLOP/s,'' 43.2 PB/s of memory bandwidth and 53.5 PB/s of fabric bandwidth per wafer, and three wafers per rack for 750 ``AI PFLOP/s,'' the same unqualified basis as the WSE-3's 125$^\dagger$. By the derivation of slide 13 that is about 31 PFLOP/s of dense FP16 per wafer (derived). The headline claim is gpt-oss-120B at over 4{,}400 tokens per second per user. The slide's ``CS-5 promised for 2027'' comes from press coverage rather than the CS-4 release, which does not mention it$^\dagger$.

The subtitle's point is the change in what Cerebras sells. The wafer was designed for training, and training paid for its first eight years through laboratories and one Gulf customer. What made it a public company was the discovery that SRAM bandwidth sells as speed per user, a product GPUs cannot match at batch one, and that one large customer, OpenAI, would pay for it at gigawatt scale. Whether speed per user stays a business as GPUs close the gap with FP4 and larger HBM is one of the open questions of slide 36.

\begin{terms}
\term{decode}{The phase of language-model inference that generates output tokens one at a time, each requiring a full read of the active weights; bandwidth-bound at small batch.}
\term{tokens per second per user}{The output speed one request sees; distinct from aggregate throughput, the tokens per second of all requests a system serves together.}
\term{speculative decoding}{Generating several candidate tokens with a cheap draft model or head and verifying them in one pass of the large model, raising tokens per second per user.}
\term{Artificial Analysis}{An independent firm that measures the speed, latency and price of public AI model endpoints by querying them repeatedly; see slide 33.}
\term{IPO}{Initial public offering: a company's first sale of shares to the public and listing on an exchange.}
\term{overallotment option}{An underwriter's option to sell up to 15 percent more shares than the base offering; also called the greenshoe.}
\term{remaining performance obligations}{Contracted revenue not yet recognized; an accounting measure of backlog.}
\term{CS-4}{Cerebras's 2026 system, Nexus: three WSE-3T wafers (WSE-3 silicon at about twice the power) per rack, aimed at inference.}
\end{terms}

\begin{intuition}
The wafer's 21 PB/s was built to train models and found its market serving them: at batch one, bandwidth is speed, and nothing else in the industry had that much of it in one place. The IPO measured the excitement about that discovery; the stock's fall by November measured the doubt about whether one customer, and a lead that GPUs are narrowing, make a lasting business.
\end{intuition}
% =====================================================================
\setstage{mlp}
\section{Dataflow and streaming machines (slides 16--20)}

\slide{16}{Groq: a deterministic tensor streaming processor}{One core, 220 MiB of SRAM, no DRAM, every cycle scheduled by the compiler; Llama 2 70B took 576 chips.}

\para{What the slide shows} The upper left draws the GroqChip die as Abts and colleagues described it at ISCA 2020$^\dagger$: a long horizontal floorplan of \emph{functional slices}, each running the full height of the die. From the outside in on each side there is a memory block (MEM, 44 slices holding 110 MiB, drawn in orange), a switch slice (SXM, purple) that shifts, permutes and transposes vectors, a matrix unit (MXM, green) labeled ``320$\times$320 MACs'' (the chip has four such planes, two per hemisphere, so each drawn block stands for 2 $\times$ 320$\times$320), and a vector unit (VXM, blue); the two halves meet at a thin center column. Arrows above and below the die show data streaming east and west, and the caption gives the mechanics: 320-element vectors travel on 32 eastward and 32 westward \emph{streams}, and 144 instruction queues issue in lock step. Below the floorplan a two-column table compares the first chip, GroqChip 1 (the TSP, 2020), with the third, the Groq 3 LPU that NVIDIA presented in 2026: process and die (14 nm, 725 mm$^2$, 26.8 billion transistors, against Samsung 4 nm), on-chip SRAM (220 MiB at 55 TiB/s, against about 500 MB at 150 TB/s), off-chip memory (none, in both), peak (750 TOPS INT8 and 188 TFLOP/s FP16, dense, against 1.2 PFLOP/s FP8, which NVIDIA does not label dense or sparse) and system scale (a node of 8 chips, a rack of 72, a network designed for 10{,}440; against an LPX rack of 256 LPUs and 315 PFLOP/s FP8). A caption reconciles the numbers and notes that Groq has no MLPerf result. The panel on the right, ``From TSP to LPU to NVIDIA's LPX'', gives the history in five steps: the design for determinism, the public demonstration of February 2024, the absence of a second Groq-branded chip, NVIDIA's licensing deal of December 24, 2025, and the Groq 3 LPX rack shown at GTC 2026.

\para{Details} Groq was founded in 2016 by Jonathan Ross, one of the engineers behind Google's first TPU (slide 9), and its chip is the most radical answer in this lecture to the question of slide 6: what does a processor look like if every piece of hardware whose job is to cope with uncertainty is removed? A GPU spends area and energy on caches (which guess what data will be reused), on warp schedulers (which pick a ready thread when another stalls), on memory controllers that reorder requests, on arbiters that decide who wins a contested bus, and on back-pressure that stalls a producer whose consumer is full. Each of these exists because the hardware does not know in advance what the program will do. The tensor streaming processor (TSP) removes all of them. It has no caches, no branch predictor, no out-of-order logic, no arbitration and no DRAM; every functional unit executes an instruction stream that the compiler wrote down cycle by cycle, and the compiler knows, for every cycle of the program, which vector is on which stream at which slice. The paper's phrase is that the hardware is \emph{deterministic}: the same program on the same input takes the same number of cycles every time, to the cycle.

The floorplan follows from that choice. Instead of many identical cores, each with its own memory and arithmetic (the organization of a GPU, an IPU or a Tenstorrent chip, slides 18 and 19), the TSP is one very wide core turned on its side. Each slice performs one kind of operation on 320-byte vectors that flow past it: the MEM slices read and write SRAM, the VXM slices do elementwise arithmetic, the MXM slices do matrix multiplication, and the SXM slices move data between lanes. A computation is a choreography in which the compiler arranges for a vector read from MEM on one cycle to arrive at an MXM a known number of cycles later, exactly when the instruction that consumes it issues there. This is the dataflow idea of slide 5 implemented in time rather than in space: the data moves, the instructions wait for it, and nothing in hardware checks that they meet, because the compiler has proved that they will.

The figures on the slide come from two sources that do not quite agree, and the caption reconciles them. The ISCA 2020 paper gives 820 TOPS of INT8 at a 1 GHz clock; the chip shipped at 900 MHz, which scales the peak to about 738 TOPS (derived), and Groq's product brief rounds that to 750 TOPS$^\dagger$. The paper gives the on-chip SRAM bandwidth as 55 TiB/s; Groq's marketing quotes 80 TB/s, which counts differently (the stream bandwidth into the functional units rather than the read bandwidth of the memory slices). FP16 runs at a quarter of the INT8 rate, 188 TFLOP/s, because a 16-bit operand occupies two byte-planes of the vector and the multiplier array is built for 8-bit products. These are dense peaks; Groq publishes no sparse figure.

The capacity is the cost. 220 MiB is about 231 MB, and Llama 2 70B in 8-bit weights is about 69 GB, so the weights alone need at least 69 / 0.23 $\approx$ 300 chips (derived) before any room is left for activations, the key--value cache or the replication that pipelining requires. Groq's production deployment used 576 chips: eight racks of nine nodes of eight cards$^\dagger$. The aggregate SRAM bandwidth of those chips, 576 $\times$ 80 TB/s $\approx$ 46 PB/s (derived), is three orders of magnitude more than any single user's decode step can consume, which is why the binding limits are the depth of the pipeline across chips and the latency of each hop rather than memory bandwidth (slide 8 gives the general argument; slide 33 the measured speeds).

The ISCA 2022 paper extended determinism from the chip to the network. A conventional cluster network is the least deterministic part of a computer: packets are routed adaptively, queues fill and drain, and a lost packet is retransmitted. Groq's multi-chip fabric schedules every link in advance, uses forward error correction instead of retry so that a corrupted packet never changes the timing, and synchronizes the chips' clocks so that the whole system runs one static program. The paper reports 722 ns per hop and a topology (a Dragonfly) designed to reach 10{,}440 chips. The result is a machine whose tail latency is a compile-time constant, which is exactly the property an interactive chat service wants.

The ``LPU moment'' came in February 2024, when Artificial Analysis, an independent benchmarking site, listed Groq's public endpoint serving Llama 2 70B at 241 output tokens per second to one user, about twice the next provider$^\dagger$. Groq renamed the chip the \emph{LPU} (language processing unit), opened GroqCloud to developers, and in 2025 received a \$1.5 billion commitment from Saudi Arabia for a data center in the kingdom$^\dagger$. What it did not do was ship a second chip of its own. In August 2023 Groq announced that its next chip would be made on Samsung's 4 nm process at the new Taylor, Texas fab; that part did not reach GroqCloud before the NVIDIA deal, but it was not abandoned: Samsung 4 nm is the process of the Groq 3 LPU that NVIDIA now builds (Korean press, March 2026$^\dagger$). Meanwhile, in September 2025, Groq raised \$750 million at a \$6.9 billion valuation on the strength of the original 14 nm part$^\dagger$.

On December 24, 2025 NVIDIA announced a non-exclusive license to Groq's inference technology and hired Jonathan Ross and most of the engineering team; the price was not disclosed and was reported at about \$20 billion$^\dagger$. Groq Inc.\ kept the GroqCloud service and in August 2026 raised again at \$3.5 billion (single secondary source; treat as reported). At GTC in March 2026 NVIDIA showed the result: the Groq 3 LPX, a rack of 256 LPUs of about 500 MB of SRAM each, 1.2 PFLOP/s of FP8 per chip and 315 PFLOP/s per rack, placed beside a Vera Rubin NVL72$^\dagger$. The division of labor is the interesting part, and it is finer than ``prefill on GPUs, decode on LPUs.'' Prefill (reading the prompt) is compute-bound and runs on the Rubin GPUs, and the key--value cache it produces stays in the GPUs' HBM. Decode is then split by layer type: the attention half of each decode step, which reads the KV cache, also runs on the GPUs, while the feed-forward (FFN, or mixture-of-experts) half, whose cost is reading weights, runs on the LPUs, whose weights sit in SRAM; the LPX also runs the draft model for speculative decoding. The slide's panel compresses this to ``decode FFN racks beside Vera Rubin, which does prefill, attention and KV''$^\dagger$. NVIDIA reported Gemma 4 31B at 3{,}431 tokens per second on a private pre-release endpoint, a number that is NVIDIA's own and was not measured on a public service$^\dagger$. The rack's 256 $\times$ 500 MB $=$ 128 GB of SRAM and 256 $\times$ 150 TB/s $\approx$ 38 PB/s (derived; ``about 40 PB/s'' on the slide) make it a weight store for the FFN layers of models up to roughly 100 billion 8-bit parameters per rack.

\begin{formulas}
Minimum chips to hold a model in SRAM, and the batch-1 bandwidth bound on decode speed:
\[
N_{\text{chips}} \ge \frac{\text{weight bytes}}{\text{SRAM per chip}},\qquad
\text{tokens/s} \le \frac{N_{\text{chips}}\times B_{\text{SRAM}}}{\text{weight bytes}} .
\]
Llama 2 70B, INT8: $69\,\text{GB}/0.23\,\text{GB}\approx 300$ chips; Groq used 576. Bound: $576\times 80\,\text{TB/s}/69\,\text{GB}\approx 670{,}000$ tokens/s, against 241 measured, so bandwidth is not the limit. FP16 peak from INT8: $750/4\approx 188$ TFLOP/s. LPX: $256\times 0.5\,\text{GB} = 128$ GB; $256\times 1.2\,\text{PF}\approx 307$ PF (derived; NVIDIA quotes 315).
\end{formulas}

\begin{terms}
\term{tensor streaming processor (TSP)}{Groq's name for its chip architecture: one wide core organized as functional slices across which vectors stream east and west on a schedule fixed by the compiler.}
\term{LPU}{Language processing unit: Groq's product name, from 2024, for its TSP chips sold for language-model inference; NVIDIA kept the name for the Groq 3 part.}
\term{functional slice}{A column of the TSP die that performs one class of operation (memory, vector arithmetic, matrix multiply or data movement) on every vector that passes it.}
\term{stream}{One of the TSP's 64 vector channels (32 in each direction) that carry 320-byte vectors between functional slices, one hop per cycle.}
\term{determinism}{The property that a program's timing is fixed and known in advance to the cycle; achieved on the TSP by removing caches, arbitration and dynamic scheduling.}
\term{static schedule}{An execution plan in which the compiler assigns every operation to a unit and a cycle before the program runs, so that no hardware decides anything at run time.}
\term{back-pressure}{A hardware mechanism by which a full consumer stalls its producer; absent on the TSP because the schedule guarantees the consumer is ready.}
\term{forward error correction}{Adding redundant bits to transmitted data so the receiver can correct errors itself, without asking for a retransmission; it keeps link timing fixed.}
\term{Dragonfly}{A network topology of fully connected groups joined by global links, giving a small number of hops at large scale; Groq's multi-chip fabric uses it.}
\term{prefill and decode}{The two phases of serving a language model: prefill processes the prompt in parallel and is compute-bound; decode generates one token at a time and, at small batch, is bound by reading the weights.}
\term{LPX}{NVIDIA's 2026 rack of 256 Groq 3 LPUs beside a Vera Rubin NVL72: the LPUs run the FFN and mixture-of-experts part of decode and speculative drafts, the GPUs prefill, attention and the KV cache.}
\term{Artificial Analysis}{An independent company that measures the speed, latency and price of public model-serving endpoints; the source of most of the speed numbers in slides 16 and 33.}
\end{terms}

\begin{intuition}
The TSP is a GPU with the guessing removed: a compiler that knows every cycle replaces the caches, schedulers and arbiters, and the area they occupied becomes SRAM and arithmetic. The price is that the model must fit in SRAM spread over hundreds of chips, which is why Groq won on speed per user and never on cost per token, and why its best use turned out to be as the weight-reading half of someone else's decode step.
\end{intuition}

% ---------------------------------------------------------------------
\slide{17}{SambaNova: reconfigurable dataflow, three memory tiers}{Compute and memory units wired per model by the compiler; SRAM, HBM and DDR hold a trillion parameters.}

\para{What the slide shows} The upper left draws the reconfigurable dataflow unit (RDU) as a checkerboard: blue \emph{pattern compute units} (PCU) alternate with orange \emph{pattern memory units} (PMU), with a small purple switch beside each, and arrows show data passing from unit to unit. The caption says that a decoder layer becomes one fused pipeline whose intermediate results never leave the chip, and that the SN40L carries 1{,}040 PCUs and 1{,}040 PMUs on two dies. To the right, three stacked bars draw the SN40L's memory tiers: SRAM, 520 MB at ``PB/s''; HBM3, 64 GB at 1.6 TB/s; DDR5, up to 1.5 TB, labeled ``model store''; the caption says that experts and whole models swap from DDR into HBM in milliseconds. Below, a table follows four generations: SN10 (2020; 7 nm, 40 billion transistors, 320 MB SRAM, no HBM, 1.5 TB DDR4, 325 TFLOP/s BF16), SN30 (2022; 7 nm, 86 billion, 640 MB, 1 TB DDR4, 688), SN40L (2023; 5 nm, 102 billion, 520 MB, 64 GB HBM3, 1.5 TB DDR5, 638) and SN50 (2026; process not published, 432 MB, 64 GB HBM2E, up to 2 TB DDR, about 1{,}600 TFLOP/s BF16; the system row reads ``8 per chassis''). The caption cites the MICRO 2024 paper, says the peaks are dense, that shipped SN40L-16 racks carry 768 GB of DDR5 per RDU, and that SambaNova has no MLPerf result. The panel, ``Composition of experts, and the business'', covers the architectural idea, the Samba-CoE system and the cloud service, the customers, and the company's finances.

\para{Details} SambaNova was founded in 2017 by Kunle Olukotun and Christopher R\'e of Stanford and Rodrigo Liang, and its architecture comes directly from Olukotun's group: \emph{Plasticine} (ISCA 2017), a coarse-grained reconfigurable array (CGRA) designed to execute \emph{parallel patterns}, the map, reduce, filter and group-by operations from which most data-parallel programs are composed. An FPGA is reconfigurable at the level of single bits and lookup tables; a CGRA is reconfigurable at the level of whole arithmetic units and memories, which costs less area and runs faster but can express less. The RDU is a CGRA: an array of PCUs, each a SIMD pipeline of arithmetic stages, interleaved with PMUs, each a bank of scratchpad SRAM with its own address generators, joined by a mesh of programmable switches. The compiler, SambaFlow, takes the model's computation graph, assigns each operator to a group of PCUs, each tensor to PMUs, and each edge to a route through the switches, and then the whole configuration runs as a pipeline.

The difference from a GPU is the difference between running a graph \emph{in time} and running it \emph{in space}. A GPU executes a model as a sequence of kernels: a matrix multiply, then a normalization, then an activation, each reading its input from HBM and writing its output back (or, in a well-fused kernel, to on-chip memory for a short while). The RDU lays the operators out side by side on the chip, so the output of one is fed directly into the next through the switches, and a whole decoder layer, or several, becomes one pipeline through which tokens flow. SambaNova calls this \emph{dataflow fusion}; slide 5 classifies it as the dataflow organization. Its benefit is that intermediate tensors never touch HBM, which saves both bandwidth and energy (slide 8 gives the energy ratio); its cost is that the compiler must find a layout that fits the chip, and a layer that does not fit must be split across configurations, which are swapped at a cost.

The three memory tiers of the SN40L are the company's answer to the memory wall (slide 4). The SN10 and SN30 had large SRAM and very large DDR4 but no HBM: a model too large for SRAM was streamed from DDR at roughly 100 GB/s per chip, which is fine for training a convolutional network and slow for decoding a language model. The SN40L, described by Prabhakar and colleagues at MICRO 2024, added 64 GB of HBM3 between the two: 520 MB of SRAM holds the working set of the current layers, 64 GB of HBM holds the active model, and up to 1.5 TB of DDR5 holds many models or many experts at once (the racks that shipped carry 768 GB per RDU$^\dagger$). The design point is not one huge model but many medium ones that can be switched quickly.

That is the idea of \emph{composition of experts} (CoE). Instead of one mixture-of-experts network trained end to end, Samba-CoE (2024) is a router in front of about 150 separately trained expert models, adding up to a trillion parameters, all resident in the DDR tier of eight sockets; a request is routed to an expert, whose weights are copied from DDR into HBM in milliseconds if they are not already there. On that workload the MICRO paper reports 3.7$\times$ the throughput of a DGX H100, whose 640 GB of HBM cannot hold the whole collection and must reload from host memory over PCIe. The comparison is real but narrow: it measures a workload chosen because it needs capacity, and the paper's authors are SambaNova's engineers.

The SambaNova Cloud service (2024) turned the hardware to the speed-per-user market that Groq and Cerebras had opened. The slide quotes Llama 3.1 405B at 132 tokens per second in 16-bit precision on 16 RDUs$^\dagger$. It is worth checking such a number against the bandwidth bound. 405 billion parameters in BF16 are about 810 GB; 16 RDUs have $16\times 64 = 1{,}024$ GB of HBM, so the weights fit, but $16\times 1.6 = 25.6$ TB/s of HBM bandwidth streams them at most $25.6/0.81\approx 32$ times per second (derived). A rate four times higher than that bound means the service reuses each weight read across several tokens: speculative decoding, in which a small model drafts tokens that the large one verifies in a batch, is the usual technique, though the sources used here do not say what SambaNova does. The general lesson, which applies to every speed figure in slide 33, is that a service's tokens per second measures the whole serving stack, not the chip.

The slide's table needs two remarks. First, the SN40L's BF16 peak (638 TFLOP/s) is a little \emph{lower} than the SN30's (688): the generation traded arithmetic for the HBM tier and moved the SRAM from 640 MB to 520 MB, which shows what the company thought the bottleneck was. Second, the SN50 figures are announced, not shipped: SambaNova describes it as ``5$\times$'' the SN40L in compute with FP8 support (3.2 PFLOP/s FP8), and the slide's ``about 1{,}600'' TFLOP/s of BF16, which the slide's caption labels derived, is in fact SambaNova's published 16-bit figure of 1.6 PFLOP/s$^\dagger$. The slide's ``8 per chassis'' could not be confirmed; The Register gives 16 RDUs per air-cooled rack (SN50 details from the February 2026 announcement and Next Platform: two chiplets of 216 MB SRAM each, 64 GB of HBM2E at 1.84 TB/s, 256 GB to 2 TB of DDR5$^\dagger$).

The business history explains why the SN50 matters. SambaNova raised \$676 million at a \$5.1 billion valuation in 2021, led by SoftBank's Vision Fund 2, and sold whole systems (DataScale) to laboratories and enterprises: Argonne's ALCF AI Testbed (SN10, then SN30: 8 nodes of 8 RDUs), Lawrence Livermore, RIKEN, Saudi Aramco, SoftBank and JPMorgan Chase$^\dagger$. In 2025 it cut staff and pivoted from training to inference; in December 2025 Intel was reported to be negotiating to buy it, and the talks ended as a partnership and a \$35 million stake$^\dagger$. In July 2026 it raised \$1 billion at \$11 billion to build and ship the SN50 in the second half of 2026$^\dagger$. Like Groq and Cerebras, it has never submitted to MLPerf (slide 32).

\begin{formulas}
Batch-1 decode bound from HBM, and the capacity of the three tiers on $n$ sockets:
\[
\text{tokens/s}\le\frac{n\,B_{\text{HBM}}}{\text{weight bytes}},\qquad
C_{\text{tier}} = n\times c_{\text{tier}} .
\]
Llama 3.1 405B, BF16, 16 SN40L: $16\times1.6\,\text{TB/s}/810\,\text{GB}\approx 32$ tokens/s (derived) against 132 served. Samba-CoE on 8 sockets: SRAM $8\times0.52\approx4.2$ GB, HBM $8\times64=512$ GB, DDR $8\times1.5=12$ TB, enough for a trillion BF16 parameters (2 TB) six times over.
\end{formulas}

\begin{terms}
\term{RDU}{Reconfigurable dataflow unit: SambaNova's chip, an array of compute and memory units whose connections the compiler configures for each model.}
\term{CGRA}{Coarse-grained reconfigurable array: a chip of arithmetic units, memories and switches whose wiring is set by configuration, at the granularity of words rather than the bits of an FPGA.}
\term{PCU and PMU}{Pattern compute unit and pattern memory unit: the RDU's two tile types, a pipelined SIMD arithmetic unit and a bank of scratchpad SRAM with address generators.}
\term{Plasticine}{The 2017 Stanford CGRA for parallel patterns (Prabhakar, Olukotun and others) on which SambaNova's RDU is based.}
\term{parallel pattern}{A data-parallel building block such as map, reduce, filter or group-by; Plasticine and the RDU are designed to execute compositions of them.}
\term{dataflow fusion}{Executing several operators of a model as one spatial pipeline on the chip, so that intermediate tensors pass between units without being written to off-chip memory.}
\term{SambaFlow}{SambaNova's compiler, which maps a PyTorch model's graph onto the RDU's compute units, memories and switches.}
\term{three-tier memory}{The SN40L's hierarchy of on-chip SRAM, on-package HBM and board-level DDR5, each larger and slower than the last.}
\term{composition of experts (CoE)}{Serving many separately trained models behind a router, as one system; distinct from a mixture-of-experts network, whose experts are trained jointly inside one model.}
\term{speculative decoding}{Generating several draft tokens with a small model and verifying them with the large model in one batched step, so that each read of the large model's weights produces more than one token.}
\term{DataScale}{SambaNova's packaged system product: racks of RDUs with host CPUs, memory and software.}
\end{terms}

\begin{intuition}
A GPU runs a model as a sequence of kernels through a cache; an RDU lays the model out on the chip and pours data through it. SambaNova's distinctive bet was not the dataflow, which several companies tried, but the memory hierarchy: a terabyte and a half of cheap DDR per chip, so that capacity, not FLOPS, is what it sells. When a speed figure looks too good, divide the HBM bandwidth by the weight bytes; the gap tells you how much the software is doing.
\end{intuition}

% ---------------------------------------------------------------------
\slide{18}{Graphcore: the IPU and bulk-synchronous parallelism}{1,472 MIMD tiles with 900 MB of SRAM exchange data in compiler-scheduled supersteps; sold to SoftBank in 2024.}

\para{What the slide shows} The upper left draws an IPU as a grid of small tiles, each a pale square with a darker strip for its core, captioned ``one core, six threads, 624 KB of SRAM each (GC200)''. To its right a timeline draws the execution model: a block labeled ``compute on local SRAM'', a thin purple bar labeled ``sync'', then a block labeled ``exchange all-to-all'', repeated, with an arrow for time; the caption says the compiler schedules every exchange so that none collide, with 133 ns tile to tile and 7.7 TB/s of measured aggregate exchange bandwidth. Below, a table compares three generations: GC2 (2018; 16 nm, 1{,}216 tiles, 304 MiB SRAM, 1.6 GHz, 125 TFLOP/s FP16, no DRAM except the host's, sold as the C2 card), GC200 (2020; 7 nm, 59.4 billion transistors, 1{,}472 tiles, 900 MB, 1.325 GHz, 250 TFLOP/s, streaming memory in DDR4, M2000 machines to POD256) and Bow (2022; the same 7 nm die on a power-delivery die, 1{,}472 tiles, 900 MB, 1.85 GHz, 350 TFLOP/s, Pod16 to Pod1024). The caption notes the GC200 die area (823 mm$^2$), that the 45 TB/s SRAM figure is nominal (a Citadel team measured 30.7), and that IPUs publish no sparse peak, so they should be compared with a GPU's dense FP16. The panel, ``What it proved, and what it did not'', describes the model, the wafer-on-wafer Bow, the MLPerf results and the ending.

\para{Details} Graphcore was founded in Bristol in 2016 by Nigel Toon and Simon Knowles and was for several years the best-funded AI chip startup outside the United States. Its intelligence processing unit (IPU) is the clearest commercial implementation of a model of parallel computation that Leslie Valiant proposed in 1990, \emph{bulk-synchronous parallelism} (BSP). In BSP a computation is a sequence of \emph{supersteps}. In each superstep every processor computes using only its own local memory; then all processors reach a barrier; then they exchange the data the next superstep needs; then the next superstep begins. Because computation and communication never overlap and every processor waits at the barrier, the model is easy to reason about: the time of a superstep is the time of the slowest processor's computation plus the time of the exchange.

The IPU makes the model physical. Each of the GC200's 1{,}472 tiles is an independent processor (the chip is MIMD, multiple instruction streams on multiple data, unlike a GPU's SIMT warps) with six hardware threads and 624 KB of SRAM that only it can address; 1{,}472 $\times$ 624 KB is about 918 MB, the ``900 MB'' of the slide (derived). There is no cache and no shared memory. During the compute phase each tile works on its own data. At the sync the tiles stop. During the exchange phase data moves between tiles over an on-chip all-to-all network, and because the Poplar compiler has computed the whole exchange in advance, it can schedule every transfer so that no two collide, without any arbitration hardware. Jia and colleagues at Citadel, who reverse-engineered the GC2 in 2019, measured tile-to-tile latency in the low hundreds of nanoseconds and aggregate exchange bandwidth of 7.7 TB/s on the GC2; they measured the on-chip SRAM bandwidth below Graphcore's nominal figure (30.7 against 45 TB/s, the caption's numbers$^\dagger$).

Two comparisons place the IPU among the other designs. Against the Groq TSP (slide 16), both rely on a compiler that schedules everything and both keep weights in SRAM, but Groq has one wide core and moves data in a single synchronized pipeline, while Graphcore has 1{,}472 independent cores that compute asynchronously and synchronize only at barriers. Against a GPU, the IPU is much better at fine-grained irregular work, such as sparse operations, graph networks or small batches, where each tile can do something different, and much worse at holding a large model: on the GC2 the entire model and its activations had to fit in 304 MiB, because there was no DRAM at all. The GC200 added \emph{streaming memory}, DDR4 on the M2000 board from which weights and activations could be loaded between phases, but at a small fraction of the SRAM bandwidth.

The Bow IPU (2022) is a manufacturing first worth knowing. TSMC bonded a second wafer, carrying deep-trench capacitors and power delivery but no logic, face to face to the logic wafer (wafer-on-wafer hybrid bonding, TSMC's SoIC-WoW). The capacitors smooth the supply voltage so the unchanged logic die can run faster: the clock rose from 1.325 to 1.85 GHz, about 40\%, and the FP16 peak from 250 to 350 TFLOP/s (250 $\times$ 1.85/1.325 $\approx$ 349, derived), with no redesign of the processor. It was the first production use of wafer-on-wafer stacking, a technique that later appeared in other products.

Graphcore was, with Google and Habana, one of the few non-NVIDIA companies to submit to MLPerf Training (slide 32): rounds v1.0 through v2.0 (June 2021 to June 2022). The last IPU results, Graphcore's own and Baidu's IPU-POD16 and POD64 entries, are in v2.0; v2.1 (November 2022) has none. In v1.1 an IPU-POD16 trained ResNet-50 in 28.3 minutes against 29.1 for a DGX A100. Graphcore had priced the comparison at its first MLPerf round in June 2021, before Bow existed: the IPU-POD16 (four M2000 machines of GC200 IPUs) listed at \$149{,}995, about half the price of the NVIDIA system$^\dagger$. The pitch was price per result, not peak, and it was honest; the trouble was that by 2022 the workloads customers cared about were transformers of tens of billions of parameters, for which 900 MB of SRAM per chip was the wrong shape and price per ResNet-50 result was the wrong metric.

The ending was quick. In March 2022 Graphcore announced the ``Good computer'', a machine of 8{,}192 next-generation IPUs, 10 exaflops of AI compute and support for 500-trillion-parameter models, for about \$120 million; it was never built. Its 2022 accounts showed revenue of \$2.7 million against a loss of \$205 million$^\dagger$. In July 2024 SoftBank bought the company for a reported \$500--600 million, against a 2020 valuation of \$2.77 billion (the slide rounds to \$2.8B)$^\dagger$; no new silicon has been announced since. Slide 35 uses Graphcore as the example of a company that was right about a mechanism, compiler-scheduled communication on SRAM-rich MIMD silicon, and wrong about the size of the models.

\begin{terms}
\term{IPU}{Intelligence processing unit: Graphcore's processor, an array of 1{,}216 or 1{,}472 independent tiles each with private SRAM, executing in bulk-synchronous supersteps.}
\term{bulk-synchronous parallelism (BSP)}{Valiant's 1990 model of parallel computation in which processors alternate local computation, a global barrier and a data exchange; each round is a superstep.}
\term{superstep}{One round of a BSP program: compute on local data, synchronize, exchange.}
\term{tile}{On the IPU, one processor core with its own SRAM and six hardware threads; more generally a repeated unit of a chip's floorplan.}
\term{MIMD}{Multiple instruction, multiple data: each processor runs its own instruction stream; contrasted with SIMD and with the GPU's SIMT, where many lanes share one stream.}
\term{barrier}{A synchronization point that every participating processor must reach before any may continue.}
\term{Poplar}{Graphcore's compiler and runtime, which partitions a graph across tiles and schedules the compute and exchange phases; PopTorch and PopART are its PyTorch and ONNX front ends.}
\term{streaming memory}{Graphcore's term for the DDR4 attached to an IPU machine, from which data is loaded into tile SRAM between supersteps.}
\term{wafer-on-wafer bonding}{Bonding two whole wafers face to face before dicing; Bow used it to put a power-delivery wafer with capacitors under the logic wafer.}
\term{deep-trench capacitor}{A capacitor etched deep into silicon, giving high capacitance per area; used to stabilize a chip's supply voltage so it can run at a higher clock.}
\end{terms}

\begin{intuition}
The IPU is 1{,}472 small computers that agree to stop at the same moment and trade data on a timetable the compiler wrote. It was elegant and fast for models that fit in its SRAM, and models stopped fitting. The lesson for every SRAM-only design is to ask where the model lives when it is a hundred times larger than the chip.
\end{intuition}

% ---------------------------------------------------------------------
\slide{19}{Tenstorrent: Tensix cores, RISC-V and an open stack}{Clusters of five small RISC-V cores around a matrix engine and 1.5 MB of SRAM, meshed; Ethernet on the die.}

\para{What the slide shows} The upper left draws one Tensix ``core'' of the Blackhole chip as Tenstorrent presented it at Hot Chips 2024$^\dagger$: across the top, five small blue boxes, five RISC-V processors labeled ``NoC in'', ``unpack'', ``math'', ``pack'' and ``NoC out''; below them a green matrix engine (``FPU'', block floating point, FP8 and FP16) and a purple vector engine (``SFPU'', activations and reductions); along the bottom an orange bar, 1.5 MB of software-managed L1 SRAM with two networks-on-chip. The caption says the chip is a mesh of 120 to 140 of these, plus 16 ``big'' 64-bit RISC-V cores, with no caches and no shared memory. A stat tile gives the price of a Blackhole p150 card with 32 GB: \$1{,}399. The table compares three cards: Grayskull e150 (cards shipped 2024; 12 nm, 120 Tensix at 1.2 GHz, 120 MB SRAM, 8 GB LPDDR4 at 118 GB/s, 332/83 TFLOP/s FP8/FP16, no Ethernet, 200 W), Wormhole n300 (2024; 12 nm, two dies, 128 Tensix at 1.0 GHz, 192 MB, 24 GB GDDR6 at 576 GB/s, 466/131, 16 $\times$ 100 Gb Ethernet per die, 300 W) and Blackhole p150 (2025; TSMC 6 nm, 140 Tensix designed at 1.35 GHz, 180--210 MB, 32 GB GDDR6 at 512 GB/s, 745/372, 10--12 $\times$ 400 Gb Ethernet, 300 W). The caption says the peaks are dense datasheet figures, that shipping Blackhole cards enable 120 of the 140 Tensix, that the Blackhole column is from Hot Chips 2024, and that the p150 product page lists 664 TFLOP/s ``BLOCKFP8'' at 120 cores. The panel, ``Not a core, not a GPU, not a closed stack'', explains what a Tensix is, the Ethernet scale-out, the open software, and the money.

\para{Details} Tenstorrent was founded in Toronto in 2016 by Ljubisa Bajic and became widely known when Jim Keller, the designer of AMD's Zen and Apple's early A-series processors, joined as chief technology officer and then, in 2023, chief executive. It is the most open company in this lecture: its kernel-level programming environment, its operator library and its compiler are published as open source, its cards are sold on a website with a price, and its RISC-V CPU designs are licensed to other companies. Its architecture is a mesh of identical compute clusters, like a GPU or an IPU, but each cluster is built differently from both.

The word ``core'' needs care here. A Tensix is a \emph{cluster}: five small in-order RISC-V processors (Tenstorrent calls them ``baby'' cores) share one matrix engine, one vector engine and 1.5 MB of SRAM. Two of the five are data movers: one reads tiles of data from the network-on-chip (NoC) into the local SRAM, the other writes results out. The other three form a software pipeline around the math engines: one \emph{unpacks} data from SRAM into the engines' source registers, one issues the math instructions, one \emph{packs} results back into SRAM. The programmer, through TT-Metalium, writes three kinds of small kernels, a reader, a compute kernel and a writer, that run concurrently on these processors and synchronize through circular buffers in the shared SRAM. The data unit is a 32 $\times$ 32 tile. This is an explicit, software-managed dataflow, in the taxonomy of slide 5 closer to a MIMD mesh with scratchpads than to a GPU's cached SIMT. The count of ``752 RISC-V cores'' on Blackhole adds up these baby cores across the clusters and the chip's other controllers; only the 16 SiFive X280 cores are application-class processors able to run Linux, which lets a Blackhole card run without a host CPU.

The table's numbers show the company's priorities. Each generation keeps roughly the same number of Tensix clusters (120 to 140) and the same 1.5 MB per cluster, so SRAM grows from 120 MB to 180--210 MB (120 $\times$ 1.5 $=$ 180; 140 $\times$ 1.5 $=$ 210, derived); compute grows mainly with clock, format and process. Memory is GDDR6, the graphics-card DRAM, not HBM: 512 GB/s on a Blackhole is a sixth of an H100's 3.35 TB/s, but GDDR6 costs a fraction as much and needs no advanced packaging, which is how a 32 GB card sells for \$1{,}399. The peak figures need two cautions. Tenstorrent publishes FP8 and ``block FP8'' (BFP8, a format in which a block of values shares one exponent) and has labeled them inconsistently between documents: Hot Chips 2024 gave 745 FP8 and 372 FP16 for the 140-core design, which is what the slide's table and the side-by-side table of slide 31 use; a press report gives 774 ``BLOCKFP8'' for 140 cores and 664 for the 120-core shipping part (664 is 774 $\times$ 120/140, derived), and the product page lists the 664. The 745 and 774 are both 140-core figures that differ in how the format is counted; the slides stay with the Hot Chips pair so that FP8 and FP16 come from one source. And shipping Blackhole cards enable 120 of 140 clusters, after a firmware change reported in the press$^\dagger$, so per-card peaks should be scaled by 120/140 unless the datasheet already has.

The Ethernet ports are the most distinctive system decision. A GPU scales up through a proprietary fabric and switches (NVLink, slide 2's reference) and scales out through network cards; a Tenstorrent chip has Ethernet controllers and routers on the die and connects directly to its neighbors. Thirty-two chips are wired into a 4 $\times$ 8 mesh in a Galaxy server: 5.2 PFLOP/s of block FP8 and 384 GB of GDDR6 for Wormhole, about 23.8 PFLOP/s of FP8 and 1 TB for Blackhole$^\dagger$ (32 $\times$ 745 TFLOP/s $\approx$ 23.8 PFLOP/s and 32 $\times$ 32 GB $=$ 1{,}024 GB, derived; at the 120 cores of shipping cards it is about 32 $\times$ 664 $\approx$ 21 PFLOP/s), and Galaxies connect to each other the same way. There is no switch, no HBM and no NVLink-class fabric, so the system is cheap and the network bandwidth per chip is modest. The design suits inference of models that can be partitioned over a mesh and leaves the largest training runs to others.

The open stack is the company's other bet. TT-Metalium is the bare-metal C++ environment for writing kernels for the five processors of a Tensix; TT-NN is an operator library on top; TT-Forge is an MLIR-based compiler that lowers PyTorch, JAX and ONNX models to TT-NN and TT-Metalium. Tenstorrent argues that the CUDA lock-in that slide 35 blames for most startups' failures is broken only by software anyone can read and modify. It also licenses its Ascalon RISC-V CPU as intellectual property (to LG and BOS Semiconductor, and in a 2 nm edge-AI project with Rapidus$^\dagger$), which gives it revenue that does not depend on selling accelerators.

The money has followed. A Series D of \$693 million in December 2024 valued the company at about \$2.6 billion after the investment (led by Samsung Securities and AFW Partners, with Bezos Expeditions, Fidelity, LG and Hyundai), and a Series E of \$1.5 billion in October 2026 doubled the valuation to \$6.2 billion, with Galaxy servers described as ``shipping commercially''$^\dagger$. Tenstorrent has never submitted to MLPerf (slide 32); its performance claims are its own.

\begin{terms}
\term{Tensix}{Tenstorrent's compute cluster: five small RISC-V processors, a matrix engine, a vector engine and 1.5 MB of SRAM, repeated 120 to 140 times in a mesh on each chip.}
\term{RISC-V}{An open instruction-set architecture that anyone may implement without a license fee; Tenstorrent uses small RISC-V cores to control its engines and designs large ones as licensable CPUs.}
\term{network-on-chip (NoC)}{The on-die network of routers and links that carries data between clusters, memory controllers and I/O; Tensix meshes have two.}
\term{software-managed SRAM}{On-chip memory with no cache controller: programs place and move data explicitly, so access time is predictable.}
\term{block floating point}{A number format in which a block of values shares a single exponent and each stores its own mantissa, cheaper than full floating point; Tenstorrent's BFP8 and Microsoft's MX formats are examples.}
\term{GDDR6}{Graphics double-data-rate DRAM: soldered beside the chip on the board, cheaper and lower in bandwidth than HBM, which sits on the package.}
\term{Galaxy}{Tenstorrent's 32-chip server, its chips wired into a 4 $\times$ 8 mesh by on-die Ethernet with no switch.}
\term{TT-Metalium}{Tenstorrent's open-source low-level programming environment: C++ reader, compute and writer kernels for the processors of each Tensix.}
\term{TT-Forge}{Tenstorrent's open-source MLIR-based compiler from PyTorch, JAX and ONNX to its hardware.}
\term{MLIR}{Multi-Level Intermediate Representation: a compiler framework (from the LLVM project) for building domain-specific compilers out of reusable dialects; most new accelerator compilers use it.}
\term{Ascalon}{Tenstorrent's high-performance RISC-V CPU design, sold as licensable intellectual property.}
\end{terms}

\begin{intuition}
Tenstorrent sells the opposite of a GPU's lock-in: cheap memory, Ethernet instead of a proprietary fabric, published prices and published source. Each Tensix is a tiny dataflow machine whose five processors move tiles in, compute and move tiles out, and the programmer can see all of it. Whether openness can substitute for CUDA's head start is the question its 2026 valuation is a bet on.
\end{intuition}

% ---------------------------------------------------------------------
\slide{20}{Korea's accelerators: Rebellions and FuriosaAI}{Latency-first inference chips on Samsung and TSMC 5 nm, then chiplets and HBM; the national fund's picks of 2026.}

\para{What the slide shows} Two small diagrams sit at the top left. The first, ``REBEL-Quad: 4 chiplets, 4 HBM3E'', draws four identical blue compute chiplets (Samsung SF4X) in a row, each with an orange HBM3E stack above it and small purple die-to-die links between neighbors; the caption gives UCIe-Advanced links of about 4 TB/s aggregate (unconfirmed; see below), 144 GB of HBM3E and one 2 PFLOP/s FP8 device. The second, ``RNGD: 8 processing elements of 64 slices'', draws eight processing elements (PEs) in a row, ``fusable; 32 MB SRAM each'', and below them one \emph{slice}: a green contraction engine of eight dot-product units with 512-bit operands, a purple vector engine and an orange transpose engine; its caption says that tensor contraction, not matrix multiplication, is the primitive (ISCA 2024). The table compares three chips: Rebellions ATOM (2023; Samsung 5 nm; 32 TFLOP/s FP16 and 128 TOPS INT8; 64 MB SRAM; 16 GB GDDR6 at 256 GB/s; 60--150 W; MLPerf Inference v3.0), REBEL-Quad (2025--26; four SF4X chiplets; 1{,}024 TFLOP/s FP16 and 2{,}048 FP8; SRAM not published; 144 GB HBM3E at 4{,}800 GB/s; 350 W, then 600 W; no MLPerf) and FuriosaAI RNGD (2024--26; TSMC 5 nm; 512 TFLOP/s FP8 and 256 BF16; 256 MB SRAM; 48 GB HBM3 at 1{,}500 GB/s; 150 W, then 180 W; no official MLPerf). The caption notes that both chips' power rose between product brief and silicon and that Furiosa's first chip, Warboy (2021, 64 TOPS, an INT8 figure), ran MLPerf Inference v1.1. The panel, ``Two paths out of the same market'', tells the two companies' stories.

\para{Details} South Korea has two of the world's three memory makers (Samsung and SK hynix), the world's second-largest foundry (Samsung) and large telecommunications and internet companies (KT, SK Telecom, Naver) that run inference at scale for Korean-language services. It also has a government that decided in 2025--26 to build a domestic AI-chip industry with public money. The two leading startups, Rebellions and FuriosaAI, began with the same target, low-latency inference for those domestic customers, and by 2026 had taken two different technical paths to the same problem: their first data-center chips did not have enough memory for large language models.

Rebellions, founded in September 2020 in Seongnam, designed ATOM as a single 5 nm die for ``latency-critical'' inference (the title of its ISSCC 2024 paper): 64 MB of SRAM, 16 GB of GDDR6 rather than HBM, hardware partitioning into as many as 16 isolated instances so that one card can serve several tenants, and a configurable 60 to 150 W envelope. It went into KT Cloud's public NPU service, the first in Korea, and into SK Telecom's call-summarization product$^\dagger$. It is also one of the few startup chips with an MLPerf result: Rebellions submitted to Inference v3.0 (April 2023), single-stream latency entries, and never again (slide 32). The arithmetic of its memory shows the limit: 256 GB/s of GDDR6 streams the weights of an 8-billion-parameter model in 8-bit precision about 32 times per second (256/8, derived), adequate for a small model and nowhere near enough for a 70-billion-parameter one.

REBEL is Rebellions' answer, and it answers by packaging. Four identical compute chiplets of about 320 mm$^2$ each on Samsung's SF4X (a 4 nm-class process) sit on a silicon interposer, each with its own 12-high HBM3E stack of 36 GB, so the package carries 144 GB at about 4.8 TB/s (four stacks at 1.2 TB/s, derived) and peaks at 1{,}024 TFLOP/s of dense FP16 and 2{,}048 of FP8. The chiplets are joined by UCIe-Advanced, the industry-standard die-to-die interface, with under 12 ns of latency (ISSCC 2026)$^\dagger$; Rebellions' brochure gives 1 TB/s per UCIe-A channel, and the slide's ``about 4 TB/s aggregate'' is a sum that the sources consulted do not print, so treat it as an estimate. Making the chiplets identical means one mask set yields a one-, two- or four-chiplet product. The power figure is the warning: the November 2024 brief said 350 W, the ISSCC 2026 paper 600 W$^\dagger$. A vendor's first number for a chip that has not shipped describes a target. At Hot Chips 2025 Rebellions demonstrated Llama 3.3 70B on a development board at 35.5 ms per token, about 28 tokens per second$^\dagger$.

The company's finances are a national project. Rebellions merged with SK Telecom's chip unit Sapeon in December 2024, with Rebellions' management keeping control; raised \$250 million at \$1.4 billion in September 2025 with Arm and Samsung as strategic investors; and in March 2026 raised about \$400 million at about \$2.3 billion in a round led by Mirae Asset and the Korea National Growth Fund, the fund's first ``K-NVIDIA'' investment$^\dagger$. The chief executive has said the company plans to list on the Korean exchange (KOSPI) in 2027.

FuriosaAI, founded in 2017 in Seoul by June Paik, took the other path. Its first chip, Warboy, a vision-inference part, submitted to MLPerf Inference v1.1 (September 2021) and v2.0 (April 2022); Furiosa reported about 1.5$\times$ an NVIDIA T4 in single-stream image tasks, its own framing$^\dagger$. Its second, RNGD (``Renegade''), presented at Hot Chips 2024 and at the industry track of ISCA 2024, is a single large TSMC 5 nm die with two HBM3 stacks (48 GB, 1.5 TB/s) and 256 MB of SRAM, at 150 W in the 2024 description and 180 W on the current specification page$^\dagger$. Its distinctive idea is in the slice diagram. A matrix multiply is a special case of a \emph{tensor contraction}, a sum over shared indices of two multi-dimensional arrays; attention, convolution and batched matrix products are all contractions with different index patterns. A GPU or a TPU reshapes each of them into matrix multiplies, often by copying data into a new layout. RNGD's compiler maps the contraction itself onto slices whose contraction engines take 512-bit operand vectors (32 BF16 or 64 FP8 values) from software-managed SRAM, and a separate transpose engine does the layout changes so they do not consume arithmetic cycles. Eight PEs of 64 slices each give the peak: 8 $\times$ 64 $=$ 512 slices, and 512 TFLOP/s of FP8 implies 1 TFLOP/s per slice at the 1 GHz clock (derived). The eight PEs can be fused into larger virtual devices so one card serves a model of about 10 billion parameters or several smaller ones.

RNGD's customers and deals tell the rest. LG AI Research adopted it in July 2025 to serve its EXAONE models, reporting up to 2.25$\times$ the performance per watt of A100-class GPUs (LG's figure) and EXAONE 32B at 50--60 tokens per second at batch 1 on four cards$^\dagger$. Earlier in 2025 Meta approached Furiosa with an acquisition offer reported at about \$800 million; Furiosa declined, reportedly over strategic direction rather than price$^\dagger$. Mass production began in January 2026, and on October 30, 2026 Furiosa raised \$620 million at about \$2.9 billion, \$280 million of it from the National Growth Fund, with a third-generation accelerator to be co-developed with Broadcom$^\dagger$. The chief executive's own verdict on RNGD, given in July 2025, is the slide's last line: it needed four HBM stacks and more than one compute die. In other words, the next Furiosa chip will look like REBEL.

Neither RNGD nor REBEL has an MLPerf result, and their performance claims are vendor-run; slide 32 explains why that matters.

\begin{formulas}
Package memory from stacks, and the per-unit peak:
\[
C = n_{\text{stacks}}\times c_{\text{stack}},\qquad B = n_{\text{stacks}}\times b_{\text{stack}},\qquad
P = N_{\text{units}}\times p_{\text{unit}} .
\]
REBEL-Quad: $4\times 36 = 144$ GB, $4\times 1.2\approx 4.8$ TB/s. RNGD: $2\times 24 = 48$ GB; $512\ \text{slices}\times 1\ \text{TF} = 512$ TF FP8. ATOM decode bound for an 8B INT8 model: $256\,\text{GB/s}/8\,\text{GB}\approx 32$ tokens/s (derived).
\end{formulas}

\begin{terms}
\term{chiplet}{A die designed to be combined with others in one package; REBEL-Quad uses four identical compute chiplets.}
\term{UCIe}{Universal Chiplet Interconnect Express: an industry standard for die-to-die links inside a package; the ``Advanced'' variant targets silicon interposers and bridges with dense wiring.}
\term{interposer}{A slab of silicon under several dies and memory stacks that carries the fine-pitch wiring between them; TSMC's CoWoS-S is the common form.}
\term{HBM3E}{The enhanced third generation of high-bandwidth memory, with stacks of 8 or 12 DRAM dies at roughly 1 to 1.2 TB/s per stack.}
\term{tensor contraction}{A sum over one or more shared indices of two multi-dimensional arrays; matrix multiplication, batched matrix products, convolution and attention are special cases.}
\term{processing element (PE)}{A repeated compute unit of an accelerator; on RNGD, one of eight blocks of 64 slices with its own control core and 32 MB of SRAM.}
\term{slice}{On RNGD, the smallest compute unit: a contraction engine, a vector engine and a transpose engine fed from a portion of the SRAM.}
\term{single-stream latency}{An MLPerf Inference scenario that sends one query at a time and measures the time to answer it; the measure that latency-first chips such as ATOM were built for.}
\term{multi-instance partitioning}{Dividing one accelerator into isolated hardware partitions that serve different tenants; ATOM supports up to 16.}
\term{Korea National Growth Fund}{The South Korean government's 2026 investment vehicle for strategic industries, whose first AI-chip investments were Rebellions and FuriosaAI.}
\end{terms}

\begin{intuition}
Both Korean companies built a good small-memory inference chip first and then discovered what every company in this lecture discovered: language models are sold by the gigabyte of fast memory. REBEL fixed it with four chiplets and four HBM stacks; Furiosa's next chip will do the same. The architecture differences (contraction engines, chiplet links) matter less than the number of HBM stacks on the package.
\end{intuition}
% =====================================================================
\setstage{posc}
\section{The hyperscalers' chips and the established vendors (slides 21--24)}

\slide{21}{AWS Inferentia and Trainium: captive volume, four generations}{A systolic NeuronCore, HBM and a torus; UltraServers reach 144 chips, and Anthropic's Rainier runs 500,000.}

\para{What the slide shows} The top left is a block diagram of one NeuronCore, the compute unit of AWS's chips from the second version on$^\dagger$: four engines in a row (a tensor engine labeled ``systolic array'', a vector engine, a scalar engine, and a ``GPSIMD'' engine of eight C/C++ cores on which custom kernels written in NKI run), above a software-managed on-chip SRAM scratchpad, with a box for HBM and NeuronLink at the right. The caption says that Trainium2 has eight NeuronCores on two dies with four HBM3 stacks, that 16 chips form a 2D torus inside a Trn2 instance, and that Trainium3 adds a switched fabric. The table compares four chips: Trainium (GA October 2022; 2 NeuronCore-v2; 190 TFLOP/s BF16; 32 GB HBM2e at 0.82 TB/s; NeuronLink; 16 chips per server), Inferentia2 (April 2023; the same cores and peak; 32 GB HBM at 0.82 TB/s; 384 GB/s links; 12 chips), Trainium2 (December 2024; 8 NeuronCore-v3; 1.3 PFLOP/s FP8; 96 GB HBM3 at 2.9 TB/s; 1 TB/s torus links; 64 chips in an UltraServer) and Trainium3 (December 2025; 8 cores plus communication cores; 2.5 PFLOP/s MXFP8; 144 GB HBM3e at 4.9 TB/s; 2.6 TB/s chip links into a switched fabric; 144 chips per UltraServer). The caption says the sources are AWS pages and blogs, that Trainium2's 5.2 PFLOP/s is with 4:1 sparsity, that the Trainium3 per-chip figure is derived from 362 PFLOP/s per 144 chips, that die sizes and power are not published, and that AWS has made no MLPerf submission. The panel, ``Why AWS builds its own'', lists price-performance, the growth of scale-up, Project Rainier, Trainium4 and the software.

\para{Details} Amazon's chip business began with the 2015 purchase of Annapurna Labs, an Israeli designer of networking and storage processors whose Nitro cards now sit in every EC2 server, and whose Graviton Arm CPUs (from 2018) showed that AWS could design a competitive processor and sell it only as cloud capacity. Inferentia (announced November 2018, instances from December 2019) and Trainium (announced December 2020, generally available October 2022) apply the same model to machine learning: AWS designs the chip, Marvell has been its reported design partner (slide 23), TSMC makes it, and customers never buy it; they rent instances. That is what the slide's title means by \emph{captive volume}. A merchant startup must win each customer against NVIDIA; AWS's chips have a guaranteed buyer in AWS itself and in the customers it steers to them with price.

The NeuronCore is a conventional design for a training accelerator, closer to a TPU core (slide 9) than to anything else in this section. Its tensor engine is a systolic array (slide 7) fed from an on-chip scratchpad; the vector engine handles normalizations and activations; the scalar engine handles elementwise functions and control; and the GPSIMD engine, eight general-purpose cores programmable in C and C++, gives developers a place to run operations that the fixed engines do not support. The Neuron Kernel Interface (NKI, beta in September 2024) is a Python-embedded kernel language that exposes these engines and the scratchpad directly, the same role Triton plays on GPUs and Pallas on TPUs. Above it, the Neuron SDK provides PyTorch and JAX support (JAX through the PJRT plugin interface), a distributed-training library and a vLLM-based serving path.

The table's peaks are dense and come from AWS's own pages. Inferentia2 and Trainium share the NeuronCore-v2 and the 190 TFLOP/s BF16 peak; Inferentia2 is the inference variant with cheaper links. Trainium2 moved to two compute dies and four HBM3 stacks per package: 1.3 PFLOP/s of dense FP8 (AWS's headline 5.2 PFLOP/s assumes 4:1 sparsity, four times dense), and, per the Neuron documentation, 667 TFLOP/s of dense BF16$^\dagger$. Trainium3, AWS's first 3 nm chip, is described only at the system level: a Trn3 UltraServer of 144 chips delivers 362 PFLOP/s of MXFP8 with 20.7 TB of HBM3e at 706 TB/s, from which the per-chip numbers follow, 362/144 $\approx$ 2.5 PFLOP/s, 20.7 TB/144 $\approx$ 144 GB and 706/144 $\approx$ 4.9 TB/s (derived). MXFP8 is the microscaling 8-bit format of the Open Compute Project, in which blocks of 32 values share a scale factor (slide 22 describes the family). AWS publishes no power figure and no die size for any of the chips.

The ``9$\times$ in a year'' of the panel is the growth of the scale-up domain, the set of chips joined by the fast chip-to-chip fabric. A Trn2 instance has 16 chips in a 2D torus (20.8 PFLOP/s FP8 dense, 16 $\times$ 1.3, derived); a Trn2 UltraServer joins four instances into 64 chips (83 PFLOP/s, 6 TiB of HBM, 64 $\times$ 96 GiB); a Trn3 UltraServer joins 144 chips through NeuronSwitch, AWS's first switched scale-up fabric, so that every chip reaches every other in one hop instead of several torus hops; each Trainium3 drives 2.56 TB/s of NeuronLink-v4 (the table's 2.6 TB/s) into it, per the Neuron documentation$^\dagger$. That is the same move NVIDIA made from the 8-GPU HGX board to the 72-GPU NVL72 rack, and for the same reason: a mixture-of-experts model's all-to-all communication runs badly on a torus.

Project Rainier is the deployment that makes Trainium matter. On October 29, 2025 AWS opened an \$11 billion campus in Indiana with ``nearly 500{,}000'' Trainium2 chips running for Anthropic and said Anthropic expected to be on more than one million Trainium2 by the end of the year, which the slide quotes as a target$^\dagger$. What was reported in December 2025 was different in scope: Andy Jassy said more than a million Trainium2 were in production across all of AWS, and AWS said Anthropic was using over 500{,}000 of them$^\dagger$. The million is a fleet-wide count, not Rainier's or Anthropic's. Half a million Trainium2 at 1.3 PFLOP/s each is about 650 EFLOP/s of dense FP8 (derived), in one company's service, without a single MLPerf result. Anthropic trains and serves Claude on these chips, as well as on Google's TPUs (slide 12) and NVIDIA GPUs. AWS claims 30--40\% better price-performance for Trn2 instances than for its H200 (P5e) instances$^\dagger$; the claim is AWS's own and compares rental prices, which AWS sets.

Trainium4, announced at re:Invent in December 2025, is described mainly by one decision: it will use NVIDIA's NVLink Fusion, the program under which NVIDIA licenses its scale-up fabric to other chips, so Trainium4 racks can share a fabric with NVIDIA GPUs and Grace CPUs$^\dagger$. For the first time a hyperscaler's own accelerator will sit on its competitor's interconnect; slide 36 counts it among the ways the GPU absorbs its rivals.

\begin{formulas}
System figures from chip figures (and back):
\[
P_{\text{server}} = N\times P_{\text{chip}},\qquad
\text{ridge} = \frac{P_{\text{chip}}}{B_{\text{chip}}} .
\]
Trn2: $16\times1.3 = 20.8$ PF; $64\times1.3\approx83$ PF; ridge $1.3\,\text{PF}/2.9\,\text{TB/s}\approx 450$ FLOP/byte. Trn3: $362/144\approx2.5$ PF; ridge $2.5/4.9\approx 510$ FLOP/byte. Rainier: $5\times10^{5}\times1.3\,\text{PF}\approx 650$ EF FP8 dense (derived).
\end{formulas}

\begin{terms}
\term{NeuronCore}{The compute unit of AWS's Inferentia and Trainium chips: a systolic tensor engine, a vector engine, a scalar engine and (from version 2) a general-purpose GPSIMD engine around an on-chip scratchpad.}
\term{Inferentia}{AWS's inference accelerator line (2019, 2023), sold only as EC2 Inf instances.}
\term{Trainium}{AWS's training accelerator line (2022, 2024, 2025, with Trainium4 announced), sold only as EC2 Trn instances and UltraServers.}
\term{captive volume}{Demand guaranteed by the chip designer's own business; a hyperscaler's chip has its own cloud as a buyer, which a merchant startup lacks.}
\term{GPSIMD engine}{The NeuronCore's block of eight general-purpose cores that run custom operators written in C, C++ or NKI.}
\term{NKI}{Neuron Kernel Interface: AWS's Python-based language for writing kernels that program the NeuronCore engines and scratchpad directly.}
\term{2D torus}{A grid network whose rows and columns wrap around, so every chip has four neighbors and no edge; Trn2 instances and TPU v5e pods use it.}
\term{scale-up domain}{The set of accelerators joined by the fast, memory-semantic chip-to-chip fabric, as opposed to the slower scale-out network between such domains.}
\term{UltraServer}{AWS's multi-instance scale-up system: 64 Trainium2 or 144 Trainium3 chips joined by NeuronLink.}
\term{MXFP8}{Microscaling 8-bit floating point, an Open Compute Project format in which each block of 32 FP8 values shares one 8-bit scale factor.}
\term{NVLink Fusion}{NVIDIA's program, from 2025, that licenses NVLink scale-up connectivity to third-party CPUs and accelerators.}
\end{terms}

\begin{intuition}
Trainium is not trying to beat NVIDIA's chip; it is trying to beat NVIDIA's margin. A conventional systolic design, sold only as rented capacity to customers the cloud already has, wins if it is good enough and much cheaper to make, and Rainier's half-million chips show that it has been. The engineering story of the generations is the fabric: from a 16-chip torus to a 144-chip switch to NVIDIA's own NVLink.
\end{intuition}

% ---------------------------------------------------------------------
\slide{22}{Microsoft Maia and Meta MTIA: inference first, then everything}{Maia runs on Ethernet; MTIA grew from a 25 W recommendation chip into a 1.7 kW HBM part in three years.}

\para{What the slide shows} The top left draws the Maia 100 die as Microsoft presented it at Hot Chips 2024$^\dagger$: a reticle-sized square divided into 16 clusters, each of four tiles (a tensor unit for the MX formats plus a vector unit) with an orange strip of cluster SRAM; beside it three boxes give the chip's external connections: 12 $\times$ 400 Gb Ethernet (4.8 Tb/s per chip), 64 GB of HBM2e at 1.8 TB/s, and three switch planes. The table compares five chips: Maia 100 (announced November 2023; TSMC N5; 820 mm$^2$; 0.8 PFLOP/s BF16; 500 MB SRAM; 64 GB HBM2e at 1.8 TB/s; 700 W; 4.8 Tb/s Ethernet), Maia 200 (January 2026; N3P; die size not published; 5.1 PFLOP/s FP8; 272 MB; 216 GB HBM3e at 7 TB/s; 750 W; 2.8 TB/s Ethernet), MTIA v1 (May 2023; 7 nm; 373 mm$^2$; 102 TOPS INT8; 128 MB; 64 GB LPDDR5 at 176 GB/s; 25 W; PCIe 4), MTIA v2 (April 2024; 5 nm; 421 mm$^2$; 354 TOPS INT8; 256 MB; 128 GB LPDDR5 at 205 GB/s; 90 W; PCIe 5) and MTIA 400 (March 2026; chiplets; 6 PFLOP/s FP8, a figure consistent with Meta's MTIA 450 numbers but not printed by Meta; SRAM not published; 288 GB HBM at 9.2 TB/s; 1{,}200 W; 72 chips per rack). The caption gives the sources and notes that Maia 100 has also been quoted at 3 POPS for 6-bit MX. The panel, ``Two companies, two reasons'', explains what each company wanted and what the two have in common.

\para{Details} Microsoft and Meta are the two largest buyers of NVIDIA GPUs after the cloud providers that resell them, and both started building their own accelerators for the same reason AWS did (slide 21), but with different workloads in mind. Microsoft's is serving OpenAI's and its own models in Azure, where the measure is tokens per dollar. Meta's was, for years, ranking and recommendation, the deep-learning recommendation models (DLRMs) that choose what appears in every Facebook and Instagram feed, whose mix of huge embedding tables and modest dense layers fits a GPU poorly.

Maia 100 was announced at Ignite in November 2023 and described at Hot Chips 2024: a reticle-limited TSMC N5 die of 820 mm$^2$ and about 105 billion transistors, 64 tiles in 16 clusters, about 500 MB of on-chip SRAM, four HBM2e stacks, and two unusual choices$^\dagger$. The first was its number formats. Maia computed in Microsoft's \emph{MX} (microscaling) formats, block floating-point types in which a block of values shares a scale, at 4, 6 and 9 bits; Microsoft, with AMD, Arm, Intel, Meta, NVIDIA and Qualcomm, then standardized the idea through the Open Compute Project as MXFP8, MXFP6, MXFP4 and MXINT8 (2023). The second was its fabric. Instead of a proprietary scale-up link, every Maia 100 has twelve 400 Gb Ethernet ports carrying a custom RDMA protocol through three planes of Ethernet switches, so the scale-up network uses the same commodity switches as the data-center network. The peak figures for Maia 100 are a known mess: the Hot Chips slides gave 3 POPS at 6-bit, 1.5 at 9-bit and 0.8 PFLOP/s at BF16, while 2023 press reports gave 1.6 PFLOP/s MXINT8 and 3.2 MXFP4; both sets cannot be dense peaks of the same silicon, and the slide uses the Hot Chips BF16 figure. Maia 100 saw little use beyond internal testing, and its successor, Maia 200 itself (codename Braga), slipped a year, from 2025 to 2026, according to press reports, through design changes and staff turnover$^\dagger$.

Maia 200, announced on January 26, 2026, is the part that went into production. It moves to TSMC N3P, drops the bespoke MX6 and MX9 formats for standard FP8 and FP4 (10.15 PFLOP/s FP4 and 5.07 FP8, dense; Microsoft publishes no BF16 figure), carries 216 GB of HBM3e at 7 TB/s and 272 MB of SRAM at 750 W, and keeps Ethernet scale-up at 2.8 TB/s per chip, in domains of up to 6{,}144 chips$^\dagger$. Microsoft says it serves GPT-5.2, and it runs only inside Azure, deployed first in Iowa and then Arizona, with no external sale (slide 34); a blade carries four Maia 200 and one Cobalt 200 Arm CPU. Its memory-to-compute ratio is the clearest sign of the purpose: 7 TB/s against 5.07 PFLOP/s of FP8 is a ridge point of about 720 FLOP per byte (derived); against Maia 100's 0.8 PFLOP/s BF16 on 1.8 TB/s (about 440) the compute grew faster than the bandwidth, but the capacity grew most, from 64 to 216 GB, which is what serving large models with long contexts needs (slides 4 and 8).

Meta's MTIA (Meta Training and Inference Accelerator) has the best-documented first generation in this section, because Meta published it at ISCA 2023 (Firoozshahian et al.). MTIA v1 is a 7 nm, 373 mm$^2$ chip of 64 processing elements in an 8 $\times$ 8 grid, each with two RISC-V cores and fixed-function dot-product and SIMD units, sharing 128 MB of SRAM, with 64 GB of LPDDR5 rather than HBM, at 25 W and 102.4 TOPS INT8. That is 4.1 TOPS per watt (102.4/25, derived), and the LPDDR5 choice is the point: a recommendation model's embedding tables are large and accessed sparsely, so capacity per dollar matters more than bandwidth, and LPDDR gives far more of it. MTIA v2 (April 2024) kept the grid on 5 nm with 256 MB of SRAM and 128 GB of LPDDR5 at 90 W; its headline 708 TOPS is with sparsity, 354 dense$^\dagger$. Seventy-two v2 chips fit a rack of three chassis.

The 2026 family, announced on March 11, 2026 with Broadcom as the design partner, is a different design point altogether: HBM chiplet parts (MTIA 300, 400, 450 and 500) of 1.2 to 1.7 kW, 6 to 10 PFLOP/s of FP8, 288 to 512 GB of HBM, in 72-chip racks with a switched scale-up backplane, aimed at generative AI as well as ranking$^\dagger$. The slide's MTIA 400 column (288 GB at 9.2 TB/s, 1{,}200 W) rests on press summaries of Meta's post, which disagree about which numbers belong to which model; treat the per-model attribution as uncertain. Going from 25 W to 1.7 kW (68$\times$, derived) is the subtitle's point, with one qualification: the 1.7 kW part is the MTIA 500, due for mass deployment in 2027 and not yet shipped, so ``three years'' counts from MTIA v1 (2023) to the announcement, and it is four years to deployment: a recommendation chip became a general accelerator because Meta's workload became language models. Meta hedges all the same, renting Google TPUs and signing a 6 GW agreement with AMD in 2026 (slide 34).

What the two programs share is listed at the end of the panel. Both put PyTorch and Triton front ends on their chips, so models written for NVIDIA GPUs can be compiled for them; neither has ever submitted to MLPerf (slide 32); and both run the owner's own workloads, which is why neither needs to.

\begin{terms}
\term{Maia}{Microsoft's in-house AI accelerator line: Maia 100 (2023--24, internal), Maia 200 (2026, in production for inference, only inside Azure).}
\term{MTIA}{Meta Training and Inference Accelerator: Meta's in-house line, from the 25 W recommendation-inference v1 (2023) to the 1.2--1.7 kW HBM chiplet parts of 2026.}
\term{MX formats}{Microscaling formats: block floating-point types (MXFP8, MXFP6, MXFP4, MXINT8) in which each block of 32 values shares an 8-bit scale; standardized through the Open Compute Project in 2023.}
\term{DLRM}{Deep-learning recommendation model: a network that combines large sparse embedding-table lookups with dense layers to rank items; Meta's dominant inference workload before generative AI.}
\term{embedding table}{A large table of learned vectors indexed by a categorical feature (a user, a page, an ad); recommendation models hold billions of rows and read few of them per query.}
\term{LPDDR5}{Low-power DDR5 DRAM, designed for phones: cheaper and lower in bandwidth than HBM, with more capacity per dollar.}
\term{RDMA}{Remote direct memory access: one machine reading or writing another's memory over the network without involving its CPU; Maia's Ethernet fabric carries a custom RDMA protocol.}
\term{ridge point}{The arithmetic intensity, peak FLOP/s divided by memory bandwidth, above which a kernel is compute-bound and below which it is bandwidth-bound.}
\term{Cobalt}{Microsoft's Arm server CPU; a Cobalt 200 hosts four Maia 200 on a blade.}
\end{terms}

\begin{intuition}
Microsoft and Meta built chips for the work they already had, which is the advantage no startup can buy. Both first versions were tuned to a past workload (Maia's private formats, MTIA's recommendation models) and both second versions converged on what everyone else converged on: standard FP8 and FP4, lots of HBM, and a switched rack.
\end{intuition}

% ---------------------------------------------------------------------
\slide{23}{OpenAI--Broadcom and the custom-XPU model}{The buyer owns the architecture, Broadcom the silicon and SerDes, TSMC the wafers; 10 GW for OpenAI by 2029.}

\para{What the slide shows} The upper left is a flow diagram of how a custom accelerator gets built, 2023--26$^\dagger$. Three boxes are joined left to right: the buyer's architects (Google, Meta, OpenAI, Apple, ByteDance), passing ``RTL'' to an ASIC house (Broadcom, Marvell, MediaTek, Alchip: physical design, SerDes, packaging, test), passing ``GDS'' to the foundry and memory makers (TSMC N3/N5 and CoWoS; HBM from SK hynix, Samsung and Micron). An arrow runs down to a purple bar for the rack: Ethernet scale-up (Broadcom's Tomahawk switches) or the buyer's own fabric, liquid cooling, the buyer's cloud. The caption: the buyer keeps the architecture and the software; the ASIC house turns it into a chip in 12--18 months and sells it back at a margin far below NVIDIA's. Below, a bar chart shows Broadcom's AI semiconductor revenue (custom accelerators and Ethernet): \$3.8 billion in fiscal 2023, \$12.2 billion in fiscal 2024, about \$20 billion in fiscal 2025 (labeled ``(d)'' on the chart, though Broadcom stated the figure), and \$8.2 billion guided for the first quarter of fiscal 2026 alone. The panel, ``The deals, and what they buy'', describes the OpenAI--Broadcom agreement, Broadcom's customer book, and the benefits and costs of the model.

\para{Details} Between the hyperscalers' fully in-house programs (slides 9--12, 21, 22) and the merchant startups (slides 13--20) sits a third way to get a non-GPU accelerator, which by 2026 had become the largest by money: a company that knows its workload designs the architecture and hires a specialist to turn it into a chip. Broadcom is widely reported to have done this for Google since the first TPU (slide 9), though the sources consulted do not document the v1 relationship; for the later generations Google's TPU is, in manufacturing terms, a Broadcom product. The industry's name for such a part is an \emph{XPU}, an accelerator of unspecified kind built to one customer's design.

The division of labor in the diagram follows the stages of chip design. The buyer writes the architecture and its register-transfer-level description (RTL), the hardware description in a language such as Verilog or SystemVerilog that says what each block does every clock. It also writes the compiler and software, which is most of the long-term effort. The ASIC house does \emph{physical design}: synthesizing the RTL into gates, placing and routing billions of transistors to meet timing and power, and integrating its own intellectual property, above all the high-speed serializer--deserializer (SerDes) circuits that drive the chip's network and memory links, where Broadcom has the industry's best portfolio. It produces the GDS (the geometric layout file sent to the foundry), arranges TSMC capacity, buys HBM, does the 2.5D packaging (TSMC's CoWoS), tests the parts and often builds the boards and racks. The buyer pays for engineering and then per chip, at a gross margin that analysts put well below the roughly 75\% gross margin NVIDIA has reported company-wide since 2024$^\dagger$; the slide's ``far below NVIDIA's'' is a qualitative statement, not a disclosed figure.

The OpenAI--Broadcom agreement, announced on October 13, 2025, is the largest of these by declared scale: 10 gigawatts of custom accelerators and racks, with OpenAI designing the accelerator and the system and Broadcom developing and deploying them over Ethernet, starting in the second half of 2026 and finishing by the end of 2029$^\dagger$. The two companies said they had worked together for more than 18 months. Reuters had reported in February 2025 that the design was a TSMC 3 nm part with a systolic array and HBM, built by a team of about 40 led by Richard Ho, who had worked on Google's TPUs$^\dagger$. As of November 2026 no shipped rack had been reported. The agreement sits beside two others OpenAI signed in the same weeks, a 10 GW letter of intent with NVIDIA (September 22, 2025) and 6 GW with AMD (October 6, 2025), so the custom chip is one supply line among three rather than a replacement.

Gigawatts are a power budget, not a chip count, and the conversion is worth doing once. If a deployed accelerator with its share of CPUs, networking, cooling and facility overhead draws on the order of 1.5 to 2 kW, 10 GW corresponds to roughly 5 to 7 million accelerators (an illustrative assumption, not a disclosed figure), and at any plausible price per installed system it is a program of hundreds of billions of dollars over four years (again illustrative). That is why such announcements are read as commitments of intent, and why slide 34 separates commitments from deployments.

Broadcom's own disclosures show the model working. Its AI semiconductor revenue was \$12.2 billion in fiscal 2024, up 220\% (so about \$3.8 billion in fiscal 2023, derived); Broadcom stated fiscal 2025 AI revenue as \$20 billion, up 65\%, on its fourth-quarter call, so the chart's ``(d)'' for derived can be dropped; and it guided to \$8.2 billion for the first quarter of fiscal 2026 alone, double the year before$^\dagger$. Broadcom says it has five XPU customers. Google's TPU is the first. Meta's MTIA (slide 22) is another. In September 2025 Broadcom booked a \$10 billion order from an unnamed customer and in December an additional \$11 billion from the same customer for delivery in late 2026; the customer was reported to be Anthropic, buying TPUs (slide 12), and Broadcom said it was not OpenAI$^\dagger$. Apple (a part reported under the codename Baltra) and ByteDance are reported customers that Broadcom has not confirmed. Marvell plays the same role for AWS (Trainium) and Microsoft, MediaTek is reported for an inference TPU variant, and Alchip and others serve smaller buyers.

The panel's last two lines are the trade. The custom model gives the buyer an accelerator with its workload built in, memory and fabric sized to its models, and a chip bought near cost, which matters when the alternative's margin is three-quarters of the price. It costs two to three years from the first architecture to a deployed rack, a compiler and runtime team of hundreds, and all the risk: if the design is wrong for the models of 2028, the buyer owns the mistake. Only a company spending tens of billions a year on accelerators can amortize that, which is why the list of XPU customers is a list of the largest AI spenders.

\begin{formulas}
Illustrative conversion of a power commitment into accelerators, and growth from reported revenue:
\[
N_{\text{acc}}\approx\frac{P_{\text{total}}}{p_{\text{per accelerator, all-in}}},\qquad
R_{t-1}=\frac{R_t}{1+g} .
\]
$10\,\text{GW}/(1.5\text{--}2\,\text{kW})\approx 5$--$7$ million accelerators (illustrative assumption). Broadcom FY2023 AI revenue: $12.2/(1+2.2)\approx 3.8$ billion dollars (derived).
\end{formulas}

\begin{terms}
\term{XPU}{Industry shorthand for a custom accelerator built to one customer's architecture by an ASIC design partner; Google's TPU, Meta's MTIA and OpenAI's planned chip are XPUs.}
\term{ASIC house}{A company that turns a customer's chip design into manufactured silicon: physical design, intellectual-property blocks, foundry and packaging management, test; Broadcom, Marvell, MediaTek and Alchip are examples.}
\term{RTL}{Register-transfer level: a hardware description, in a language such as Verilog, of what each block of a chip does on each clock cycle; the hand-off from architects to physical designers.}
\term{physical design}{Turning RTL into a manufacturable layout: synthesis to gates, placement, routing, timing and power closure.}
\term{GDS}{The layout database (GDSII or its successor OASIS) describing every shape on every mask layer; the file a foundry manufactures from.}
\term{SerDes}{Serializer--deserializer: the analog circuit that sends and receives very high-rate serial data over a link; it limits every chip-to-chip network and is a core Broadcom competence.}
\term{CoWoS}{Chip-on-wafer-on-substrate: TSMC's 2.5D packaging that places logic dies and HBM stacks on a silicon interposer; its capacity has been a bottleneck for every accelerator.}
\term{Tomahawk}{Broadcom's line of high-radix Ethernet switch chips, used for both scale-out and Ethernet-based scale-up networks.}
\term{gigawatt commitment}{An agreement denominated in the electrical power of the data centers to be built, a proxy for the number of accelerators that is not a chip count.}
\end{terms}

\begin{intuition}
The custom-XPU model separates knowing what to build from knowing how to build it. The buyer supplies the workload, the architecture and the software; Broadcom supplies everything that makes a design into a few million working packages. It is how the TPU was always made, and from 2025 it became how every company large enough to afford a compiler team tried to escape NVIDIA's margin.
\end{intuition}

% ---------------------------------------------------------------------
\slide{24}{The established vendors: Gaudi, Cloud AI, Spyre, Dojo}{Intel bought one line and killed another; Qualcomm bets on LPDDR; IBM sits beside the mainframe; Tesla paused.}

\para{What the slide shows} The upper left is a table of five chips from four established companies$^\dagger$: Intel's Gaudi 2 (2022; 7 nm; 24 tensor processor cores (TPC) and matrix engines (MME); 432 TFLOP/s BF16; 48 MB SRAM; 96 GB HBM2e at 2.46 TB/s; 600 W; 24 $\times$ 100 Gb Ethernet) and Gaudi 3 (2024; 5 nm, two dies; 64 TPC and 8 MME; 1{,}835 TFLOP/s; 96 MB; 128 GB HBM2e at 3.67 TB/s; 900 or 1{,}200 W; 24 $\times$ 200 GbE); Qualcomm's Cloud AI 100 Ultra (2023; 7 nm, four AI 100 SoCs on one card; 64 AI cores; 870 TOPS; 576 MB SRAM; 128 GB LPDDR4x at 548 GB/s; 150 W; a PCIe card); IBM's Spyre (2025; 5 nm; 32 cores; 300 TOPS; 64 MB; 128 GB LPDDR5, bandwidth not published; 75 W; ``8 per drawer''); and Tesla's D1 (2021; 7 nm; 354 nodes; 362 TFLOP/s BF16; 440 MB SRAM; no DRAM on the die; 2 TB/s per edge; 400 W; 25 dies per training tile). The caption gives the sources and their quirks (Gaudi 3 gives one figure for BF16 and FP8; IBM gives no LPDDR bandwidth; Tesla's 2021 slides gave 4 TB/s per edge, its Hot Chips 2022 slides 2). Two stat tiles give \$2B (Intel's purchase of Habana, December 2019) and 200 MW (Humain's commitment to Qualcomm's AI200, 2026). The panel, ``Four fates'', gives each line's outcome.

\para{Details} This slide collects four companies that were already large before they built AI accelerators, and whose chips illustrate four different outcomes. None of them is a startup, so none of them failed for lack of money; their stories are about strategy.

\emph{Intel} entered twice. In August 2016 it bought Nervana Systems, a deep-learning startup, for a price reported at \$350--400 million$^\dagger$, and in 2019 showed two Nervana chips at Hot Chips: NNP-T for training (TSMC 16 nm, 24 tensor processing clusters, 32 GB of HBM2) and NNP-I for inference. In December 2019 it bought Habana Labs, an Israeli startup whose Gaudi training chip already worked, for about \$2 billion, and in January 2020 cancelled NNP-T within weeks of its launch$^\dagger$: Intel ``bought one line and killed another''. Gaudi 2 (2022) is a heterogeneous design: matrix multiply engines (MMEs) for the dense products and programmable VLIW tensor processor cores (TPCs) for everything else, with 96 GB of HBM2e and, distinctively, 24 ports of 100 Gb Ethernet with RDMA on the die, so that Gaudi servers scale out over standard Ethernet with no separate network card. Gaudi was one of the few non-NVIDIA chips in MLPerf Training, from v1.0 (2021) through v4.0 (June 2024), when 1{,}024 Gaudi 2 trained the GPT-3 175B benchmark in 66.9 minutes$^\dagger$. Google's TPUs also ran GPT-3 in those rounds, so the slide's claim is that Gaudi 2 was the only \emph{merchant} non-NVIDIA part, one sold to anyone, with GPT-3 training results (slide 32). Gaudi 3 (2024) doubled almost everything, two 5 nm dies, 1{,}835 TFLOP/s of dense BF16 or FP8 (Intel gives one figure for both), 128 GB, but it never appeared in MLPerf and Intel said it missed its \$500 million revenue target for 2024. On January 30, 2025 Intel cancelled Falcon Shores, the planned successor that was to merge Gaudi with its GPU line, as a product, leaving Gaudi 3 as the last Gaudi and a rack-scale GPU, ``Jaguar Shores'', as a plan$^\dagger$. What ended in 2025 was the roadmap; Gaudi 3 itself kept shipping.

\emph{Qualcomm} took its smartphone neural processing unit (the Hexagon NPU lineage), scaled it up and put it on a PCIe card. The Cloud AI 100 (announced 2019, in production 2021) and the Ultra (2023), a card carrying four AI 100 SoCs of 16 cores and 144 MB of SRAM each (4 $\times$ 16 $=$ 64 cores, 4 $\times$ 144 $=$ 576 MB of SRAM, derived; 870 TOPS INT8 dense, 128 GB of LPDDR4x at 548 GB/s, 150 W$^\dagger$) are designed for performance per watt, and they were the one non-NVIDIA accelerator that submitted MLPerf Inference results, including power measurements, round after round: Qualcomm from v1.0 (April 2021) through v4.0 (March 2024), and the independent submitter Krai in later rounds$^\dagger$. The efficiency is easy to compute: 870 TOPS at 150 W is 5.8 TOPS per watt (derived), against roughly 2 to 3 for a contemporary data-center GPU at INT8. The bet on LPDDR rather than HBM is the same as Meta's MTIA v1 (slide 22): more capacity per dollar and per watt, much less bandwidth. In October 2025 Qualcomm announced the AI200 (2026), with 768 GB of LPDDR per card in 160 kW liquid-cooled racks, and the AI250 (2027), which adds near-memory computing (slide 28), and named Humain, the Saudi AI company, as the first customer with a 200 MW commitment$^\dagger$. Qualcomm has published no TOPS or FLOP/s figures for either, and as of November 2026 no shipment had been reported.

\emph{IBM} builds accelerators to sit beside its mainframes. Its Telum processor (2021) put an AI accelerator on the CPU die so that a bank could score every card transaction for fraud inside the transaction itself, without sending data off the machine. Spyre (Hot Chips 2024) is the external companion: a Samsung 5 nm chip of 32 cores, 300 TOPS, 128 GB of LPDDR5 and 75 W on a PCIe card, generally available on October 28, 2025 for the z17 mainframe and LinuxONE and in early December 2025 for Power11 servers$^\dagger$. IBM's release allows up to 48 cards in a z17 and 16 in a Power server; the slide's ``8 per drawer'', like its 64 MB and 300 TOPS (precision not stated), could not be confirmed in IBM's own material. It runs language-model inference next to the transactions and the data that regulation keeps on premises. It is not trying to compete for training or for cloud serving; it is the clearest example in this lecture of an accelerator defined by where it sits rather than by how fast it is. (IBM's research chip NorthPole is a different design; slide 30.)

\emph{Tesla} built Dojo to train its driving networks on video. The D1 die (AI Day 2021; Hot Chips 34, 2022) is a 645 mm$^2$ TSMC 7 nm mesh of 354 nodes, each a full superscalar CPU core with four-way multithreading, an 8 $\times$ 8 $\times$ 4 matrix unit and 1.25 MB of SRAM: 354 $\times$ 1.25 $\approx$ 440 MB (derived), with no DRAM on the die at all$^\dagger$. Twenty-five D1 dies are bonded in a 5 $\times$ 5 array onto a carrier to form a 15 kW ``training tile'', with high-bandwidth links at every edge; DRAM sits only in separate interface processors at the edges of a tile. The design is closer to Cerebras's wafer (slide 13) assembled from tested dies than to a GPU. The bandwidth figures are an example of how to read conference slides: AI Day 2021 gave 4 TB/s per die edge and Hot Chips 2022 gave 2 TB/s, almost certainly the same link counted in both directions and in one. Dojo went into service in 2023; in August 2025 Bloomberg reported that the team had been disbanded, and Elon Musk called the next version ``an evolutionary dead end'', saying Tesla's effort would converge on its AI5 and AI6 inference chips$^\dagger$. That was not the end: on January 20, 2026 Musk said Tesla had restarted work on a Dojo3, built on AI5-family silicon and intended for space-based AI compute$^\dagger$. The slide's summary, ``shipped 2023, cut Aug 2025, restarted Jan 2026'', and its subtitle's ``Tesla paused'' reflect that; nothing about a Dojo3 machine had shipped by November 2026.

The four fates, then: Intel's line ended after a product that worked but did not sell; Qualcomm's continues on the strength of power efficiency and a single large customer; IBM's survives by not competing; Tesla's paused because buying GPUs was easier than keeping a unique architecture's software alive, and its restart in 2026 is for a different purpose and a different chip. Slide 35 generalizes the pattern.

\begin{formulas}
Efficiency and on-die memory:
\[
\eta = \frac{P_{\text{peak}}}{\text{power}},\qquad
\text{SRAM}_{\text{die}} = N_{\text{nodes}}\times s_{\text{node}} .
\]
Cloud AI 100 Ultra: $870/150\approx5.8$ TOPS/W (INT8); Spyre: $300/75 = 4$ TOPS/W; Gaudi 2: $432/600\approx0.72$ TFLOP/W BF16; Gaudi 3 at 900 W: $1{,}835/900\approx2.0$ (derived, all vendor peaks). D1: $354\times1.25\,\text{MB}\approx440$ MB.
\end{formulas}

\begin{terms}
\term{Gaudi}{Habana Labs' training accelerator line (Gaudi 2019, Gaudi 2 2022, Gaudi 3 2024), owned by Intel since December 2019.}
\term{MME}{Matrix multiply engine: Gaudi's fixed-function block for dense matrix products.}
\term{TPC}{Tensor processor core: Gaudi's programmable VLIW SIMD core for the operations the MME does not handle.}
\term{RoCE}{RDMA over Converged Ethernet: remote direct memory access carried on standard Ethernet; Gaudi integrates RoCE ports on the die.}
\term{merchant silicon}{A chip sold to any customer, as opposed to one designed for and used only by its maker (captive) or by one buyer (custom).}
\term{NPU}{Neural processing unit: a generic name for an accelerator for neural networks, especially the blocks in phone and PC processors from which Qualcomm's Cloud AI cards descend.}
\term{Spyre}{IBM's PCIe AI accelerator for the z17 mainframe and Power11 servers (2025), companion to the AI unit on the Telum II processor.}
\term{Telum}{IBM's mainframe processor family (2021 on) with an on-chip AI accelerator for in-transaction inference.}
\term{Dojo}{Tesla's in-house training supercomputer built from D1 dies in 25-die training tiles; in service from 2023, cut in August 2025, with a Dojo3 restarted in January 2026.}
\term{D1}{Tesla's Dojo die: 354 CPU-like nodes with matrix units and 440 MB of SRAM in total, no DRAM.}
\term{training tile}{Tesla's module of 25 D1 dies bonded in a 5 $\times$ 5 array with power, cooling and edge links, about 15 kW.}
\end{terms}

\begin{intuition}
Size did not save any of these programs; fit did. Intel had the money and a working chip and still could not make customers switch from CUDA; Qualcomm and IBM picked niches (watts, and the mainframe) where NVIDIA was not trying; Tesla discovered that owning a unique supercomputer means owning its software forever, and set it aside until it found a new use for it. When a large company's accelerator survives, look for the customer it serves that NVIDIA does not.
\end{intuition}
% =====================================================================
\setstage{trn}
\section{China: Ascend, CloudMatrix and the policy backdrop (slides 25--27)}

\slide{25}{Huawei Ascend: Da Vinci cores to the 950/960/970 roadmap}{Cube-MAC cores on 7 nm; cut off from TSMC in 2020, Huawei designs its own HBM and scales by count.}

\para{What the slide shows} The upper left of the slide draws one Da Vinci AI core as Huawei presented it at Hot Chips 31 in August 2019. Data enters through a software-managed L1 buffer, passes into two small L0 buffers that hold the A and B operands of a matrix multiply, and feeds the \emph{cube unit}, drawn as three stacked squares and labeled $16\times16\times16$, 4{,}096 FP16 multiply-accumulates per cycle. Results land in an L0 C buffer of accumulators and from there go to a vector unit; a scalar unit issues instructions. A caption under the core gives the die-level facts: 30 to 32 cores per Ascend 910 die (456 mm$^2$, TSMC 7 nm), 20 to 25 active cores on the SMIC-built 910B, and two 910B dies to make a 910C. Below the core a table runs across six generations, Ascend 910 (August 2019), 910B (2023), 910C (2025), 950PR (first quarter 2026), 950DT (fourth quarter 2026) and 960 (first quarter 2027), with rows for process, FP16 throughput in TFLOP/s, FP8 and FP4 throughput in PFLOP/s, HBM, and interconnect bandwidth per chip in TB/s. A footnote explains the provenance of each column (the 950 and 960 figures from Huawei Connect 2025 and 2026) and adds the Ascend 970 for 2028 (3.6 PFLOP/s FP8, 14 PFLOP/s FP4). The panel on the right, ``From launch to the roadmap'', tells the story in four steps: the Da Vinci core and its software, the cut-off from TSMC, the failed attempt to train DeepSeek's R2 on Ascend, and the roadmap Eric Xu announced in September 2025 and Huawei revised a year later.

\para{Details} The Da Vinci core is the closest thing China has to the TPU's tensor core, and the comparison is worth making precise. A TPU v4 TensorCore holds four $128\times128$ systolic arrays, each a two-dimensional grid of multiply-accumulators through which operands march in lockstep (slide 7). The Da Vinci cube unit is a three-dimensional arrangement: in each cycle it multiplies a $16\times16$ tile of A by a $16\times16$ tile of B and adds the result into a $16\times16$ tile of C, which is $16\times16\times16 = 4{,}096$ multiply-accumulates. The difference is in how operands are delivered. A systolic array passes partial sums between neighbors and reads each operand once per row or column; the cube is a small, wide matrix engine fed directly from the L0 buffers, closer in spirit to a GPU tensor core than to a weight-stationary systolic array. Around the cube sit the two units every tensor core needs: a vector unit for element-wise work (activation functions, normalization, the softmax inside attention) and a scalar unit for control and addressing. The L1 and L0 buffers are \emph{software-managed}: the compiler, not a cache controller, decides what is resident, as in the TPU and unlike the GPU (slide 6).

The arithmetic of a whole chip follows. At launch Huawei quoted 256 TFLOP/s of FP16 and 512 TOPS of INT8 for the 910, at about 310 W$^\dagger$; CSET lists the fastest SKU at about 320 TFLOP/s FP16. With 32 cores, 256 TFLOP/s implies $256\times10^{12}/(32\times4{,}096\times2)\approx 1.0$ GHz (derived), a plausible clock for a 7 nm part of 2019. That put the 910 in the same class as NVIDIA's V100 (125 TFLOP/s of dense FP16 tensor math) and, a year later, the A100 (312). It was a competitive first chip. The software arrived with it: CANN (Compute Architecture for Neural Networks), Huawei's operator library, compiler and runtime, the layer that plays the role CUDA and cuDNN play for NVIDIA; and MindSpore, Huawei's own training framework, open-sourced in March 2020. The 910 made one appearance in MLPerf: 512 chips trained ResNet-50 in the Training v0.7 round of July 2020, in the closed division, plus four of five open-division results (R4 §3.1). Huawei never submitted again (slide 32).

The second paragraph of the panel is where the story turns. Huawei was added to the US Entity List in May 2019, which barred US firms from selling to it without a license; in 2020 the Commerce Department extended the \emph{foreign direct product rule} to Huawei, which forbade any foundry anywhere that uses US tools and software from making chips to Huawei's designs. TSMC, which made the 910, stopped. The 910B, first visible in documents in October 2022 and shipping in volume in 2023, moved to SMIC's N$+$2 process: a 7 nm-class node built with deep-ultraviolet (DUV) lithography and multiple patterning, because the extreme-ultraviolet (EUV) tools TSMC uses at 7 nm and below cannot be exported to China. Multipatterning exposes each critical layer two or more times with offset masks to print features finer than one exposure allows; it works, but every extra exposure adds cost, time and a chance of misalignment, so yields are lower (CSET 2024). The 910B shows the cost: CSET reports 20 to 25 active Da Vinci cores rather than 30 to 32, consistent with disabling defective ones, and up to about 400 TFLOP/s of FP16 at higher clocks, 64 GB of HBM2e at up to 1.6 TB/s$^\dagger$. SemiAnalysis adds an uncomfortable detail, which the deck quotes as theirs: most 910B and 910C dies that shipped in 2024--25 came from TSMC after all, ordered through the front company Sophgo before TSMC cut it off (about \$500 million of wafers, roughly 2.9 million dies), and most of the HBM was a stockpile of Samsung stacks bought before the December 2024 HBM controls (slide 27)$^\dagger$. The 910C (2025) puts two 910B-class dies in one package with 128 GB of HBM; its dense BF16 rate, about 780--800 TFLOP/s, is not published and is inferred from the CloudMatrix total of about 300 PFLOP/s across 384 chips (slide 26). It has no native FP8.

The third paragraph records a test. After DeepSeek's R1 in January 2025, Chinese authorities encouraged the company to train its successor on Ascend. The Financial Times reported in August 2025 that, with Huawei engineers on site, the R2 effort ``never completed a single successful training run'' on Ascend: instability, slow chip-to-chip links and immature software sent training back to NVIDIA's H20, with Ascend kept for inference$^\dagger$. A week later DeepSeek released V3.1 with a note that its FP8 scaling format (UE8M0) was designed ``for the next generation of domestically produced chips''. The lesson the slide draws is not that Ascend cannot train models (Huawei says its own Pangu models are trained on it) but that a frontier training run is a test of the whole stack, interconnect and compiler included, and the chip is the easy part.

The fourth paragraph is the roadmap Eric Xu, Huawei's rotating chairman, announced at Huawei Connect on September 18, 2025, and it is unusually candid. Xu said the mainland would ``lag in process nodes for a relatively long time'', and he described a strategy built around that fact rather than against it: Huawei's own HBM, narrower number formats, and scale-out by chip count. The 950PR (first quarter 2026) is aimed at prefill and recommendation: 1 PFLOP/s of FP8 (and Huawei's own HiF8), 2 PFLOP/s of MXFP4, 2 TB/s of interconnect per chip (2.5 times the 910C; the table's 0.8 TB/s for the 910C is derived from that ratio), and a first in-house HBM, ``HiBL 1.0'', whose capacity and bandwidth were not stated. The 950DT (fourth quarter 2026) is the same die with a second HBM design, ``HiZQ 2.0'', 144 GB at 4 TB/s, aimed at decode and training. The split mirrors the industry's: prefill is compute-bound and decode bandwidth-bound (slide 4), so one die with two memory options serves both. As announced in 2025, the 960 (fourth quarter 2027) doubled to 2 PFLOP/s FP8 and 4 PFLOP/s FP4 with about twice the memory, and the 970 (fourth quarter 2028) to 4 PFLOP/s FP8 and 8 PFLOP/s FP4 with a 4 TB/s link$^\dagger$. At Huawei Connect on September 18, 2026 Huawei revised the plan, and the slide follows the revision where it can: the 960 now comes as a 960DT in the first quarter of 2027 (the table's ``Q1 2027'') and a 960PR in the third quarter of 2027, with a 2.2 TB/s chip link (the table's 2.2; the 4 TB/s that some summaries attached to the 960 was the 970's 2025 figure), and press coverage gives the 960PR 2 PFLOP/s of FP8 and 8 PFLOP/s of FP4, twice the FP4 of the 2025 plan (the table's FP4 cell gives the 960PR's 8); the 970 moves to 2028 at 3.6 PFLOP/s FP8, 14 PFLOP/s FP4 and a 4.4 TB/s link; and a 980 is planned for 2029$^\dagger$. All of these are announcements. In November 2026 the 950PR was reported shipping and the 950DT was due, and neither has published die details, process or power.

Two cautions apply to the whole table. First, every figure after the 910 is either an analyst estimate (the 910B and 910C have no public datasheet) or a keynote number, which is why the table carries a dagger and the ``n.p.'' cells are left as they are. Second, the FP8 and FP4 rows begin where the FP16 row ends: comparing the 910C's 800 TFLOP/s of BF16 with the 950's 1 PFLOP/s of FP8 is comparing different number formats, and on a like-for-like basis the 950 is not a large step in raw arithmetic. The step is in memory and links, which is where the next slide goes.

\begin{terms}
\term{Ascend}{Huawei's line of AI accelerators: the 310 for inference at the edge, the 910 family (910, 910B, 910C) for the data center, and the 950, 960 and 970 announced in 2025 and rescheduled in 2026.}
\term{Da Vinci core}{Huawei's AI core: a cube unit of $16\times16\times16$ FP16 multiply-accumulators per cycle with vector and scalar units and software-managed L1 and L0 buffers; replicated 30 or more times per Ascend 910 die.}
\term{cube unit}{The matrix engine of a Da Vinci core: it multiplies two $16\times16$ tiles and accumulates into a third each cycle, 4{,}096 multiply-accumulates.}
\term{software-managed buffer}{On-chip memory whose contents are placed and evicted by the compiler or program, not by a cache controller; used by the TPU, Ascend and most ASICs in this lecture.}
\term{CANN}{Compute Architecture for Neural Networks: Huawei's operator library, graph compiler and runtime for Ascend, the counterpart of CUDA and its libraries; Huawei announced in August 2025 that it would open the source, a process to be completed by the end of 2025.}
\term{MindSpore}{Huawei's deep-learning framework, open-sourced in 2020; PyTorch reaches Ascend through a separate plug-in (torch\_npu).}
\term{Entity List}{A US Commerce Department list of foreign organizations to which US exports need a license; Huawei was added in May 2019.}
\term{foreign direct product rule}{A US rule that extends export controls to foreign-made items produced with US technology or tools; applied to Huawei in 2020, it ended TSMC's manufacturing for Huawei.}
\term{SMIC}{Semiconductor Manufacturing International Corporation, China's largest foundry; its N$+$2 process is a 7 nm-class node built without EUV lithography.}
\term{DUV multipatterning}{Printing features finer than one deep-ultraviolet exposure allows by exposing a layer two or more times with offset masks; slower, costlier and lower-yielding than EUV.}
\term{EUV lithography}{Extreme-ultraviolet (13.5 nm wavelength) lithography, used by TSMC, Samsung and Intel for 7 nm-class and smaller nodes; the tools cannot be exported to China.}
\term{prefill and decode}{The two phases of LLM inference: prefill processes the prompt in parallel and is compute-bound; decode generates one token at a time and is bandwidth-bound. Huawei's 950PR and 950DT target one each.}
\term{HiF8}{Huawei's own 8-bit floating-point format, supported from the Ascend 950 alongside the standard FP8 and MXFP8 formats.}
\end{terms}

\begin{intuition}
Ascend began as a competent TPU-like chip on the world's best process and then lost the process. Everything since is substitution: SMIC for TSMC with fewer working cores, stockpiled and then home-made HBM for Korean HBM, two dies for one, narrower formats for more transistors. When a roadmap says ``we will lag in process'', expect it to win on count, links and memory rather than on per-chip arithmetic, and check the next slide for the bill.
\end{intuition}

% ---------------------------------------------------------------------
\slide{26}{CloudMatrix 384 and the SuperPoDs: scale, not per-chip speed}{384 Ascend 910C on an all-optical fabric out-total a GB200 NVL72 at about four times the power.}

\para{What the slide shows} The top of the slide draws the CloudMatrix 384, sold as the Atlas 900 A3 SuperPoD: twelve compute racks, each with four drawers, and four narrower switch racks on the right, all joined underneath by an orange line for the optical scale-up fabric. The caption states the arithmetic: 12 compute racks of 32 Ascend 910C make 384 NPUs, the four switch racks carry an all-to-all optical network built on 400G linear-drive transceivers and Huawei's UnifiedBus 1.0 protocol, and the title of the drawing gives the totals, 16 racks and 6{,}912 optical transceivers. Below, a table compares the CloudMatrix 384 with NVIDIA's GB200 NVL72, row by row: accelerators (384 against 72, a ratio of 5.3), dense BF16 (about 300 PFLOP/s against 180, 1.7 times), HBM capacity (49 TB against 13.4, 3.6 times), HBM bandwidth (about 1{,}230 TB/s against 576, 2.1 times), the scale-up fabric (optical all-to-all against copper NVLink 5) and power, all-in (about 559 kW against about 145 kW, 3.9 times; SemiAnalysis's estimates, including each system's share of networking and cooling). A footnote gives the sources and SemiAnalysis's efficiency ratios: per FLOP the CloudMatrix uses 2.3 times the power, per TB/s of HBM 1.8 times, per TB of HBM 1.1 times. The panel on the right reads the comparison in four steps and ends with the SuperPoDs Huawei announced for 2026 and 2027.

\para{Details} The NVL72 (covered in the previous two lectures) is NVIDIA's answer to a specific problem: the largest models are too big for one GPU and their parallel parts exchange data too often for an ordinary network, so NVIDIA built a 72-GPU domain in which every GPU reaches every other at 1.8 TB/s over copper NVLink, in one rack, which draws on the order of 120--130 kW itself and about 145 kW all-in by SemiAnalysis's count. That domain is called the \emph{scale-up} domain, as opposed to \emph{scale-out}, the slower network (InfiniBand or Ethernet) that joins many such domains. The size of the scale-up domain bounds how far tensor and expert parallelism can go before traffic falls to the slow network, and for mixture-of-experts models, whose all-to-all exchanges route every token to the GPUs that hold its experts, a bigger scale-up domain is worth a great deal.

Huawei could not match a B200 chip for chip, so it built a scale-up domain five times larger. The CloudMatrix 384 joins 384 910Cs in one all-to-all fabric. Copper reaches a meter or two at these data rates, enough for one rack; 16 racks need light, so every chip-to-switch link is optical. SemiAnalysis counts 6{,}912 transceivers of 400 Gb/s, of the linear-drive kind that omits the digital signal processor of a conventional optical module to save power and latency$^\dagger$. The transceivers are the expensive and power-hungry part, and they are also the failure point: an optical module fails far more often than a copper cable, so a system with almost 7{,}000 of them must route around failures as a matter of routine.

The table then reads in pairs. On aggregate arithmetic, 5.3 times as many chips give 1.7 times the dense BF16, so one 910C is about $1.7/5.3 = 0.32$ of a B200 (derived): about 780--800 TFLOP/s against 2{,}500. On memory the trade is better: each 910C has 128 GB, so 384 chips have 49 TB against the NVL72's 13.4 TB (186 GB per B200 in the GB200 superchip), 3.6 times, and 2.1 times the bandwidth. For serving a large mixture-of-experts model, whose weights and KV cache are what limit the batch, capacity in one domain is the quantity that matters, and this is the CloudMatrix's real advantage. The bill is power: SemiAnalysis estimates about 559 kW for the whole CloudMatrix against about 145 kW for the NVL72, both all-in, which is $559/145 \approx 3.9$ times (derived), and 2.3 times the power per FLOP, 1.8 times per TB/s of bandwidth and 1.1 times per TB of capacity (SemiAnalysis's ratios). Ratios computed against other NVL72 figures in circulation (a 120 kW rack figure from Dell, for example) mix an all-in estimate with a rack-only one and are not comparable (559/120 would give 4.7); the slide uses SemiAnalysis's matched pair. Sixteen racks also take more floor than one.

Why accept that bill? Because in China the alternative is not an NVL72. With H100s and B200s unavailable and H20 sales on and off (slide 27), a Chinese cloud's choice is between a domestic system that is large, hot and available and an imported one that cannot be bought. Power is cheaper and more plentiful in China than in the US, and the government is building data-center capacity in its western provinces; SemiAnalysis argued in April 2025 that power is not the binding constraint for Huawei, chips are. The comparison on the slide is therefore one of \emph{exchange rate}: Huawei pays in watts, transceivers and floor area for what NVIDIA gets from a better process and NVLink copper. It is NVIDIA's own NVL72 idea (one big scale-up domain) executed with worse parts at a larger scale.

The third paragraph of the panel records the deployment: deliveries from March 2025, and over 300 systems by September 2025 according to Huawei$^\dagger$; at 384 chips each, that is more than 115{,}000 910Cs (derived). UnifiedBus is Huawei's protocol for the scale-up fabric, comparable to NVLink in role: version 1.0 runs in the Atlas 900 A3, and version 2.0, whose specification Huawei published in 2025, claims more than 200 m of optical reach, 100 ns failover from a broken link and 2.1 $\mu$s latency between NPUs in different cabinets$^\dagger$. Those are vendor numbers, and the 2.1 $\mu$s is a cabinet-to-cabinet figure, to be compared with latencies inside a copper NVLink rack rather than with a network hop.

The last paragraph lists the announced successors, and they push the same strategy further. The Atlas 950 SuperPoD (fourth quarter 2026) joins up to 8{,}192 Ascend 950DT in 160 cabinets (128 for compute, 32 for communication) over about 1{,}000 m$^2$, with 8 EFLOP/s of FP8, 16 of FP4, 1{,}152 TB of memory and 16 PB/s of interconnect. The Atlas 960 SuperPoD (fourth quarter 2027, as announced in 2025, before the 2026 revision moved the 960 itself to 2027) has 15{,}488 Ascend 960 in 220 cabinets, 30 EFLOP/s FP8. Above them, SuperClusters of 64 SuperPoDs would reach more than half a million chips in 2026 and a million in 2027, joined by UnifiedBus over Ethernet$^\dagger$. As a check, $8{,}192\times1$ PFLOP/s $= 8.2$ EFLOP/s of FP8 and $8{,}192\times 144$ GB $= 1{,}180$ TB (derived), consistent with the announced totals. None of these had shipped in November 2026; the slide marks them ``announced, not shipped''.

\begin{terms}
\term{CloudMatrix 384}{Huawei's 384-chip Ascend 910C system of 2025 (sold as the Atlas 900 A3 SuperPoD): 12 compute racks and 4 switch racks joined by an all-to-all optical scale-up fabric.}
\term{SuperPoD}{Huawei's name for a scale-up domain of many cabinets: the Atlas 900 A3 (384 chips), Atlas 950 (8{,}192) and Atlas 960 (15{,}488).}
\term{scale-up domain}{The set of accelerators that reach each other over the fast fabric (NVLink, UnifiedBus, the TPU's ICI) rather than over the data-center network; 72 GPUs in an NVL72, 384 chips in a CloudMatrix 384.}
\term{scale-out}{The slower network (InfiniBand, Ethernet) that joins many scale-up domains into a cluster.}
\term{all-to-all fabric}{A network in which every endpoint can reach every other in one hop through switches, at full rate; contrasted with a torus, where traffic passes through intermediate chips.}
\term{linear-drive optics}{Optical transceivers (LPO) that drive the laser directly from the host's electrical signal without a digital signal processor, saving power and latency at the cost of reach and margin.}
\term{UnifiedBus}{Huawei's interconnect protocol for its SuperPoDs (UB 1.0 in the Atlas 900 A3, UB 2.0 published in 2025), the counterpart of NVLink; UBoE carries it over Ethernet.}
\term{GB200 NVL72}{NVIDIA's 2024--25 rack of 36 Grace CPUs and 72 B200 GPUs joined by copper NVLink 5 at 1.8 TB/s per GPU, about 145 kW all-in by SemiAnalysis's estimate.}
\term{mixture of experts}{A model whose feed-forward layers are split into many experts, of which each token uses a few; it needs large memory for all experts and all-to-all communication to route tokens.}
\end{terms}

\begin{intuition}
If you cannot build a faster chip, build a bigger room: the CloudMatrix trades a third of a B200 per chip for five times the chips, light instead of copper and four times the power, and comes out ahead on total memory and arithmetic. The question to ask of any such comparison is which resource is scarce. Where chips are scarce and power is not, the trade is rational; where power is the limit, as in most of the US, it is not.
\end{intuition}

% ---------------------------------------------------------------------
\slide{27}{Kunlun, Hanguang, Cambricon and the policy backdrop}{China's ASIC lines, its GPU-like rivals' IPO wave, and the export-control calendar behind both.}

\para{What the slide shows} The upper left of the slide is a table of seven Chinese ASICs with the year, peak throughput and a note. Baidu's Kunlun 2 (2021, 128 TFLOP/s FP16, 7 nm, about 120 W); the Kunlun P800 (2025, about 345 TFLOP/s FP16, deployed in a 30{,}000-card cluster in April 2025); the Kunlun M100 and M300 (2026--27, not published, for inference and training, with a 512-card node); Alibaba's Hanguang 800 (2019, 820 TOPS INT8, 12 nm, 17 billion transistors, inference); Alibaba T-Head's PPU (2025, ``H20 class'', 96 GB of HBM2e, a 16{,}384-card order); Cambricon's MLU370 (2021, 96 TFLOP/s FP16, the last part before sanctions) and MLU590 (2023, about 345 TFLOP/s FP16, adds FP8, with the next part, the 690, slipped to 2026). A caption names the sources (Caixin, Recode China AI, TrendForce) and explains that the Chinese GPU-like designs, Biren, Moore Threads, MetaX, Hygon and Iluvatar, are left out of the table because they are GPUs. The panel ``What the controls made'' describes three effects of US export controls: demand for domestic chips, a ceiling on what they can do, and an IPO wave. Across the bottom of the slide runs a timeline of eight events from October 2022 to September 2025, alternating above and below the line, ending in bold with Beijing's instruction to Chinese firms to stop buying NVIDIA's RTX Pro 6000D.

\para{Details} Three Chinese companies besides Huawei have designed data-center AI ASICs for long enough to have generations. Baidu's chip unit, now spun out as Kunlunxin, began with the Kunlun 1 (2020, Samsung 14 nm) and the Kunlun 2 (2021, 7 nm, 256 TOPS INT8 and 128 TFLOP/s FP16 at about 120 W). The third generation, the P800, is the one that matters: about 345 TFLOP/s of FP16 according to Guosen Securities, a 10{,}000-card cluster in February 2025 and a 30{,}000-card cluster by April 2025 that runs most of Baidu's own inference$^\dagger$. In November 2025 Baidu published a roadmap: the M100 for inference in 2026, the M300 for training and inference in 2027, a 512-card supernode in 2026 and a million-card cluster by 2030; Kunlunxin filed for a Hong Kong listing in January 2026. Alibaba's chip unit, T-Head, built the Hanguang 800 in 2019 as an inference chip for its own image and search workloads (820 TOPS INT8 as reported, 12 nm, about 17 billion transistors$^\dagger$) and in 2025 the PPU, a general training-and-inference part on a domestic 7 nm-class process with 96 GB of HBM2e and 700 GB/s of chip-to-chip links, pitched against NVIDIA's H20. China Unicom's Sanjiangyuan data center ordered 16{,}384 PPU cards in 1{,}024 servers, the largest domestic order at that site$^\dagger$. Cambricon, a 2016 spin-out of the Chinese Academy of Sciences that once licensed the NPU in Huawei's Kirin phone chips, makes the MLU (machine learning unit) line: the MLU370 of 2021 (96 TFLOP/s FP16) was its last part before the company was added to the Entity List in December 2022; the MLU590 of 2023, on a domestic 7 nm process, reaches about 345 TFLOP/s FP16 and adds FP8; the Siyuan 690 slipped from the end of 2025 to the second half of 2026$^\dagger$.

The FP16 column is the one to compare, and it shows the ceiling the panel names. Domestic chips in 2025 delivered roughly 100 to 400 TFLOP/s of dense FP16 against 1{,}000 to 2{,}500 for NVIDIA's current parts (H100 989, B200 2{,}250 dense)$^\dagger$. The reasons are the ones slide 25 gave for Ascend: SMIC's 7 nm-class process by DUV multipatterning with low yields and limited capacity (SemiAnalysis estimated about 50{,}000 advanced wafers a month in 2025), and HBM from stockpiles rather than current production. Hence the conclusion the panel draws and slide 26 illustrates: China scales by count.

The panel's first paragraph describes what the controls did to demand, and the effect is the opposite of the one the controls intended for the domestic industry. Each rule that removed an NVIDIA part from the Chinese market moved buyers to domestic parts. Cambricon is the clearest case: a loss of about RMB 450 million in 2024 became a net profit of RMB 2.06 billion in 2025 on revenue of RMB 6.5 billion, up 453\% (TrendForce, March 2026)$^\dagger$. Baidu's 30{,}000 P800s and Alibaba's 16{,}384-card PPU order are the same story at the cloud providers.

The third paragraph is the IPO wave of December 2025 and January 2026, when China's GPU-like designers went public within five weeks. Moore Threads listed on Shanghai's STAR Market on December 5, 2025, at RMB 114.28 a share and closed its first day up about 425\% (the often-quoted 468\% was the intraday peak); MetaX followed on December 17 (up about 643\% on the day); Biren listed in Hong Kong on January 2, 2026, raising HK\$5.58 billion; Iluvatar CoreX on January 8; Kunlunxin filed in January and T-Head's listing plan was cleared$^\dagger$. The deck leaves these GPU-like companies out of the table because this lecture is about designs that are not GPUs, but their listings belong on the slide because they measure the same thing as Cambricon's profit: the value the market places on being the domestic alternative.

The timeline at the bottom gives the dates. On October 7, 2022 the Bureau of Industry and Security (BIS) imposed the first controls on advanced computing chips to China, covering the A100 and H100 class. NVIDIA designed reduced versions, the A800 and H800, with slower chip-to-chip links to fall under the thresholds; on October 17, 2023 BIS replaced the bandwidth test with a test on \emph{total processing performance} (TPP) and performance density, which closed that gap. On December 2, 2024 BIS added controls on HBM itself and about 140 Entity List additions. On January 13, 2025 the outgoing administration published the AI Diffusion rule, which sorted the world into three tiers with country caps on accelerator imports. In April 2025 the US told NVIDIA that its China-specific H20 would need a license; NVIDIA took a charge of up to \$5.5 billion in its quarter. On May 13, 2025 BIS rescinded the Diffusion rule and issued guidance that using Huawei's Ascend 910B, 910C or 910D ``anywhere in the world'' could violate US controls. In July and August H20 and AMD MI308 sales were allowed to resume, reportedly on condition that the companies pay 15\% of the revenue to the US government. On September 17, 2025, according to the Financial Times, the Cyberspace Administration of China told ByteDance, Alibaba and others to stop buying NVIDIA's RTX Pro 6000D$^\dagger$. The last event, set in bold on the slide, is the reversal: by late 2025 the restriction was coming from Beijing as well as Washington.

Some of these dates come from law-firm summaries and press reports rather than the Federal Register, and the research behind the deck marks the 2022 and 2023 days as recalled rather than re-fetched; a student citing them should check the BIS rule texts. The direction of the story is not in doubt.

\begin{terms}
\term{Kunlun}{Baidu's AI chip line (Kunlun 1, 2, P800, M100, M300), now made by its spun-out unit Kunlunxin.}
\term{Hanguang 800}{Alibaba's 2019 inference chip from its T-Head unit: 12 nm, 820 TOPS INT8.}
\term{PPU}{Alibaba T-Head's 2025 training-and-inference accelerator, 96 GB HBM2e, pitched against NVIDIA's H20.}
\term{Cambricon MLU}{Cambricon's machine-learning-unit accelerators (MLU290, MLU370, MLU590, Siyuan 690); Cambricon was added to the Entity List in December 2022.}
\term{BIS}{The Bureau of Industry and Security of the US Department of Commerce, which writes and enforces export controls.}
\term{total processing performance}{TPP: the BIS metric for export controls on chips, roughly the peak dense operations per second times the bit width of the operands; with performance density (TPP per mm$^2$) it sets the thresholds from October 2023.}
\term{AI Diffusion rule}{A BIS rule of January 2025 that capped accelerator exports by country in three tiers; rescinded in May 2025.}
\term{H20}{NVIDIA's China-specific Hopper part, designed under the 2023 thresholds with full HBM but reduced compute; its sales were stopped and resumed in 2025.}
\term{STAR Market}{The Shanghai Stock Exchange's Science and Technology Innovation Board, where Moore Threads and MetaX listed in December 2025.}
\end{terms}

\begin{intuition}
Export controls are a tax on one side's chips and a subsidy to the other's. Each rule lowered the ceiling on what Chinese buyers could get from NVIDIA and raised the floor under what they would pay a domestic supplier, and the profits and IPOs of 2025--26 are the subsidy showing up in the accounts. The chips themselves remain a few times slower each; whether that gap closes is a question about SMIC and HBM, not about design.
\end{intuition}
% =====================================================================
\setstage{outc}
\section{Emerging designs: in-memory, hardwired, photonic, neuromorphic (slides 28--30)}

\slide{28}{Compute in the memory: d-Matrix, Untether, Mythic, EnCharge}{Put the multiply in the SRAM array and the bytes stop moving; capacity and precision are the price.}

\para{What the slide shows} The upper left of the slide draws ``three places to put the multiply-accumulate'' as three small panels. The first, \emph{conventional}, has a box for SRAM or HBM holding the weights, an arrow labeled ``bus, registers'' and a MAC array below it; its caption says that every weight travels across the chip, or from DRAM, on every use. The second, \emph{digital in-memory}, is a grid of twelve SRAM banks, each with a dark strip for the multiply-accumulate logic built into it; the caption says the weights stay put and the math is exact integer arithmetic, as in d-Matrix, Untether and IBM's NorthPole. The third, \emph{analog in-memory}, is a crossbar of horizontal and vertical wires with a dot at each crossing; the caption says Kirchhoff's current law sums the dot product, and that drift and the converters at the edges bound the precision. Below the three panels a log-scale bar chart gives the energy to move one bit, from Dally, Turakhia and Han (CACM 2020): 0.05 pJ from a local 8 KB SRAM, 0.7 pJ across 100 MB of on-chip memory, 4 pJ from LPDDR4, 10 pJ over a SerDes link and 20 pJ from DDR4. The panel on the right, ``The parts, and where they stand'', describes d-Matrix's Corsair, Untether's speedAI240, Mythic's analog chips and EnCharge's EN100, and ends with the trade between the two kinds of in-memory compute.

\para{Details} Slide 8 made the energy argument with Horowitz's 45 nm table: a DRAM read costs hundreds of times a multiply. The chart on this slide is Dally's update for 14 to 28 nm, expressed per bit moved rather than per operation, and the ratios are larger: a bit read from the 8 KB SRAM next to the arithmetic costs 0.05 pJ; the same bit carried across a chip with 100 MB of SRAM costs 0.7 pJ, fourteen times more, because the wires are long; from LPDDR4 it costs 4 pJ, from DDR4 20 pJ, 400 times the local read. Dally's paper puts an 8-bit integer add at about 0.01 pJ. So on a modern chip the multiply is nearly free and the cost is distance. The SRAM machines of section 5 (Groq, Cerebras, Graphcore) attack the far end of the chart by removing DRAM; in-memory compute attacks the near end by removing the wire between the SRAM and the multiplier.

There are two ways to do it, and the slide separates them. \emph{Digital in-memory compute} (also called at-memory or near-memory compute) puts ordinary digital multiply-accumulate logic inside or right beside each SRAM bank, so a weight is read a few micrometers from where it is used and never crosses the chip. The arithmetic is exact integer or block-floating-point math, so the results match a conventional processor's bit for bit. The cost is density: an SRAM bank with logic beside it stores fewer bits per mm$^2$ than a plain SRAM, and SRAM itself stores far fewer than DRAM, so a chip holds at most a few gigabytes. \emph{Analog in-memory compute} goes further: it stores each weight as a conductance (the resistance of a flash cell, a resistive memory cell or a capacitor's charge) at the crossing of two wires, applies each input as a voltage on a row, and lets Ohm's law do the multiplies and Kirchhoff's current law do the sums, so the current flowing out of each column is a dot product. A $1{,}000\times1{,}000$ crossbar computes a million multiply-accumulates in one step with almost no data movement. The cost is precision: the inputs must be converted from digital to analog (DACs) and the outputs back (ADCs), the converters eat much of the energy saved, the cells drift with temperature and age, and effective precision rarely exceeds 8 bits.

d-Matrix's Corsair, launched in November 2024, is the most developed digital design. As d-Matrix described it at Hot Chips 2025, each card carries two chips of four chiplets each, eight TSMC 6 nm chiplets in all, with 2 GB of SRAM across the card, in which some SRAM arrays are repurposed as matrix-vector multipliers; the card's ``performance memory'' delivers 150 TB/s, and 256 GB of LPDDR5X behind it is ``capacity memory'' for weights that do not fit, at 275 to 550 W per card$^\dagger$. d-Matrix quotes ``9.6 POPS'' (9{,}600 trillion operations per second) at MXINT4, a 4-bit block floating-point format; as a peak in an unusual format it is quoted in quotation marks and not compared with dense FP8 or BF16 figures elsewhere, and it does not appear in the Hot Chips 2025 coverage. A SquadRack of eight servers of eight cards, 64 cards with $64\times2 = 128$ GB of SRAM (derived), would hold a 100-billion-parameter model at 8 bits entirely in SRAM, which is d-Matrix's claim rather than a published deployment$^\dagger$. The next generation, Raptor (2026), is to stack DRAM on the compute in three dimensions to raise capacity; the Hot Chips 2025 coverage names the 3D-DRAM test vehicle Pavehawk, and Raptor's details are d-Matrix's announcement$^\dagger$. d-Matrix raised \$275 million at a \$2 billion valuation in November 2025$^\dagger$.

Untether AI's speedAI240 (``Boqueria'') also computed inside its SRAM arrays, with 1{,}458 RISC-V cores and LPDDR5 behind them; it added FP8 and BF16 to its predecessor's INT8. Untether was one of the few companies to submit MLPerf \emph{power} results, in Inference v4.1 (August 2024), on ResNet-50 only (slide 32). In June 2025 it stopped supplying and supporting its products and AMD hired its engineering team$^\dagger$. That outcome, a team acquired and a product line ended, is the commonest exit in this lecture (slide 35).

Mythic is the analog case. Its M1076 stored weights as the conductances of flash cells on a 40 nm process and was sold for vision at the edge, quoted at about 25 TOPS in about 3 W$^\dagger$ (vendor figures the research could not re-fetch). After raising more than \$170 million it ran out of money in November 2022; in March 2023 it was revived with \$13 million and a defense focus, and in December 2025 it raised an oversubscribed \$125 million, the figure the slide gives$^\dagger$. Analog in-memory compute is therefore not dead as a business, though Mythic's market remains the edge rather than the data center. EnCharge AI's EN100 (2025) computes in the charge domain, summing charges on capacitors rather than currents, which is less sensitive to drift; it is aimed at laptops and other client devices, and the slide notes there is no independent data on it. Qualcomm's AI250 (2027) is announced as ``near-memory computing'' with no architectural detail (slide 24).

The last paragraph of the panel is the summary. Digital in-memory compute keeps exact arithmetic and caps the model at gigabytes per chip, so a large model needs many chips or a DRAM tier behind them, which brings back some of the data movement it set out to remove. Analog in-memory compute saves more energy and loses precision, which suited 8-bit vision networks at the edge and has not yet suited large language models.

\begin{terms}
\term{in-memory compute}{Doing multiply-accumulate operations inside or beside the memory arrays that hold the weights, so the weights do not travel to a separate arithmetic unit.}
\term{digital in-memory compute}{In-memory compute with ordinary digital logic placed in or next to SRAM banks; exact integer or block-floating-point results, limited capacity.}
\term{analog in-memory compute}{Storing weights as conductances or charges at the crossings of a wire grid and computing dot products by summing currents or charges; very low data movement, limited precision.}
\term{crossbar}{A grid of row and column wires with a memory cell at each crossing; with inputs as row voltages, each column current is a dot product.}
\term{Kirchhoff's current law}{The currents flowing into a node sum to the currents flowing out; in a crossbar it adds the cell currents of a column, performing the sum of a dot product.}
\term{DAC and ADC}{Digital-to-analog and analog-to-digital converters; in analog in-memory compute they convert inputs in and results out and consume much of the energy budget.}
\term{block floating point}{A number format in which a block of values shares one exponent and each value keeps a small mantissa (MXINT8, MXINT4, MXFP4); cheaper than full floating point at similar accuracy.}
\term{SerDes}{Serializer-deserializer: the circuit that sends parallel data over a fast serial link between chips; about 10 pJ per bit (Dally 2020).}
\term{Corsair}{d-Matrix's 2024 inference card: digital in-memory compute in eight 6 nm SRAM chiplets on two chips (2 GB per card, 150 TB/s) with 256 GB of LPDDR5X behind.}
\term{speedAI240}{Untether AI's at-memory inference chip with 1{,}458 RISC-V cores; discontinued in 2025.}
\end{terms}

\begin{intuition}
On a modern chip a multiply costs less than moving its operand a millimeter, so the frontier of efficiency is about geometry: put the arithmetic where the weights already are. Digital designs do this exactly and run out of room; analog designs do it with physics and run out of precision. A design that claims both capacity and exactness should be asked where its weights live when the model does not fit.
\end{intuition}

% ---------------------------------------------------------------------
\slide{29}{The specialization spectrum: from GPU to hardwired model}{Each step right trades freedom for speed; Taalas etched a model into metal and was bought in six months.}

\para{What the slide shows} An arrow runs across the top left of the slide under the label ``what is fixed at design time, left to right'', with five boxes along it. From the left: the \emph{GPU} (any kernel; NVIDIA, AMD), the \emph{TPU and Trainium} (tensor operations and a compiler; Google, AWS), the \emph{LPU and dataflow} machines (a static schedule and SRAM; Groq, Cerebras), the \emph{transformer ASIC} (attention and MLP datapaths; Etched, MatX) and the \emph{hardwired model} (one network in metal; Taalas). Below the arrow a table lists five products with what each is, the claim made for it, and its status in November 2026: Etched's Sohu, MatX's MatX One, Positron's Atlas and Asimov, and Taalas's HC1. A caption says the figures are the vendors' own and for one user, except Etched's, which is a server's aggregate throughput; that none of these companies has submitted to MLPerf (slide 32); and that the Asimov figures are simulation. The panel on the right, ``Claims against shipments'', sets the money raised against what has shipped, explains what a transformer ASIC removes, describes the hardwired endpoint, and lists the acquisitions of 2025--26.

\para{Details} The arrow is the lecture's taxonomy (slide 5) turned into a single axis. A GPU fixes almost nothing: its cores run any kernel a programmer writes, and a new model needs new software, not new silicon. A TPU or Trainium fixes the kind of operation, dense tensor algebra in systolic arrays, and leaves the model to the compiler. A Groq LPU or a Cerebras wafer fixes more: the compiler must schedule every operation in advance and the weights must fit in SRAM, so models that do not suit a static schedule or exceed the SRAM lose much of the advantage. A transformer ASIC fixes the model family: its datapaths are shaped for attention and the feed-forward (MLP) blocks of a transformer, and other networks may not run at all. A hardwired model fixes the model itself, weights included. Each step right removes hardware the earlier designs needed for generality (schedulers, caches, general matrix units, weight memory) and spends the area and power on the one thing that remains. Each step also shortens the useful life of the chip, because each fixes something that the field may change.

Etched's Sohu is the best-funded transformer ASIC. Announced in June 2024 with a \$120 million round, it was claimed to run Llama 70B at more than 500{,}000 tokens per second in an eight-chip server, which Etched said replaced 160 H100s; it is built on TSMC N4P with 144 GB of HBM3E per chip$^\dagger$. Etched raised \$500 million at a \$5 billion valuation in December 2025, then \$300 million at \$10.3 billion in late July 2026 and \$700 million at \$21 billion, led by Jane Street, on August 18, 2026: the slide's ``\$1B at \$10--21B (Jul--Aug 2026)''$^\dagger$. It claims more than \$1 billion of orders. It announced first-pass (A0) silicon in June 2026. In November 2026 no independent benchmark of Sohu had been published, and the 500{,}000 tokens per second, from Etched's June 2024 launch post, is the aggregate throughput of an eight-chip server across many users, not the speed one user sees; it is the one entry in the table that is not a per-user figure. MatX, founded in 2022 by former Google engineers, describes the MatX One as a split systolic array with weights in SRAM and the KV cache in HBM, and claims more than 2{,}000 tokens per second on a large mixture-of-experts model; it raised \$500 million in February 2026 and had shown no silicon$^\dagger$.

Positron takes a different route to the same market. Its shipping product, Atlas, is built on Altera Agilex 7 FPGAs with HBM and DDR5 and is sold on price per token rather than peak speed; more than 50 racks were at Oracle Cloud by September 2026. Its ASIC, Asimov, on TSMC N3P, is designed around capacity: up to 2.3 TB of LPDDR5X per chip rather than HBM, at about 3 TB/s, with tape-out at the end of 2026 and volume in 2027. Positron raised \$875 million at a \$5 billion valuation in September 2026, before Asimov taped out; all of Asimov's figures are simulation$^\dagger$. That an FPGA product is the only one in the table that ships in volume is itself a lesson about how long silicon takes.

Taalas is the endpoint. Its HC1, shown in February 2026, has the weights of one model, Llama 3.1 8B, built into the metal layers of the chip: 6 nm, 815 mm$^2$, 53 billion transistors, with weights in what Taalas calls a custom 3-bit base type that combines 3-bit and 6-bit parameters, so ``3-bit'' alone understates it; the slide gives 3/6-bit, precision 3--6 bits$^\dagger$. Because the weights are wiring, not data, there is no weight memory to read and no bandwidth limit on decode, and Taalas demonstrated about 17{,}000 tokens per second for a single user$^\dagger$. The price is total inflexibility: a new model is a new mask set, at least the metal layers of a new chip, and the precision is fixed at manufacture. An 8-billion-parameter model is also small; whether the method scales to models of hundreds of billions is unshown. In early August 2026, six months after the HC1 demonstration, AMD announced an agreement to acquire Taalas, subject to customary closing conditions; the slide's ``AMD buys'' and ``bought'' should be read as ``agreed to buy'', since closing had not been confirmed in the sources consulted$^\dagger$.

The panel's first paragraph puts the money in order: \$1 billion in July and August 2026 alone for Etched, at valuations rising from \$10.3 to \$21 billion, \$500 million for MatX, \$875 million for Positron, all before any of their ASICs had an independent measurement, and only Positron's FPGA product shipping. The second paragraph names the bet a transformer ASIC makes: that the architecture stops changing. Transformers have been stable since 2017 in outline, but the details that a fixed datapath depends on have not: mixture-of-experts routing, grouped and multi-latent attention, state-space and linear-attention hybrids, and longer contexts each change the shapes and the ratio of compute to memory traffic. The last paragraph lists the exits that cluster at this end of the spectrum: NVIDIA's license of Groq's technology (slide 16), AMD's hiring of Untether's team (slide 28) and agreed purchase of Taalas, and Marvell's purchase of Celestial AI (slide 30). The large companies are buying the ideas and the people at the specialized end rather than competing with them.

\begin{terms}
\term{specialization spectrum}{The ordering of accelerator designs by how much is fixed at design time, from the GPU (nothing but the instruction set) to the hardwired model (the weights).}
\term{transformer ASIC}{A chip whose datapaths are built specifically for the attention and feed-forward blocks of transformer models, giving up general matrix flexibility; Etched's Sohu is the example.}
\term{hardwired model}{A chip in which one model's weights are fixed in the wiring (the metal layers) rather than stored in memory; Taalas's HC1.}
\term{metal layers}{The upper layers of a chip that carry the wiring between transistors; changing only them is cheaper than a full new design, and Taalas encodes weights in them.}
\term{A0 silicon}{The first fabricated version of a chip design, used for bring-up and debugging before revisions (A1, B0) go to production.}
\term{FPGA}{Field-programmable gate array: a chip whose logic and wiring are configured after manufacture; slower and less dense than an ASIC but available without a tape-out.}
\term{tape-out}{The point at which a chip design is sent to the foundry for its masks to be made; volume production typically follows a year or more later.}
\term{tokens per second per user}{The speed one user sees in decode; distinct from system throughput, which sums over all users in a batch.}
\end{terms}

\begin{intuition}
Moving right on the spectrum is borrowing against the future: you get speed now in return for betting that what you fixed will not change. The GPU never makes the bet and never wins by much; a hardwired model makes the biggest bet and is obsolete with the next model release. Read a claim from the right end by asking what it fixes, how long that thing has been stable, and whether anyone outside the company has measured it.
\end{intuition}

% ---------------------------------------------------------------------
\slide{30}{Light and spikes: photonics and neuromorphic computing}{Photonic compute was deferred and money went to interconnect; neuromorphic machines stay research tools.}

\para{What the slide shows} The upper left of the slide draws a photonic interposer of the kind Lightmatter's Passage M1000 and Celestial AI's Photonic Fabric provide: a pale blue base labeled ``waveguides, modulators, detectors'' carrying four boxes marked ``XPU die, HBM beside it'', joined to each other by light paths and down to a common bus that leaves on the right ``to fibers''; a note gives the M1000's figures, 256 fibers and 114 Tb/s. Below, a table lists six projects with what each is and its status in November 2026: Lightmatter's Envise (photonic matrix-vector multiply, deferred to 2030 or later), Lightmatter's Passage (the M1000 interposer, 114 Tb/s over 4{,}000 mm$^2$, and the L200 co-packaged optics, up to 64 Tb/s per chip; near-package optics in 2026, co-packaged optics in 2027), Celestial AI (the Photonic Fabric chiplet, 16 Tb/s, stackable on hot dies; Marvell bought it in December 2025 for \$3.25 billion and more), Ayar Labs (TeraPHY optical I/O chiplets; valued above \$1 billion in December 2024), Intel's Hala Point (1{,}152 Loihi 2 chips, 1.15 billion neurons, 20 POPS, 2.6 kW; a Sandia research system since April 2024) and IBM's NorthPole (12 nm, 22 billion transistors, 256 cores, 224 MB, no DRAM; a prototype that runs a 3-billion-parameter LLM on 16 cards). The panel on the right explains why photonic compute and neuromorphic machines stay niche, why light won in interconnect, and why NorthPole is the near-memory limit case.

\para{Details} Light was proposed for computing long before deep learning, and neural networks seemed to offer it the perfect workload: a matrix-vector multiply can be done by passing light through a mesh of interferometers whose phase settings encode the matrix, at the speed of light and with almost no energy spent in the multiply itself. Lightmatter's Envise was the commercial attempt, with claims of 262 TOPS/W and 200 ps per multiply$^\dagger$. It ran into three problems the panel names. Precision: optical multiplication is analog, with the same noise and drift problems as the analog crossbars of slide 28, and practical precision is a few bits. Conversion energy: every input must be converted from electrical to optical and every output back, and the converters cost more than the multiply saves. Memory: there is no optical memory, so weights and activations live in electronic SRAM and DRAM and must be converted on every use. Lightmatter moved Envise to a 2030-or-later roadmap and turned the company to interconnect$^\dagger$.

Interconnect is where light wins, for the reason in Dally's chart on slide 28: moving a bit over an electrical SerDes link costs about 10 pJ, and the cost and the reach both get worse as data rates rise, so a copper link at today's rates reaches a meter or two. Light in fiber goes hundreds of meters at little extra cost per meter. The industry is therefore moving the electrical-to-optical conversion ever closer to the chip: from pluggable transceivers at the front of a switch (the 6{,}912 in Huawei's CloudMatrix, slide 26), to near-package optics, to co-packaged optics in the same package as the switch or accelerator, to optical interposers under the dies. Lightmatter's Passage M1000 is the last: a 4{,}000 mm$^2$ photonic interposer with 256 fibers and 114 Tb/s, on which accelerator dies sit; its L200 is a co-packaged optical engine of 32 to 34 Tb/s per chip, and the L200X of 64 Tb/s; Lightmatter's ``more than 200 Tb/s'' is the total I/O of a multi-chip XPU package built with several of them, not one chip's rate$^\dagger$. Celestial AI's Photonic Fabric puts an optical chiplet (16 Tb/s in its first version) beside or under hot compute dies. Marvell agreed on December 2, 2025 to buy Celestial for \$3.25 billion up front (\$1.0 billion in cash and about 27.2 million shares) plus up to \$2.25 billion more if revenue targets are met by the end of fiscal 2029, which is the ``up to \$5.5 billion'' in the panel (Marvell's 8-K filing). Ayar Labs, maker of TeraPHY optical I/O chiplets and SuperNova lasers, raised \$155 million at a valuation above \$1 billion in December 2024 with NVIDIA, AMD and Intel as strategic investors$^\dagger$. NVIDIA's own co-packaged-optics products are its Spectrum-X and Quantum-X Photonics switches, shown in 2025 and in production in 2026, which is what the panel's ``NVIDIA's 2026 switches use co-packaged optics'' refers to; they are scale-out switches, and at GTC 2026 NVIDIA said that copper still has a role in the scale-up fabric of its Feynman generation, with co-packaged optics reserved for NVLink domains larger than about 1{,}152 GPUs$^\dagger$. The common thread is that every scale-up domain larger than one rack (slide 26) needs light, and the money followed.

Neuromorphic computing takes its inspiration from the brain: neurons that communicate by discrete spikes, compute only when a spike arrives, and keep their state locally. Intel's Loihi 2 (2021, Intel 4 process) implements spiking neurons and programmable synapses in silicon; Hala Point, delivered to Sandia National Laboratories in April 2024, joins 1{,}152 of them into a 6U system of 140{,}544 neuromorphic cores, 1.15 billion neurons and 128 billion synapses that draws up to about 2.6 kW and is quoted at 20 POPS and more than 15 TOPS/W at 8 bits$^\dagger$. The energy of such a machine scales with the number of spikes, not with the number of multiply-adds, so it is efficient when activity is sparse: sensor processing, optimization problems, some scientific simulation. Dense LLM inference is the opposite case, with every weight used for every token, and the slide's conclusion is that no spiking machine had shown competitive dense-LLM throughput by November 2026. SpiNNaker2 (TU Dresden), with 152 Arm cores per chip and a machine of about five million cores, is the other large neuromorphic system and is likewise a research instrument.

IBM's NorthPole, published in \emph{Science} in October 2023 (Modha et al.), belongs on the slide because of its lineage (it descends from IBM's TrueNorth neuromorphic chip) but it is not spiking: it is the purest example of near-memory compute. A 12 nm die of 22 billion transistors carries 256 cores, each doing 2{,}048 8-bit operations per cycle, and 224 MB of on-chip memory, with no off-chip DRAM at all; to the host it ``appears as an active memory''. A model must fit in 224 MB per chip, so IBM's 2024 demonstration split a 3-billion-parameter LLM across 16 cards in a 2U server, with INT4 weights, and measured 28{,}356 tokens per second across the system, under 1 ms per token per user, and 72.7 times the energy efficiency (tokens per joule) of ``the next lowest-latency GPU'' in a low-latency setting. These are IBM's own measurements; IBM's blog does not name the GPU, and the slide's ``an H100's'' is the deck's reading of it, not IBM's wording. The comparison is at low batch, where GPUs are least efficient (slide 4). Still, the system makes the argument of this lecture in its most extreme form: with no DRAM, decode is never waiting on memory bandwidth, and the price is that model size is set by chip count.

\begin{terms}
\term{photonic compute}{Performing matrix-vector multiplies with light, typically in a mesh of interferometers whose phases encode the matrix; analog, with limited precision.}
\term{optical interconnect}{Carrying data between chips or racks as light in waveguides or fiber; lower energy per bit and much longer reach than copper at high data rates.}
\term{co-packaged optics}{CPO: optical transceivers placed in the same package as a switch or accelerator chip, so the electrical path to the optics is millimeters long.}
\term{near-package optics}{NPO: optical engines placed on the board next to the package rather than inside it; an intermediate step toward co-packaged optics.}
\term{photonic interposer}{A silicon base carrying optical waveguides, modulators and detectors on which electronic dies are mounted, so dies communicate by light; Lightmatter's Passage M1000.}
\term{Photonic Fabric}{Celestial AI's optical interconnect chiplet technology, bought by Marvell in 2025--26.}
\term{neuromorphic computing}{Hardware modeled on biological neurons that communicate by spikes and compute only when spikes arrive; Intel's Loihi 2 and SpiNNaker2 are examples.}
\term{spiking neural network}{A neural network whose units emit discrete spikes over time rather than continuous activations; energy scales with spike count.}
\term{Hala Point}{Intel's 2024 neuromorphic system at Sandia: 1{,}152 Loihi 2 chips, 1.15 billion neurons, up to about 2.6 kW.}
\term{NorthPole}{IBM's 2023 inference chip: 12 nm, 256 cores, 224 MB of on-chip memory and no off-chip DRAM; the limit case of near-memory compute.}
\end{terms}

\begin{intuition}
Light is excellent at carrying bits and poor at computing with them, so the photonics money went where light is excellent. Spikes are excellent when little is happening and poor when everything is, which is what a dense LLM does every token. Both lessons are the same as the rest of the lecture: a new physical mechanism wins where it matches the dominant cost, and in 2026 the dominant cost is moving data.
\end{intuition}
% =====================================================================
\setstage{vc}
\section{Measured: specifications, MLPerf, speed and deals (slides 31--34)}

\slide{31}{The chips side by side: what each vendor published}{Dense peaks, memory and fabric for eleven non-GPU parts against one GPU; the gaps (n.p.) are part of the picture.}

\para{What the slide shows} The slide is one table, the widest in the deck, with thirteen columns: a label column and twelve chips. The first chip column is the reference GPU, NVIDIA's B200 in its HGX module; the other eleven are the non-GPU parts the lecture has discussed, ordered roughly from the hyperscalers' captive chips to the merchant startups and then China: Google's TPU v7 (Ironwood), AWS Trainium2, Microsoft Maia 200, Meta MTIA v2, the Cerebras WSE-3, the SambaNova SN40L, the Groq LPU (the GroqChip of 2020), the Graphcore Bow IPU, the Tenstorrent Blackhole, the FuriosaAI RNGD and the Huawei Ascend 910C. Sixteen rows follow: vendor, year first shipped, process, how the chip is organized (its count of cores, tiles or processing elements), four rows of dense peak throughput at BF16, FP8, FP4 and INT8, on-chip SRAM, the kind of off-chip memory, its capacity and its bandwidth, the per-chip scale-up bandwidth, the size of the largest coherent domain that bandwidth reaches, the power, and the software stack. A long footnote under the table records how each awkward entry was obtained. Dashes mean the chip does not support the format or has no such memory; ``n.p.'' means the vendor has not published the number as of November 2026.

\para{Details} The table is the lecture's attempt to put eleven very different machines on one page, and its first lesson is how hard that is. Every number in it is a vendor figure from a datasheet, a product page, a Hot Chips talk or a paper; none is a measurement, which is why the header carries a dagger. The second lesson is in the gaps: a quarter of the cells are dashes or n.p., and the pattern of what is missing says as much as what is there.

\emph{First shipped and process.} The dates run from 2020 (the GroqChip, a 2019 design on GlobalFoundries 14 nm, the oldest process in the table) to 2026 (Maia 200 and the RNGD). Google and AWS do not state the process of Ironwood or Trainium2; analysts assume TSMC's 5 nm or 3 nm class, but the lecture does not print a guess. Maia 200 is the only part on TSMC N3P, the newest node in the table, and the 910C is on SMIC's N+2, a 7 nm-class process made without extreme-ultraviolet lithography (slide 25). The Bow is the odd one: a TSMC N7 compute die bonded wafer-on-wafer (WoW) to a second wafer that carries deep-trench capacitors for power delivery, the first commercial use of that TSMC technique.

\emph{Organization.} The ``organization'' row is the closest the table comes to comparing architectures, and it shows that the word ``core'' means nothing across vendors. The B200 has 148 streaming multiprocessors; Ironwood has two TensorCores per chip (two dies, one TensorCore each, each with several matrix units); Trainium2 has eight NeuronCores; Maia 200 is listed at about 96 cores, a figure Microsoft has not published; MTIA v2 has 64 processing elements in an 8$\times$8 grid; the WSE-3 has about 900{,}000 cores, each a small in-order engine with 48 KB of SRAM; the SN40L has 1{,}040 pattern compute units (PCUs) paired with as many memory units; the GroqChip is one core, a single very wide VLIW machine; the Bow has 1{,}472 independent tiles; the Blackhole has 140 Tensix cores (120 enabled on shipping cards); the RNGD has eight processing elements; the 910C is two dies, each a 910B. A ``core'' ranges from a few square millimeters of a wafer to half a chip.

\emph{Peaks.} The four throughput rows are all \emph{dense}: no 2:4 structured sparsity is counted, so NVIDIA's B200 appears at 2{,}250 TFLOP/s BF16, 4{,}500 FP8 and 9{,}000 FP4 rather than at the doubled headline figures. The INT8 entry for the B200 (4{,}500 TOPS) is the same as its FP8 rate, because the tensor cores process both 8-bit formats at one speed. Ironwood (2{,}307 BF16, 4{,}614 FP8) matches the B200 at 16 and 8 bits and has no FP4; Maia 200 (1{,}270 BF16, which is derived rather than published, since Microsoft gives only FP8 and FP4; 5{,}070 FP8; 10{,}150 FP4) is the only non-GPU part that supports 4-bit floating point and the only one whose FP8 and FP4 exceed the B200's. Trainium2 is 667 BF16 and 1{,}300 FP8. The startups' peaks are lower by a factor of three to ten: the SN40L 638 BF16, the GroqChip 188 FP16 and 750 TOPS INT8, the Bow 350, the Blackhole 372 BF16 and 745 FP8 (Hot Chips 2024 figures for the 140-core design; the shipping 120-core p150 is listed at 664 BLOCKFP8), the RNGD 256 BF16 and 512 FP8, MTIA v2 177 BF16 and 354 TOPS INT8. Two entries carry a tilde. The 910C's ${\sim}$800 TFLOP/s BF16 is an analyst estimate from two 910B dies, because Huawei does not publish a 910C datasheet. The WSE-3's ${\sim}$15{,}000 is \emph{derived}: Cerebras publishes only ``125 PF AI compute'' for the WSE-3, a figure that counts sparsity and does not name a precision, so the lecture starts instead from the 7.5 PFLOP/s of FP16 quoted for the WSE-2 (a figure that carries no dense/sparse qualifier at its source, slide 13) and doubles it for the WSE-3's wider (8-wide instead of 4-wide) FP16 datapath at about the same core count. The result, about 15 PFLOP/s, is still the largest number in the row by a factor of six, but it belongs to a whole wafer drawing 23 kW, not a chip.

\emph{On-chip SRAM.} This row is where the designs divide. The B200, Ironwood and Trainium2 are n.p.: NVIDIA, Google and AWS describe their caches and scratchpads in general terms but do not publish a total, and for a GPU-like design the figure would be tens to a few hundred megabytes. Every chip that \emph{does} publish an SRAM total publishes it because SRAM is the selling point: 272 MB on Maia 200, 256 MB on MTIA v2 and the RNGD, 192 MB on the 910C (an analyst figure, not Huawei's), 210 MB on the Blackhole, 230 MB (220 MiB) on the GroqChip, 520 MB on the SN40L, 900 MB on the Bow, and 44 GB on the WSE-3, two hundred times the others.

\emph{Off-chip memory.} The ``memory'' row is the mirror image. The hyperscaler parts and the B200 use HBM: HBM3e on the B200 (180 GB at 7.7 TB/s in the HGX module), Ironwood (192 GB, 7.4 TB/s) and Maia 200 (216 GB, 7.0 TB/s), HBM3 on Trainium2 (96 GB, 2.9 TB/s) and the RNGD (48 GB, 1.5 TB/s), HBM2e on the 910C (128 GB, 3.2 TB/s, an analyst estimate). The SN40L is the only chip with two off-chip tiers, 64 GB of HBM3 at 1.6 TB/s and up to 1.5 TB of DDR5 behind it (slide 17). Three chips have \emph{no} off-chip DRAM at all: the WSE-3, which streams weights from external MemoryX servers in training (slide 14); the GroqChip, which holds weights only in SRAM and so needs hundreds of chips per large model (slide 16); and the Bow, whose weights beyond 900 MB stream from host DDR through the IPU's links. MTIA v2 uses LPDDR5 (128 GB at 0.2 TB/s), a low-cost, low-bandwidth choice that only makes sense for its first job, recommendation models whose embedding tables are big but whose per-query bytes are small. The Blackhole uses GDDR6 (32 GB at 0.5 TB/s), the graphics-card memory, chosen for price.

\emph{Scale-up and domain.} The scale-up row gives the bandwidth one chip has into its largest directly connected group, and the next row gives that group's size. The B200 has 1{,}800 GB/s of NVLink 5; as an HGX B200 module, the form the table's footnote specifies, its NVLink domain is the 8 GPUs on the baseboard, which is the table's ``8''. The 72-GPU domain belongs to the GB200 NVL72 rack, a different product built from the same GPU. Ironwood has 1{,}200 GB/s of inter-chip interconnect into a 9{,}216-chip pod joined by optical circuit switches, by far the largest coherent domain in the table (slide 11). Trainium2 has 1{,}000 GB/s of NeuronLink into a 64-chip UltraServer; Maia 200 has 2{,}800 GB/s over its Ethernet-based fabric into a 6{,}144-chip domain. MTIA v2 uses PCIe and its ``domain'' of 72 is a rack of cards. The WSE-3's 150 GB/s is the CS-3 system's host I/O, and its domain of 2{,}048 is the largest cluster Cerebras says its MemoryX/SwarmX fabric supports; that is a data-parallel cluster, not a coherent scale-up domain in the footnote's sense, so the entry is not comparable with its neighbors. The GroqChip's 480 GB/s of chip-to-chip links (a figure the sources consulted do not print; Trainium2's 1{,}000 GB/s is likewise unconfirmed) reach, in Groq's ISCA 2022 design, 10{,}440 chips in a deterministic Dragonfly. The Bow's IPU-Links run at 64 GB/s each, into a Bow Pod of up to 1{,}024 IPUs; the Blackhole's Ethernet ports give 400 GB/s into a 32-chip Galaxy server (ten to twelve 400 Gb/s ports would be 500--600 GB/s, so the table's 400 is a conservative or reduced-port figure); the RNGD's 128 GB/s is PCIe 5 into an 8-card server; the 910C reaches 384 chips in a CloudMatrix (slide 26). The SN40L's per-chip link bandwidth and the 910C's are n.p.

\emph{Power.} Five vendors do not publish chip power: Google (Ironwood), AWS (Trainium2), SambaNova (SN40L), Graphcore (Bow) and Huawei (910C). The B200 is 1{,}000 W in the HGX module; Maia 200 750 W; Blackhole 300 W; RNGD 180 W; MTIA v2 90 W. The GroqChip's ${\sim}$300 W is from third-party listings of the GroqCard, not a Groq datasheet, and the WSE-3's 23{,}000 W is the whole CS-3 system including cooling and power supplies. Power per chip is, for most of these parts, the number the buyer most needs and the vendor least wants to compare.

\emph{Software.} The last row names the stack each chip is programmed through: CUDA; XLA and JAX (with PyTorch/XLA) for the TPU; AWS Neuron; the Maia SDK (Triton-based); PyTorch directly for MTIA; Cerebras's CSoft (PyTorch front end, CSL for kernels); SambaFlow; GroqWare; Poplar; TT-Metalium (``TT-Metal'') with TT-NN and TT-Forge above it; Furiosa's SDK; and CANN for the Ascend. Eleven different compilers is the cost of eleven different architectures, and slide 35 lists it as the first reason most of the programs did not survive.

\begin{formulas}
Per-chip ridge point (the arithmetic intensity at which compute and memory balance) is peak over bandwidth:
\[ I^\ast = \frac{P_{\text{dense}}}{B_{\text{mem}}} . \]
B200 at FP8: $4{,}500\times10^{12} / 7.7\times10^{12} \approx 584$ FLOP/byte; Ironwood at FP8: $4{,}614/7.4 \approx 624$; Trainium2 at FP8: $1{,}300/2.9 \approx 448$; RNGD at FP8: $512/1.5 \approx 341$ (all derived). A chip whose weights sit in SRAM has no meaningful $I^\ast$ against DRAM; its ridge is set by on-chip bandwidth (80 TB/s for the GroqChip, 21 PB/s for the WSE-3).
\end{formulas}

\begin{terms}
\term{n.p.}{Not published: a figure the vendor had not disclosed as of November 2026. The lecture prints n.p. rather than an estimate, because a missing number is itself information about how the part is sold.}
\term{dense peak}{The vendor's maximum matrix throughput counted without structured sparsity; the only peak this lecture compares across chips.}
\term{HGX module}{NVIDIA's eight-GPU baseboard form of a data-center GPU; the B200 in HGX is rated at 1,000 W and slightly lower peaks than the B200 in a GB200 superchip.}
\term{wafer-on-wafer bonding}{Bonding two whole wafers face to face before dicing; Graphcore's Bow used it to add a power-delivery wafer under its compute die.}
\term{coherent domain}{The largest group of chips joined by the scale-up fabric so that each can load from and store to the others' memory directly; 72 GPUs in an NVL72, 9,216 chips in an Ironwood pod.}
\term{GDDR6}{Graphics DRAM soldered on the board beside the chip; cheaper and lower in bandwidth than HBM. Tenstorrent's Blackhole uses it.}
\term{LPDDR5}{Low-power DRAM designed for phones and laptops; high capacity per dollar and per watt, low bandwidth. Used by MTIA v2, Qualcomm's AI cards and d-Matrix.}
\term{ridge point}{Peak FLOP/s divided by memory bandwidth: the arithmetic intensity, in FLOP per byte, above which a kernel is compute-bound and below which it is memory-bound.}
\end{terms}

\begin{intuition}
Read the table by columns for what each chip is, and by rows for what each design gave up. The SRAM row and the memory row are complements: chips that publish a large SRAM number have little or no DRAM, and chips with the most HBM do not bother to publish their SRAM. The n.p. cells cluster on the captive chips (no buyer to convince) and on power (no vendor wants that comparison).
\end{intuition}

% ---------------------------------------------------------------------
\slide{32}{MLPerf: who shows up, and who never did}{The neutral chip benchmark: TPUs, Gaudi, the IPU and a few inference parts came; Cerebras and Groq never did.}

\para{What the slide shows} The upper left of the slide is a participation matrix. Its columns are the years 2018 to 2026; its rows are non-GPU accelerator lines: Google TPU, Intel Gaudi, Graphcore IPU, Qualcomm Cloud AI 100, Huawei Ascend, Rebellions ATOM, Furiosa Warboy, Untether speedAI, and a last row for Cerebras, Groq and SambaNova together. A filled dot marks a year in which the vendor's part appeared in an MLPerf Training round; a hollow circle marks an MLPerf Inference round. Google's row has Training dots in every year from 2018 to 2025 but Inference circles only in 2019, 2023, 2024 and 2025; Gaudi has Training dots from 2021 to 2024 and Inference circles in 2023 and 2024; Graphcore has Training dots in 2021 and 2022 only; Qualcomm has Inference circles from 2021 to 2024; Huawei has a single Training dot in 2020; Rebellions a single Inference circle in 2023; Furiosa Inference circles in 2021 and 2022; Untether one in 2024. The last row says ``no submission in any round, 2018--2026''. Below the matrix, a five-row table lists one notable result for each of five vendors. The right panel, ``Reading the matrix'', explains who submits, who left, who never came, and what that means for the rest of the lecture.

\para{Details} MLPerf is the benchmark suite run by MLCommons, an industry consortium founded as MLPerf in May 2018 and incorporated in December 2020. It has two main families. \emph{Training} measures time to train a fixed set of models (ResNet-50, BERT, and from 2023 GPT-3 175B, later Llama 2 70B LoRA fine-tuning, Llama 3.1 and others) to a fixed quality target; \emph{Inference} measures throughput and latency for fixed models under four scenarios (offline, server, single-stream, multi-stream). Results are submitted by vendors or their partners, peer-reviewed by the other submitters before publication, and split into a \emph{closed} division (the same model and rules for everyone, comparable across hardware) and an \emph{open} division (submitters may change the model, for instance to a sparse or quantized one, so results are not directly comparable). Training rounds have run twice a year from v0.5 (December 2018) to v5.1 (November 2025) and v6.0 (June 2026); Inference from v0.5 (November 2019) to v6.0 (April 2026). It is the only cross-vendor benchmark with that kind of review, which is why participation itself is a signal.

The matrix follows the participation table compiled for this lecture from the MLCommons result pages and press releases$^\dagger$. \emph{Google} submitted TPUs to Training from v0.5 (TPU v2 and v3) through v5.0 (June 2025, Trillium listed as available), except v2.1 (November 2022) and v3.0 (June 2023), from both of which it was absent; every year from 2018 to 2025 still has at least one Google Training round, which is why the matrix row is full. Its TPU Inference entries are fewer than its presence suggests: TPU v3 in v0.5 (2019), TPU v5e in v3.1 (2023) and v4.0 (2024), and Trillium in v5.0 (2025). Google submitted no inference results at all from v0.7 to v3.0, and its entries in Inference v4.1, v5.1 and v6.0 are NVIDIA GPU systems, not TPUs. It was not among the submitters of Training v5.1 (November 2025), its Training v6.0 entry is a GB200 NVL72, and Ironwood, generally available since November 2025, has not appeared in any round. \emph{Habana}, then Intel after the December 2019 acquisition, entered Training in v1.0 (June 2021) with the first Gaudi, debuted Gaudi 2 in v2.0 (June 2022), ran GPT-3 on 384 Gaudi 2 in v3.0 (June 2023) and on 1{,}024 in v4.0 (June 2024); it submitted Inference in v3.1 (September 2023) and v4.0 (March 2024). Intel said in June 2024 that it planned to submit Gaudi 3; it never did. \emph{Graphcore} entered Training in v1.0 (June 2021) and last appeared in v2.0 (June 2022), where Baidu also entered IPU-POD16 and POD64 systems; there are no IPU entries in v2.1, and Graphcore never submitted Inference. \emph{Qualcomm} submitted the Cloud AI 100 to Inference from v1.0 (April 2021) through v4.0 (March 2024), never to Training; the independent submitter Krai has carried AI 100 systems in later rounds. \emph{Huawei} appears once, in Training v0.7 (July 2020), with 512 Ascend 910 chips on ResNet-50 in the closed division and four of the five open-division results, the last Chinese accelerator submission. \emph{Rebellions} submitted the ATOM to Inference v3.0 (April 2023) only. \emph{FuriosaAI} submitted the Warboy to Inference v1.1 (September 2021) and v2.0 (April 2022); its second-generation RNGD has not been submitted. (Habana's Goya inference chip in v0.5, 2019, and Furiosa's own v0.5 entry are not drawn in the matrix.) \emph{Untether AI} submitted the speedAI240 to Inference v4.1 (August 2024), as a first-time submitter with notably good performance per watt, and folded ten months later (slide 28).

The bottom row is the striking one. Cerebras, Groq and SambaNova, the three best-funded non-GPU startups and the three whose inference speed is the subject of slide 33, have never submitted to any MLPerf round. Neither have Tenstorrent, the hyperscalers' captive chips (Trainium and Inferentia, Maia, MTIA), d-Matrix, Etched or Positron. AWS's Inferentia appeared once indirectly, in Inference v2.0 (April 2022), only as a backend of NVIDIA's Triton server submission.

The notable-results table picks one result per vendor to show what a submission bought. Huawei's v0.7 entry put 512 Ascend 910 on ResNet-50 and established that the Ascend could train at scale, the year before the sanctions cut it off from TSMC. Graphcore's first entry (v1.0) ran BERT in about nine minutes and ResNet-50 in 14.5 minutes on an IPU-POD64, and Graphcore framed it on price: an IPU-POD16 at a list price of \$149{,}995 against roughly \$300{,}000 for a DGX A100, for a claimed 1.6$\times$ (ResNet) and 1.3$\times$ (BERT) better performance per dollar$^\dagger$. Qualcomm's v3.0 entry put eighteen Cloud AI 100 Pro cards in one 2U server for over 430{,}000 ResNet-50 inferences per second offline, 241 per second per watt. Intel's v4.0 entry trained GPT-3 175B (to the benchmark's checkpoint target) on 1{,}024 Gaudi 2 in 66.9 minutes, the only merchant non-NVIDIA part with GPT-3 results. Google's v4.1 entry (November 2024) showed a Trillium preview and v5p pods on the same GPT-3 benchmark with 99\% weak-scaling efficiency, Google's own figure.

The panel's reading of the pattern is the lecture's. \emph{Who submits} are vendors that sell (or rent) a general-purpose part against a GPU that already has a number in the same table: a buyer comparing Gaudi 2 with an H100 wants the same benchmark run on both, and an absent result is read as a bad one. \emph{Who left} did so when the comparison stopped flattering them: Graphcore after v2.0, as the H100 arrived; Intel after Gaudi 2, rather than show Gaudi 3 against Blackwell; Qualcomm after v4.0; Untether after one round; and Google after Training v5.0, with Ironwood not entered. \emph{Who never came} falls into two groups. The subtitle names Cerebras and Groq, and the panel adds SambaNova: ``these three sell \emph{latency per user}'', tokens per second for one stream, which neither time-to-train nor MLPerf's throughput-oriented inference scenarios reward (a design that needs 576 chips to hold a 70B model scores badly on per-chip throughput even when it delivers one user 300 tokens per second). Being SRAM-heavy is not by itself the reason, since Graphcore's SRAM-only IPU submitted Training results in 2021--22 and Untether's SRAM-based part submitted to Inference v4.1; what the three non-submitters share is the per-user-latency product. The captive chips have no buyer to convince: AWS, Microsoft and Meta deploy them to their own services, and a public benchmark offers them only a chance to look worse than NVIDIA. Huawei, Kunlun and Cambricon stopped after 2020 (or never came) for both reasons and a third, export controls; the Shenzhen sparse-inference company Moffett did submit open-division Inference results in 2022--23, so ``the Chinese parts'' in the panel should be read as those three.

\emph{What it means} for the rest of the lecture: almost every performance number in it is either the vendor's own or a third party's measurement of a \emph{service} (slide 33), not a peer-reviewed measurement of a chip. A peak without a benchmark is a claim. Even SemiAnalysis's open InferenceMAX suite (October 2025) and its successor InferenceX (February 2026) covered NVIDIA and AMD GPUs only, with TPU and Trainium backends promised and not delivered within the year$^\dagger$.

\begin{terms}
\term{MLPerf}{The benchmark suites of the MLCommons consortium: Training (time to train fixed models to fixed quality) and Inference (throughput and latency for fixed models under fixed scenarios), with peer review of every result before publication.}
\term{MLCommons}{The nonprofit industry consortium, incorporated in December 2020, that runs MLPerf; its members include the chip vendors, the cloud providers and academic groups.}
\term{closed division}{The MLPerf category in which every submitter runs the same model under the same rules, so results are comparable across hardware.}
\term{open division}{The MLPerf category in which submitters may change the model (sparsify it, quantize it, retrain it), so results show what a system can do but are not directly comparable.}
\term{time to train}{MLPerf Training's metric: wall-clock minutes from start to reaching the benchmark's quality target, for a stated number of chips.}
\term{weak scaling}{Growing the problem with the number of chips so that the work per chip stays constant; 99\% weak-scaling efficiency means the time per step barely rises as chips are added.}
\term{submitter}{The organization that runs and submits an MLPerf result; usually the chip vendor or a system builder, sometimes an independent party such as Krai.}
\term{InferenceMAX}{SemiAnalysis's open-source, nightly-run inference benchmark (October 2025), renamed InferenceX in 2026; it has measured only GPUs so far.}
\end{terms}

\begin{intuition}
Treat MLPerf participation as a revealed preference. A vendor that expects to look good on a neutral, peer-reviewed comparison enters it; a vendor that leaves, or never comes, has decided the comparison does not help it sell. That the fastest inference services in the world have never entered says that what they sell, latency for one user, is not what the benchmark measures, and also that they have chosen not to be measured.
\end{intuition}

% ---------------------------------------------------------------------
\slide{33}{Inference speed and price: what the SRAM machines deliver}{Thousands of tokens per second for one user, third-party measured; the catch is the cost per token.}

\para{What the slide shows} A horizontal bar chart on a logarithmic axis, from 20 to 6{,}000 output tokens per second for one user, titled ``Artificial Analysis, as vendors cited it''. Nine bars, from top to bottom: GPU clouds serving a 70B model in 2024 (a range of 30--100); Groq on Llama 2 70B in February 2024 (241); Cerebras on Llama 3.1 70B at its August 2024 launch (450); SambaNova on Llama 3.1 70B in September 2024 (461); Cerebras on Llama 3.1 70B in October 2024 (2{,}100); SambaNova on Llama 3.1 405B in September 2024 (132); Cerebras on Llama 3.1 405B in November 2024 (969); Cerebras on gpt-oss-120B on the CS-4 in August 2026 (4{,}400); and NVIDIA's LPX rack on Gemma 4 31B in August 2026 (3{,}431). A caption explains the conditions. The right panel, ``Speed, price and the batch'', gives the reason the SRAM machines are fast, what that speed costs, the published prices, and the caveat that a service is measured, not a chip.

\para{Details} The quantity on the axis is \emph{output tokens per second for one user}: how fast text appears for a single request, after the first token. It is the metric the SRAM-resident designs were built for, and the one MLPerf does not report (slide 32). The source of every bar is Artificial Analysis, an independent company that sends standard prompts to public API endpoints and records the speed, the time to first token and the price, but the bars are labeled ``as vendors cited it'' because each figure was taken from the vendor's announcement quoting Artificial Analysis at a moment when it flattered the vendor$^\dagger$. They are therefore best-case snapshots, not averages.

The bars tell a short history. In February 2024 Groq became the first non-GPU chip on Artificial Analysis's API leaderboard, at 241 tokens per second on Llama 2 70B (Groq's own figure was about 300); GPU clouds at the time served 70B models at roughly 30 to 100 tokens per second per user, because they batch many users together and each user's decode step waits on the batch (the range on the slide is ``typical 2024 cloud serving at batch''). Cerebras launched its inference service on August 27, 2024 at 450 tokens per second on Llama 3.1 70B and 1{,}800 on the 8B, with 16-bit weights, and two months later, on October 24, 2024, announced 2{,}100 on the same 70B model on the same hardware, a 4.7$\times$ improvement (derived) from software alone: better placement of layers across the wafers, better pipelining, and (as the industry soon did generally) speculative decoding. SambaNova, which had launched its cloud on the SN40L in September 2024, reported 461 tokens per second on the 70B and 132 on Llama 3.1 405B; Cerebras answered with 969 on the 405B in November. In August 2026 Cerebras's CS-4 ran gpt-oss-120B at 4{,}400 tokens per second per user (Cerebras's own figure, not Artificial Analysis's), and Artificial Analysis measured NVIDIA's Groq-derived LPX rack, on a private pre-release endpoint at concurrency one, at 3{,}431 tokens per second on Gemma 4 31B. That last bar is the first third-party inference number for the LPU since NVIDIA licensed Groq's technology (slide 16), and it is on a smaller model than the Cerebras bar beside it.

\emph{Why they are fast.} At batch size one, generating a token requires reading every active weight once, so the decode rate per user is bounded by weight-memory bandwidth divided by weight bytes (slide 8). A 70B model in 16 bits is 140 GB; an H100's 3.35 TB/s of HBM bounds one GPU at about 24 tokens per second for its share, and tensor parallelism across eight GPUs raises the bound to roughly 190 tokens per second before NVLink and kernel overheads (all derived). Four CS-3 wafers holding the same model in 176 GB of SRAM have $4\times21$ PB/s of on-wafer bandwidth, a bound of hundreds of thousands of tokens per second; 576 GroqChips have $576\times80$ TB/s. The bound is so far above the observed speeds that it is not the constraint any more: the limit moves to the pipeline depth (the number of chips or wafers a token passes through) and to chip-to-chip latency. That is why the same Cerebras hardware could go from 450 to 2{,}100 tokens per second on software.

\begin{formulas}
Batch-1 decode bound and cost per million tokens:
\[ \text{tokens/s per user} \;\le\; \frac{B_{\text{weights}}}{\text{bytes per token}}, \qquad \text{bytes per token} \approx \text{active weight bytes} \;(+\;\text{KV-cache bytes}), \]
\[ \text{\$ per million tokens} = \frac{\text{system price per hour}}{\text{aggregate tokens/s}\times 3{,}600}\times 10^{6}. \]
At batch $b$, one weight read serves $b$ users, so aggregate tokens/s rise nearly in proportion to $b$ until compute saturates, and the cost per token falls by the same factor; the per-user speed falls.
\end{formulas}

\emph{What it costs.} Capacity is bought in chips. Groq served Llama 2 70B from 576 GroqChips (8 racks of 9 nodes of 8 cards), because the weights alone need over 300 chips at 230 MB each and the rest hold the KV cache, activations and pipeline buffers. SemiAnalysis's February 2024 analysis of Groq's economics concluded that Groq had an advantage in silicon cost per token for latency-optimized serving (the operating point Groq sells), but that in throughput-optimized serving, where a GPU batches many users so each weight read is shared, NVIDIA was about ten times better in performance per dollar on a bill-of-materials basis (the slide's paraphrase of SemiAnalysis, not a quoted figure)$^\dagger$. The two claims are consistent: the batch amortizes the weight read, and an SRAM machine that does not need to batch for speed also does not benefit from batching as much, because it has already paid for the bandwidth in chips.

\emph{The prices.} Groq priced Mixtral 8x7B at \$0.27 per million tokens in 2024. Cerebras priced Llama 3.1 405B at \$6 per million input tokens and \$12 per million output tokens, several times the GPU clouds' prices for the same model. On the Artificial Analysis leaderboard for gpt-oss-120B in November 2026, Cerebras charged a blended \$0.39 per million tokens and delivered the first answer token (including the model's reasoning) in 2.1 seconds; the panel's comparison range of 6 to 69 seconds includes SambaNova at about 6 seconds, which is not a GPU cloud; the GPU clouds at \$0.04--0.05 took about 20 to 56 seconds$^\dagger$. Eight to ten times the price for about ten to twenty-five times the speed to the first answer (derived) is the product, and whether that is a business depends on how many users will pay for speed (slide 36's second open question).

\emph{The caveat.} Every bar is a measurement of a \emph{service}: the provider chooses the quantization (8-bit or 4-bit weights are faster than 16-bit), whether to use speculative decoding (a small draft model proposes several tokens that the large model verifies in one pass), the batching policy, and the hardware allocation behind the endpoint. A faster bar may be a better chip, or a more aggressive quantization, or a quieter hour. The lecture uses these numbers because they are the only third-party measurements that include non-GPU parts, not because they isolate the chip.

\begin{terms}
\term{output tokens per second}{The rate at which generated tokens are returned to one user after the first; the decode speed a reader experiences.}
\term{time to first token}{The delay from sending a request to receiving the first output token; dominated by prefill (processing the prompt) and, for reasoning models, by the hidden reasoning tokens.}
\term{Artificial Analysis}{An independent company that benchmarks public AI model APIs for speed, latency and price on standard prompts; the source of most third-party non-GPU inference numbers.}
\term{concurrency}{The number of simultaneous requests a benchmark sends to an endpoint; concurrency one measures single-user speed.}
\term{batching}{Serving several users' requests in one pass through the model so that each weight read is shared; raises throughput and lowers cost per token at the expense of per-user speed.}
\term{speculative decoding}{A serving technique in which a small, fast draft model proposes several tokens and the large model verifies them in one forward pass, accepting those it would have generated; raises tokens per second without changing the output distribution.}
\term{blended price}{A single price per million tokens that averages input and output prices at an assumed ratio (Artificial Analysis uses three input tokens per output token).}
\term{tensor parallelism}{Splitting each weight matrix across several chips so that every chip reads only its share of the weights for each token; divides the bytes per chip per token by the number of chips.}
\end{terms}

\begin{intuition}
At batch one, speed is a memory-bandwidth problem, and SRAM machines bought so much bandwidth that it stopped being their limit. The price of that bandwidth is the number of chips per copy of the model, which is why the fastest service in the chart is also among the most expensive per token. A GPU cloud is a bus; Cerebras and Groq are taxis.
\end{intuition}

% ---------------------------------------------------------------------
\slide{34}{Deals and share, 2025--26: commitments are not deployments}{The non-GPU chips won gigawatt-scale orders for the first time; most of that power is not yet installed.}

\para{What the slide shows} On the left, a table of eight commitments with three columns: who and when, what was committed, and the status in November 2026. Anthropic--Google (October 2025 and April 2026): up to a million TPUs and over a gigawatt in 2026, with Broadcom booking \$21 billion of TPU orders and multiple gigawatts more from 2027; status ``Ironwood GA; racks delivering''. Anthropic--AWS Rainier (2025): an \$11 billion Indiana campus, about 500{,}000 Trainium2 live in October 2025 (``1M$+$ AWS-wide''), 2.2 GW planned; ``live''. OpenAI--Broadcom (October 2025): 10 GW of custom accelerators and racks from the second half of 2026 to 2029; ``no racks reported''. OpenAI--Cerebras (January and April 2026): over \$20 billion, 750 MW through 2028 with an option to 2 GW; ``CS-3/CS-4 shipping''. Meta--Google (February 2026): a multibillion-dollar TPU rental through Google Cloud, purchases discussed for 2027; ``rental''. Saudi Arabia (``KSA, Humain''): a \$1.5 billion Groq LPU expansion (February 2025) and 200 MW of Qualcomm AI200/AI250 from 2026; ``Groq live; Qualcomm announced''. Microsoft Maia 200 (January 2026): inference in two Azure regions, no external sale; ``in production''. Huawei CloudMatrix 384: over 300 systems delivered by September 2025, about 115{,}000 910C (derived); ``shipping''. The right panel, ``How much of the market'', gives market share, Broadcom's revenue as a proxy for custom silicon, the startup financing record and exits, and how to read the table.

\para{Details} Until 2025 no non-GPU accelerator had been bought in gigawatt quantities by anyone other than its designer. The table records the year that changed, and its third column records how little of the change was installed by November 2026. All entries are from company announcements and press reports$^\dagger$, and the commitments are stated in the units the parties used: chips, dollars or megawatts of power. A rough conversion between them is that a gigawatt of data-center power holds on the order of half a million to a million accelerators, depending on the chip's power and the facility's overhead; the lecture does not convert, because the parties did not.

\emph{Anthropic and Google.} In October 2025 Anthropic and Google announced access to ``up to one million TPUs'', bringing ``well over a gigawatt'' of capacity online in 2026, in a deal reported as worth tens of billions of dollars; a second agreement in April 2026 added multiple gigawatts from 2027. Broadcom, which co-designs and builds the TPU, disclosed a \$10 billion order in its third fiscal quarter of 2025 and an additional \$11 billion order ``from this same customer'' for delivery in late 2026, \$21 billion in all, and the customer was reported to be Anthropic. Two cautions apply: the deal is for capacity, not a purchase of a million chips, and \$21 billion is what Broadcom booked, not Anthropic's total spend. Ironwood became generally available in November 2025, and racks are being delivered.

\emph{Anthropic and AWS.} Project Rainier is the \$11 billion, seven-building campus in New Carlisle, Indiana, that AWS built for Anthropic. On October 29, 2025 AWS said nearly 500{,}000 Trainium2 chips were live there and that Anthropic expected to be using more than a million by the end of 2025; 23 more buildings are planned to bring the site above 2.2 GW. In December 2025 AWS reported more than a million Trainium2 in production, but that count is AWS-wide, across all customers; Anthropic's own figure at the time was ``over 500{,}000''. The slide therefore puts the million in parentheses as ``1M$+$ AWS-wide'': the half-million is Rainier and Anthropic, the million is the fleet (slides 3 and 21).

\emph{OpenAI and Broadcom.} The October 13, 2025 announcement is the largest number in the table and the least advanced: 10 GW of accelerators designed by OpenAI and built and deployed by Broadcom, with Ethernet networking, deployment starting in the second half of 2026 and finishing by the end of 2029 (slide 23). By November 2026 no racks had been reported. The same autumn OpenAI also signed a 10 GW letter of intent with NVIDIA (September 22) and a 6 GW agreement with AMD (October 6), so the Broadcom program is one of three.

\emph{OpenAI and Cerebras.} On January 14, 2026 Cerebras and OpenAI announced a multiyear deal for 750 MW of Cerebras capacity through 2028, reported at over \$10 billion; in April it was expanded to over \$20 billion over three years, and Cerebras's 2026 S-1 describes an option to 2 GW by 2030, a \$1 billion working-capital loan from OpenAI and warrants for OpenAI to acquire Cerebras shares if its spending reaches agreed levels. It is the deal that changed Cerebras's customer mix: G42 and MBZUAI had provided 86\% of its 2025 revenue (slide 15). CS-3 and CS-4 systems are shipping against it.

\emph{Meta and Google.} Meta, which builds MTIA for recommendation and inference (slide 22), signed a multibillion-dollar agreement on February 26, 2026 to rent TPUs through Google Cloud, with talks reported on buying TPUs for its own data centers from 2027. It is the first time a hyperscaler with its own accelerator program rented another's.

\emph{Saudi Arabia and Humain.} The Kingdom of Saudi Arabia announced a \$1.5 billion commitment to expand Groq LPU capacity at the LEAP conference on February 10, 2025, extending a cluster Groq and Aramco Digital had built in Dammam in 2024$^\dagger$. That commitment predates Humain, the Public Investment Fund's AI company, which was launched in May 2025, so the \$1.5 billion is the kingdom's, not Humain's; the slide's ``KSA, Humain'' row groups them because Humain has since become the kingdom's vehicle for AI capacity. In October 2025 Humain committed to 200 MW of Qualcomm's AI200 and AI250 racks from 2026 (slide 24); no shipments had been reported by November 2026.

\emph{Maia 200 and CloudMatrix.} Microsoft announced the Maia 200 on January 26, 2026 as running inference in two Azure regions; it is not sold outside Microsoft. Huawei said over 300 CloudMatrix 384 systems had been delivered by September 2025; at 384 Ascend 910C each that is about 115{,}000 chips ($300\times384 = 115{,}200$, derived), an order of magnitude below the Trainium2 count at Rainier, but in a market NVIDIA cannot legally serve with its best parts (slides 26--27).

The right panel puts the table in proportion. \emph{Share:} TrendForce estimated in October 2025 that NVIDIA held about 70\% of the 2025 AI-chip market and forecast that custom ASICs would pass GPUs in \emph{unit} shipments in 2026, driven by the North American cloud providers and Chinese firms$^\dagger$. By value GPUs still dominate, and the comparison is partly an accounting artifact: a TPU or a Trainium is never sold as a chip, only as cloud capacity, so its revenue appears as Google's or Amazon's cloud revenue. \emph{Broadcom as the proxy:} the one public company whose revenue tracks custom accelerators most directly is Broadcom, which reported \$12.2 billion of AI semiconductor revenue in fiscal 2024, up 220\%, about \$20 billion in fiscal 2025 (the slide marks it derived, but Broadcom stated it on its fourth-quarter call), and guided \$8.2 billion for the first quarter of fiscal 2026 alone, with five XPU customers. \emph{Startups:} US semiconductor startups raised a record \$6.2 billion of venture funding in 2025, up 85\%, about half of it in three rounds (Cerebras's \$1.1 billion Series G, PsiQuantum's \$1 billion, Groq's \$750 million Series E)$^\dagger$. Then came the exits: NVIDIA's non-exclusive license of Groq's technology and hiring of its leaders for a reported \$20 billion (December 24, 2025); Cerebras's IPO at \$185 a share, \$6.4 billion gross including the over-allotment (May 14, 2026); SambaNova's \$1 billion first close at an \$11 billion valuation (July 2026); and Tenstorrent's \$1.5 billion Series E at a \$6.2 billion valuation (October 2026).

\emph{Read the table as} three kinds of number: power contracted, chips live, and racks reported. Only Rainier's Trainium2, Google's Ironwood, Microsoft's Maia 200 and Huawei's CloudMatrix are deployments at scale in November 2026; the OpenAI--Broadcom 10 GW is a plan, and the Cerebras and Humain contracts are being filled.

\begin{terms}
\term{gigawatt-scale deal}{A capacity agreement denominated in electrical power (1 GW = 1,000 MW), the unit data-center builders now use because power, not chips, is the binding constraint on new AI capacity.}
\term{XPU}{Broadcom's name for a custom accelerator it designs and builds for one customer, such as Google's TPU or OpenAI's planned chip.}
\term{capacity agreement}{A contract for the use of accelerators in someone else's data center, measured in chips, power or dollars over years, as opposed to a purchase of the chips.}
\term{over-allotment}{The option (greenshoe) that lets an IPO's underwriters sell up to 15\% more shares than the base offering; Cerebras's \$6.4 billion gross includes it.}
\term{non-exclusive license}{A grant of rights to use technology that leaves the licensor free to license it to others; the form of NVIDIA's December 2025 deal with Groq, which left Groq an independent company.}
\end{terms}

\begin{intuition}
Gigawatts announced, chips running and racks delivered are three different numbers, and the press often prints the first as if it were the third. In November 2026 the non-GPU chips that are actually running at large scale are the ones whose owners also own the workload: Google's TPUs, Amazon's Trainium, Microsoft's Maia and Huawei's Ascend.
\end{intuition}
% =====================================================================
\setstage{ink}
\section{Wrap-up (slides 35--37)}

\slide{35}{Why most of them failed, and what the survivors share}{A few found captive volume, a niche or a buyer for their IP; the rest ran out of software or money.}

\para{What the slide shows} On the left, a timeline of exits from 2020 to 2026 drawn along a horizontal line with five dated points (2020, 2022, 2024, 2025, 2026). Above the line, in red, the failures and distressed sales: Nervana cancelled and Wave Computing bankrupt (2020); Mythic out of cash in November 2022; Graphcore sold for about \$600 million against a \$2.77 billion peak (2024); Untether folding, Esperanto winding down and Tesla's Dojo ending (2025; Tesla restarted a Dojo3 in January 2026, for a different purpose); Taalas agreeing to sell itself to AMD and SambaNova's sale talks stalling (2026). Below the line, in green, the successes in the same years: TPU v4 deployed and Trainium announced (2020); Trn1 generally available and the Ascend 910B on SMIC (2022); Cerebras Inference launched and Tenstorrent raising \$693 million (2024); Groq's technology licensed to NVIDIA for about \$20 billion and Celestial AI sold to Marvell (2025); Cerebras's IPO, SambaNova's \$11 billion valuation, Maia 200 and TPU 8 (2026). Under the timeline a paragraph states the pattern. The right panel, ``Five reasons, and three ways out'', lists software, memory, scale-out, the moving target and capital as the reasons, and captive volume, a niche, and becoming IP as the ways out.

\para{Details} The timeline is selective: it shows dated events for which the lecture has a source (on the section slides) and leaves out the many smaller start-ups that closed quietly. Even so it makes the lecture's central empirical point. Of the merchant startups that set out between 2016 and 2019 to sell a better training or inference chip than NVIDIA's, by November 2026 one had gone public on its own terms (Cerebras), one had been valued as a going concern at a premium (SambaNova, though only after failed sale talks), one had raised steadily as an IP and RISC-V company (Tenstorrent), and the rest had been sold, licensed, folded into a larger company, or closed.

The failures are worth reading individually. \emph{Intel Nervana} was cancelled in early 2020, weeks after Intel bought Habana for about \$2 billion; Intel judged Habana's Gaudi closer to customers' needs, and five years later cancelled Gaudi's successor as well (slide 24). \emph{Wave Computing}, a dataflow-processor company that had bought MIPS, entered Chapter 11 in April 2020. \emph{Mythic}, the analog in-memory design, ran out of cash in November 2022 and survived on a \$13 million rescue in 2023, before raising \$125 million in December 2025 for edge parts (slide 28). \emph{Graphcore} reached a \$2.77 billion valuation in December 2020, reported revenue of \$2.7 million for 2022, and was sold to SoftBank in July 2024 for a reported \$600 million (slide 18). \emph{Untether AI} entered MLPerf in August 2024 with a notably efficient inference card and folded in June 2025, its team hired by AMD. \emph{Esperanto} wound down its AI chip business in 2025 (reported). \emph{Tesla's Dojo}, the in-house training computer with the wafer-scale training tile, was cut in August 2025, its team disbanded; Musk restarted a Dojo3 in January 2026, aimed at space-based compute rather than Tesla's video training, so the program is paused rather than over (slide 24). \emph{Taalas}, which etched a model's weights into a chip's metal layers, agreed in August 2026 to be acquired by AMD, subject to closing (slide 29). SambaNova's 2025 talks to sell itself to Intel for a reported \$1.6 billion stalled before it raised at \$11 billion.

The successes divide into the same three groups the panel names. The hyperscalers (Google's TPU v4 in 2020 through TPU 8 announced in 2026; Amazon's Trainium announced in 2020 and generally available in 2022; Microsoft's Maia 200 in 2026) and Huawei (the 910B on SMIC's process in 2022) succeeded by owning both the chip and the workload. Cerebras succeeded by finding a niche, speed per user, and then a customer large enough to fill it (OpenAI). Groq and Celestial AI succeeded as \emph{intellectual property}: NVIDIA paid about \$20 billion for a non-exclusive license to Groq's inference technology and hired its leaders, and Marvell paid \$3.25 billion upfront for Celestial's photonic fabric.

The paragraph under the timeline states the pattern: a chip two to five times better on one axis (energy, latency, memory capacity per dollar) arrives a year late against a GPU whose software already runs the newest model; a buyer that owns the workload and the compiler (Google, Amazon, Meta, Huawei) can afford to wait out that year, and a merchant startup, which must sell each chip to a customer who can simply buy more GPUs, cannot. Celestial, Taalas and Untether ended as engineering teams inside larger companies, and so did most of Groq's, though Groq Inc.\ itself continues to run GroqCloud and raised again at \$3.5 billion in 2026.

The panel's five reasons are the lecture's diagnosis.
\begin{enumerate}
\item \emph{Software.} Every design here needed its own compiler, and most needed a compiler that produces a \emph{static} schedule (slide 5) plus the libraries users expect (attention kernels, collective communication, quantization, serving). CUDA had a fifteen-year start (slide 6). The startups' compilers handled the models of their launch year well and the next year's late; the convergence on PyTorch 2's compiler interface, OpenXLA and vLLM's hardware plugins (from 2023--25) came too late for most of them.
\item \emph{Memory.} The SRAM-only parts (Graphcore, Groq, the first Cerebras systems) capped the model per chip at hundreds of megabytes or tens of gigabytes just as models grew to hundreds of gigabytes; the HBM parts queued behind NVIDIA for HBM and advanced packaging (slides 8 and 31).
\item \emph{Scale-out.} Training moved to clusters of $10^4$ chips joined by a fast scale-up fabric and a mature collective library; NVIDIA had NVLink and NCCL, Google had its torus and optical switches (slide 10), Huawei built CloudMatrix (slide 26), and the startups sold nodes or racks.
\item \emph{The moving target.} The workload changed every two years: convolutional networks (2015--18), transformers (from 2018), mixture-of-experts and long context (from 2023), reasoning models with long decode (from 2024). A fixed datapath fits the workload it was designed for, and a chip taped out in year $t$ ships into the workload of year $t+2$ (slides 29 and 32).
\item \emph{Capital.} Each generation needs a new tape-out at an advanced node every 18 to 24 months, at nine figures each for masks and engineering, before any revenue from that generation (slide 34).
\end{enumerate}

The three ways out are the survivors' common feature. \emph{Captive volume}: a designer that is also the user (TPU, Trainium, MTIA, Maia, Ascend) does not need to win a sale, only to beat the price of the GPUs it would otherwise buy, and its compiler only has to run its own models. \emph{A niche}: a product that GPUs serve badly, speed per user for interactive and agentic inference, which Cerebras and Groq turned into services rather than chips. \emph{Becoming IP}: selling the design or the team to a company with the volume (Groq to NVIDIA, Celestial to Marvell, Taalas to AMD) or licensing it from the start (Tenstorrent's RISC-V and chiplet licensing).

\begin{terms}
\term{merchant chip}{An accelerator sold to outside customers, as opposed to a captive chip used only by its designer.}
\term{captive chip}{An accelerator designed by a company for its own data centers and services and not sold as a chip: the TPU, Trainium, MTIA, Maia.}
\term{tape-out}{The release of a finished chip design to the foundry for mask making; at advanced nodes a tape-out costs tens to hundreds of millions of dollars before the first wafer.}
\term{acqui-hire}{An acquisition, or a hiring agreement, whose main purpose is the target company's engineering team rather than its products.}
\term{static schedule}{A plan, fixed by the compiler before the program runs, of which operation executes on which unit in which cycle; the alternative to scheduling by hardware at run time.}
\end{terms}

\begin{intuition}
Being two to five times better on one axis was not enough, because the GPU was also improving two to three times per generation on every axis at once and already ran next year's model. The designs that survived changed the question from ``is the chip better?'' to ``who needs this chip?'': its owner, a niche the GPU serves badly, or a larger company buying the ideas.
\end{intuition}

% ---------------------------------------------------------------------
\slide{36}{Summary: what to take from this lecture}{The rivals won where someone else owned the workload; the 2027 questions are about scale, not cores.}

\para{What the slide shows} Two full-height panels. The left, ``Eight things to remember'', is a numbered list, each item with the slides that support it. The right, ``Open questions for 2027'', poses five questions, each with the evidence that will answer it.

\para{Details} The eight items compress the lecture, and each can be expanded from the section it cites.

\emph{1. Decode is bandwidth-bound; every design here answers that (slide 4).} Peak compute has grown about 3$\times$ every two years, DRAM bandwidth 1.6$\times$ and interconnect 1.4$\times$ (Gholami et al.). Generating a token at batch one reads every active weight once, so it runs at memory speed, not compute speed. Every non-GPU design is one answer to that: keep the weights in SRAM (Groq, Cerebras, Graphcore, d-Matrix), reuse each operand many times in a systolic array (TPU, Trainium), add memory tiers (SambaNova), or compute inside the memory (Mythic, EnCharge, NorthPole).

\emph{2. Five organizations: systolic, dataflow, streaming, MIMD mesh, GPU (slide 5).} The lecture's taxonomy asks where operands wait and who decides when they move. A systolic array moves them in lockstep through a grid of multipliers (TPU); a reconfigurable dataflow machine wires compute and memory units per model (SambaNova); a streaming processor moves them through functional slices on a compiler-fixed schedule (Groq); a MIMD mesh gives each tile its own program and memory (Graphcore, Tenstorrent, Cerebras); the GPU schedules threads in hardware.

\emph{3. The TPU is the proof (slides 9--12).} Eight generations of systolic matrix units, HBM from v2, liquid cooling from v3, an optically switched torus from v4, FP8 and 9{,}216-chip pods with Ironwood, and a compiler (XLA) that owns the whole graph. It is the non-GPU line with the longest record of training frontier models, and in 2025--26 it became the first ASIC sold to others at gigawatt scale (slide 12).

\emph{4. The wafer works, inference first (slides 13--15).} Cerebras solved yield with spare cores, power with vertical delivery and heat with a cold plate, and put 44 GB of SRAM and 21 PB/s on one 46{,}225 mm$^2$ die. It was sold for training, but it succeeded commercially as an inference engine selling speed per user.

\emph{5. SRAM streaming buys latency, costs capacity (slides 16 and 33).} Groq's deterministic chip and Cerebras's wafer deliver thousands of tokens per second per user because weight reads never wait on DRAM; the price is hundreds of chips (or several wafers) per copy of a large model, and a cost per token several times a batched GPU's.

\emph{6. Hyperscaler chips are the volume (slides 21--23).} More than a million Trainium2 in production across AWS (over 500{,}000 of them for Anthropic), TPUs contracted by the gigawatt, Maia in production, OpenAI's 10 GW program with Broadcom: by unit count, the custom accelerators of the cloud companies are where the non-GPU volume is, and analysts expect them to pass GPUs in units in 2026.

\emph{7. China scales out what it cannot scale up (slides 25--27).} Cut off from TSMC and from NVIDIA's best parts, Huawei built 384-chip all-optical SuperPoDs whose aggregate exceeds a GB200 NVL72 at about four times the power, developed its own HBM, and published a roadmap to 8{,}192-chip systems.

\emph{8. Distrust ``AI FLOPS'' (slides 31--33).} Ask which precision, dense or sparse, and measured by whom. Most of the numbers in this lecture are vendor peaks, the benchmark that would test them is skipped by most of the vendors, and the third-party speed measurements are of services, not chips.

The right panel's open questions are each tied to evidence that should arrive in 2027.

\emph{Do the custom chips ship at the promised scale?} The OpenAI--Broadcom 10 GW, Trainium 3 and 4, Maia 200 and TPU 8 are commitments and roadmaps; deployments, not announcements, will show whether ASICs pass GPUs in units as TrendForce forecast. The test is racks reported in service, the third column of slide 34.

\emph{Does inference speed stay a business?} Cerebras and NVIDIA's Groq-derived LPX rack sell speed per user. GPUs with FP4, more HBM per chip and speculative decoding close the gap each year; if a batched GPU serves a reasoning model fast enough for most users, the premium for speed shrinks to the users who need the very fastest.

\emph{Does Ascend reach its roadmap?} Huawei promised the 950DT in the fourth quarter of 2026, the 960 in the fourth quarter of 2027 and 8{,}192-chip Atlas 950 SuperPoDs, on SMIC's 7 nm-class processes and its own HBM. Each step is a test of whether scale-out can keep substituting for process technology.

\emph{Which startups are left?} SambaNova's SN50, Rebellions' REBEL systems and d-Matrix's Raptor are due to ship, or slip, in 2027. The lecture's base rate (slide 35) is not encouraging; each has a specific niche argument.

\emph{Does the GPU absorb the ideas?} NVIDIA's Blackwell added tensor memory beside the tensor cores, close to the on-chip operand storage the TPU and the startups argued for, and its 2026 racks include an LPU, Groq's design. If the GPU keeps absorbing the rivals' ideas, the lasting effect of a decade of non-GPU accelerators may be on the GPU.

\begin{terms}
\term{base rate}{The fraction of past cases with a given outcome, used as the starting estimate for a new case; here, the share of AI-chip startups that survived independently.}
\term{tensor memory}{A per-SM on-chip memory introduced with NVIDIA's Blackwell to hold tensor-core operands and accumulators close to the multipliers; an idea the specialized accelerators had used for years.}
\end{terms}

\begin{intuition}
The lecture's one-sentence summary is in the subtitle: the rivals won where someone else owned the workload. The next year's questions are not about whose core is cleverer but about who installs gigawatts, whose software runs the next model, and whether speed for one user remains worth paying for.
\end{intuition}

% ---------------------------------------------------------------------
\slide{37}{Sources}{Papers and talks, vendor documentation, filings, benchmark tables and reporting; $^\dagger$ marks vendor or press material.}

\para{What the slide shows} Two columns of compact references in six groups: Concepts and the TPU; Cerebras, Groq, SambaNova; Graphcore, Tenstorrent, Rebellions, Furiosa; Hyperscalers; China; Emerging; and Measured and money. Groups whose material is mostly vendor or press sources carry a dagger in their heading.

\para{Details} The sources for this lecture are of five kinds, and the slide's grouping makes it possible to see which kind supports which part.

\emph{Peer-reviewed papers and conference talks} carry the principles and the best-documented designs. The concepts section rests on Horowitz's ISSCC 2014 energy table, Dally, Turakhia and Han's 2020 CACM update, Gholami and colleagues' memory-wall analysis (IEEE Micro 2024) and Valiant's 1990 bulk-synchronous parallel model. The TPU is the best-documented accelerator in the lecture because Google published it: Jouppi et al. at ISCA 2017 (TPU v1), Norrie et al. at Hot Chips 32 in 2020 and in IEEE Micro 2021 (v2 and v3; the talk is the source of the utilization argument for 128 $\times$ 128 arrays on slide 7), Jouppi et al. at ISCA 2021 (TPUv4i) and ISCA 2023 (TPU v4). The startups published less but not nothing: Groq's two ISCA papers (2020, 2022), SambaNova's SN40L at MICRO 2024, Cerebras's architecture paper in IEEE Micro (2023), the Graphcore IPU microbenchmark study (arXiv 1912.03413), Meta's MTIA at ISCA 2023, Huawei's DaVinci at Hot Chips 31 and HPCA 2021, IBM's NorthPole in Science (2023). Hot Chips talks, which are reviewed for inclusion but are vendor presentations, supply most second-generation details (Ironwood, Maia, Blackhole, RNGD, Gaudi 3).

\emph{Vendor documentation} (datasheets, product pages, cloud documentation such as Google's TPU pages and AWS's Neuron documentation, press releases) supplies most specifications for parts that have no paper. These numbers are daggered on the slides: they are what the vendor states, not what anyone has measured.

\emph{Filings} are the most reliable financial sources: Cerebras's S-1 registration statements (2024 and 2026) and quarterly 8-Ks for revenue and customer concentration, Broadcom's 8-Ks for AI revenue. Government documents (the US Commerce Department's export rules of 2022--25) anchor the China section's policy calendar.

\emph{Benchmark tables} are the MLCommons results for Training v0.5 to 6.0 and Inference v0.5 to 6.0, the source of slide 32, and Artificial Analysis's provider leaderboards, the source of slide 33. The ALCF AI Testbed evaluations (Emani et al., PMBS 2022 and arXiv 2310.04607) are among the few published side-by-side measurements of Cerebras, SambaNova, Graphcore, Groq and Habana systems on the same workloads.

\emph{Analysis and reporting} fills in the business story: SemiAnalysis on Groq's economics and on CloudMatrix 384, CSET on China's semiconductor position, TrendForce on market share, Crunchbase News on venture funding, and trade press (The Next Platform, The Register, ServeTheHome, Tom's Hardware, HPCwire) for dated events. These are secondary sources, daggered, and wherever two disagreed the lecture says so or uses the more conservative figure.

The full list, with one line per source on why it matters and a link, is Appendix~D of these notes and the paper kit that accompanies them; the kit's script fetches the papers that are freely available and lists the rest with their landing pages.

\begin{terms}
\term{Hot Chips}{The annual IEEE symposium on high-performance chips, held at Stanford each August, where vendors present the architecture of new processors; talks are selected by a committee but are vendor material.}
\term{S-1}{The registration statement a company files with the US Securities and Exchange Commission before an initial public offering; it discloses audited revenue, customers and risks.}
\end{terms}

\begin{intuition}
Trust the sources in the order the lecture uses them: papers and filings first, benchmark tables second, vendor documentation third, reporting last. The dagger on a number is a reminder of which tier it came from.
\end{intuition}
% =====================================================================
\clearpage
\appendix
\setstage{ink}
\section{Notation}

\begin{center}\scriptsize\setlength{\tabcolsep}{4pt}\renewcommand{\arraystretch}{1.05}
\begin{tabular}{@{}lp{4.5cm}p{6.9cm}@{}}
\toprule
symbol or term & meaning & notes\\
\midrule
MAC, FMA & one multiply-accumulate (fused multiply-add) & counted as two operations in every peak: peak $= 2\times$ MACs per cycle $\times f$\\
MXU & matrix unit: the TPU's systolic array & $256\times256$ 8-bit MACs on TPU v1; $128\times128$ BF16 from v2; 4 per TensorCore on v4\\
$N$, $f$ & systolic array side, clock & one $N\times N$ array does $N^2$ MACs per cycle; TPU v1 $65{,}536$ at 0.7 GHz\\
PE, core, tile & the replicated processing element & 900{,}000 cores on a WSE-3; 1{,}472 tiles on a GC200 IPU; 140 Tensix cores on Blackhole (120 enabled on the p150)\\
TensorCore, NeuronCore, Da Vinci core & the per-chip compute cores of Google, AWS and Huawei & 2 TensorCores on TPU v4; 8 NeuronCores on Trainium2; 30--32 Da Vinci cores per Ascend 910 die\\
SRAM, HBM, DRAM & on-chip static RAM; stacked DRAM on the package; DDR or LPDDR off the package & always stated: Groq 220 MiB SRAM and no DRAM; WSE-3 44 GB SRAM; SN40L 520 MB SRAM $+$ 64 GB HBM $+$ 1.5 TB DDR\\
MiB, MB & binary ($2^{20}$ B) and decimal ($10^6$ B) megabytes & 220 MiB $= 230.7$ MB (Groq); TPU v1's 28 MiB\\
$W$ & weight bytes of a model & parameters $\times$ bytes per parameter: 70B at 8 bits $= 70$ GB, at 16 bits 140 GB\\
$B$ & memory bandwidth that serves the weights & 3.35 TB/s (H100 HBM3), 7.4 (Ironwood), 21 PB/s (WSE-3 SRAM), 80 TB/s (Groq SRAM)\\
$I^\ast$ & ridge point $P/B$, FLOP per byte & above it a kernel is compute-bound, below it bandwidth-bound; decode at batch 1 is about 2 FLOP/B\\
OCS & optical circuit switch & 48 in a TPU v4 pod; mirrors steer light between 64-chip cubes\\
torus & a network in which each chip links to its neighbors in each dimension, with wraparound & TPU v2/v3 2D; v4 and later 3D ($4\times4\times4$ cubes)\\
ICI & inter-chip interconnect & the TPU's scale-up links\\
BSP superstep & compute, synchronize, exchange & Graphcore's execution model (Valiant 1990)\\
PCU, PMU & pattern compute and memory units & SambaNova's RDU tiles (Plasticine lineage)\\
Tensix & Tenstorrent's core & five RISC-V cores, a matrix and a vector engine, 1.5 MB SRAM\\
LPU, TSP & Groq's language processing unit, tensor streaming processor & the same chip: one core, 220 MiB SRAM, compiler-scheduled\\
XPU & a custom accelerator built for one buyer & Broadcom's term for the TPU, MTIA, OpenAI's chip\\
SuperPoD & Huawei's multi-cabinet scale-up domain & 384 (Atlas 900 A3), 8{,}192 (Atlas 950), 15{,}488 (Atlas 960)\\
dense, sparse & peaks without and with 2:4 structured sparsity & sparse $= 2\times$ dense; all deck peaks are dense with the precision stated\\
``AI FLOPS'' & a vendor peak with no precision or density stated & quoted, daggered, never compared: WSE-3 ``125 PF'', CS-4 ``750 AI PFLOP/s''\\
tok/s per user & decode speed one user sees & distinct from system throughput summed over a batch\\
n.p. & not published & a figure the vendor had not disclosed as of November 2026\\
$^\dagger$ & vendor, press or analyst source & datasheet, product page, release, analyst note, press report\\
\bottomrule
\end{tabular}
\end{center}

\clearpage
\section{Formulas at a glance}

\begin{description}[leftmargin=0pt,labelindent=0pt,itemsep=6pt,font=\normalfont\bfseries\sffamily]\sloppy
\item[Systolic peak.] An $N\times N$ array does $N^2$ MACs per cycle, so peak $= A\times N^2\times 2\times f$ for $A$ arrays per chip at clock $f$. TPU v1: $1\times256^2\times2\times0.7\ \mathrm{GHz} = 91.8 \approx 92$ TOPS INT8. TPU v4: $2\ \text{TensorCores}\times4\ \text{MXUs}\times128^2\times2\times1.05\ \mathrm{GHz} = 275$ TFLOP/s BF16. Da Vinci core: $4{,}096$ FP16 MACs per cycle; 32 cores at about 1.0 GHz give $32\times4{,}096\times2\times10^9 = 262$ TFLOP/s, consistent with the 910's 256 (derived).
\item[Operand reuse.] In a weight-stationary $N\times N$ array each weight is loaded once and used for every activation vector that passes; each input entering from the edge is used $N$ times as it crosses. Operand reads per MAC fall as $1/N$: $1/256$ on TPU v1, against one or more register reads per operation on a scalar core.
\item[Bytes per token at batch 1.] Every generated token reads every active weight once: bytes per token $= W$ (plus the KV cache read for the context). A dense 70B model: $W = 70$ GB at 8 bits, 140 GB at 16. A mixture-of-experts model reads only its active parameters.
\item[Decode bound.] Tokens per second per copy $\le B/W$. A 70B model at 8 bits on one H100 (3.35 TB/s): $3.35/0.070 = 48$ tok/s; on a B200 (7.7 TB/s) 110; on an Ironwood (7.4 TB/s) 106 (derived). With $n$ chips each reading its own slice the bound is $nB/W$ only if the per-chip slices are read in parallel; in a pipeline the latency is the sum of the stages.
\item[Ridge point.] $I^\ast = P/B$. H100 FP16 $989/3.35 = 295$ FLOP/B; Ironwood FP8 $4{,}614/7.4 = 624$; Groq FP16 $188/80 = 2.4$ FLOP/B; WSE-3 dense FP16 $15{,}000/21{,}000 = 0.7$ FLOP/B (all derived). An SRAM machine's ridge point is near the intensity of decode itself, so decode runs near peak arithmetic.
\item[Chips per model.] $n \ge W/S$ for $S$ bytes of on-chip memory per chip. 70 GB at 8 bits: Groq ($S = 230$ MB) $\ge 305$; Graphcore Bow (900 MB) $\ge 78$; WSE-3 (44 GB) $\ge 2$; at 16 bits double. KV cache, activations and pipeline duplication raise the count: Groq used 576 for Llama 2 70B.
\item[Energy per bit moved.] $E = \text{bits}\times e$, with $e$ from Dally, Turakhia and Han (CACM 2020): local 8 KB SRAM 0.05 pJ/bit, 100 MB on-chip 0.7, LPDDR4 4, SerDes 10, DDR4 20. Horowitz (45 nm): FP32 multiply 3.7 pJ, 8 KB SRAM read of 32 bits 5 pJ, DRAM read of 32 bits 640 pJ; a 32-bit DRAM read is $640/3.7 = 173$ FP32 multiplies, and a 64-bit read (1.3--2.6 nJ, about 2{,}000 pJ) about $2{,}000/3.7 \approx 540$, the ``about 500'' of slide 8.
\item[Growth rates.] Over $t$ years at a rate $r$ per two years, the factor is $r^{t/2}$. Gholami et al.: compute $3.0^{t/2}$, DRAM bandwidth $1.6^{t/2}$, interconnect $1.4^{t/2}$. From 2012 to 2026 ($t = 14$): $3^7 = 2{,}187$, $1.6^7 = 26.8$, $1.4^7 = 10.5$; compute per byte of bandwidth grew $2{,}187/26.8 = 82$ times (derived).
\item[Pod and rack totals.] Pod peak $=$ chips $\times$ chip peak. TPU v4: $4{,}096\times275 = 1.13$ EFLOP/s BF16; Ironwood: $9{,}216\times4{,}614 = 42.5$ EFLOP/s FP8; Atlas 950 SuperPoD: $8{,}192\times1 = 8.2$ EFLOP/s FP8; CloudMatrix 384: $384\times0.78 \approx 300$ PFLOP/s BF16 (all derived from per-chip figures).
\item[Per-chip from totals.] Chip value $=$ system total $/$ chips: 910C BF16 $300/384 = 0.78$ PFLOP/s; HBM $49\ \mathrm{TB}/384 = 128$ GB; HBM bandwidth $1{,}230/384 = 3.2$ TB/s (derived, as the deck's table does).
\item[Efficiency ratios.] Perf/W $= P/\text{power}$; system power ratio $\times$ inverse performance ratio gives energy per unit work. CloudMatrix against NVL72, all-in: $559/145 = 3.9$ times the power; $3.9/1.7 = 2.3$ times the power per FLOP; $3.9/2.1 = 1.9$ per TB/s (SemiAnalysis, from unrounded totals: 1.8); $3.9/3.6 = 1.1$ per TB (all derived; SemiAnalysis gives 2.3, 1.8 and 1.1).
\item[Achieved bandwidth from a measured speed.] Bandwidth used $=$ tokens per second $\times$ bytes per token. Cerebras at 2{,}100 tok/s on Llama 3.1 70B at 16 bits: $2{,}100\times140\ \mathrm{GB} = 294$ TB/s, about $0.35\%$ of four wafers' $84$ PB/s of SRAM bandwidth (derived): the limit is latency, not bandwidth.
\item[Cost per token.] Cost per million tokens $=$ price per hour $/(\text{tokens per second}\times3{,}600)\times10^6$. A service's price per million tokens is the provider's choice and includes its margin; it is not a chip's cost.
\item[Facility power to chips.] Chips $= P_{\text{facility}}/p_{\text{all-in}}$. Anthropic's ``up to 1 million TPUs'' at ``well over a gigawatt'' implies about 1 kW or more per chip all-in (derived).
\item[MLPerf.] A training result is time to train to a target quality; compare two systems at matched scale as $t_A/t_B$, and across scales only with the chip counts stated.
\end{description}

\clearpage
\section{Exercises}

\begin{enumerate}[leftmargin=1.6em,itemsep=6pt]
\item \textbf{Systolic peaks.} TPU v1 has one $256\times256$ array of 8-bit MACs at 700 MHz; TPU v4 has two TensorCores of four $128\times128$ BF16 arrays at 1.05 GHz. Compute both peaks and the number of MACs per cycle. How many times as many MACs per cycle does v4 have, and how much of the 3$\times$ rise in peak comes from the clock? \emph{Answer: v1 $65{,}536$ MACs per cycle, $65{,}536\times2\times0.7\times10^9 = 91.8$ TOPS; v4 $8\times16{,}384 = 131{,}072$ MACs per cycle, 275 TFLOP/s. Twice the MACs and 1.5 times the clock give the factor of 3; the switch from INT8 to BF16 means v4's operations are much more expensive, so the comparison understates the change.}
\item \textbf{The memory wall.} Using Gholami et al.'s rates (compute $3.0\times$, DRAM bandwidth $1.6\times$, interconnect $1.4\times$ per two years), compute the growth of each from 2012 to 2026 and the growth in FLOPs available per byte of DRAM bandwidth. What does the last number mean for a kernel whose intensity is fixed, like batch-1 decode? \emph{Answer: $3^7 = 2{,}187$, $1.6^7 = 26.8$, $1.4^7 = 10.5$; FLOPs per byte rose 82 times. A kernel at fixed intensity uses an ever smaller fraction of peak: decode at about 2 FLOP/B was always bandwidth-bound and is now bound by a factor about 80 larger.}
\item \textbf{Decode bounds.} A dense 70B model at 8 bits is served at batch 1. Compute the upper bound on tokens per second on one H100 (3.35 TB/s HBM3), one B200 (7.7 TB/s), one Ironwood (7.4 TB/s) and one Groq LPU's slice (80 TB/s, if the model fit). Why can no single one of these hold the model except the GPUs and Ironwood? \emph{Answer: 48, 110, 106 and 1{,}143 tok/s. A Groq chip holds 230 MB, so the 70 GB model needs at least 305 chips; the H100 (80 GB), B200 (180 GB) and Ironwood (192 GB) hold 70 GB in HBM.}
\item \textbf{Chips per model.} At 8 and 16 bits, how many Groq LPUs (230 MB SRAM), Graphcore Bow IPUs (900 MB), Cerebras WSE-3 wafers (44 GB) and Ascend 910C (128 GB HBM) are needed just to hold a 70B model's weights? Why did Groq's 2024 deployment of Llama 2 70B use 576 chips rather than the minimum? \emph{Answer: 8 bits: 305, 78, 2, 1; 16 bits: 609, 156, 4, 2. Groq's 576 chips (8 racks of 9 nodes of 8 cards) also hold the KV cache and activations and divide the layers into a pipeline whose stage count is set by the network topology, not by weight bytes alone.}
\item \textbf{SRAM against DRAM energy.} Using Dally's figures, compute the energy to move the 70 GB of weights of a 70B model at 8 bits once (one decode step at batch 1) from DDR4 (20 pJ/bit), LPDDR4 (4), 100 MB-scale on-chip SRAM (0.7) and local 8 KB SRAM (0.05). At 50 tokens per second, what power does each represent? \emph{Answer: $5.6\times10^{11}$ bits; 11.2 J, 2.24 J, 0.39 J and 0.028 J per token; at 50 tok/s, 560 W, 112 W, 20 W and 1.4 W. Batching amortizes the DRAM figure over many users; at batch 1 it does not.}
\item \textbf{Ridge points.} Compute $I^\ast = P/B$ for an H100 (989 TFLOP/s FP16, 3.35 TB/s), an Ironwood (4{,}614 FP8, 7.4 TB/s), a Groq LPU (188 FP16, 80 TB/s SRAM) and a WSE-3 (about 15{,}000 dense FP16, 21 PB/s SRAM). For which of them is batch-1 decode (about 2 FLOP/B) compute-bound? \emph{Answer: 295, 624, 2.4 and 0.7 FLOP/B. Decode is bandwidth-bound on the H100 and Ironwood by factors of about 150 and 300; on the LPU it is near the ridge and on the wafer compute-bound, which is why the SRAM machines are fast per user and why their limit becomes latency.}
\item \textbf{TPU pods.} From the per-chip peaks (v2 46, v3 123, v4 275 TFLOP/s BF16; Ironwood 4{,}614 FP8) compute the pod peaks for 256, 1{,}024, 4{,}096 and 9{,}216 chips, and compare with the deck's 11.5 PF, 100$+$ PF, 1.1 EF and 42.5 EF. Which differs, and why? \emph{Answer: 11.8 PF, 126 PF, 1.13 EF and 42.5 EF. The v3 entry is Google's own ``100+ petaflops'', a lower bound rather than a rounding: the derived peak is 1{,}024 $\times$ 123 $\approx$ 126 PF; the others agree to rounding. Ironwood's 42.5 EF is FP8; in BF16 it is 21.3 EF, so the 37$\times$ rise from v4 is about 19$\times$ on a like-for-like format.}
\item \textbf{CloudMatrix arithmetic.} From the totals (384 chips, about 300 PFLOP/s dense BF16, 49 TB HBM, about 1{,}230 TB/s) derive the 910C's per-chip values and compare with a B200 (2{,}500 TFLOP/s dense BF16, 186 GB, 8 TB/s in the GB200). Using SemiAnalysis's all-in estimates, about 145 kW for the NVL72 and about 559 kW for the CloudMatrix, what is the power ratio and each system's power per accelerator? \emph{Answer: 0.78 PFLOP/s, 128 GB and 3.2 TB/s per 910C: 0.31, 0.69 and 0.40 of a B200. Power ratio $559/145 \approx 3.9$; about $145/72 \approx 2.0$ kW per B200 slot and $559/384 \approx 1.46$ kW per 910C, all-in (derived); the 910C draws less per chip but about 2.3 times as much per FLOP.}
\item \textbf{Bandwidth actually used.} Cerebras served Llama 3.1 70B at 16 bits at 450 tok/s (August 2024) and 2{,}100 tok/s (October 2024) on four CS-3 systems. Compute the weight bandwidth each speed implies and compare with four wafers' SRAM bandwidth (4 $\times$ 21 PB/s). What changed between August and October, and what does the exercise say about the limit? \emph{Answer: 63 TB/s and 294 TB/s, 0.08\% and 0.35\% of 84 PB/s. The hardware did not change; software did (the deck: ``on software''). The limit is pipeline latency across wafers and the time per layer, not bandwidth, so software can move it.}
\item \textbf{Groq's 576 chips.} Compute the aggregate SRAM bandwidth of 576 Groq chips at 80 TB/s each and the decode bound it implies for a 70 GB model, and compare with the 241 tok/s Artificial Analysis measured in February 2024. What fraction of the bound is that, and what binds instead? \emph{Answer: 46 PB/s; $46{,}080/70 \approx 660{,}000$ tok/s; 241 is about 0.04\% of the bound. The binding constraints are the pipeline's depth across hundreds of chips and chip-to-chip latency; the bandwidth argument explains why the design is fast, not how fast.}
\item \textbf{``AI FLOPS''.} Cerebras quotes ``125 PFLOP/s peak AI'' for the WSE-3 and the deck derives about 15 PFLOP/s dense FP16. NVIDIA quotes 9 PFLOP/s dense FP4 and 18 sparse for the B200, and 2.25 PFLOP/s dense FP16. Compute the ratio of each headline to its dense FP16 figure. Why does the deck refuse to compare the 125 with anything? \emph{Answer: $125/15 = 8.3$ for the WSE-3; $18/2.25 = 8$ for the B200's sparse FP4 against dense FP16. A headline that states neither precision nor sparsity could be any of these multiples; only dense figures at a stated precision compare.}
\item \textbf{Captive volume.} Project Rainier was reported at about 500{,}000 Trainium2 chips (1.3 PFLOP/s dense FP8, 667 TFLOP/s BF16 each). Compute its aggregate FP8 and BF16 peaks and express the FP8 figure in Ironwood pods (42.5 EFLOP/s FP8). \emph{Answer: 650 EFLOP/s FP8 and 333 EFLOP/s BF16 (derived peaks); about 15 Ironwood pods. Peaks summed over half a million chips say nothing about how much of them one training job uses.}
\item \textbf{Three memory tiers.} An SN40L has 520 MB of SRAM, 64 GB of HBM and 1.5 TB of DDR5. For eight chips, compute each tier's total and say which tier can hold a one-trillion-parameter composition of experts at 16 bits, and which holds the experts in use. \emph{Answer: 4.2 GB, 512 GB and 12 TB. Two trillion bytes fit only in DDR; the active experts (tens to hundreds of GB) are copied to HBM, and the tiles of the layer being computed live in SRAM. The design bets that switching experts from DDR is rare enough to hide.}
\item \textbf{Hardwired bandwidth.} Taalas's HC1 holds Llama 3.1 8B in a custom 3-bit base type that mixes 3-bit and 6-bit parameters, and generates about 17{,}000 tokens per second for one user. Taking 3 bits per weight as a lower bound, compute the weight bytes and the bandwidth a conventional chip would need to read them at that speed, and express it in B200s (7.7 TB/s). \emph{Answer: $8\times10^9\times3/8 = 3$ GB; $3\ \mathrm{GB}\times17{,}000 = 51$ TB/s, about 6.6 B200s of HBM bandwidth for one user (derived); the 6-bit parameters push both up, to at most 6 GB and 102 TB/s if all were 6-bit. The weights are wiring, so nothing is read.}
\item \textbf{NorthPole.} A NorthPole chip has 224 MB of on-chip memory and no DRAM. A 3-billion-parameter model with INT4 weights is served on 16 cards. Compute the weight bytes, the minimum cards for weights alone, and the per-card share. What else must fit? \emph{Answer: 1.5 GB; at least 7 cards; 94 MB per card on 16. The KV cache (also INT4 in IBM's run), activations and the pipeline's buffers take the rest, which is why 16 cards were used.}
\item \textbf{Price and speed for one user.} On slide 33, gpt-oss-120B is priced at \$0.39 per million tokens on Cerebras and \$0.04--0.05 on GPU clouds. A user asks for 10{,}000-token answers. Compute the cost per answer on each and, assuming 4{,}400 tok/s on Cerebras and an assumed 300 tok/s on a GPU cloud, the generation time. When is the faster service worth its price? \emph{Answer: \$0.0039 against \$0.0004--0.0005, about 8--10 times; 2.3 s against 33 s (the GPU speed is an assumption, not a deck figure). The SRAM service is worth it when the user's waiting time, or an agent's chain of many calls, costs more than a fraction of a cent per answer; for bulk batch work it is not.}
\end{enumerate}
%=====
\clearpage
\setstage{ink}
\section{Reading list}
\label{app:reading}
The papers, talks, datasheets, filings, benchmark tables and reports behind each part of the deck, in the order they are discussed, with one line on why each is worth reading. Entries with an arXiv identifier are included in the paper archive of the kit; entries without one are listed with a URL. Items marked \emph{secondary source} are press coverage, analyst notes or commentary; vendor material is named as such.

\subsection*{1. Overview: a decade of alternatives (slides 1-3)}
\begin{itemize}[leftmargin=1.2em,itemsep=3pt]\sloppy
\item \textbf{Reuther, Michaleas, Jones, Gadepally, Samsi and Kepner (MIT Lincoln Laboratory) (2022).} \textit{AI and ML accelerator survey and trends (IEEE HPEC 2022)}. arXiv:2210.04055. Survey paper. The fourth annual scatter of peak operations per second against power for published accelerators, with release dates and the neuromorphic, photonic and memristive entries; the model for the landscape chart. Slides 2 and 3.
\item \textbf{Reuther, Michaleas, Jones, Gadepally, Samsi and Kepner (MIT Lincoln Laboratory) (2021).} \textit{AI accelerator survey and trends (IEEE HPEC 2021)}. arXiv:2109.08957. Survey paper. The previous year's edition, with the 2016 to 2021 wave of startups (Graphcore, Groq, SambaNova, Cerebras) on one chart; read it to see which entries disappeared by 2026. Slides 2 and 3.
\item \textbf{Hennessy and Patterson (2019).} \textit{A new golden age for computer architecture (Communications of the ACM 62(2), Turing lecture)}. \url{https://cacm.acm.org/?p=7368}. Primary paper (lecture). The end of Dennard scaling and Moore's law as the case for domain-specific architectures, with the TPU as the worked example; the thesis behind every program on the who's-who table. Slides 1 and 3.
\item \textbf{Sastry and Emani (Argonne) (2026).} \textit{ALCF AI Testbed (ATPESC 2026 slides, August 4, 2026)}. \url{https://extremecomputingtraining.anl.gov/sites/atpesc/files/2026-08/AI%20Testbeds.pdf}. Laboratory slides. The non-GPU machines one DOE facility still runs in 2026: 4 CS-3, SN40L nodes with 16 RDUs each, 2 Tenstorrent Galaxy nodes, a 72-chip GroqRack and a Bow Pod64. Slides 3 and 35.
\item \textbf{Argonne Leadership Computing Facility (2026).} \textit{ALCF AI Testbed (system page)}. \url{https://ai.alcf.anl.gov}. Facility page. The testbed inventory: GroqRack (72 chips), Cerebras CS-2, SambaNova DataScale (64 RDUs), Graphcore Bow Pod64, Habana Gaudi; the systems behind the ALCF evaluations cited later. Slides 3 and 32.
\item \textbf{Hot Chips (2026).} \textit{Hot Chips 2026 conference program (August 23-25, 2026)}. \url{https://hc2026.hotchips.org/program/conference/}. Conference program. Where the 2026 accelerators were presented, including the Samsung LPDDR5X-PIM talk; the dates behind the timeline. Slides 3 and 28.
\item \textbf{ServeTheHome (2025).} \textit{Hot Chips 2025 preliminary schedule released}. \url{https://www.servethehome.com/hot-chips-2025-preliminary-schedule-released/}. Press, secondary source. The 2025 machine-learning sessions: Google Ironwood, d-Matrix Corsair, AMD MI350 and NVIDIA GB10 on one day; who was presenting non-GPU silicon that year. Slides 3, 11 and 28.
\item \textbf{The Register (2026).} \textit{AI chip startups soak up \$1.1B in VC funding this week (February 25, 2026)}. \url{https://www.theregister.com/2026/02/25/ai_chips_vc_funding_1point1billion/}. Press, secondary source. MatX (\$500M Series B), Axelera (over \$250M) and SambaNova (\$350M with an Intel collaboration) in one week; the 2026 funding status on the who's-who table. Slides 3 and 34.
\end{itemize}
\subsection*{2. Principles: the memory wall, dataflow and SRAM (slides 4-8)}
\begin{itemize}[leftmargin=1.2em,itemsep=3pt]\sloppy
\item \textbf{Kung (1982).} \textit{Why systolic architectures? (IEEE Computer 15(1))}. \url{https://www.cs.virginia.edu/~smk9u/CS6501F19/kung_-_1982_-_why_systolic_architectures.pdf}. Primary paper (course mirror). The origin of the systolic-array vocabulary: data pulsing through a grid of cells so each operand fetched is used many times; the idea behind the MXU. Slides 5 and 7.
\item \textbf{Horowitz (2014).} \textit{Computing's energy problem (and what we can do about it) (ISSCC 2014, paper 1.1)}. \url{https://www.Gwern.net/doc/cs/hardware/2014-horowitz-2.pdf}. Primary paper (mirror). The 45 nm energy table: adds and multiplies in fractions of a pJ to a few pJ, cache reads of 10 to 100 pJ, a DRAM access of 1 to 2 nJ; the bar chart on the SRAM slide. Slides 4 and 8.
\item \textbf{Dally (NVIDIA and Stanford) (2016).} \textit{High-performance hardware for machine learning (UC Berkeley CS294 lecture, October 19, 2016)}. \url{https://berkeley-deep-learning.github.io/cs294-dl-f16/slides/DL_HW_Berkeley_0916.pdf}. Lecture slides. A legible reproduction of the Horowitz table (32-bit FP multiply 3.7 pJ, 8 KB SRAM read 5 pJ, 32-bit DRAM read 640 pJ) and the argument for keeping data on chip. Slide 8.
\item \textbf{Dally, Turakhia and Han (2020).} \textit{Domain-specific hardware accelerators (Communications of the ACM 63(7))}. \url{https://www.doc.ic.ac.uk/~wl/teachlocal/arch/papers/cacm20dsa.pdf}. Primary paper (course mirror). The four sources of accelerator gain and the updated energy costs: an add instruction about 4,000 times the add, 8 KB SRAM 50 fJ/bit, 100 MB on chip 0.7 pJ/bit, DDR4 about 20 pJ/bit, SerDes about 10 pJ/bit. Slides 6, 8 and 30.
\item \textbf{Gholami, Yao, Kim, Hooper, Mahoney and Keutzer (2024).} \textit{AI and memory wall (IEEE Micro 2024)}. arXiv:2403.14123. Primary paper. Peak server FLOP/s grew 3.0 times every two years, DRAM bandwidth 1.6 and interconnect 1.4; the three exponentials on the memory-wall slide. Slide 4.
\item \textbf{Valiant (1990).} \textit{A bridging model for parallel computation (Communications of the ACM 33(8))}. \url{https://www.osti.gov/biblio/6502724}. Bibliographic record (paper paywalled). The bulk-synchronous parallel model of compute, communicate and synchronize supersteps, which Graphcore's IPU implements literally. Slides 5 and 18.
\item \textbf{Chen, Emer and Sze (2016).} \textit{Eyeriss: a spatial architecture for energy-efficient dataflow for convolutional neural networks (ISCA 2016)}. \url{https://sands.kaust.edu.sa/classes/CS290E/F19/papers/eyeriss.pdf}. Primary paper (course mirror). Defines the weight-stationary, output-stationary, no-local-reuse and row-stationary dataflows, the vocabulary for where operands wait in each machine. Slides 5 and 7.
\item \textbf{Sze, Chen, Yang and Emer (2017).} \textit{Efficient processing of deep neural networks: a tutorial and survey (Proceedings of the IEEE)}. arXiv:1703.09039. Survey paper. Dataflows, memory hierarchy and reduced precision organized in one place, with the energy-per-access hierarchy from registers to DRAM; the best single introduction to this section. Slides 4, 5 and 8.
\item \textbf{Prabhakar, Zhang, Koeplinger, Feldman, Zhao, Hadjis, Pedram, Kozyrakis and Olukotun (2017).} \textit{Plasticine: a reconfigurable architecture for parallel patterns (ISCA 2017)}. \url{https://sands.kaust.edu.sa/classes/CS290E/F19/papers/plasticine.pdf}. Primary paper (course mirror). The coarse-grained reconfigurable array of pattern compute and memory units, 113 mm2 at 28 nm, 49 W, up to 76.9 times an FPGA's perf/W; the academic ancestor of SambaNova's RDU. Slides 5 and 17.
\item \textbf{Fowers, Ovtcharov, Papamichael, ... Chung and Burger (Microsoft) (2018).} \textit{A configurable cloud-scale DNN processor for real-time AI (ISCA 2018)}. \url{https://www.microsoft.com/en-us/research/publication/a-configurable-cloud-scale-dnn-processor-for-real-time-ai/}. Primary paper (publication page). Project Brainwave: an FPGA soft NPU with weights pinned in on-chip memory, 10 to over 35 TFLOP/s at batch 1 on Stratix 10; the first data-center SRAM-resident serving design and Microsoft's step before Maia. Slides 8 and 22.
\item \textbf{Shilov (Tom's Hardware) (2022).} \textit{TSMC's 3nm node: no SRAM scaling implies more expensive CPUs and GPUs (December 16, 2022)}. \url{https://tomshardware.com/news/no-sram-scaling-implies-on-more-expensive-cpus-and-gpus}. Press, secondary source. The SRAM bitcell stays at about 0.021 um2 from N5 to N3E, so on-chip capacity no longer gets cheaper with each node; why SRAM-only designs hit a capacity wall. Slide 8.
\item \textbf{Somala (Epoch AI) (2026).} \textit{Memory has grown to nearly two-thirds of AI chip component costs (May 21, 2026)}. \url{https://epoch.ai/data-insights/ai-chip-component-cost-shares}. Analyst estimate, secondary source. HBM rose from 52 to 63 percent of accelerator component spending between Q1 2024 and Q4 2025, logic dies about 13 percent; the cost side of the SRAM-versus-HBM trade. Slides 4 and 8.
\end{itemize}
\subsection*{3. The Google TPU (slides 9-12)}
\begin{itemize}[leftmargin=1.2em,itemsep=3pt]\sloppy
\item \textbf{Jouppi, Young, Patil, Patterson et al. (Google) (2017).} \textit{In-datacenter performance analysis of a tensor processing unit (ISCA 2017)}. arXiv:1704.04760. Primary paper. TPU v1: a 256x256 8-bit MXU (65,536 MACs), 92 TOPS at 700 MHz, 24 MiB unified buffer, 8 GiB DDR3 at 34 GB/s, 75 W, 28 nm; 15 to 30 times a K80 and 30 to 80 times its perf/W, with the roofline showing DDR3 as the limit. Slides 6, 7 and 9.
\item \textbf{Norrie, Patil, Yoon, Kurian, Li, Laudon, Young, Jouppi and Patterson (Google) (2020).} \textit{Google's training chips revealed: TPUv2 and TPUv3 (Hot Chips 32 slides)}. \url{https://www.hc32.hotchips.org/assets/program/conference/day2/HotChips2020_ML_Training_Google_Norrie_Patil.v01.pdf}. Primary conference slides (the talk behind Norrie et al., IEEE Micro 41(2), 2021). Why 128x128: about 1.6 times the utilization of one 256x256 array at constant FLOP/s; 256-chip 11.5 PF and 1,024-chip 100+ PF pods, 4 ICI links of 500 Gb/s in a 2D torus. Slides 7 and 9.
\item \textbf{Norrie, Patil, Yoon, Kurian, Li, Laudon, Young, Jouppi and Patterson (Google) (2021).} \textit{The design process for Google's training chips: TPUv2 and TPUv3 (IEEE Micro 41(2))}. \url{https://www.Gwern.net/doc/ai/scaling/hardware/2021-norrie.pdf}. Primary paper (mirror). TPU v2 as a different machine from v1: two TensorCores with VLIW scalar units driving vector units and 128 x 128 MXUs (32,768 operations per cycle at 700 MHz), BF16, HBM at 700 GB/s (20 times v1's DDR3), and the ICI torus. Slides 7 and 9.
\item \textbf{Jouppi, Yoon, Ashcraft, Gottscho et al. (Google) (2021).} \textit{Ten lessons from three generations shaped Google's TPUv4i (ISCA 2021, industry track)}. \url{https://www.cs.cmu.edu/~18742/papers/Jouppi2021.pdf}. Primary paper (course mirror). TPUv4i: 7 nm, 138 TFLOP/s BF16, 128 MB of CMEM chosen as the knee of the curve, 175 W chip TDP; and the ten design lessons, including why TPU v2 went to two cores. Slides 7 and 10.
\item \textbf{Jouppi, Kurian, Li, Ma, Nagarajan, Nai, Patil, Subramanian, Swing, Towles, Young, Zhou, Zhou and Patterson (Google) (2023).} \textit{TPU v4: an optically reconfigurable supercomputer for machine learning with hardware support for embeddings (ISCA 2023)}. arXiv:2304.01433. Primary paper. 4,096 chips in 64 cubes joined by optical circuit switches (under 5 percent of cost and under 3 percent of power), SparseCores at 5 to 7 times on embeddings, 1.2 to 1.7 times an A100. Slide 10.
\item \textbf{Jouppi, Lakshmanamurthy, Young and Patterson (Google) (2026).} \textit{Google's training supercomputers from TPU v2 to Ironwood: architectural stability, scale, resilience, power efficiency, and sustainability across five generations (IEEE Micro)}. arXiv:2606.15870. Primary paper. Peak per chip up about 100 times (46 TFLOP/s BF16 to 4,614 FP8), HBM about 10 times, pods from 256 to 9,216 chips, peak perf per TDP watt about 30 times (Ironwood 6 times v5p), 93 percent goodput training Gemini 2.5. Slides 9 to 11.
\item \textbf{Jouppi and Lakshmanamurthy (Google) (2025).} \textit{Ironwood: delivering best in class perf, perf/TCO and perf/Watt for reasoning model training and serving (Hot Chips 2025 slides)}. \url{https://hc2025.hotchips.org/assets/program/conference/day2/61_Google_Ironwood-Final.pdf}. Primary conference slides. The two-die TPU: 4,614 TFLOP/s FP8, 192 GiB HBM3E in 8 stacks at 7.3 TB/s, 1.2 TB/s ICI, 9,216 chips through OCS at 42.5 EFLOP/s FP8 with 1.77 PB of shared HBM. Slide 11.
\item \textbf{Google Cloud (2026).} \textit{TPU7x (Ironwood) system architecture (documentation)}. \url{https://docs.cloud.google.com/tpu/docs/tpu7x}. Vendor documentation. The official per-chip row: 2,307 TFLOP/s BF16, 4,614 FP8, 192 GiB at 7,380 GB/s, 1,200 GB/s bidirectional ICI, 4x4x4 cubes up to 9,216 chips. Slides 11 and 31.
\item \textbf{Google Cloud (2025).} \textit{TPU v6e (Trillium) documentation}. \url{https://docs.cloud.google.com/tpu/docs/v6e}. Vendor documentation. 918 TFLOP/s BF16, 1,836 TOPS INT8, 32 GB at 1,638 GB/s, 800 GB/s ICI, 256-chip 2D torus. Slide 11.
\item \textbf{Google Cloud (2025).} \textit{TPU v5p documentation}. \url{https://docs.cloud.google.com/tpu/docs/v5p}. Vendor documentation. 459 TFLOP/s BF16, 95 GiB usable of 96 GiB HBM at 2,765 GB/s, 1,200 GB/s ICI, 8,960-chip 3D torus. Slide 11.
\item \textbf{Google Cloud (2025).} \textit{TPU v5e documentation}. \url{https://docs.cloud.google.com/tpu/docs/v5e}. Vendor documentation. 197 TFLOP/s BF16, 393 TOPS INT8, 16 GB at 800 GiB/s, 400 GB/s ICI, 256-chip 2D torus; the first inference-line TPU. Slide 11.
\item \textbf{Lohmeyer (Google Cloud) (2024).} \textit{Announcing the general availability of Trillium, our sixth-generation TPU (December 11, 2024)}. \url{https://cloud.google.com/blog/products/compute/trillium-tpu-is-ga}. Vendor blog. 4.7 times peak per chip, double the HBM and ICI, 67 percent better energy efficiency; Gemini 2.0 trained on it; GA date on the business timeline. Slides 11 and 12.
\item \textbf{Pichika and Singh (Google Cloud) (2024).} \textit{Trillium MLPerf 4.1 training benchmarks (November 14, 2024)}. \url{https://cloud.google.com/blog/products/compute/trillium-mlperf-41-training-benchmarks}. Vendor blog. GPT-3 175B in MLPerf Training 4.1 (preview): 99 percent scaling efficiency across Trillium pods with Multislice and up to 1.8 times the perf per dollar of v5p. Slides 11, 12 and 32.
\item \textbf{Google Cloud (2025).} \textit{Announcing Ironwood TPUs general availability and new Axion VMs (November 6, 2025)}. \url{https://cloud.google.com/blog/products/compute/ironwood-tpus-and-new-axion-based-vms-for-your-ai-workloads/}. Vendor blog. Ironwood GA in the coming weeks, 10 times v5p peak, 9,216-chip superpods, and Anthropic's plan to access up to one million TPUs. Slides 11 and 12.
\item \textbf{StorageReview (2026).} \textit{Google announces TPU 8t Sunfish and TPU 8i Zebrafish (May 1, 2026)}. \url{https://www.storagereview.com/news/google-announces-tpu-8t-sunfish-and-tpu-8i-zebrafish}. Press, secondary source. The TPU 8 columns: 8t at 12.6 PF FP4, 216 GB at 6.5 TB/s, 9,600-chip superpods; 8i at 10.1 PF FP4, 288 GB at 8.6 TB/s, 384 MB SRAM, the Boardfly network. Slides 11 and 31.
\end{itemize}
\subsection*{4. Cerebras and the wafer (slides 13-15)}
\begin{itemize}[leftmargin=1.2em,itemsep=3pt]\sloppy
\item \textbf{Lie (Cerebras) (2023).} \textit{Cerebras architecture deep dive: first look inside the hardware/software co-design for deep learning (IEEE Micro 43(3); author's blog version)}. \url{https://www.cerebras.ai/blog/cerebras-architecture-deep-dive-first-look-inside-the-hw-sw-co-design-for-deep-learning}. Primary paper (author's web version, with a link to the IEEE Micro issue). The core: 48 KB of SRAM in eight banks, a four-FMAC FP16 datapath, dataflow triggering; weight streaming. The WSE-2's 7.5 PF FP16 used on slides 2, 13 and 31 is the figure analysts quote; the fact-check found no dense/sparse qualifier at its source. Slides 13 and 14.
\item \textbf{Rocki, Van Essendelft, Sharapov, Schreiber, Morrison, Kibardin, Portnoy, Dietiker, Syamlal and James (2020).} \textit{Fast stencil-code computation on a wafer-scale processor (SC20)}. arXiv:2010.03660. Primary paper. A bandwidth-bound 3D stencil from NETL runs near peak on a CS-1; the evidence that the wafer behaves as one device with PB/s of local bandwidth, and the reason labs bought the first units. Slides 13 and 14.
\item \textbf{Cerebras Systems (2024).} \textit{Cerebras Systems unveils world's fastest AI chip with 4 trillion transistors (WSE-3 press release, March 13, 2024)}. \url{https://www.cerebras.ai/press-release/cerebras-announces-third-generation-wafer-scale-engine}. Vendor press release. WSE-3: 4 trillion transistors, 900,000 cores, 44 GB SRAM, TSMC 5 nm, '125 petaflops' with no precision or sparsity stated, 1.5 TB to 1.2 PB external memory, 2,048-system clusters, 24-trillion-parameter models. Slides 13 and 14.
\item \textbf{Morgan (The Next Platform) (2024).} \textit{Cerebras goes hyperscale with third-gen waferscale supercomputers (March 14, 2024)}. \url{https://www.nextplatform.com/ai/2024/03/14/cerebras-goes-hyperscale-with-third-gen-waferscale-supercomputers/1642584}. Press, secondary source. The 8-wide FP16 SIMD of WSE-3 (4-wide before), the sparse qualifier on the 256-exaflop cluster figure, MemoryX tiers from 24 TB to 1,200 TB, and CG-1 at \$100M for 32 nodes. Slides 13 and 14.
\item \textbf{Cerebras Systems (2024).} \textit{Cerebras and G42 announce Condor Galaxy 3 (March 13, 2024)}. \url{https://cerebras.ai/press-release/cerebras-g42-announce-condor-galaxy-3}. Vendor press release. CG-3 in Dallas with 64 CS-3 at 8 'AI exaflops'; with CG-1 and CG-2 the network totals 16; the Condor Galaxy table. Slide 14.
\item \textbf{Cerebras Systems (SEC Form S-1) (2024).} \textit{Form S-1 registration statement (filed September 30, 2024)}. \url{https://www.sec.gov/Archives/edgar/data/2021728/000162828024041596/cerebras-sx1.htm}. Filing (SEC). Revenue of \$24.6M (2022), \$78.7M (2023) and \$136.4M (H1 2024); G42 was 83 percent of 2023 revenue and 87 percent of H1 2024 revenue. Slides 14 and 15.
\item \textbf{Cerebras Systems (2024).} \textit{Introducing Cerebras Inference: AI at instant speed (August 27, 2024)}. \url{https://cerebras.ai/blog/introducing-cerebras-inference-ai-at-instant-speed}. Vendor blog. The launch: Llama 3.1 8B at 1,800 tokens/s and 70B at 450 tokens/s with 16-bit weights, claimed 20 times GPU clouds; the first bar of the speed chart. Slides 15 and 33.
\item \textbf{Cerebras Systems (2024).} \textit{Cerebras triples its industry-leading inference performance (October 25, 2024)}. \url{https://www.cerebras.ai/press-release/cerebras-triples-its-industry-leading-inference-performance-setting-new-all-time-record}. Vendor press release. 2,100 tokens/s on the 70B Llama model as measured by Artificial Analysis; the second bar of the speed chart. Slides 15 and 33.
\item \textbf{Futurum Group (2026).} \textit{Cerebras S-1 teardown: is the \$23B wafer-scale IPO the end of GPU homogeneity? (April 22, 2026)}. \url{https://futurumgroup.com/insights/cerebras-s-1-teardown-is-the-23b-wafer-scale-ipo-the-end-of-gpu-homogeneity/}. Analyst note, secondary source. The April 2026 S-1: 2025 revenue \$510M, MBZUAI 62 and G42 24 percent (86 together), backlog \$24.6B, GAAP profit only from a \$363.3M one-time gain, OpenAI terms of 750 MW with an option to 2 GW. Slides 14 and 15.
\item \textbf{TechCrunch (2026).} \textit{Cerebras raises \$5.5B, kicking off 2026's IPO season with a bang (May 14, 2026)}. \url{https://techcrunch.com/2026/05/14/cerebras-raises-5-5b-kicking-off-2026s-ipo-season-with-a-bang}. Press, secondary source. Priced at \$185 (range raised twice), \$5.5B raised on 30 million shares, opened at \$385 and closed at \$311, about \$66B at the close per Yahoo Finance. Slides 15 and 34.
\item \textbf{Cerebras Systems (SEC Form 8-K) (2026).} \textit{Cerebras announces second quarter 2026 financial results (August 12, 2026)}. \url{https://www.sec.gov/Archives/edgar/data/0002021728/000162828026056186/cbrsannouncesfinancialresu.htm}. Filing (SEC). Q2 2026 GAAP revenue \$180.1M, RPO \$25.4B, \$8.6B cash, full-year core revenue guide of \$880M to \$890M. Slides 15 and 34.
\item \textbf{Cerebras Systems (2026).} \textit{Cerebras unveils CS-4: up to 30 times faster than GPU-based solutions (August 18, 2026)}. \url{https://investors.cerebras.ai/node/7401/pdf}. Vendor press release. Three WSE-3T wafers per Nexus rack, 250 'PFLOPS' and 43.2 PB/s per wafer (750 per rack), 53.5 PB/s fabric, 2 us wafer-to-wafer; over 4,400 tokens/s per user on gpt-oss-120B. Slides 15 and 33.
\item \textbf{Morgan (The Next Platform) (2026).} \textit{Cerebras overclocks WSE-3 waferscale engine to boost inference oomph in Nexus CS-4 (August 19, 2026)}. \url{https://www.nextplatform.com/compute/2026/08/19/cerebras-overclocks-wse-3-waferscale-engine-to-boost-inference-oomph-in-nexus-cs-4/5289400}. Press, secondary source. WSE-3T is the same 900,000-core die at about twice the power (clock estimated at 1.4 to 2.8 GHz, the author's estimate), three backpacks per rack, GA later in Q3 2026. Slide 15.
\end{itemize}
\subsection*{5. Dataflow and streaming machines (slides 16-20)}
\begin{itemize}[leftmargin=1.2em,itemsep=3pt]\sloppy
\item \textbf{Abts, Ross, Sparling et al. (Groq) (2020).} \textit{Think fast: a tensor streaming processor (TSP) for accelerating deep learning workloads (ISCA 2020)}. \url{https://groq.com/wp-content/uploads/2024/02/2020-Isca.pdf}. Primary paper. One core of functional slices on a 25 x 29 mm 14 nm die, 26.8 billion transistors, 220 MiB of SRAM at about 55 TiB/s, 820 TOPS at 1 GHz (900 MHz nominal), no caches or arbiters; ResNet-50 at 20.4K images/s at batch 1. Slides 5, 6 and 16.
\item \textbf{Abts, Kimmell, Ling, Kim et al. (Groq) (2022).} \textit{A software-defined tensor streaming multiprocessor for large-scale machine learning (ISCA 2022)}. \url{https://groq.com/wp-content/uploads/2024/02/GroqISCAPaper2022_ASoftwareDefinedTensorStreamingMultiprocessorForLargeScaleMachineLearning.pdf}. Primary paper. Determinism extended to the network: 8-TSP nodes, 72-TSP racks, a Dragonfly to 10,440 TSPs, 722 ns per hop, forward error correction and software replay instead of link retry. Slides 16 and 31.
\item \textbf{Groq (2024).} \textit{Groq LPU inference engine leads in first independent LLM benchmark (February 13, 2024)}. \url{https://groq.com/newsroom/groq-lpu-inference-engine-leads-in-first-independent-llm-benchmark}. Vendor press release. Llama 2 Chat 70B at 241 tokens/s on Artificial Analysis, the first independent LLM API measurement of a non-GPU chip. Slides 16 and 33.
\item \textbf{DatacenterDynamics (2025).} \textit{Nvidia to license tech from AI inference chip company Groq, hire its leadership (December 2025)}. \url{https://www.datacenterdynamics.com/en/news/nvidia-to-license-tech-from-ai-inference-chip-company-groq-hire-its-leadership/}. Press, secondary source. A non-exclusive license plus the hiring of Jonathan Ross and Sunny Madra; the \$20B figure reported by CNBC and not confirmed; Groq stays independent with GroqCloud. Slides 16, 34 and 35.
\item \textbf{Mann (The Register) (2026).} \textit{Nvidia slaps \$20B Groq tech into massive new LPX racks to speed AI response time (March 16, 2026)}. \url{https://www.theregister.com/2026/03/16/nvidia_lpx_groq_3/}. Press, secondary source. The GTC 2026 LPX rack: 256 LPUs, about 500 MB of SRAM and 150 TB/s per LPU, 1.2 PF FP8 each, shipping later in 2026; in NVIDIA's split the LPX runs the FFN and MoE half of decode and speculative drafts, while Vera Rubin does prefill, attention and the KV cache. Slides 16 and 31.
\item \textbf{Tom's Hardware (2026).} \textit{Nvidia presents Groq 3 LPX architecture and unveils its first third-party inference benchmark (August 26, 2026)}. \url{https://tomshardware.com/tech-industry/semiconductors/nvidia-presents-groq-3-lpx-architecture-and-unveils-its-first-third-party-inference-benchmark}. Press, secondary source. Hot Chips 2026: LP30 with no HBM, 350 ns chip to chip; 256 chips give 128 GB, 40 PB/s and 315 PF FP8; Gemma 4 31B at 3,431 tokens/s on a private endpoint at concurrency 1; in production. Slides 16 and 33.
\item \textbf{Prabhakar, Sivaramakrishnan, Gandhi, Du, ... Olukotun (SambaNova) (2024).} \textit{SambaNova SN40L: scaling the AI memory wall with dataflow and composition of experts (MICRO 2024)}. arXiv:2405.07518. Primary paper. The three-tier SRAM, HBM and DDR design; Samba-CoE with 150 experts and about a trillion parameters on one 8-socket node, 15 to 31 times faster model switching, 3.7 times a DGX H100. Slides 8 and 17.
\item \textbf{Prabhakar (SambaNova) (2024).} \textit{SambaNova SN40L RDU: breaking the barrier of trillion+ parameter scale gen AI computing (Hot Chips 2024 slides)}. \url{https://www.hc2024.hotchips.org/assets/program/conference/day1/48_HC2024.Sambanova.Prabhakar.final-withoutvideo.pdf}. Primary conference slides. TSMC 5 nm, 102 billion transistors, 1,040 PCUs and PMUs, 520 MB SRAM, 64 GB HBM at 1.6 TB/s, 1.5 TB DDR, 638 TFLOP/s BF16; 16 chips serve Llama 3.1 405B at 114 tokens/s. Slides 17 and 31.
\item \textbf{SambaNova Systems (2026).} \textit{SambaNova unveils fastest chip for agentic AI, collaborates with Intel, and raises \$350M+ (February 24, 2026)}. \url{https://www.businesswire.com/news/home/20260224971025/en/}. Vendor press release. SN50: five times the SN40L's compute and four times its network, 256-chip domains, Llama 3.3 70B FP8 at 895 tokens/s per user against 184 on a B200 (vendor test); Series E and the Intel collaboration. Slides 17 and 34.
\item \textbf{Jia, Tillman, Maggioni and Scarpazza (Citadel) (2019).} \textit{Dissecting the Graphcore IPU architecture via microbenchmarking}. arXiv:1912.03413. Independent measurement paper. GC2: 1,216 tiles, 304 MiB of SRAM, 30.7 TB/s measured against 45 nominal, 133 ns exchange latency, 64 GB/s IPU-Links; the only outside look at the IPU. Slides 18 and 31.
\item \textbf{Ward-Foxton (EE Times Asia) (2022).} \textit{Graphcore supercharges IPU with wafer-on-wafer (March 9, 2022)}. \url{https://www.eetasia.com/?p=36263}. Press, secondary source. Bow: a power-delivery wafer with deep-trench capacitors hybrid-bonded under the unchanged 7 nm IPU, 1.325 to 1.85 GHz, 350 TFLOP/s; Bow Pod systems from Pod16 to Pod1024; the planned 10-exaflop Good computer. (The \$149,995 list price on slide 18 was the GC200 IPU-POD16 of June 2021, before Bow.) Slide 18.
\item \textbf{Robinson (The Register) (2023).} \textit{AI chip biz Graphcore seeks capital to remain going concern (October 6, 2023)}. \url{https://www.theregister.com/2023/10/06/brit_ai_chip_outfit_graphcore/}. Press, secondary source. The 2022 accounts: revenue \$2.7M (down 46 percent), pre-tax loss \$204.6M, staff from 631 to 494, a material going-concern uncertainty; the year before the SoftBank sale. Slides 18 and 35.
\item \textbf{The Register (2024).} \textit{Tenstorrent's Blackhole chips boast 768 RISC-V cores and almost as many FLOPS (August 27, 2024)}. \url{https://www.theregister.com/2024/08/27/tenstorrent_ai_blackhole/}. Press, secondary source. The Hot Chips 36 Blackhole talk: 140 Tensix, 16 big and 752 baby RISC-V cores, 745 TFLOP/s FP8 and 372 FP16, 10 x 400 GbE, 32-chip Galaxy at 23.8 PF FP8 with 1 TB. Slides 19 and 31.
\item \textbf{ServeTheHome (2025).} \textit{Rebellions REBEL-Quad UCIe and 144GB HBM3E accelerator at Hot Chips 2025 (August 24, 2025)}. \url{https://www.servethehome.com/rebellions-rebel-quad-ucie-and-144gb-hbm3e-accelerator-at-hot-chips-2025/}. Press, secondary source. Four Samsung SF4X chiplets on CoWoS-S joined by UCIe-A, four HBM3E sites (144 GB), Llama 3.3 70B at 35.5 ms per token on a development board. Slide 20.
\item \textbf{Lam (Chips and Cheese) (2024).} \textit{FuriosaAI's RNGD at Hot Chips 2024: accelerating AI with a more flexible primitive (September 11, 2024)}. \url{https://chipsandcheese.com/p/furiosaais-rngd-at-hot-chips-2024-accelerating-ai-with-a-more-flexible-primitive}. Independent analysis, secondary source. The written record of the RNGD talk: TSMC 5 nm, 1 GHz, 150 W, 48 GB HBM3 at 1.5 TB/s, 64 slices plus a spare per unit, 256 MB SRAM; the tensor-contraction primitive of Furiosa's ISCA 2024 paper. Slide 20.
\end{itemize}
\subsection*{6. The hyperscalers' chips and the established vendors (slides 21-24)}
\begin{itemize}[leftmargin=1.2em,itemsep=3pt]\sloppy
\item \textbf{AWS Machine Learning Blog (2023).} \textit{AWS Inferentia2 builds on AWS Inferentia1 by delivering 4x higher throughput and 10x lower latency (June 13, 2023)}. \url{https://aws.amazon.com/blogs/machine-learning/aws-inferentia2-builds-on-aws-inferentia1-by-delivering-4x-higher-throughput-and-10x-lower-latency/}. Vendor blog. The Inferentia1 and 2 table: 128 against 380 TOPS INT8, 64 against 190 TFLOP/s BF16, 8 GB of DRAM at 50 GB/s against 32 GB of HBM at 820 GB/s, NeuronCore-v1 to v2. Slide 21.
\item \textbf{Barr (AWS) (2024).} \textit{Amazon EC2 Trn2 instances and Trn2 UltraServers for AI/ML training and inference are now available (December 3, 2024)}. \url{https://aws.amazon.com/blogs/aws/amazon-ec2-trn2-instances-and-trn2-ultraservers-for-aiml-training-and-inference-is-now-available}. Vendor blog. Trainium2: 8 NeuronCores, 96 GiB HBM at 2.9 TB/s, 1.3 PF FP8 dense (5.2 sparse); 16-chip instances on a 1 TB/s 2D torus; 64-chip UltraServers at 83 PF dense; 30 to 40 percent better price-performance than P5e (AWS claim). Slides 21 and 31.
\item \textbf{AWS (2025).} \textit{Amazon EC2 Trn3 UltraServers (product page)}. \url{https://aws.amazon.com/ec2/instance-types/trn3/}. Vendor page. Trainium3: AWS's first 3 nm chip, 144 GB HBM3e at 4.9 TB/s, 144-chip UltraServers at 362 PF MXFP8 with 20.7 TB and 706 TB/s, NeuronSwitch-v1 and NeuronLink-v4 (2.56 TB/s per chip in the Neuron documentation, the 2.6 TB/s on slide 21); 4.4 times Trn2 (AWS claim). Slide 21.
\item \textbf{TechCrunch (2025).} \textit{Amazon releases an impressive new AI chip and teases an Nvidia-friendly roadmap (December 2, 2025)}. \url{https://techcrunch.com/2025/12/02/amazon-releases-an-impressive-new-ai-chip-and-teases-a-nvidia-friendly-roadmap}. Press, secondary source. Trainium3 at re:Invent 2025 and the Trainium4 preview with NVIDIA NVLink Fusion support, no date given. Slide 21.
\item \textbf{ServeTheHome (2024).} \textit{Microsoft Maia 100 AI accelerator for Azure (Hot Chips 2024)}. \url{https://servethehome.com/microsoft-maia-100-ai-accelerator-for-azure}. Press, secondary source. The Hot Chips 2024 talk: TSMC N5 on CoWoS-S, 64 GB HBM2E, about 500 MB of on-chip SRAM, 16 clusters of 4 tiles, 12 x 400 GbE, 700 W TDP (about 500 W in service), 6- and 9-bit formats. Slide 22.
\item \textbf{Microsoft Azure Infrastructure Blog (2026).} \textit{Deep dive into the Maia 200 architecture (January 2026)}. \url{https://techcommunity.microsoft.com/blog/azureinfrastructureblog/deep-dive-into-the-maia-200-architecture/4489312}. Vendor blog. Maia 200: TSMC N3P, 216 GB HBM3e at 7 TB/s, 272 MB SRAM, 10.15 PF FP4 and 5.07 PF FP8, 750 W, a 2.8 TB/s Ethernet scale-up link and 6,144-chip domains; deployed only inside Azure, with no BF16 figure published. Slides 22 and 31.
\item \textbf{Meta AI (2024).} \textit{Our next-generation Meta Training and Inference Accelerator (April 10, 2024)}. \url{https://ai.meta.com/blog/next-generation-meta-training-inference-accelerator-AI-MTIA/}. Vendor blog. The v1 against v2 table: 354 TOPS INT8 dense (708 sparse), 177 TFLOP/s BF16 dense, 256 MB SRAM at 2.7 TB/s, 128 GB LPDDR5 at 204.8 GB/s, 90 W, 1.35 GHz. Slides 22 and 31.
\item \textbf{HPCwire (2026).} \textit{Meta packs gobs of HBM into homegrown AI accelerators (March 11, 2026)}. \url{https://www.hpcwire.com/2026/03/11/meta-packs-gobs-of-hbm-into-homegrown-ai-accelerators/}. Press, secondary source. The 2026 MTIA family: MTIA 400 at 1,200 W in 72-chip racks, MTIA 450 with 288 GB at 18.4 TB/s and 21 PF MX4, MTIA 500 with up to 512 GB, 27.6 TB/s, 30 PF MX4 and 1,700 W in 2027. Slide 22.
\item \textbf{Tom's Hardware (2025).} \textit{OpenAI, Broadcom to co-develop 10GW of custom AI chips (October 13, 2025)}. \url{https://www.tomshardware.com/openai-broadcom-to-co-develop-10gw-of-custom-ai-chips}. Press, secondary source. The deal terms: 10 GW of accelerators and racks, OpenAI designing and Broadcom building and deploying over Ethernet, from H2 2026 to the end of 2029. Slides 23 and 34.
\item \textbf{Kaplan (Intel) (2024).} \textit{Intel Gaudi 3 AI accelerator: architected for gen AI training and inference (Hot Chips 2024 slides)}. \url{https://www.hc2024.hotchips.org/assets/program/conference/day1/60_HC2024.Intel.RomanKaplan.Gaudi3-0826.pdf}. Primary conference slides. Two 5 nm dies, 8 MMEs and 64 TPCs, 96 MB SRAM at 12.8 TB/s, 128 GB HBM2e at 3.67 TB/s, 1,835 TFLOP/s for both BF16 and FP8 matrix, 900 W air or 1,200 W liquid, 24 x 200 GbE. Slide 24.
\item \textbf{Qualcomm (2023).} \textit{Qualcomm Cloud AI 100 Ultra product brief}. \url{https://www.qualcomm.com/content/dam/qcomm-martech/dm-assets/documents/Prod-Brief-QCOM-Cloud-AI-100-Ultra.pdf}. Vendor datasheet. 64 AI cores, 870 TOPS INT8, 576 MB SRAM, 128 GB LPDDR4x at 548 GB/s, 150 W on a PCIe card. Slide 24.
\item \textbf{Kennedy (ServeTheHome) (2025).} \textit{Qualcomm announces new integrated AI racks with 768GB cards and a 200MW AI deal (October 27, 2025)}. \url{https://www.servethehome.com/qualcomm-announces-new-integrated-ai-racks-with-768gb-cards-and-a-200mw-ai-deal/}. Press, secondary source. AI200 with 768 GB of LPDDR per card, AI250 with near-memory compute and a claimed 10 times effective bandwidth, about 160 kW liquid-cooled racks, Humain's 200 MW from 2026. Slides 24, 28 and 34.
\item \textbf{ServeTheHome (2024).} \textit{IBM Telum II processor and Spyre AI updates at Hot Chips 2024 (August 26, 2024)}. \url{https://servethehome.com/ibm-telum-ii-processor-and-spyre-ai-updates-at-hot-chips-2024}. Press, secondary source. Spyre: 32 cores with 2 MB scratchpads, Samsung 5 nm, 300 TOPS, 75 W, 128 GB LPDDR5 per card, eight cards per drawer as presented at Hot Chips (IBM's 2025 GA release gives up to 48 cards per z17 and 16 per Power11). Slide 24.
\item \textbf{Talpes, Williams and Das Sarma (Tesla) (2022).} \textit{The microarchitecture of Tesla's exa-scale computer (Hot Chips 34 slides)}. \url{https://hc34.hotchips.org/assets/program/conference/day2/Machine%20Learning/HotChips_tesla_dojo_uarch.pdf}. Primary conference slides. D1: TSMC 7 nm, 645 mm2, 354 nodes at 2 GHz, 362 TFLOP/s BF16/CFP8, 440 MB SRAM, 2 TB/s per die edge; a 5 x 5 tile of known-good dies. Slide 24.
\item \textbf{TechCrunch (2025).} \textit{Tesla Dojo: the rise and fall of Elon Musk's AI supercomputer (September 2, 2025)}. \url{https://techcrunch.com/2025/09/02/tesla-dojo-the-rise-and-fall-of-elon-musks-ai-supercomputer/}. Press, secondary source. Dojo shut down and its team disbanded in mid-August 2025, Dojo 2 called an evolutionary dead end, the \$16.5B Samsung AI6 contract weeks earlier. Read with the January 2026 TechCrunch report of the Dojo3 restart. Slides 24 and 35.
\item \textbf{TechCrunch (2026).} \textit{Elon Musk says Tesla's restarted Dojo3 will be for space-based AI compute (January 20, 2026)}. \url{https://techcrunch.com/2026/01/20/elon-musk-says-teslas-restarted-dojo3-will-be-for-space-based-ai-compute}. Press, secondary source. Five months after the team was disbanded, Musk says Dojo3 is restarted on AI5-family silicon, aimed at space-based AI compute rather than Tesla's video training; why slide 24 says Tesla paused rather than stopped. Slides 3 and 24.
\end{itemize}
\subsection*{7. China: Ascend, CloudMatrix and the policy backdrop (slides 25-27)}
\begin{itemize}[leftmargin=1.2em,itemsep=3pt]\sloppy
\item \textbf{Liao, Tu, Xia and Zhou (Huawei) (2019).} \textit{DaVinci: a scalable architecture for neural network computing (Hot Chips 31 slides)}. \url{https://www.old.hotchips.org/hc31/HC31_1.11_Huawei.Davinci.HengLiao_v4.0.pdf}. Primary conference slides. The Da Vinci core's 16 x 16 x 16 cube (4,096 FP16 or 8,192 INT8 MACs per cycle) with vector and scalar units; Ascend 910 at 256 TFLOP/s FP16, 350 W, 7+ nm EUV. Slide 25.
\item \textbf{Feldgoise and Dohmen (CSET) (2024).} \textit{Pushing the limits: Huawei's AI chip tests U.S. export controls (June 17, 2024)}. \url{https://cset.georgetown.edu/publication/pushing-the-limits-huaweis-ai-chip-tests-u-s-export-controls/}. Analyst note, secondary source. Ascend 910 against 910B: TSMC 7 nm against SMIC N+2, 30-32 against 20-25 active cores, 320 against 400 TFLOP/s FP16 (part of the gain a counting change), HBM2e at up to 1.6 TB/s and 64 GB; low SMIC yields inferred. Slides 25 and 27.
\item \textbf{Patel, Nishball, Xie et al. (SemiAnalysis) (2025).} \textit{Huawei AI CloudMatrix 384: China's answer to Nvidia GB200 NVL72 (April 16, 2025)}. \url{https://newsletter.semianalysis.com/p/huawei-ai-cloudmatrix-384-chinas-answer-to-nvidia-gb200-nvl72}. Analyst note, secondary source. 384 Ascend 910C in 16 racks with 6,912 400G LPO transceivers, about 300 PF dense BF16, 3.6 times the HBM and 2.1 times the bandwidth of an NVL72 at about 3.9 times the power (about 559 kW against 145 kW, all-in; 2.3 times per FLOP, 1.8 per TB/s, 1.1 per TB); most 910B/C dies from TSMC via Sophgo, HBM mainly Samsung. Slides 25, 26 and 31.
\item \textbf{Zuo, Lin, Deng et al. (Huawei Cloud) (2025).} \textit{Serving large language models on Huawei CloudMatrix384}. arXiv:2506.12708. Primary paper. The system as Huawei describes it: 384 Ascend NPUs and 192 Kunpeng CPUs on the Unified Bus; DeepSeek-R1 served at 6,688 prefill and 1,943 decode tokens/s per NPU. Slide 26.
\item \textbf{Xu (Huawei Connect 2025 keynote) (2025).} \textit{Groundbreaking SuperPoD interconnect: leading a new paradigm for AI infrastructure (September 18, 2025)}. \url{https://www.huawei.com/en/news/2025/9/hc-xu-keynote-speech}. Vendor keynote text. The Ascend 950PR (Q1 2026), 950DT (Q4 2026, HiZQ 2.0 HBM, 144 GB at 4 TB/s), 960 (Q4 2027) and 970 (Q4 2028); Atlas 950 (8,192 chips, 8 EF FP8) and 960 (15,488 chips, 30 EF FP8) SuperPoDs; UnifiedBus; the mainland 'will lag behind in process nodes'. The 960 and 970 dates and figures were revised in September 2026 (see the Tom's Hardware entry). Slides 25 and 26.
\item \textbf{Tom's Hardware (2026).} \textit{Huawei details AI accelerator roadmap, pulls in next-generation Ascend NPUs by quarters; FP4 performance of the Ascend 960PR doubles expectations (Huawei Connect 2026)}. \url{https://www.tomshardware.com/tech-industry/artificial-intelligence/huawei-details-ai-accelerator-roadmap-pulls-in-next-generation-ascend-npus-by-quarters-fp4-performance-of-the-ascend-960pr-doubles-expectations}. Press, secondary source. The September 18, 2026 revision: 960DT in Q1 2027 and 960PR in Q3 2027 (2 PF FP8, 8 PF FP4) with a 2.2 TB/s link; 970 in 2028 at 3.6 PF FP8, 14 PF FP4 and 4.4 TB/s; a 980 in 2029. Slide 25.
\item \textbf{The Register (2025).} \textit{DeepSeek's R2 training on Huawei Ascend reportedly failed (August 14, 2025)}. \url{https://www.theregister.com/2025/08/14/dodgy_huawei_deepseek/}. Press, secondary source (relaying the Financial Times). DeepSeek never completed a successful R2 training run on Ascend, citing unstable chips, slow interconnect and immature software, and went back to H20s for training. Slide 25.
\item \textbf{Recode China AI (2025).} \textit{Who will fill Nvidia's AI chip void in China?}. \url{https://recodechinaai.substack.com/p/who-will-fill-nvidias-ai-chip-void}. Analyst newsletter, secondary source. Hanguang 800 (12 nm, 17 billion transistors, 820 TOPS), Kunlun 2 (7 nm, 128 TFLOP/s FP16, about 120 W), Cambricon MLU370 (96 TF) and MLU590 (345 TF, FP8). Slide 27.
\item \textbf{Caixin Global (2025).} \textit{Baidu unveils ambitious AI chip roadmap targeting 1 million-card cluster by 2030 (November 14, 2025)}. \url{https://www.caixinglobal.com/2025-11-14/baidu-unveils-ambitious-ai-chip-roadmap-targeting-1-million-card-cluster-by-2030-102382580.html}. Press, secondary source. Kunlun M100 (2026, inference) and M300 (2027, training), the 30,000-card P800 cluster of April 2025, 512-card nodes in 2026, a million cards by 2030. Slide 27.
\item \textbf{TrendForce (2025).} \textit{Alibaba T-Head PPU lands China Unicom deal, reportedly matches Nvidia H20 performance (September 17, 2025)}. \url{https://www.trendforce.com/news/2025/09/17/news-alibaba-t-head-ppu-lands-china-unicom-deal-reportedly-matches-nvidia-h20-performance/}. Analyst press, secondary source. The PPU: 96 GB HBM2e, 700 GB/s chip to chip, domestic 7 nm; a 16,384-card China Unicom order. Slide 27.
\item \textbf{TrendForce (2026).} \textit{China GPU race intensifies: Cambricon turns profitable in 2025, Moore Threads and MetaX narrow losses (March 9, 2026)}. \url{https://www.trendforce.com/news/2026/03/09/news-china-gpu-race-intensifies-cambricon-turns-profitable-in-2025-moore-threads-metax-narrow-losses/}. Analyst press, secondary source. Cambricon's 2025 revenue of RMB 6.5B (up 453 percent) and net profit of RMB 2.06B after a 2024 loss; Moore Threads and MetaX revenue. Slide 27.
\item \textbf{NVIDIA (SEC Form 8-K) (2025).} \textit{Form 8-K: H20 export license requirement (April 9, 2025, signed April 15)}. \url{https://www.sec.gov/Archives/edgar/data/1045810/000104581025000082/nvda-20250409.htm}. Filing (SEC). The license requirement for the H20 notified April 9, made indefinite April 14, with up to about \$5.5B of charges; the April 2025 step on the policy timeline. Slide 27.
\item \textbf{Morrison Foerster (2025).} \textit{BIS issues interim final rule: framework for artificial intelligence diffusion (February 4, 2025)}. \url{https://www.mofo.com/resources/insights/250204-bis-issues-interim-final-rule}. Law-firm summary, secondary source. The January 13, 2025 rule and its three country tiers (authorized, restricted including China, all others), with the May 2025 rescission noted. Slide 27.
\item \textbf{Kirkland and Ellis (2025).} \textit{BIS rescission of the Biden administration's AI diffusion rule (May 27, 2025)}. \url{https://www.kirkland.com/publications/kirkland-alert/2025/05/bis-rescission-of-the-biden-administration}. Law-firm summary, secondary source. The May 13, 2025 rescission and the GP10 guidance naming Ascend 910B, 910C and 910D as presumptively made in violation of US controls. Slide 27.
\item \textbf{Associated Press (via The Atlanta Journal-Constitution) (2025).} \textit{Nvidia, AMD to pay 15\% of China chip sale revenue to US government (August 2025)}. \url{https://www.ajc.com/news/2025/08/nvidia-amd-to-pay-15-of-china-chip-sale-revenue-to-us-government-2}. Press, secondary source. H20 and MI308 sales resumed with 15 percent of China revenue to the US government as a license condition (FT, confirmed to AP by an official). Slide 27.
\item \textbf{Bloomberg (2025).} \textit{China agency orders firms to stop buying Nvidia AI chip, FT says (September 17, 2025)}. \url{https://www.bloomberg.com/news/articles/2025-09-17/china-agency-orders-firms-to-stop-buying-nvidia-ai-chip-ft-says}. Press, secondary source. The Cyberspace Administration of China told Alibaba and ByteDance to stop testing and cancel orders for the RTX Pro 6000D; the last point on the policy timeline. Slide 27.
\end{itemize}
\subsection*{8. Emerging designs: in-memory, hardwired, photonic, neuromorphic (slides 28-30)}
\begin{itemize}[leftmargin=1.2em,itemsep=3pt]\sloppy
\item \textbf{SiliconANGLE (2025).} \textit{Chip startup d-Matrix raises \$275M to speed up inference with memory-compute (November 12, 2025)}. \url{https://siliconangle.com/2025/11/12/chip-startup-d-matrix-raises-275m-speed-inference-memory-compute/}. Press, secondary source. The \$275M round at \$2B; Corsair with 2 GB of SRAM per card, part of it repurposed as multipliers, 256 GB LPDDR5 behind, '9,600 TOPS' in MXINT4 (vendor); Raptor (2026) to stack DRAM on 4 nm (vendor plan). For the chiplet count (two chips of four 6 nm chiplets per card) use the ServeTheHome Hot Chips 2025 entry. Slides 3 and 28.
\item \textbf{ServeTheHome (2025).} \textit{d-Matrix presents Corsair, an in-memory computing architecture for inference, at Hot Chips 2025}. \url{https://www.servethehome.com/d-matrix-presents-corsair-an-in-memory-computing-architecture-for-inference-at-hot-chips-2025/}. Press, secondary source. The Hot Chips 2025 talk: a Corsair card is two chips of four chiplets each, eight TSMC 6 nm chiplets with 2 GB of SRAM in all, 256 GB of LPDDR5X, 275 to 550 W; the Pavehawk 3D-DRAM test vehicle. Slide 28.
\item \textbf{Gomez-Luna, El Hajj, Fernandez, Giannoula, Oliveira and Mutlu (2021).} \textit{Benchmarking a new paradigm: an experimental analysis of a real processing-in-memory architecture}. arXiv:2105.03814. Primary paper. The first independent characterization of commercial processing-in-DRAM (UPMEM DIMMs with 128 DRAM processing units each): where moving compute to the memory pays and where it does not. Slide 28.
\item \textbf{TechCrunch (2023).} \textit{AI chip startup Mythic rises from the ashes with \$13M, new CEO (March 9, 2023)}. \url{https://techcrunch.com/2023/03/09/ai-chip-startup-mythic-rises-from-the-ashes-with-13m-new-ceo}. Press, secondary source. The analog flash-in-memory company nearly out of money in November 2022, revived with \$13M and a defense focus (before the \$125M round of December 2025, separate entry). Slides 28 and 35.
\item \textbf{Mythic (BusinessWire) (2025).} \textit{Mythic raises \$125 million (December 17, 2025)}. \url{https://www.businesswire.com/news/home/20251217546249/en/}. Vendor press release. An oversubscribed \$125M round three years after the near-shutdown of November 2022; the analog in-memory company back in business for edge inference. Slides 28 and 35.
\item \textbf{HPCwire (2025).} \textit{Untether AI team acquired by AMD; product support discontinued (June 5, 2025)}. \url{https://www.hpcwire.com/off-the-wire/untether-ai-team-acquired-by-amd-product-support-discontinued/}. Press release, secondary source. Untether ends supply and support of speedAI and its SDK as AMD takes the team. Slides 28, 29 and 35.
\item \textbf{CNX Software (2026).} \textit{Taalas HC1 hardwired Llama 3.1 8B AI accelerator delivers up to 17,000 tokens/s (February 22, 2026)}. \url{https://www.cnx-software.com/2026/02/22/taalas-hc1-hardwired-llama-3-1-8b-ai-accelerator-delivers-up-to-17000-tokens-s/}. Press, secondary source. HC1: Llama 3.1 8B in the metal layers, TSMC 6 nm, 815 mm2, 53 billion transistors, mixed 3- and 6-bit parameters, close to 17,000 tokens/s for one user. Slide 29.
\item \textbf{NYU Shanghai RITS (2026).} \textit{AMD acquires Taalas, the startup etching LLMs into silicon (August 2026)}. \url{https://rits.shanghai.nyu.edu/ai/amd-acquires-taalas-the-startup-etching-llms-into-silicon/}. Secondary source. AMD's definitive agreement of August 6, 2026 to buy Taalas, terms undisclosed, closing expected in Q4 2026. Slides 29 and 34.
\item \textbf{GadgetScout (2026).} \textit{Etched reaches \$10.3bn valuation with SK Hynix backing (August 2, 2026)}. \url{https://www.gadgetscout.co.uk/articles/etched-reaches-10-3bn-valuation-with-sk-hynix-backing.html}. Press, secondary source. A \$300M Series C at \$10.3B (July 23, 2026), followed in August by \$700M at \$21B (separate entry); Sohu on TSMC N4P with 144 GB, A0 silicon in June 2026, no independent benchmark; the 500,000 tokens/s claim for an eight-chip server, which is aggregate throughput, not one user's speed. Slide 29.
\item \textbf{AI Weekly (2026).} \textit{Etched doubles to \$21B in \$700M round led by Jane Street (August 18, 2026)}. \url{https://aiweekly.co/alerts/etched-doubles-to-21b-in-700m-round-led-by-jane-street}. Press, secondary source. \$700M at \$21B three weeks after the \$300M round at \$10.3B, so \$1B in July and August 2026 with still no independent benchmark of Sohu; the money-before-silicon row on slide 29. Slides 29 and 34.
\item \textbf{Implicator (2026).} \textit{Positron raises \$875 million at \$5 billion before its Asimov chip tapes out (September 10, 2026)}. \url{https://www.implicator.ai/positron-raises-875-million-at-5-billion-before-its-asimov-chip-tapes-out.md}. Press, secondary source. Atlas on Agilex 7 FPGAs with HBM and DDR5 (50+ racks at Oracle); Asimov on N3P with 288 GB to 2.3 TB of LPDDR5X, tape-out end 2026, volume H2 2027, figures from simulation. Slide 29.
\item \textbf{Marvell Technology (SEC Form 8-K, Ex. 99.1) (2025).} \textit{Marvell to acquire Celestial AI (December 2, 2025)}. \url{https://www.sec.gov/Archives/edgar/data/1835632/000119312525305289/d34367dex991.htm}. Filing (SEC). \$1.0B cash and \$2.25B in stock up front, up to \$2.25B more on revenue milestones; the Photonic Fabric chiplet at 16 Tb/s, 3D co-packaged with XPUs. Slides 29 and 30.
\item \textbf{Lightmatter (2026).} \textit{Lightmatter product summary (llms.txt)}. \url{https://lightmatter.co/llms.txt}. Vendor page. Passage M1000 (114 Tb/s on a 4,000 mm2 photonic interposer) and the L200 co-packaged optics; Envise placed in a 2030-plus full photonic stack, not cancelled. For L200 rates use the DCD entry: 32 to 34 Tb/s for the L200 and 64 Tb/s for the L200X per chip, more than 200 Tb/s only per multi-chip package. Slide 30.
\item \textbf{DatacenterDynamics (2025).} \textit{Lightmatter unveils M1000 and L200 Passage photonic interconnects}. \url{https://datacenterdynamics.com/en/news/lightmatter-unveils-m1000-and-l200-passage-photonic-interconnects}. Press, secondary source. The M1000 interposer at 114 Tb/s over 256 fibers; the L200 co-packaged optical engine at 32 to 34 Tb/s and the L200X at 64 Tb/s per chip, with more than 200 Tb/s claimed only per multi-chip XPU package. Slide 30.
\item \textbf{TechPowerUp (2024).} \textit{Intel builds world's largest neuromorphic system to enable more sustainable AI (April 17, 2024)}. \url{https://Techpowerup.com/321645/intel-builds-worlds-largest-neuromorphic-system-to-enable-more-sustainable-ai}. Vendor release (republished). Hala Point at Sandia: 1,152 Loihi 2 chips, up to 1.15 billion neurons, up to 20 petaops, 2,600 W maximum. Slide 30.
\item \textbf{Gonzalez, Huang, Kelber, Nazeer et al. (TU Dresden) (2024).} \textit{SpiNNaker2: a large-scale neuromorphic system for event-based and asynchronous machine learning}. arXiv:2401.04491. Primary paper. 152 Arm cores per chip with local SRAM, 48-chip boards, a multi-million-core machine: the MIMD counterpoint to Loihi 2 among neuromorphic research systems. Slide 30.
\item \textbf{Modha (IBM) (2023).} \textit{NorthPole: neural inference at the frontier of energy, space, and time (October 19, 2023; companion to Modha et al., Science 382(6668))}. \url{https://modha.org/2023/10/northpole-neural-inference-at-the-frontier-of-energy-space-and-time}. Author's page for the Science paper, with a free full-text link. 12 nm, 22 billion transistors, 800 mm2, 256 cores, 224 MB on chip and no off-chip memory; the near-memory limit case. Slides 28 and 30.
\item \textbf{Modha (IBM) (2024).} \textit{Breakthrough low-latency, high-energy-efficiency LLM inference performance using NorthPole (September 26, 2024, IEEE HPEC)}. \url{https://modha.org/2024/09/breakthrough-low-latency-high-energy-efficiency-llm-inference-performance-using-northpole}. Author's summary of an HPEC 2024 paper. A 3B model on 16 NorthPole cards in 2U at 28,356 tokens/s, under 1 ms per token per user, 72.7 times the energy efficiency of 'the next lowest-latency GPU' (IBM's measurement; IBM's blog does not name the GPU, which slide 30 reads as an H100). Slide 30.
\item \textbf{DeBole, Appuswamy, McGlohon et al. (IBM) (2025).} \textit{A scalable NorthPole system with end-to-end vertical integration for low-latency and energy-efficient LLM inference}. arXiv:2511.15950. Primary paper. 288 NorthPole cards in 18 2U servers: 115 peta-ops at INT4, 3.7 PB/s, 30 kW in one rack, serving 8B models to 28 users at 2.8 ms per token or a 70B model; the step from prototype to rack. Slide 30.
\end{itemize}
\subsection*{9. Measured: specifications, MLPerf, speed and deals (slides 31-34)}
\begin{itemize}[leftmargin=1.2em,itemsep=3pt]\sloppy
\item \textbf{Peng, Choudhury, Ding, Barker, Geng and Li (UConn, PNNL, Rochester) (2023).} \textit{Evaluating emerging AI/ML accelerators: IPU, RDU, and NVIDIA/AMD GPUs}. arXiv:2311.04417. Independent evaluation paper. A side-by-side specification table (SN10: 40 billion transistors, 640 PCUs and PMUs, 320 MB SRAM, 325 TFLOP/s, 1.5 TB DRAM) and measured kernels on the IPU, the RDU and GPUs. Slides 17 and 31.
\item \textbf{Emani, Foreman, Sastry, Xie et al. (Argonne) (2023).} \textit{A comprehensive performance study of large language models on novel AI accelerators}. arXiv:2310.04607. Independent evaluation paper. GPT-2 XL, GenSLM and sparse GPT-3 runs on A100, SN30, CS-2, Bow Pod64, Gaudi2 and MI250 at the ALCF; the closest thing to a neutral LLM benchmark for the dataflow machines, which skipped MLPerf. Slides 32 and 35.
\item \textbf{Mattson et al. (2020).} \textit{MLPerf Training benchmark (MLSys 2020)}. arXiv:1910.01500. Primary paper. The time-to-train method, the reference models and the closed and open divisions behind every Training round on the participation chart. Slide 32.
\item \textbf{MLCommons (2022).} \textit{MLPerf Training results repositories (training\_results\_v0.5 to v6.0; this link is v2.1, November 2022)}. \url{https://github.com/mlcommons/training_results_v2.1}. Benchmark consortium, open data. One GitHub repository per round, with a directory per submitter and systems/*.json per system; the primary record for the Training dots on slide 32. v2.1 is the round with neither Google nor any IPU entry; change the version in the URL for other rounds. Slides 12, 18 and 32.
\item \textbf{MLCommons (2024).} \textit{MLPerf Inference results repositories (inference\_results\_v0.5 to v6.0; this link is v4.0, March 2024)}. \url{https://github.com/mlcommons/inference_results_v4.0}. Benchmark consortium, open data. The primary record for the Inference circles on slide 32: TPU entries only in v0.5, v3.1, v4.0 and v5.0 (Google's v4.1, v5.1 and v6.0 entries are NVIDIA systems); v4.0 is Qualcomm's last round and has TPU v5e and Gaudi 2. Slides 24 and 32.
\item \textbf{MLCommons (2025).} \textit{MLPerf Training v5.0 results (June 4, 2025)}. \url{https://mlcommons.org/2025/06/mlperf-training-v5-0-results}. Benchmark consortium, open. Twenty submitters with AMD new; TPU Trillium the one non-GPU processor among the twelve; Google's last Training round before it skipped 5.1. Slide 32.
\item \textbf{MLCommons (2026).} \textit{Supplemental results discussion for MLPerf Training v6.0 (June 2026)}. \url{https://mlcommons.org/wp-content/uploads/2026/06/MLPerf-Training-v6.0-Supplemental-Discussion-UNDER-EMBARGO-UNTIL-6_16_26-8_00-AM-PT-1.pdf}. Benchmark consortium, open. The submitters' own notes for Training 6.0: Google is listed as a submitter but only AMD and NVIDIA accelerators are named; the basis for 'no TPU result in 6.0'. Slides 12 and 32.
\item \textbf{Artificial Analysis (2026).} \textit{gpt-oss-120B (high): API provider benchmarks (live page, read November 2026)}. \url{https://artificialanalysis.ai/models/gpt-oss-120b/providers}. Independent benchmark service. Sixteen providers of one model: Cerebras first answer in about 2.1 s at \$0.39 per million tokens, the highest price listed; Groq and SambaNova slower to first answer. Read it with a date; it changes. Slide 33.
\item \textbf{Patel and Nishball (SemiAnalysis) (2024).} \textit{Groq inference tokenomics: speed, but at what cost? (February 21, 2024)}. \url{https://newsletter.semianalysis.com/p/groq-inference-tokenomics-speed-but}. Analyst note, secondary source. 576 chips (8 racks of 9 servers of 8) to serve Mixtral, 230 MB SRAM and about 725 mm2 on GlobalFoundries 14 nm per chip, against an H100 with about \$1,150 of HBM; the cost-per-token argument. Slides 8, 16 and 33.
\item \textbf{South China Morning Post (2025).} \textit{Google, Anthropic unveil cloud deal worth tens of billions of dollars (October 24, 2025)}. \url{https://www.scmp.com/tech/big-tech/article/3330143/google-anthropic-unveil-cloud-deal-worth-tens-billions-dollars}. Press, secondary source. Up to one million TPUs for Anthropic and over a gigawatt online in 2026; the first row of the deals table. Slides 12 and 34.
\item \textbf{Analytics India Magazine (2025).} \textit{Broadcom reveals \$21 billion Google TPU order from Anthropic (December 12, 2025)}. \url{https://analyticsindiamag.com/ai-news-updates/broadcom-reveals-21-billion-google-tpus-order-from-anthropic/}. Press, secondary source. From Broadcom's Q4 FY2025 call: about \$10B of TPU orders in one quarter and \$11B more for delivery in late 2026, the customer named as Anthropic. Slides 12, 23 and 34.
\item \textbf{Dataconomy (2026).} \textit{Meta signs multibillion-dollar deal to rent Google TPUs for AI training (February 27, 2026)}. \url{https://dataconomy.com/2026/02/27/meta-signs-multibillion-dollar-deal-to-rent-google-tpus-for-ai-training/}. Press, secondary source (relaying Reuters, February 26, 2026). Meta leasing TPUs through Google Cloud, with outright purchases discussed for the following year. Slides 12 and 34.
\item \textbf{SiliconANGLE (2025).} \textit{AWS opens \$11B Project Rainier data center campus built for Anthropic (October 29, 2025)}. \url{https://siliconangle.com/2025/10/29/aws-opens-11b-project-rainier-data-center-campus-built-anthropic/}. Press, secondary source. Nearly 500,000 Trainium2 in seven Indiana buildings, over a million targeted by the end of 2025 (the 1M+ AWS reported in December 2025 was fleet-wide, with Anthropic on over 500,000), more than 2.2 GW planned. Slides 21 and 34.
\item \textbf{Broadcom (SEC Form 8-K, Ex. 99) (2025).} \textit{Broadcom announces fourth quarter and fiscal year 2025 financial results (December 11, 2025)}. \url{https://www.sec.gov/Archives/edgar/data/1730168/000173016825000116/avgo-11022025x8kxex99.htm}. Filing (SEC). AI semiconductor revenue up 74 percent in Q4 FY2025 and \$8.2B guided for Q1 FY2026, driven by custom accelerators and Ethernet switches; the XPU business in one number. Slides 23 and 34.
\item \textbf{TrendForce (2025).} \textit{Strong demand from CSPs and sovereign cloud to drive over 20\% growth in AI server shipments by 2026 (October 30, 2025)}. \url{https://www.trendforce.com/presscenter/news/20251030-12762.html}. Analyst press release, secondary source. NVIDIA at about 70 percent of the 2025 AI-chip market and custom-ASIC shipments expected to pass GPU shipments in 2026. Slides 34 and 36.
\item \textbf{Crunchbase News (2026).} \textit{Record-high US semiconductor startup funding in 2025 (January 6, 2026)}. \url{https://news.crunchbase.com/venture/record-high-us-semiconductor-startup-funding-eoy-2025}. Press (data provider), secondary source. About \$6.2B of US semiconductor startup equity in 2025, up 85 percent, nearly half in three rounds: Cerebras \$1.1B, PsiQuantum \$1B, Groq \$750M. Slide 34.
\end{itemize}
\subsection*{10. Wrap-up (slides 35-37)}
\begin{itemize}[leftmargin=1.2em,itemsep=3pt]\sloppy
\item \textbf{Emani, Xie, Raskar, Sastry, ... Papka (Argonne) with Cerebras, SambaNova, Graphcore and Groq (2022).} \textit{A comprehensive evaluation of novel AI accelerators for deep learning workloads (PMBS 2022, SC22 workshops)}. \url{https://sc22.supercomputing.org/proceedings/workshops/workshop_pages/ws_pmbsf120.html}. Workshop paper (proceedings page). The first peer-reviewed side-by-side of the four dataflow machines on the same science workloads at the ALCF testbed; what the programming effort looked like from a user's side. Slides 35 and 37.
\item \textbf{Lattner, Amini, Bondhugula, Cohen, Davis, Pienaar, Riddle, Shpeisman, Vasilache and Zinenko (2021).} \textit{MLIR: scaling compiler infrastructure for domain specific computation (CGO 2021)}. arXiv:2002.11054. Primary paper. The compiler substrate under StableHLO, IREE, TT-Forge and most startup stacks; why a new chip in 2026 needs less compiler work than one in 2017. Slides 19 and 35.
\item \textbf{Chen, Moreau, Jiang, Zheng et al. (2018).} \textit{TVM: an automated end-to-end optimizing compiler for deep learning (OSDI 2018)}. arXiv:1802.04799. Primary paper. The auto-tuning tensor compiler that several NPU vendors built on; the alternative to writing every kernel by hand. Slide 35.
\item \textbf{Google Open Source Blog (2023).} \textit{OpenXLA is ready to accelerate and simplify ML development (March 8, 2023)}. \url{https://opensource.googleblog.com/2023/03/openxla-is-ready-to-accelerate-and-simplify-ml-development.html}. Vendor blog. XLA, StableHLO and IREE opened up, with AWS, Cerebras, Graphcore, Intel, Meta and NVIDIA among the co-developers; when the TPU's compiler became a shared target. Slides 12 and 35.
\item \textbf{vLLM project (2025).} \textit{Introducing vLLM hardware plugin, best practice from Ascend NPU (May 12, 2025)}. \url{https://vllm.ai/blog/2025-05-12-hardware-plugin}. Open-source project blog. The platform-plugin mechanism (RFC 11162) used by vllm-ascend and vllm-spyre; how non-NVIDIA silicon reaches production LLM serving without forking. Slides 25 and 35.
\item \textbf{Quach (The Register) (2020).} \textit{Two startups enter, one leaves: Intel kills off much-delayed Nervana AI training chip, pushes on with Habana (January 31, 2020)}. \url{https://www.theregister.com/2020/01/31/intel_kills_spring_crest/}. Press, secondary source. Intel ends the NNP-T (Spring Crest) weeks after buying Habana; Nervana had cost \$350M in 2016; an early example of an acquirer choosing between two accelerator lines. Slides 24 and 35.
\item \textbf{MIPS (PR Newswire) (2021).} \textit{Wave Computing and MIPS emerge from Chapter 11 bankruptcy (March 1, 2021)}. \url{https://prnewswire.com/news-releases/wave-computing-and-mips-emerge-from-chapter-11-bankruptcy-301237051.html}. Vendor press release. The dataflow-chip company Wave leaves bankruptcy as MIPS under Tallwood (\$61M winning bid), turning to RISC-V; the first of the dataflow startups to fail. Slides 5 and 35.
\item \textbf{corsix.org (2025).} \textit{Esperanto lives on at AINekko (October 29, 2025)}. \url{https://www.corsix.org/content/esperanto-lives-on}. Independent analysis, secondary source. Esperanto's July 2025 shutdown (staff poached at up to four times the pay) and the ET-SoC-1's 1,000-plus RISC-V cores and 32 GB of LPDDR4x, now open-sourced by AINekko. Slide 35.
\end{itemize}
%=====
\clearpage
\printindex
\end{document}
